\documentclass[10pt]{article}
\usepackage[verbose, a4paper, hmargin=2.5cm, vmargin=2.5cm]{geometry}

\usepackage{fontspec}
\usepackage{ctex}
\usepackage{paratype}

\usepackage{amsmath}
\usepackage{amssymb}
\usepackage{amsfonts}
\usepackage{esint}
\usepackage{stmaryrd}

\usepackage[mathbf=sym]{unicode-math}
\setmainfont{STIX Two Text}
\setmathfont{STIX Two Math}

\setsansfont{Arial}
\setmonofont{Courier New}

\usepackage{makecell}
\usepackage{wasysym}
\usepackage{pifont}
\usepackage[export]{adjustbox}
\usepackage{graphicx}
\usepackage{mdframed}
\usepackage{booktabs,array,multirow}
\usepackage{tabularx}
\usepackage{hyperref}
\hypersetup{colorlinks=true, linkcolor=blue, filecolor=magenta, urlcolor=cyan,}
\urlstyle{same}
\usepackage[most]{tcolorbox}
\definecolor{mygray}{RGB}{240,240,240}
\tcbset{
  colback=mygray,
  boxrule=0pt,
}
\graphicspath{ {./images/} }
\newcommand{\HRule}{\begin{center}\rule{0.9\linewidth}{0.2mm}\end{center}}
\newcommand{\customfootnote}[1]{
  \let\thefootnote\relax\footnotetext{#1}
}
\newcommand{\lt}{\mathrel{<}}
\newcommand{\gt}{\mathrel{>}}
\newcommand*\circled[1]{\tikz[baseline=(char.base)]{\node[shape=circle,draw,inner sep=1.5pt] (char) {\small #1};}}
\setcounter{MaxMatrixCols}{256}
\begin{document}
J.-P. Serre

\section*{Truiter initiancing}

线性

表示

有限群

Springer-Verlag 纽约 海德堡 柏林

数学研究生教材

\begin{center}
\includegraphics[max width=0.2\textwidth]{images/bo_d6kn08k601uc73e1e9r0_1_978_306_140_117_0.jpg}
\end{center}
\hspace*{3em} 

编辑委员会

F. W. Gehring

P. R. Halmos

主编

C. C. Moore

\section*{Jean-Pierre Serre 线性表示论(有限群)}

法文原著翻译

Leonard L. Scott

\begin{center}
\includegraphics[max width=0.2\textwidth]{images/bo_d6kn08k601uc73e1e9r0_2_1091_1865_140_188_0.jpg}
\end{center}
\hspace*{3em} 

Springer-Verlag

纽约 海德堡 柏林

让-皮埃尔·塞雷 伦纳德·L·斯科特

法兰西学院 弗吉尼亚大学

代数与几何讲席 数学系

法国巴黎 弗吉尼亚州夏洛茨维尔 22903

\section*{编辑委员会}

P. R. 哈尔莫斯 C. C. 摩尔 F. W. 盖尔林

常务编辑 加州大学伯克利分校 密歇根大学

加州大学 数学系 数学系

数学系 加州伯克利 94720 密歇根安娜堡 48104

加州圣塔芭芭拉 93106

美国数学学会主题分类:主要 20Cxx

次要 20Dxx, 12B15, 16A50, 16A54

国会图书馆编目资料

塞雷，让-皮埃尔。

有限群的线性表示。

(数学研究生教材；42)

法文版 Représentations linéaires des groupes finis 第2版 之译本

含参考书目和索引。

1. 群的表示。2. 有限群。

I. 书名。II. 丛书。

QA171.S5313 \({512}^{\prime }{.2}\;{76} - {12585}\)

法文版之译本

Représentations linéaires des groupes finis，巴黎:Hermann 1971

版权所有。

 \textcopyright  1977 纽约施普林格出版社(Springer-Verlag, New York Inc.)

1977年精装第一版之平装重印本

987654321

ISBN 978-1-4684-9460-0 ISBN 978-1-4684-9458-7 (电子书)

DOI 10.1007/978-1-4684-9458-7

\section*{前言}

本书由三部分组成，层次与目的各异:

第一部分最初为量子化学家而作。它描述了由 Frobenius 建立的线性表示与特征标之间的对应关系。这是一个基础性结果，在数学以及量子化学或物理中经常使用。我力求给出尽可能初等的证明，仅用到群的定义和线性代数的基本知识。例子(第5章)选自对化学家有用的内容。

第二部分是1966年为高师二年级学生开设的课程。它在以下方面补充了第一部分:

(a)表示的次数与特征值的整性性质(第6章)；

(b)诱导表示、Artin 与 Brauer 定理及其应用(第7–11章)；

(c)有理性问题(第12章与第13章)。

所用方法是线性代数(比第一部分含义更广):群代数、模、非交换张量积、半单代数。

第三部分是对Brauer理论的导论:从特征0到特征 \(p\) (及反向)的转换。我自由地使用了阿贝尔范畴的语言(投射模、Grothendieck群)，这很适合此类问题。主要结果有:

(a)分解同态是满的:所有特征 \(p\) 下的不可约表示都可以“虚拟地”(即在适当的Grothendieck群中)提升到特征0。

(b)Fong–Swan 定理，它使前述陈述中的“虚拟地”一词可被去除，前提是所考虑的群是 \(p\) -可解的。

我还给出了一些关于Artin表示的应用。

我很高兴致谢:

Gaston Berthier 与 Josiane Serre，授权我重录第一部分，该部分最初作为他们著作《量子化学》的附录撰写；

Yves Balasko，他根据一些讲义起草了第二部分的初稿；

Alexandre Grothendieck，授权我重录第三部分，该部分最早见于他1965/66年度的I.H.E.S.代数几何研讨会讲义。

\section*{目录}

第一部分

表示与特征 1

1 线性表示概论 3

1.1 定义 3

1.2 基本例子 4

1.3 子表示 5

1.4 不可约表示 7

1.5 两个表示的张量积 7

1.6 对称平方与交替平方 9

2 表示的特征 10

2.1 表示的特征值(character) 10

2.2 舒尔引理；基本应用 13

2.3 特征的正交关系 15

2.4 正则表示的分解 17

2.5 不可约表示的个数 18

2.6 表示的规范分解 21

2.7 表示的显式分解 23

3 子群、直积、诱导表示 25

3.1 阿贝尔子群 25

3.2 两群的直积 26

3.3 诱导表示 28

4 紧群 32

4.1 紧群 32

4.2 紧群上的不变测度 32

4.3 紧群的线性表示 33

vii

5 例子 35

5.1 循环群 \({C}_{n}\) 35

5.2 群 \({C}_{\infty }\) 36

5.3 二面体群 \({\mathrm{D}}_{n}\) 36

5.4 群 \({D}_{nh}\) 38

5.5 群 \({D}_{\infty }\) 39

5.6 群 \({D}_{\omega h}\) 40

5.7 交错群 \({\mathfrak{A}}_{4}\) 41

5.8 对称群 \({\mathfrak{S}}_{4}\) 42

5.9 正方体的群 43

参考书目:第一部分 44

第二部分

零特征下的表示论 45

6 群代数 47

6.1 表示与模 47

6.2 C[G] 的分解 48

6.3 \(\mathbf{C}\left\lbrack  \mathrm{G}\right\rbrack\) 的中心 50

6.4 整数的基本性质 50

6.5 特征的整性性质。应用 52

7 诱导表示；马基不等式准则 54

7.1 诱导 54

7.2 诱导表示的特征；互易公式 55

7.3 对子群的限制 58

7.4 马基的不可约性准则 59

诱导表示的 8 个例子 61

8.1 正规子群；对不可约表示度数的应用 61

不可约表示

8.2 被阿贝尔群的半直积 62

8.3 若干有限群类的回顾 63

8.4 西洛定理 65

8.5 超可解群的线性表示 66

9 阿廷定理 68

9.1 环 R(G) 68

9.2 阿廷定理的陈述 70

9.3 第一种证明 70

9.4 (i) 的第二种证明 \(\Rightarrow\) (ii) 72

10 布劳尔的一条定理 74

10.1 \(p\) -正规元素； \(p\) -基本子群 74

10.2 源于 \(p\) -基本子群的诱导特征 75

子群

10.3 构造特征 76

10.4 定理 18 与 \({18}^{\prime }\) 的证明 78

10.5 布劳尔定理 78

目录

11 布劳尔定理的应用 81

11.1 特征的刻画 81

11.2 弗罗贝nius 的一个定理 83

11.3 布劳尔定理的一个逆命题 85

11.4 \(A \otimes  R\left( G\right)\) 的谱 86

12 有理性问题 90

12.1 环 \({R}_{K}\left( G\right)\) 与 \({R}_{K}\left( G\right)\) 90

12.2 舒尔指标 92

12.3 在圆分域上的实现性 94

12.4 \({R}_{K}\left( G\right)\) 的秩 95

12.5 阿廷定理的推广 96

12.6 布劳尔定理的推广 97

12.7 定理28的证明 99

13 有理性问题:例子 102

13.1 域 \(\mathbf{Q}\) 102

13.2 域 \(\mathbf{R}\) 106

参考文献:第二部分 111

第三部分

布劳尔理论导论 113

14 群 \({R}_{K}\left( G\right) ,{R}_{k}\left( G\right)\) 和 \({P}_{k}\left( G\right)\) 115

14.1 环 \({R}_{K}\left( G\right)\) 与 \({R}_{k}\left( G\right)\) 115

14.2 群 \({P}_{k}\left( G\right)\) 和 \({P}_{A}\left( G\right)\) 116

14.3 \({P}_{k}\left( G\right)\) 的结构 116

14.4 \({\mathrm{P}}_{\mathrm{A}}\left( \mathrm{G}\right)\) 的结构 118

14.5 对偶性 120

14.6 标量扩张 122

15 cde 三角形 124

15.1 \(c : {\mathrm{P}}_{k}\left( \mathrm{G}\right)  \rightarrow  {\mathrm{R}}_{k}\left( \mathrm{G}\right)\) 的定义 124

15.2 \(d : {\mathrm{R}}_{\mathrm{K}}\left( \mathrm{G}\right)  \rightarrow  {\mathrm{R}}_{k}\left( \mathrm{G}\right)\) 的定义 125

15.3 \(e : {\mathrm{P}}_{k}\left( \mathrm{G}\right)  \rightarrow  {\mathrm{R}}_{\mathrm{K}}\left( \mathrm{G}\right)\) 的定义 127

15.4 cde 三角形的基本性质 127

15.5 示例: \({p}^{\prime }\) -群 128

15.6 示例: \(p\) -群 129

15.7 示例: \({p}^{\prime }\) -群与 \(p\) -群的直积 129

16 定理 131

16.1 cde 三角形的性质 131

16.2 \(e\) 像的刻画 133

16.3 投影 A [G]-模的刻画 134

按其特征

16.4 投影 A [G]-模的例子:缺陷为零的 136

不可约表示

17 证明 138

17.1 群的变换 138

17.2 模数情形下的布劳尔定理 139

17.3 定理33的证明 140

17.4 定理35的证明 142

17.5 定理37的证明 143

17.6 定理38的证明 144

18 模块字符 147

18.1 表示的模字符 147

18.2 模字符的独立性 149

18.3 等价表述 151

18.4 为 \(d\) 的截面 152

18.5 例:对称群 \({\mathfrak{S}}_{4}\) 的模字符 153

18.6 例:交错群 \({\mathfrak{A}}_{5}\) 的模字符 156

19 对阿廷表示的应用 159

19.1 阿廷与斯旺表示 159

19.2 阿廷与斯旺表示的有理性 161

19.3 一个不变量 162

附录 163

参考书目:第三部分 165

符号索引 167

术语索引 169

我 表示与特征

\section*{第1章 线性表示的一般性}

\section*{1.1 定义}

令 \(\mathbf{V}\) 为复数域 \(\mathbf{C}\) 上的向量空间，令 GL(V) 为把 V 映为自身的同构全体。GL(V) 的一个元素 \(a\) 按定义是将 \(V\) 映到 \(V\) 的线性映射，且具有逆映射 \({a}^{-1}\) ；此逆映射也是线性的。当 \(\mathrm{V}\) 有一个由 \(\left( {e}_{i}\right)\) 个元素构成的有限基 \(n\) 时，每一线性映射 \(a : \mathrm{V} \rightarrow  \mathrm{V}\) 由一阶为 \(n\) 的方阵 \(\left( {a}_{ij}\right)\) 决定。系数 \({a}_{ij}\) 为复数；它们由将像 \(a\left( {e}_{j}\right)\) 用基 \(\left( {e}_{i}\right)\) 表示得到:

\[
a\left( {e}_{j}\right)  = \mathop{\sum }\limits_{i}{a}_{ij}{e}_{i}.
\]

说 \(a\) 是同构等价于说 \(a\) 的行列式 \(\det \left( a\right)  = \det \left( {a}_{ij}\right)\) 不为零。故群 \(\mathbf{{GL}}\left( \mathrm{V}\right)\) 可识别为一切可逆方阵(阶为 \(n\) )所构成的群。

假设现在 \(G\) 是一个有限群，恒等元为 1，运算为 \(\left( {s,t}\right)  \mapsto  {st}\) 。在 \(\mathrm{V}\) 中 \(\mathrm{G}\) 的线性表示是一个同态映射 \(\rho\) ，将群 \(G\) 映到群 \({GL}\left( V\right)\) 。换言之，我们将每个元素 \(s \in  \mathbf{G}\) 关联到 \(\rho \left( s\right)\) 的一个元素 \(\mathbf{{GL}}\left( \mathrm{V}\right)\) ，并满足等式

\[
\rho \left( {st}\right)  = \rho \left( s\right)  \cdot  \rho \left( t\right) \;\text{ for }s,t \in  \mathrm{G}.
\]

[我们也常写 \({\rho }_{s}\) 代替 \(\rho \left( s\right)\) 。] 注意前述公式蕴含下列关系:

\[
\rho \left( 1\right)  = 1,\;\rho \left( {s}^{-1}\right)  = \rho {\left( s\right) }^{-1}.
\]

给定 \(\rho\) 时，我们说 \(\mathrm{V}\) 是 \(\mathrm{G}\) 的表示空间(或简略称为 G 的表示)。下文中，我们只考虑 V 为有限维的情形。这并非严重的限制。事实上，对大多数应用，人们关心的是处理有限个元素 \({x}_{i}\) 属于 \(\mathbf{V}\) ，总能找到包含这些 \({x}_{i}\) 的有限维子表示(将在 1.3 中给出定义):只需取由这些 \({x}_{i}\) 的像 \({\rho }_{s}\left( {x}_{i}\right)\) 所生成的向量子空间。

现假定 \(\mathbf{V}\) 为有限维，令其维数为 \(n\) ；我们也称 \(n\) 为所研究表示的次数。令 \(\left( {e}_{i}\right)\) 为 \(\mathbf{V}\) 的一基，令 \({\mathbf{R}}_{s}\) 为相对于该基的 \({\rho }_{s}\) 的矩阵。我们有

\[
\det \left( {\mathbf{R}}_{s}\right)  \neq  0,\;{\mathbf{R}}_{st} = {\mathbf{R}}_{s} \cdot  {\mathbf{R}}_{t}\;\text{ if }s,t \in  \mathbf{G}.
\]

若记矩阵 \({\mathbf{R}}_{s}\) 的系数为 \({r}_{ij}\left( s\right)\) ，则第二个公式变为

\[
{r}_{ik}\left( {st}\right)  = \mathop{\sum }\limits_{j}{r}_{ij}\left( s\right)  \cdot  {r}_{jk}\left( t\right) .
\]

反之，给定满足上述恒等式的可逆矩阵组 \({\mathbf{R}}_{s} = \left( {{r}_{ij}\left( s\right) }\right)\) ，则存在相应的线性表示 \(\rho\) ，将 \(\mathbf{G}\) 表示于 \(\mathbf{V}\) ；这就是所谓以“矩阵形式”给出表示的含义。

设 \(\rho\) 与 \({\rho }^{\prime }\) 是同一群 \(\mathbf{G}\) 在向量空间 \(\mathrm{V}\) 与 \({\mathrm{V}}^{\prime }\) 中的两个表示。若存在线性同构 \(\tau  : \mathrm{V} \rightarrow  {\mathrm{V}}^{\prime }\) 将 \(\rho\) “变换”为 \({\rho }^{\prime }\) ，即满足恒等式，则称这两个表示相似(或同构)。

\[
\tau  \circ  \rho \left( s\right)  = {\rho }^{\prime }\left( s\right)  \circ  \tau \;\text{ for all }s \in  \mathrm{G}.
\]

当 \(\rho\) 与 \({\rho }^{\prime }\) 分别以矩阵形式由 \({\mathrm{R}}_{s}\) 与 \({\mathrm{R}}_{s}^{\prime }\) 给出时，这意味着存在可逆矩阵 \(T\) 使得

\[
\mathrm{T} \cdot  {\mathbf{R}}_{s} = {\mathbf{R}}_{s}^{\prime } \cdot  \mathrm{T}\text{ , for all }s \in  \mathrm{G}\text{ , }
\]

也可写作 \({\mathbf{R}}_{s}^{\prime } = \mathbf{T} \cdot  {\mathbf{R}}_{s} \cdot  {\mathbf{T}}^{-1}\) 。我们可以将两个这样的表示视为相同(通过让每个 \(x \in  \mathbf{V}\) 对应元素 \(\tau \left( x\right)  \in  {\mathbf{V}}^{\prime }\) )；特别地， \(\rho\) 与 \({\rho }^{\prime }\) 有相同的次数。

\section*{1.2 基本例子}

(a) 群 \(G\) 的一次表示是群同态 \(\rho  : G \rightarrow  {C}^{ * }\) ，其中 \({C}^{ * }\) 表示非零复数的乘法群。由于 \(G\) 的每个元素都是有限阶， \(\rho\) 的取值 \(\rho \left( s\right)\) 为单位根；尤其，我们有 \(\left| {\rho \left( s\right) }\right|  = 1\) 。

若对所有 \(s \in  \mathbf{G}\) 取 \(\rho \left( s\right)  = 1\) ，则得到群 \(\mathbf{G}\) 的表示，称为平凡(或单位)表示。

(b) 设 \(\mathbf{G}\) 的阶为 \(g\) ，取维数为 \(g\) 的向量空间 \(\mathbf{V}\) ，其基 \({\left( {e}_{t}\right) }_{t \in  \mathrm{G}}\) 以 \(\mathrm{G}\) 的元素 \(t\) 索引。对每个 \(s \in  \mathrm{G}\) ，令 \({\rho }_{s}\) 为将 \(\mathrm{V}\) 中的 \({e}_{t}\) 映为 \({e}_{st}\) 的线性映射；这定义了一个线性表示，称为 G 的正规表示，其次数等于 G 的阶。注意 \({e}_{s} = {\rho }_{s}\left( {e}_{1}\right)\) ；因此 \({e}_{1}\) 的像构成 \(\mathbf{V}\) 的一组基。反过来，若 \(\mathbf{W}\) 是包含向量 \(w\) 的 \(\mathbf{G}\) 的一个表示，且 \({\rho }_{s}\left( w\right) ,s \in  \mathrm{G}\) 构成 \(\mathrm{W}\) 的一组基，则 W 与正规表示同构(通过令 \(\tau \left( {e}_{s}\right)  = {\rho }_{s}\left( w\right)\) 定义同构 \(\tau  : \mathbf{V} \rightarrow  \mathbf{W}\) )。

(c) 更一般地，设 G 在有限集 X 上作用。这意味着对每个 \(s \in  \mathbf{G}\) ，给定一个置换 \(x \mapsto  {sx}\) 于 \(\mathbf{X}\) ，且满足以下恒等式

\[
{1x} = x,s\left( {tx}\right)  = \left( {st}\right) x\;\text{ if }s,t \in  \mathrm{G},x \in  \mathrm{X}.
\]

令 \(\mathrm{V}\) 为一向量空间，其基 \({\left( {e}_{x}\right) }_{x \in  \mathrm{X}}\) 由 X 的元素索引。对 \(s \in  \mathrm{G}\) ，令 \({\rho }_{s}\) 为将 \({e}_{x}\) 送至 \({e}_{sx}\) 的 \(\mathrm{V}\) 上的线性映射；由此得到的 \(\mathbf{G}\) 的线性表示称为与 \(\mathbf{X}\) 相关的置换表示。

\section*{1.3 子表示}

令 \(\rho  : \mathbf{G} \rightarrow  \mathbf{{GL}}\left( \mathbf{V}\right)\) 为一线性表示，且令 \(\mathbf{W}\) 为 \(\mathbf{V}\) 的向量子空间。若 \(\mathbf{W}\) 在 \(\mathbf{G}\) 的作用下不变(亦称“稳定”)，换言之，假设对任意 \(s \in  G\) ，由 \(x \in  W\) 可推出 \({\rho }_{s}x \in  W\) 。则 \({\rho }_{s}\) 限制于 \(W\) 的映射 \({\rho }_{s}^{W}\) 是由 \(W\) 到自身的同构，且我们有 \({\rho }_{st}^{W} = {\rho }_{s}^{W} \cdot  {\rho }_{t}^{W}\) 。因此 \({\rho }^{W} : G \rightarrow  {GL}\left( W\right)\) 是在 \(W;W\) 上的 \(G\) 的线性表示，称为 \(V\) 的一个子表示。

例。取 \(G\) 的正则表示为 \(V\) [参见 1.2 (b)]，并令 W 为由元素 \(x = \mathop{\sum }\limits_{{s \in  \mathrm{G}}}{e}_{s}\) 生成的 \(V\) 的一维子空间。对所有 \(s \in  \mathrm{G}\) ，我们有 \({\rho }_{s}x = x\) ；因此 \(\mathrm{W}\) 是 \(\mathrm{V}\) 的一个子表示，与平凡表示同构。(我们将在 2.4 中确定正则表示的所有子表示。)

在继续之前，我们回顾线性代数中的一些概念。设 V 为向量空间，且 \(W\) 和 \({W}^{\prime }\) 是 \(V\) 的两个子空间。如果每一个 \(x \in  V\) 都可唯一地写成 \(x = w + {w}^{\prime }\) 的形式，其中 \(w \in  \mathrm{W}\) 与 \({w}^{\prime } \in  {\mathrm{W}}^{\prime }\) ，则称空间 \(V\) 为 \(W\) 与 \({W}^{\prime }\) 的直和；这相当于说 \(W\) 与 \({W}^{\prime }\) 的交 \(W \cap  {W}^{\prime }\) 为 0 且 \(\dim \left( V\right)  = \dim \left( W\right)  + \dim \left( {W}^{\prime }\right)\) 。我们记为 \(V = W \oplus  {W}^{\prime }\) 并说 \({\mathbf{W}}^{\prime }\) 是 \(\mathbf{V}\) 中 \(\mathbf{W}\) 的补空间。将每个 \(x \in  \mathbf{V}\) 送到其分量 \(w \in  \mathbf{W}\) 的映射 \(p\) 称为与分解 \(\mathbf{V} = \mathbf{W} \oplus  {\mathbf{W}}^{\prime }\) 相关的从 \(\mathbf{V}\) 到 \(\mathbf{W}\) 的投影； \(p\) 的像是 \(\mathbf{W}\) ，且 \(p\left( x\right)  = x\) 对于 \(x \in  \mathrm{W}\) ；反之，若 \(p\) 是将 \(\mathrm{V}\) 映到自身的线性映射并满足这两个性质，则可验证 \(\mathbf{V}\) 是 \(\mathbf{W}\) 与 \(p\) 的核 \({\mathbf{W}}^{\prime }\) 的直和

(满足 \(x\) 的 \({px} = 0\) 的集合)。因此在从 \(\mathbf{V}\) 到 \(\mathbf{W}\) 的投影与 \(\mathbf{V}\) 中 \(\mathbf{W}\) 的补空间之间建立了一一对应关系。

我们现在回到子表示:

定理 1. 设 \(\rho  : \mathbf{G} \rightarrow  \mathbf{{GL}}\left( \mathbf{V}\right)\) 是在 \(\mathbf{V}\) 中的关于 \(\mathbf{G}\) 的线性表示，且 W 是在 \(\mathbf{G}\) 下稳定的 \(\mathbf{V}\) 的向量子空间。则存在在 \(\mathrm{G}\) 下稳定的 \({\mathrm{W}}^{0}\) ，它是 \(\mathrm{W}\) 在 \(\mathrm{V}\) 中的一个补空间。

取任意的 \({\mathrm{W}}^{\prime }\) 作为 \(\mathrm{V}\) 中 \(\mathrm{W}\) 的一个补，并令 \(p\) 为相应的从 \(\mathrm{V}\) 到 \(\mathrm{W}\) 的投影。对 \(\mathbf{G}\) 中元素的共轭作用取这些投影的平均，构成 \({p}^{0}\) :

\[
{p}^{0} = \frac{1}{g}\mathop{\sum }\limits_{{t \in  \mathrm{G}}}{\rho }_{t} \cdot  p \cdot  {\rho }_{t}^{-1}\;\left( {g\text{ being the order of }\mathrm{G}}\right) .
\]

由于 \(p\) 将 V 映入 W 且 \({\rho }_{t}\) 保持 \(W\) ，我们看到 \({p}^{0}\) 将 \(V\) 映入 \(W\) ；我们对于 \(x \in  W\) 有 \({\rho }_{t}^{-1}x \in  W\) ，因此

\[
p \cdot  {\rho }_{t}^{-1}x = {\rho }_{t}^{-1}x,\;{\rho }_{t} \cdot  p \cdot  {\rho }_{t}^{-1}x = x,\;\text{ and }\;{p}^{0}x = x.
\]

因此 \({p}^{0}\) 是从 \(\mathrm{V}\) 到 \(\mathrm{W}\) 的一个投影，对应于某个 \(\mathrm{W}\) 的补 \({\mathrm{W}}^{0}\) 。此外我们还有

\[
{\rho }_{s} \cdot  {p}^{0} = {p}^{0} \cdot  {\rho }_{s}\;\text{ for all }s \in  \mathrm{G}.
\]

确实，计算 \({\rho }_{s} \cdot  {p}^{0} \cdot  {\rho }_{s}^{-1}\) ，我们得到:

\[
{\rho }_{s} \cdot  {p}^{0} \cdot  {\rho }_{s}^{-1} = \frac{1}{g}\mathop{\sum }\limits_{{t \in  \mathrm{G}}}{\rho }_{s} \cdot  {\rho }_{t} \cdot  p \cdot  {\rho }_{t}^{-1} \cdot  {\rho }_{s}^{-1} = \frac{1}{g}\mathop{\sum }\limits_{{t \in  \mathrm{G}}}{\rho }_{st} \cdot  p \cdot  {\rho }_{st}^{-1} = {p}^{0}.
\]

如果现在 \(x \in  {\mathrm{W}}^{0}\) 和 \(s \in  \mathrm{G}\) 有 \({p}^{0}x = 0\) ，因此 \({p}^{0} \cdot  {\rho }_{s}x = {\rho }_{s} \cdot  {p}^{0}x \; = 0\) ，也就是 \({\rho }_{s}x \in  {\mathrm{W}}^{0}\) ，这表明 \({\mathrm{W}}^{0}\) 在 \(\mathrm{G}\) 下不变，从而完成证明。

注。设 \(\mathrm{V}\) 赋有满足通常条件的标量积 \(\left( {x \mid  y}\right)\) :在 \(x\) 上线性、在 \(y\) 上半线性，并且若 \(x \neq  0\) 则 \(\left( {x \mid  x}\right)  > 0\) 。假设该标量积对 G 不变，即 \(\left( {{\rho }_{s}x \mid  {\rho }_{s}y}\right)  = \left( {x \mid  y}\right)\) ；通过将 \(\left( {x \mid  y}\right)\) 替换为 \(\mathop{\sum }\limits_{{t \in  \mathrm{G}}}\left( {{\rho }_{t}x \mid  {\rho }_{t}y}\right)\) ，我们总可化为此情形。在这些假设下， \(\mathrm{V}\) 中 \(\mathrm{W}\) 的正交补 \({\mathrm{W}}^{0}\) 是一个关于 \(\mathrm{W}\) 的补空间且在 \(\mathrm{G}\) 下不变；由此得到定理 1 的另一证明。注意，标量积 \(\left( {x \mid  y}\right)\) 的不变性意味着，若 \(\left( {e}_{i}\right)\) 是 \(\mathrm{V}\) 的一组标准正交基，则相对于此基， \({\rho }_{s}\) 的矩阵是酉矩阵。

保持定理 1 的假设和记号，设 \(x \in  \mathbf{V}\) 并令 \(w\) 和 \({w}^{0}\) 为其在 \(\mathbf{W}\) 与 \({\mathbf{W}}^{0}\) 上的投影。我们有 \(x = w + {w}^{0}\) ，因此 \({\rho }_{s}x = {\rho }_{s}w + {\rho }_{x}{w}^{0}\) ，且由于 \(\mathbf{W}\) 与 \({\mathbf{W}}^{0}\) 在 \(\mathbf{G}\) 下不变，我们有 \({\rho }_{s}w \in  \widetilde{\mathbf{W}}\) 与 \({\rho }_{s}{w}^{0} \in  {\mathbf{W}}^{0}\) ；于是 \({\rho }_{s}w\) 与 \({\rho }_{s}{w}^{0}\) 是 \({\rho }_{s}x\) 的投影。由此可得表示 \(\mathbf{W}\) 与 \({\mathbf{W}}^{0}\) 决定表示 \(\mathbf{V}\) 。

我们称 \(V\) 是 \(W\) 与 \({W}^{0}\) 的直和，并记作 \(V = W \oplus  {W}^{0}\) 。 \(V\) 的元素被识别为一对 \(\left( {w,{w}^{0}}\right)\) ，其中 \(w \in  W\) 与 \({w}^{0} \in  {W}^{0}\) 。若以矩阵形式给出 \(W\) 与 \({W}^{0}\) 为 \({R}_{s}\) ，则 \({R}_{s}^{0},W \oplus  {W}^{0}\) 以矩阵形式表示为

\[
\left( \begin{matrix} {\mathrm{R}}_{s} & 0 \\  0 & {\mathrm{R}}_{s}^{0} \end{matrix}\right) .
\]

任意有限个表示的直和以类似方式定义。

\section*{1.4 不可约表示}

设 \(\rho  : G \rightarrow  {GL}\left( V\right)\) 为群 \(G\) 的线性表示。若 \(V\) 非 0 且 除 0 和 V 外没有向量子空间在 G 下不变，则称其为不可约或单表示。由定理 1，第二个条件等价于说 V 不能表示为两个表示的直和(当然不包括平凡分解 \(V = 0 \oplus  V\) )。次数为 1 的表示显然是不可约的。我们稍后(3.1)将看到每个非阿贝尔群至少有一个次数为 \(\geq  2\) 的不可约表示。

不可约表示可用直和来构造其他表示:

定理 2. 每个表示都可表示为不可约表示的直和。

设 \(V\) 为群 \(G\) 的线性表示。我们对 \(\dim \left( V\right)\) 归纳。若 \(\dim \left( V\right)  = 0\) ，定理显然(0 是空族不可约表示的直和)。设 \(\dim \left( V\right)  \geq  1\) 。若 \(V\) 不可约，则无需证明。否则，由定理 1， \(V\) 可分解为直和 \({V}^{\prime } \oplus  {V}^{\prime \prime }\) ，且 \(\dim \left( {V}^{\prime }\right)  < \dim \left( V\right)\) 与 \(\dim \left( {V}^{\prime \prime }\right)  < \dim \left( V\right)\) 非零。由归纳假设 \({V}^{\prime }\) 与 \({V}^{\prime \prime }\) 分别为不可约表示的直和，因此 V 亦然。

注。设 \(\mathbf{V}\) 为一个表示，且设 \(\mathbf{V} = {\mathbf{W}}_{1} \oplus  \cdots  \oplus  {\mathbf{W}}_{k}\) 为将 \(\mathrm{V}\) 分解为若干不可约表示直和的一个分解。我们可以问这个分解是否唯一。当所有 \({\rho }_{s}\) 都等于 1 的情形表明一般情况下并非如此(在这种情况下 \({\mathbf{W}}_{i}\) 是一维子空间，我们对一个向量空间分解为若干一维子空间有大量不同的分解)。尽管如此，我们将在 2.3 中看到，和给定不可约表示同构的 \({\mathrm{W}}_{i}\) 的个数不依赖于所选的分解。

\section*{1.5 两个表示的张量积}

除了具有加法形式性质的直和运算外，还有一种“乘法”:张量积，有时称为克罗内克乘积。其定义如下:

首先，设 \({V}_{1}\) 与 \({V}_{2}\) 为两个向量空间。若有一空间 \(W\) 以及将 \({V}_{1} \times  {V}_{2}\) 映入 \(W\) 的映射 \(\left( {{x}_{1},{x}_{2}}\right)  \mapsto  {x}_{1} \cdot  {x}_{2}\) ，并满足下列两条，则称 \(W\) 为 \({V}_{1}\) 与 \({V}_{2}\) 的张量积:

(i) 对变量 \({x}_{1}\) 与 \({x}_{2}\) 中的各自分量， \({x}_{1} \cdot  {x}_{2}\) 在每一项上都是线性的。

(ii) 若 \(\left( {e}_{{i}_{1}}\right)\) 为 \({V}_{1}\) 的一组基，且 \(\left( {e}_{{i}_{2}}\right)\) 为 \({V}_{2}\) 的一组基，则各乘积 \({e}_{{i}_{1}} \cdot  {e}_{{i}_{2}}\) 构成 \(\mathbf{W}\) 的一组基。

容易证明这样的空间存在且唯一(同构意义下)；记作 \({V}_{1} \otimes  {V}_{2}\) 。条件 (ii) 表明

\[
\dim \left( {{\mathrm{V}}_{1} \otimes  {\mathrm{V}}_{2}}\right)  = \dim \left( {\mathrm{V}}_{1}\right)  \cdot  \dim \left( {\mathrm{V}}_{2}\right) .
\]

现在令 \({\rho }^{1} : G \rightarrow  {GL}\left( {V}_{1}\right)\) 和 \({\rho }^{2} : G \rightarrow  {GL}\left( {V}_{2}\right)\) 为群 G 的两个线性表示。对 \(s \in  G\) ，通过下列条件在 \({GL}\left( {{V}_{1} \otimes  {V}_{2}}\right)\) 中定义元素 \({\rho }_{s}\) :

\[
{\rho }_{s}\left( {{x}_{1} \cdot  {x}_{2}}\right)  = {\rho }_{s}^{1}\left( {x}_{1}\right)  \cdot  {\rho }_{s}^{2}\left( {x}_{2}\right) \;\text{ for }{x}_{1} \in  {\mathrm{V}}_{1},{x}_{2} \in  {\mathrm{V}}_{2}.
\]

[ \({\rho }_{s}\) 的存在性与唯一性由条件 (i) 与 (ii) 易得。] 我们写作:

\[
{\rho }_{s} = {\rho }_{s}^{1} \otimes  {\rho }_{s}^{2}
\]

这些 \({\rho }_{s}\) 在 \({\mathrm{V}}_{1} \otimes  {\mathrm{V}}_{2}\) 中给出一个 \(\mathrm{G}\) 的线性表示，称为所给表示的张量积。

该定义的矩阵表述如下:令 \(\left( {e}_{{i}_{1}}\right)\) 为 \({V}_{1}\) 的一组基，令 \({r}_{{i}_{1}{j}_{1}}\left( s\right)\) 为相对于此基的 \({\rho }_{s}^{1}\) 的矩阵，并以同样方式定义 \(\left( {e}_{{i}_{2}}\right)\) 与 \({r}_{{i}_{2}{j}_{2}}\left( s\right)\) 。公式:

\[
{\rho }_{s}^{1}\left( {e}_{{j}_{1}}\right)  = \mathop{\sum }\limits_{{i}_{1}}{r}_{{i}_{1}{j}_{1}}\left( s\right)  \cdot  {e}_{{i}_{1}},\;{\rho }_{s}^{2}\left( {e}_{{j}_{2}}\right)  = \mathop{\sum }\limits_{{i}_{2}}{r}_{{i}_{2}{j}_{2}}\left( s\right)  \cdot  {e}_{{i}_{2}}
\]

蕴含:

\[
{\rho }_{s}\left( {{e}_{{j}_{1}} \cdot  {e}_{{j}_{2}}}\right)  = \mathop{\sum }\limits_{{{i}_{1},{i}_{2}}}{r}_{{i}_{1}{j}_{1}}\left( s\right)  \cdot  {r}_{{i}_{2}{j}_{2}}\left( s\right)  \cdot  {e}_{{i}_{1}} \cdot  {e}_{{i}_{2}}.
\]

因此 \({\rho }_{s}\) 的矩阵为 \(\left( {{r}_{{i}_{1}{j}_{1}}\left( s\right)  \cdot  {r}_{{i}_{2}{j}_{2}}\left( s\right) }\right)\) ；它是 \({\rho }_{s}^{1}\) 与 \({\rho }_{s}^{2}\) 矩阵的张量积。

两个不可约表示的张量积通常不再不可约。它分解为不可约表示的直和，可通过表征理论确定(参见 2.3)。

在量子化学中，张量积常如下出现: \({V}_{1}\) 与 \({V}_{2}\) 是在 \(G\) 下保持的不函数空间，分别以 \(\left( {\phi }_{{i}_{1}}\right)\) 与 \(\left( {\psi }_{{i}_{2}}\right)\) 为基， \({V}_{1} \otimes  {V}_{2}\) 为由乘积 \({\phi }_{{i}_{1}} \cdot  {\psi }_{{i}_{2}}\) 生成的向量空间，且这些乘积线性无关。最后一条条件至关重要。以下为满足该条件的两个特例:

(1) \(\phi\) 只依赖于某些变量 \(\left( {x,{x}^{\prime },\ldots }\right)\) 而 \(\psi\) 依赖于与前者无关的变量 \(\left( {y,{y}^{\prime },\ldots }\right)\) 。

(2) 空间 \({V}_{1}\) (或 \({V}_{2}\) )有以单个函数 \(\phi\) 为基，该函数在任何区域内不恒为零；则 \({V}_{1}\) 的维数为 1。

\section*{1.6 对称平方与交错平方}

假设表示 \({V}_{1}\) 与 \({V}_{2}\) 同于相同的表示 \(\mathbf{V}\) ，因此 \({\mathbf{V}}_{1} \otimes  {\mathbf{V}}_{2} = \mathbf{V} \otimes  \mathbf{V}\) 。若 \(\left( {e}_{i}\right)\) 为 \(\mathbf{V}\) 的一组基，令 \(\theta\) 为 \(\mathbf{V} \otimes  \mathbf{V}\) 的自同构，使得

\[
\theta \left( {{e}_{i} \cdot  {e}_{j}}\right)  = {e}_{j} \cdot  {e}_{i}\;\text{ for all pairs }\left( {i,j}\right) .
\]

由此可得 \(\theta \left( {x \cdot  y}\right)  = y \cdot  x\) 对于 \(x,y \in  \mathrm{V}\) ，因此 \(\theta\) 独立于所选基 \(\left( {e}_{i}\right)\) ；此外 \({\theta }^{2} = 1\) 。空间 \(\mathbf{V} \otimes  \mathbf{V}\) 随即分解为直和

\[
\mathrm{V} \otimes  \mathrm{V} = {\mathbf{{Sym}}}^{2}\left( \mathrm{\;V}\right)  \oplus  {\mathbf{{Alt}}}^{2}\left( \mathrm{\;V}\right)
\]

其中 \({\operatorname{Sym}}^{2}\left( \mathrm{\;V}\right)\) 是满足 \(\theta \left( z\right)  = z\) 的元素 \(z \in  \mathrm{V} \otimes  \mathrm{V}\) 的集合， \({\operatorname{Alt}}^{2}\left( \mathrm{\;V}\right)\) 是满足 \(\theta \left( z\right)  =  - z\) 的元素 \(z \in  \mathrm{V} \otimes  \mathrm{V}\) 的集合。元素 \({\left( {e}_{i} \cdot  {e}_{j} + {e}_{j} \cdot  {e}_{i}\right) }_{i \leq  j}\) 构成 \({\operatorname{Sym}}^{2}\left( \mathrm{\;V}\right)\) 的一组基，元素 \({\left( {e}_{i} \cdot  {e}_{j} - {e}_{j} \cdot  {e}_{i}\right) }_{i < j}\) 构成 \({\operatorname{Alt}}^{2}\left( \mathrm{\;V}\right)\) 的一组基。我们有

\[
\dim {\mathbf{{Sym}}}^{2}\left( \mathrm{\;V}\right)  = \frac{n\left( {n + 1}\right) }{2},\;\dim {\mathbf{{Alt}}}^{2}\left( \mathrm{\;V}\right)  = \frac{n\left( {n - 1}\right) }{2}
\]

若 \(\dim \mathrm{V} = n\) 。

子空间 \({\operatorname{Sym}}^{2}\left( \mathrm{V}\right)\) 与 \({\operatorname{Alt}}^{2}\left( \mathrm{V}\right)\) 在 \(\mathrm{G}\) 下不变，因而分别给出所给表示的对称平方和交替平方表示。

\section*{第2章 特征理论}

\section*{2.1 表示的特征值(character)}

设 \(\mathrm{V}\) 为一具基为 \(\left( {e}_{i}\right)\) 共 \(n\) 个元素的向量空间，且令 \(a\) 为将 \(\mathbf{V}\) 映到其自身的线性映射，其矩阵为 \(\left( {a}_{ij}\right)\) 。称 \(a\) 的迹为标量

\[
\operatorname{Tr}\left( a\right)  = \mathop{\sum }\limits_{i}{a}_{ii}
\]

它等于 a 的特征值之和(按重数计)，且不依赖于所选基 \(\left( {e}_{i}\right)\) 。

现在令 \(\rho  : \mathbf{G} \rightarrow  \mathbf{{GL}}\left( \mathrm{V}\right)\) 为有限群 \(\mathbf{G}\) 在向量空间 \(\mathrm{V}\) 上的线性表示。对每个 \(s \in  \mathrm{G}\) ，设:

\[
{\chi }_{\rho }\left( s\right)  = \operatorname{Tr}\left( {\rho }_{s}\right) .
\]

由此得的在 \(\mathrm{G}\) 上的复值函数 \({\chi }_{\rho }\) 称为表示 \(\rho\) 的特征；该函数的重要性主要在于它刻画了表示 \(\rho\) (参见 2.3)。

命题1. 若 \(\chi\) 是一表示 \(\rho\) 的特征，且其次数为 \(n\) ，则有:

(i) \(\chi \left( 1\right)  = n\) ,

(ii) 对于 \(s \in  \mathrm{G}\) 有 \(\chi \left( {s}^{-1}\right)  = \chi {\left( s\right) }^{ * }\;\) ，

(iii) 对于 \(s,t \in  \mathrm{G}\) 有 \(\chi \left( {{ts}{t}^{-1}}\right)  = \chi \left( s\right) \;\) 。

(若 \(z = x + {iy}\) 是复数，我们用 \({z}^{ * }\) 或 \(\bar{z}\) 表示其共轭 \(x - {iy}\) 。)

我们有 \(\rho \left( 1\right)  = 1\) ，且 \(\operatorname{Tr}\left( 1\right)  = n\) 因为 \(\mathrm{V}\) 的维数为 \(n\) ；因此得 (i)。

对于(ii)我们注意到 \({\rho }_{s}\) 具有有限阶；因此其特征值 \({\lambda }_{1},\ldots ,{\lambda }_{n}\) 也具有有限阶，故其模均为1(这也可从 \({\rho }_{s}\) 可由酉矩阵给出看出，参见1.3)。因此

\[
\chi {\left( s\right) }^{ * } = \operatorname{Tr}{\left( {\rho }_{s}\right) }^{ * } = \sum {\lambda }_{i}^{ * } = \sum {\lambda }_{i}^{-1} = \operatorname{Tr}\left( {\rho }_{s}^{-1}\right)  = \operatorname{Tr}\left( {\rho }_{{s}^{-1}}\right)  = \chi \left( {s}^{-1}\right) .
\]

公式(iii)也可写为 \(\chi \left( {vu}\right)  = \chi \left( {uv}\right)\) ，将 \(u = {ts},v = {t}^{-1}\) 代入；因此它可由众所周知的公式推出

\[
\operatorname{Tr}\left( {ab}\right)  = \operatorname{Tr}\left( {ba}\right)
\]

该公式对任意两个将 \(\mathbf{V}\) 映入自身的线性映射 \(a\) 和 \(b\) 成立。

注. 在 \(\mathbf{G}\) 上满足恒等式 (iii) 的函数 \(f\) ，或者等价地说 \(f\left( {uv}\right)  = f\left( {vu}\right)\) ，称为类函数。我们将在 2.5 中看到每个类函数都是若干特征的线性组合。

命题2. 设 \({\rho }^{1} : \mathrm{G} \rightarrow  \mathbf{{GL}}\left( {\mathrm{V}}_{1}\right)\) 和 \({\rho }^{2} : \mathrm{G} \rightarrow  \mathbf{{GL}}\left( {\mathrm{V}}_{2}\right)\) 为 \(\mathrm{G}\) 的两个线性表示，令 \({\chi }_{1}\) 和 \({\chi }_{2}\) 为它们的特征。则:

(i) 直和表示 \({\mathrm{V}}_{1} \oplus  {\mathrm{V}}_{2}\) 的特征 \(\chi\) 等于 \({\chi }_{1} + {\chi }_{2}\)

(ii) 张量积表示 \({\mathrm{V}}_{1} \otimes  {\mathrm{V}}_{2}\) 的特征 \(\psi\) 等于 \({\chi }_{1} \cdot  {\chi }_{2}\) 。

设 \({\rho }^{1}\) 和 \({\rho }^{2}\) 以矩阵形式给出: \({\mathrm{R}}_{s}^{1},{\mathrm{R}}_{s}^{2}\) 。表示 \({\mathrm{V}}_{1} \oplus  {\mathrm{V}}_{2}\) 由下式给出

\[
{\mathrm{R}}_{s} = \left( \begin{matrix} {\mathrm{R}}_{s}^{1} & 0 \\  0 & {\mathrm{R}}_{s}^{2} \end{matrix}\right)
\]

从而 \(\operatorname{Tr}\left( {\mathrm{R}}_{s}\right)  = \operatorname{Tr}\left( {\mathrm{R}}_{s}^{1}\right)  + \operatorname{Tr}\left( {\mathrm{R}}_{s}^{2}\right)\) ，即 \(\chi \left( s\right)  = {\chi }_{1}\left( s\right)  + {\chi }_{2}\left( s\right)\) 。

我们对(ii)以同样方式处理:按1.5的记法，我们有

\[
{\chi }_{1}\left( s\right)  = \mathop{\sum }\limits_{{i}_{1}}{r}_{{i}_{1}{i}_{1}}\left( s\right) ,\;{\chi }_{2}\left( s\right)  = \mathop{\sum }\limits_{{i}_{2}}{r}_{{i}_{2}{i}_{2}}\left( s\right) ,
\]

\[
\psi \left( s\right)  = \mathop{\sum }\limits_{{{i}_{1},{i}_{2}}}{r}_{{i}_{1}{i}_{1}}\left( s\right) {r}_{{i}_{2}{i}_{2}}\left( s\right)  = {\chi }_{1}\left( s\right)  \cdot  {\chi }_{2}\left( s\right) .
\]

命题3. 设 \(\rho  : \mathbf{G} \rightarrow  \mathbf{{GL}}\left( \mathrm{V}\right)\) 为 \(\mathbf{G}\) 的线性表示， \(\chi\) 为其特征。令 \({\chi }_{\sigma }^{2}\) 为 \(\mathrm{V}\) 的对称平方 \({\operatorname{Sym}}^{2}\left( \mathrm{V}\right)\) (参见1.5)的特征，令 \({\chi }_{\alpha }^{2}\) 为 \({\operatorname{Alt}}^{2}\left( \mathrm{V}\right)\) 的特征。对每个 \(s \in  \mathrm{G}\) ，有

\[
{\chi }_{\sigma }^{2}\left( s\right)  = \frac{1}{2}\left( {\chi {\left( s\right) }^{2} + \chi \left( {s}^{2}\right) }\right)
\]

\[
{\chi }_{\alpha }^{2}\left( s\right)  = \frac{1}{2}\left( {\chi {\left( s\right) }^{2} - \chi \left( {s}^{2}\right) }\right)
\]

并且 \({\chi }_{\sigma }^{2} + {\chi }_{\alpha }^{2} = {\chi }^{2}\) 。

设 \(s \in  \mathbf{G}\) 。可以选取由 \(\left( {e}_{i}\right)\) 的特征向量组成的基 \(\mathbf{V}\) ，例如这是因为 \({\rho }_{s}\) 可由酉矩阵表示，参见 1.3。于是我们有 \({\rho }_{s}{e}_{i} = {\lambda }_{i}{e}_{i}\) 且 \({\lambda }_{i} \in  \mathbf{C}\) ，因此

\[
\chi \left( s\right)  = \sum {\lambda }_{i},\;\chi \left( {s}^{2}\right)  = \sum {\lambda }_{i}^{2}.
\]

另一方面，我们有

\[
\left( {{\rho }_{s} \otimes  {\rho }_{s}}\right) \left( {{e}_{i} \cdot  {e}_{j} + {e}_{j} \cdot  {e}_{i}}\right)  = {\lambda }_{i}{\lambda }_{j} \cdot  \left( {{e}_{i} \cdot  {e}_{j} + {e}_{j} \cdot  {e}_{i}}\right) ,
\]

\[
\left( {{\rho }_{s} \otimes  {\rho }_{s}}\right) \left( {{e}_{i} \cdot  {e}_{j} - {e}_{j} \cdot  {e}_{i}}\right)  = {\lambda }_{i}{\lambda }_{j} \cdot  \left( {{e}_{i} \cdot  {e}_{j} - {e}_{j} \cdot  {e}_{i}}\right) ,
\]

因此

\[
{\chi }_{\sigma }^{2}\left( s\right)  = \mathop{\sum }\limits_{{i < j}}{\lambda }_{i}{\lambda }_{j} = \sum {\lambda }_{i}^{2} + \mathop{\sum }\limits_{{i < j}}{\lambda }_{i}{\lambda }_{j} = \frac{1}{2}{\left( \sum {\lambda }_{i}\right) }^{2} + \frac{1}{2}\sum {\lambda }_{i}^{2}
\]

\[
{\chi }_{\alpha }^{2}\left( s\right)  = \mathop{\sum }\limits_{{i < j}}{\lambda }_{i}{\lambda }_{j} = \frac{1}{2}{\left( \sum {\lambda }_{i}\right) }^{2} - \frac{1}{2}\sum {\lambda }_{i}^{2}.
\]

命题得证。

(注意等式 \({\chi }_{\sigma }^{2} + {\chi }_{\alpha }^{2} = {\chi }^{2}\) ，它反映了 \(\mathrm{V} \otimes  \mathrm{V}\) 为 \({\operatorname{Sym}}^{2}\left( \mathrm{V}\right)\) 与 \({\operatorname{Alt}}^{2}\left( \mathrm{V}\right)\) 的直和。)

练习

2.1. 设 \(\chi\) 和 \({\chi }^{\prime }\) 是两个表示的特征。证明下列公式:

\[
{\left( \chi  + {\chi }^{\prime }\right) }_{\sigma }^{2} = {\chi }_{\sigma }^{2} + {\chi }_{\sigma }^{\prime 2} + \chi {\chi }^{\prime }
\]

\[
{\left( \chi  + {\chi }^{\prime }\right) }_{\alpha }^{2} = {\chi }_{\alpha }^{2} + {\chi }_{\alpha }^{\prime 2} + \chi {\chi }^{\prime }
\]

2.2. 设 \(\mathrm{X}\) 为一个有限集合， \(\mathrm{G}\) 在其上作用，令 \(\rho\) 为相应的置换表示［参见 1.2，例(c)］， \({\chi }_{\mathrm{X}}\) 为 \(\rho\) 的特征。设 \(s \in  \mathbf{G}\) ；证明 \({\chi }_{\mathrm{X}}\left( s\right)\) 为被 \(s\) 固定的 \(\mathrm{X}\) 元素个数。

2.3. 设 \(\rho  : G \rightarrow  {GL}\left( V\right)\) 为一线性表示，其特征为 \(\chi\) ，令 \({V}^{\prime }\) 为 \(V\) 的对偶，即 \(V\) 上线性函数的空间。对 \(x \in  V,{x}^{\prime } \in  {V}^{\prime }\) ，记 \(\left\langle  {x,{x}^{\prime }}\right\rangle\) 为线性泛函 \({x}^{\prime }\) 在 \(x\) 处的值。证明存在唯一的线性表示 \({\rho }^{\prime } : \mathrm{G} \rightarrow  \mathrm{{GL}}\left( {\mathrm{V}}^{\prime }\right)\) ，使得

\[
\left\langle  {{\rho }_{s}x,{\rho }_{s}^{\prime }{x}^{\prime }}\right\rangle   = \left\langle  {x,{x}^{\prime }}\right\rangle  \;\text{ for }s \in  \mathrm{G},x \in  \mathrm{V},{x}^{\prime } \in  {\mathrm{V}}^{\prime }.
\]

这称为 \(\rho\) 的对偶(或反合同)表示；其特征为 \({\chi }^{ * }\) 。

2.4. 设 \({\rho }_{1} : G \rightarrow  \operatorname{GL}\left( {V}_{1}\right)\) 和 \({\rho }_{2} : G \rightarrow  \operatorname{GL}\left( {V}_{2}\right)\) 为两个线性表示，特征分别为 \({\chi }_{1}\) 和 \({\chi }_{2}\) 。令 \(W = \operatorname{Hom}\left( {{V}_{1},{V}_{2}}\right)\) 为线性映射空间 \(f : {V}_{1} \rightarrow  {V}_{2}\) 。对 \(s \in  G\) 与 \(f \in  W\) 令 \({\rho }_{s}f = {\rho }_{2,s} \circ  f \circ  {\rho }_{1,s}^{-1}\) ，于是 \({\rho }_{s}f \in  \mathbf{W}\) 。证明这给出了一个线性表示 \(\rho  : \mathbf{G} \rightarrow  \mathbf{{GL}}\left( \mathbf{W}\right)\) ，且其特征为 \({\chi }_{1}^{ * } \cdot  {\chi }_{2}\) 。该表示与 \({\rho }_{1}^{\prime } \otimes  {\rho }_{2}\) 同构，

其中 \({\rho }_{1}^{\prime }\) 为 \({\rho }_{1}\) 的对偶，参见练习 2.3。

\section*{2.2 舒尔引理；基本应用}

命题 4(舒尔引理)。设 \({\rho }^{1} : \mathrm{G} \rightarrow  \mathbf{{GL}}\left( {\mathrm{V}}_{1}\right)\) 与 \({\rho }^{2} : \mathrm{G} \rightarrow  \mathbf{{GL}}\left( {\mathrm{V}}_{2}\right)\) 为 \(\mathrm{G}\) 的两个不可约表示，且令 \(f\) 为从 \({\mathrm{V}}_{1}\) 到 \({\mathrm{V}}_{2}\) 的线性映射，满足对所有 \(s \in  \mathrm{G}\) 有 \({\rho }_{s}^{2} \circ  f = f \circ  {\rho }_{s}^{1}\) 。则:

(1) 若 \({\rho }^{1}\) 与 \({\rho }^{2}\) 非同构，则有 \(f = 0\) 。

(2) 如果 \({\mathbf{V}}_{1} = {\mathbf{V}}_{2}\) 和 \({\rho }^{1} = {\rho }^{2},f\) 是中心相似(即标量乘以恒等映射)。

情况 \(f = 0\) 是平凡的。现在假设 \(f \neq  0\) 并令 \({\mathrm{W}}_{1}\) 为其核(即满足 \(x \in  {\mathrm{V}}_{1}\) 的所有 \({fx} = 0\) 的集合)。对于 \(x \in  {\mathrm{W}}_{1}\) 我们有 \(f{\rho }_{s}^{1}x = {\rho }_{s}^{2}{fx} \; = 0\) ，于是 \({\rho }_{s}^{1}\chi  \in  {\mathrm{W}}_{1}\) ，且 \({\mathrm{W}}_{1}\) 在 G 下不变。由于 \({\mathrm{V}}_{1}\) 是不可约的， \({W}_{1}\) 等于 \({V}_{1}\) 或 0；第一种情况被排除，因为它蕴含 \(f = 0\) 。同样的论证表明映射 \(f\) 的像 \({W}_{2}\) (即满足 \({fx}\) 的 \(x \in  {V}_{1}\) 的集合)等于 \({V}_{2}\) 。性质 \({W}_{1} = 0\) 和 \({W}_{2} = {V}_{2}\) 显示 \(f\) 是从 \({V}_{1}\) 到 \({V}_{2}\) 的同构，从而证明了命题 (1)。

现在假设 \({V}_{1} = {V}_{2},{\rho }^{1} = {\rho }^{2}\) ，并令 \(\lambda\) 为 \(f\) 的一个特征值:至少存在一个，因为标量域是复数域。设 \({f}^{\prime } = f - \lambda\) 。既然 \(\lambda\) 是 \(f\) 的特征值，则算子 \({f}^{\prime }\) 的核为 \(\neq  0\) ；另一方面，我们有 \({\rho }_{s}^{2} \circ  {f}^{\prime } = {f}^{\prime } \circ  {\rho }_{s}^{1}\) 。证明的第一部分表明这些性质只可能在 \({f}^{\prime } = 0\) 的情形，即当 \(f\) 等于 \(\lambda\) 时成立。

保留假设 \({V}_{1}\) 和 \({V}_{2}\) 是不可约的，并记 \(\mathbf{G}\) 的阶为 \(g\) 。

推论 1. 设 \(h\) 是从 \({\mathrm{V}}_{1}\) 到 \({\mathrm{V}}_{2}\) 的线性映射，并设:

\[
{h}^{0} = \frac{1}{g}\mathop{\sum }\limits_{{t \in  \mathrm{G}}}{\left( {\rho }_{t}^{2}\right) }^{-1}h{\rho }_{t}^{1}.
\]

则有:

(1) 若 \({\rho }^{1}\) 和 \({\rho }^{2}\) 不同构，则有 \({h}^{0} = 0\) 。

(2) 若 \({\mathrm{V}}_{1} = {\mathrm{V}}_{2}\) 和 \({\rho }^{1} = {\rho }^{2},{h}^{0}\) 是比为 \(\left( {1/n}\right) \operatorname{Tr}\left( h\right)\) 的同构变换，且 \(n = \dim \left( {\mathrm{V}}_{1}\right) .\)

我们有 \({\rho }_{s}^{2}{h}^{0} = {h}^{0}{\rho }_{s}^{1}\) 。确实:

\[
{\left( {\rho }_{s}^{2}\right) }^{-1}{h}^{0}{\rho }_{s}^{1} = \frac{1}{g}\mathop{\sum }\limits_{{t \in  \mathrm{G}}}{\left( {\rho }_{s}^{2}\right) }^{-1}{\left( {\rho }_{t}^{2}\right) }^{-1}h{\rho }_{t}^{1}{\rho }_{s}^{1}
\]

\[
= \frac{1}{g}\mathop{\sum }\limits_{{t \in  \mathrm{G}}}{\left( {\rho }_{ts}^{2}\right) }^{-1}h{\rho }_{ts}^{1} = {h}^{0}.
\]

将命题 4 应用于 \(f = {h}^{0}\) ，我们在情况 (1) 中看到 \({h}^{0} = 0\) ，在情况 (2) 中看到 \({h}^{0}\) 等于某个标量 \(\lambda\) 。此外，在后一种情况下，我们有:

\[
\operatorname{Tr}\left( {h}^{0}\right)  = \frac{1}{g}\mathop{\sum }\limits_{{t \in  \mathrm{G}}}\operatorname{Tr}\left( {{\left( {\rho }_{t}^{1}\right) }^{-1}h{\rho }_{t}^{1}}\right)  = \operatorname{Tr}\left( h\right)
\]

并且由于 \(\operatorname{Tr}\left( \lambda \right)  = n \cdot  \lambda\) ，我们得到 \(\lambda  = \left( {1/n}\right) \operatorname{Tr}\left( h\right)\) 。

现在我们在假设 \({\rho }^{1}\) 与 \({\rho }^{2}\) 以矩阵形式给定的情况下重写推论 1:

\[
{\rho }_{t}^{1} = \left( {{r}_{{i}_{1}{j}_{1}}\left( t\right) }\right) ,{\rho }_{t}^{2} = \left( {{r}_{{i}_{2}{j}_{2}}\left( t\right) }\right) .
\]

线性映射 \(h\) 由矩阵 \(\left( {x}_{{i}_{2}{i}_{1}}\right)\) 定义，同样 \({h}^{0}\) 由 \(\left( {x}_{{i}_{2}{i}_{1}}^{0}\right)\) 定义。由 \({h}^{0}\) 的定义我们有:

\[
{x}_{{i}_{2}{i}_{1}}^{0} = \frac{1}{g}\mathop{\sum }\limits_{{t,{j}_{1},{j}_{2}}}{r}_{{i}_{2}{j}_{2}}\left( {t}^{-1}\right) {x}_{{j}_{2}{j}_{1}}{r}_{{j}_{1}{i}_{1}}\left( t\right) .
\]

右边对于 \({x}_{{j}_{2}{j}_{1}}\) 是一个线性型；在情况 (1) 中，该型对所有 \({x}_{{j}_{2}{j}_{1}}\) 的取值系统都为零；因此其系数为零。由此得:

推论 2。 在情况 (1) 中，我们有:

\[
\frac{1}{g}\mathop{\sum }\limits_{{t \in  \mathrm{G}}}{r}_{{i}_{2}{j}_{2}}\left( {t}^{-1}\right) {r}_{{j}_{1}{i}_{1}}\left( t\right)  = 0
\]

对任意 \({i}_{1},{i}_{2},{j}_{1},{j}_{2}\) 。

在情况 (2) 中我们类似地有 \({h}^{0} = \lambda\) ，即 \({x}_{{i}_{2}{i}_{1}}^{0} = \lambda {\delta }_{{i}_{2}{i}_{1}}\left( {\delta }_{{i}_{2}{i}_{1}}\right.\) 表示克罗内克符号(当 \({i}_{1} = {i}_{2}\) 时等于 1，否则为 0)，并且有 \(\lambda  = \left( {1/n}\right) \operatorname{Tr}\left( h\right)\) ，也就是 \(\lambda  = \left( {1/n}\right) \sum {\delta }_{{j}_{2}{j}_{1}}{\chi }_{{j}_{2}{j}_{1}}\) 。因此有等式:

\[
\frac{1}{g}\mathop{\sum }\limits_{{t,{j}_{1},{j}_{2}}}{r}_{{i}_{2}{j}_{2}}\left( {t}^{-1}\right) {x}_{{j}_{2}{j}_{1}}{r}_{{j}_{1}{i}_{1}}\left( t\right)  = \frac{1}{n}\mathop{\sum }\limits_{{{j}_{1},{j}_{2}}}{\delta }_{{i}_{2}{i}_{1}}{\delta }_{{j}_{2}{j}_{1}}{x}_{{j}_{2}{j}_{1}}.
\]

比较 \({x}_{{j}_{2}{j}_{1}}\) 的系数，我们如上得到:

推论 3。 在情况 (2) 中我们有:

\[
\frac{1}{g}\mathop{\sum }\limits_{{t \in  \mathrm{G}}}{r}_{{i}_{2}{j}_{2}}\left( {t}^{-1}\right) {r}_{{j}_{1}{i}_{1}}\left( t\right)  = \frac{1}{n}{\delta }_{{i}_{2}{i}_{1}}{\delta }_{{j}_{2}{j}_{1}} = \left\{  \begin{array}{ll} 1/n & \text{ if }{i}_{1} = {i}_{2}\text{ and }{j}_{1} = {j}_{2} \\  0 & \text{ otherwise } \end{array}\right.
\]

注

(1) 若 \(\phi\) 与 \(\psi\) 是作用在 \(G\) 上的函数，设

\[
\langle \phi ,\psi \rangle  = \frac{1}{g}\mathop{\sum }\limits_{{t \in  \mathrm{G}}}\phi \left( {t}^{-1}\right) \psi \left( t\right)  = \frac{1}{g}\mathop{\sum }\limits_{{t \in  \mathrm{G}}}\phi \left( t\right) \psi \left( {t}^{-1}\right) .
\]

我们有 \(\langle \phi ,\psi \rangle  = \langle \psi ,\phi \rangle\) 。此外 \(\langle \phi ,\psi \rangle\) 在 \(\phi\) 与 \(\psi\) 上是线性的。用此记法，推论 2 与 3 分别变为

\[
\left\langle  {{r}_{{i}_{2}{j}_{2}},{r}_{{j}_{1}{i}_{1}}}\right\rangle   = 0\;\text{ and }\;\left\langle  {{r}_{{i}_{2}{j}_{2}},{r}_{{j}_{1}{i}_{1}}}\right\rangle   = \frac{1}{n}{\delta }_{{i}_{2}{i}_{1}}{\delta }_{{j}_{2}{j}_{1}}.
\]

(2) 设矩阵 \(\left( {{r}_{ij}\left( t\right) }\right)\) 为酉矩阵(可通过适当选取基实现，参见 1.3)。此时我们有 \({r}_{ij}\left( {t}^{-1}\right)  = {r}_{ji}{\left( t\right) }^{ * }\) ，且推论 2 与 3 就是下一节所定义的标量积 \(\left( {\phi  \mid  \psi }\right)\) 的正交关系。

\section*{2.3 特征的正交关系}

先给出一个记号。若 \(\phi\) 与 \(\chi\) 是群 G 上的两个复值函数，令

\[
\left( {\phi  \mid  \psi }\right)  = \frac{1}{g}\mathop{\sum }\limits_{{t \in  \mathrm{G}}}\phi \left( t\right) \psi {\left( t\right) }^{ * },\;g\text{ being the order of }\mathrm{G}.
\]

这是一个标量积:对 \(\phi\) 是线性的，对 \(\psi\) 是半线性的，并且对所有 \(\phi  \neq  0\) 我们有 \(\left( {\phi  \mid  \phi }\right)  > 0\) 。

如果 \(\check{\psi }\) 是由公式 \(\check{\psi }\left( t\right)  = \psi {\left( {t}^{-1}\right) }^{ * }\) 定义的函数，我们有

\[
\left( {\phi  \mid  \psi }\right)  = \frac{1}{g}\mathop{\sum }\limits_{{t \in  \mathrm{G}}}\phi \left( t\right) \check{\psi }\left( {t}^{-1}\right)  = \langle \phi ,\check{\psi }\rangle ,
\]

参见 2.2，注 1。特别地，若 \(\chi\) 是 G 的表示的迹字符，则我们有 \(\dot{\chi } = \chi\) (命题 1)，因此对 G 上任意函数 \(\phi\) 有 \(\left( {\phi  \mid  \chi }\right)  = \langle \phi ,\chi \rangle\) 。因此只要关心迹字符，就可以随意使用 \(\left( {\phi  \mid  \chi }\right)\) 或 \(\langle \phi ,\chi \rangle\) 。

\section*{定理 3}

(i) 若 \(\chi\) 是不可约表示的迹字符，则有 \(\left( {\chi  \mid  \chi }\right)  = 1\) (即 \(\chi\) “范数为 1”)。

ii) 若 \(\chi\) 与 \({\chi }^{\prime }\) 是两个互非同构的不可约表示的迹字符，则有 \(\left( {\chi  \mid  {\chi }^{\prime }}\right)  = 0\) (即 \(\chi\) 与 \({\chi }^{\prime }\) 正交)。

设 \(\rho\) 为一个不可约表示，迹字符为 \(\chi\) ，以矩阵形式表示为 \({\rho }_{t} = \left( {{r}_{ij}\left( t\right) }\right)\) 。我们有 \(\chi \left( t\right)  = \sum {r}_{ii}\left( t\right)\) ，因此

\[
\left( {\chi  \mid  \chi }\right)  = \langle \chi ,\chi \rangle  = \mathop{\sum }\limits_{{i,j}}\left\langle  {{r}_{ii},{r}_{jj}}\right\rangle  .
\]

但根据命题 4 的推论 3，我们有 \(\left\langle  {{r}_{ii},{r}_{jj}}\right\rangle   = {\delta }_{ij}/n\) ，其中 \(n\) 为 \(\rho\) 的次数。因此

\[
\left( {\chi  \mid  \chi }\right)  = \left( {\mathop{\sum }\limits_{{i,j}}{\delta }_{ij}}\right) /n = n/n = 1
\]

因为指标 \(i,j\) 各自取 \(n\) 个值。(ii) 的证明相同，改用推论 2 代替推论 3。

备注。不可约表示的迹字符称为不可约迹字符。定理 3 表明不可约迹字符构成一组正交归一系统；这个结果将在后文补全(2.5，定理 6)。

定理 4。设 \(\mathrm{V}\) 为群 \(\mathrm{G}\) 的线性表示，迹字符为 \(\phi\) ，并假定 \(\mathrm{V}\) 分解为若干不可约表示的直和:

\[
\mathrm{V} = {\mathrm{W}}_{1} \oplus  \cdots  \oplus  {\mathrm{W}}_{k}.
\]

那么，若 \(\mathrm{W}\) 为迹字符为 \(\chi\) 的不可约表示，和 \({\mathrm{W}}_{i}\) 同构的个数等于标量积 \(\left( {\phi  \mid  \chi }\right)  = \langle \phi ,\chi \rangle\) 。

设 \({\chi }_{i}\) 为 \({\mathrm{W}}_{i}\) 的迹字符。由命题 2，我们有

\[
\phi  = {\chi }_{1} + \cdots  + {\chi }_{k}
\]

因此 \(\left( {\phi  \mid  \chi }\right)  = \left( {{\chi }_{1} \mid  \chi }\right)  + \cdots  + \left( {{\chi }_{k} \mid  \chi }\right)\) 。但根据前述定理， \(\left( {{\chi }_{i} \mid  \chi }\right)\) 等于 1 或 0，取决于 \({\mathrm{W}}_{i}\) 是否与 W 同构。结论随之成立。

推论 1。与 \(\mathrm{W}\) 同构的 \({\mathrm{W}}_{i}\) 的个数不依赖于所选的分解。

(该数称为“W 在 V 中出现的次数”或“W 被包含在 V 中的次数”。)

确实， \(\left( {\phi  \mid  \chi }\right)\) 不依赖于分解。

备注。正是在这个意义上可以说表示分解为不可约表示时具有唯一性。我们将在2.6节回到此处。

推论2. 具有相同特征的两个表示是同构的。

事实上，推论1表明它们对每一个给定的不可约表示包含相同次数。

上述结果将表示的研究归结为其特征的研究。若 \({\chi }_{1},\ldots ,{\chi }_{h}\) 为 \(G\) 的不同不可约特征，且 \({\mathrm{W}}_{1},\ldots ,{\mathrm{W}}_{k}\) 表示对应的表示，则每个表示 \(\mathrm{V}\) 同构于直和

\[
\mathrm{V} = {m}_{1}{\mathrm{\;W}}_{1} \oplus  \cdots  \oplus  {m}_{h}{\mathrm{\;W}}_{h}\;{m}_{i}\text{ integers } \neq  0.
\]

\(\mathrm{V}\) 的特征 \(\phi\) 等于 \({m}_{1}{\chi }_{1} + \cdots  + {m}_{h}{\chi }_{h}\) ，并且我们有 \({m}_{i} = \left( {\phi  \mid  {\chi }_{i}}\right)\) 。[这对两个不可约表示的张量积 \({\mathrm{W}}_{i} \otimes  {\mathrm{W}}_{j}\) 尤为适用，并表明乘积 \({\chi }_{i} \cdot  {\chi }_{i}\) 分解为 \({\chi }_{i}{\chi }_{j} = \sum {m}_{ij}^{k}{\chi }_{k}\) ，其中 \({m}_{ij}^{k}\) 为整数 \(\geq  0\) 。] \({\chi }_{i}\) 之间的正交关系还意味着:

\[
\left( {\phi  \mid  \phi }\right)  = \mathop{\sum }\limits_{{i = 1}}^{{i = h}}{m}_{i}^{2}
\]

因此:

定理5. 若 \(\phi\) 是表示的特征， \(\mathrm{V},\left( {\phi  \mid  \phi }\right)\) 为正整数，并且我们有 \(\left( {\phi  \mid  \phi }\right)  = 1\) 当且仅当 \(\mathrm{V}\) 是不可约的。

确实，只有当某一 \({m}_{i}\) 等于1而其余等于0时， \(\sum {m}_{i}^{2}\) 才等于1，即当 \(\mathbf{V}\) 同构于某一 \({\mathbf{W}}_{i}\) 时。

由此得到一个非常方便的不可约性判据。

\section*{练习}

2.5. 设 \(\rho\) 为特征为 \(\chi\) 的线性表示。证明 \(\rho\) 包含单位表示的次数等于 \(\left( {\chi  \mid  1}\right) \; = \left( {1/g}\right) \mathop{\sum }\limits_{{s \in  \mathrm{G}}}\chi \left( s\right)\) 。

2.6. 设 \(\mathrm{X}\) 为有限集， \(\mathrm{G}\) 在其上作用， \(\rho\) 为相应的置换表示(1.2)， \(\chi\) 为其特征。

(a) 元素 \(x \in  \mathbf{X}\) 在 \(\mathbf{G}\) 下的像的集合 \(\mathbf{G}x\) 称为一个轨道。设不同轨道的数目为 \(c\) 。证明 \(c\) 等于 \(\rho\) 中包含平凡表示 1 的次数；由此推出 \(\left( {\chi  \mid  1}\right)  = c\) 。特别地，如果 \(G\) 是传递的(即如果 \(c = 1\) )，则 \(\rho\) 可以分解为 \(1 \oplus  \theta\) ，且 \(\theta\) 不包含平凡表示。如果 \(\psi\) 是 \(\theta\) 的特征，则有 \(\chi  = 1 + \psi\) 和 \(\left( {\psi  \mid  1}\right)  = 0\) 。

(b) 令 \(G\) 通过公式 \(s\left( {x,y}\right)  = \left( {{sx},{sy}}\right)\) 在 \(X\) 的积 \(X \times  X\) 上以自身作用。证明相应置换表示的迹为 \({\chi }^{2}\) 。

(c) 假设 \(\mathbf{G}\) 在 \(\mathbf{X}\) 上传递，且 \(\mathbf{X}\) 至少有两个元素。若对任意满足 \(x \neq  y\) 和 \({x}^{\prime } \neq  {y}^{\prime }\) 的 \(x,y,{x}^{\prime },{y}^{\prime } \in  \mathrm{X}\) ，存在 \(s \in  \mathrm{G}\) 使得 \({x}^{\prime } = {sx}\) 与 \({y}^{\prime } = {sy}\) 成立，则称 \(\mathbf{G}\) 为双传递。证明下列性质等价:

(i) \(G\) 为双传递。

(ii) \(G\) 在 \(X \times  X\) 上的作用有两个轨道，即对角线及其补集。

(iii) \(\left( {{\chi }^{2} \mid  1}\right)  = 2\) 。

(iv) 在 (a) 中定义的表示 \(\theta\) 是不可约的。

[等价 (i) \(\Leftrightarrow\) (ii) 显然；(ii) \(\Leftrightarrow\) (iii) 从 (a) 与 (b) 推出。若 \(\psi\) 是 \(\theta\) 的迹，则有 \(1 + \psi  = \chi\) 与 \(\left( {1 \mid  1}\right)  = 1,\left( {\psi  \mid  1}\right)  = 0\) ，由此可见 (iii) 等价于 \(\left( {{\psi }^{2} \mid  1}\right)  = 1\) ，即等价于 \(\left( {1/g}\right) \mathop{\sum }\limits_{{s \in  \mathrm{G}}}\psi {\left( s\right) }^{2} = 1\) ；由于 \(\psi\) 为实值，这确实意味着 \(\theta\) 不可约，参见定理 5。]

\section*{2.4 正则表示的分解}

记号。于第 2 章余下部分，G 的不可约特征记为 \({\chi }_{1},\ldots ,{\chi }_{h}\) ；其次数记为 \({n}_{1},\ldots ,{n}_{k}\) ；我们有 \({n}_{i} \; = {\chi }_{i}\left( 1\right)\) ，参见命题 1。

设 \(\mathrm{R}\) 为 \(\mathrm{G}\) 的正则表示。回想(参见 1.2)它有一组基 \({\left( {e}_{t}\right) }_{t \in  \mathrm{G}}\) 使得 \({\rho }_{s}{e}_{t} = {e}_{st}\) 。若 \(s \neq  1\) ，则对所有 \(t\) 有 \({st} \neq  t\) ，这表明 \({\rho }_{s}\) 的矩阵的对角项为零；特别地我们有 \(\operatorname{Tr}\left( {\rho }_{s}\right)  = 0\) 。另一方面，对 \(s = 1\) ，我们有

\[
\operatorname{Tr}\left( {\rho }_{s}\right)  = \operatorname{Tr}\left( 1\right)  = \dim \left( \mathrm{R}\right)  = g.
\]

因此:

命题 5. 正则表示的特征 \({r}_{\mathrm{G}}\) 由下列公式给出:

\[
{r}_{\mathrm{G}}\left( 1\right)  = g,\;\text{ order of }\mathrm{G},
\]

\[
{r}_{\mathrm{G}}\left( s\right)  = 0\;\text{ if }s \neq  1.
\]

推论 1. 每个不可约表示 \({\mathbf{W}}_{i}\) 都包含在正则表示中，其重数等于其次数 \({n}_{i}\) 。

根据定理 4，此数等于 \(\left\langle  {{r}_{\mathrm{G}},{\chi }_{i}}\right\rangle\) ，并且我们有

\[
\left\langle  {{r}_{\mathrm{G}},{\chi }_{i}}\right\rangle   = \frac{1}{g}\mathop{\sum }\limits_{{s \in  \mathrm{G}}}{r}_{\mathrm{G}}\left( {s}^{-1}\right) {\chi }_{i}\left( s\right)  = \frac{1}{g}g \cdot  {\chi }_{i}\left( 1\right)  = {\chi }_{i}\left( 1\right)  = {n}_{i}.
\]

推论 2。

(a) 次数 \({n}_{i}\) 满足关系式 \(\mathop{\sum }\limits_{{i = 1}}^{{i = h}}{n}_{i}^{2} = g\) 。

(b) 若 \(s \in  \mathbf{G}\) 不等于 1，则有 \(\mathop{\sum }\limits_{{i = 1}}^{{i = h}}{n}_{i}{\chi }_{i}\left( s\right)  = 0\) 。

由推论 1，可得对所有 \(s \in  \mathrm{G}\) 都有 \({r}_{\mathrm{G}}\left( s\right)  = \sum {n}_{i}{\chi }_{i}\left( s\right)\) 。取 \(s = 1\) 得到 (a)，取 \(s \neq  1\) 得到 (b)。

\section*{备注}

(1) 上述结果可用于确定群 G 的不可约表示:设已构造出若干两两非同构的不可约表示，其次数为 \({n}_{1},\ldots ,{n}_{k}\) ；要使它们成为 \(G\) 的所有不可约表示(同构意义下)，充要条件是 \({n}_{1}^{2} + \cdots  + {n}_{k}^{2} = g\) 。

(2) 我们将在后文(第二部分 6.5)看到次数 \({n}_{i}\) 的另一条性质:它们整除 \(g\) 的阶数 \(\mathrm{G}\) 。

\section*{练习}

2.7. 证明:若一切在所有 \(s \neq  1\) 上为零的 \(G\) 的特征，是正规表示的特征 \({r}_{\mathrm{G}}\) 的整数倍。

\section*{2.5 不可约表示的个数}

回忆(参见 2.1):若对所有 \(s,t \in  \mathrm{G}.\) 有 \(f\left( {{ts}{t}^{-1}}\right)  = f\left( s\right)\) ，则定义在 \(\mathbf{G}\) 上的函数 \(f\) 称为类函数。

命题 6. 令 \(f\) 为定义在 \(\mathbf{G}\) 上的类函数，令 \(\rho  : \mathbf{G} \rightarrow  \mathbf{{GL}}\left( \mathbf{V}\right)\) 为 \(\mathbf{G}\) 的一个线性表示。令 \({\rho }_{f}\) 为将 \(\mathbf{V}\) 映到自身的线性映射，定义如下:

\[
{\rho }_{f} = \mathop{\sum }\limits_{{t \in  \mathrm{G}}}f\left( t\right) {\rho }_{t}
\]

若 \(\mathrm{V}\) 为次数为 \(n\) 、特征为 \(\chi\) 的不可约表示，则 \({\rho }_{f}\) 为比率为 \(\lambda\) 的相似变换，其给出为:

\[
\lambda  = \frac{1}{n}\mathop{\sum }\limits_{{t \in  \mathrm{G}}}f\left( t\right) \chi \left( t\right)  = \frac{g}{n}\left( {f \mid  {\chi }^{ * }}\right) .
\]

让我们计算 \({\rho }_{s}^{-1}{\rho }_{f}{\rho }_{s}\) 。我们有:

\[
{\rho }_{s}^{-1}{\rho }_{f}{\rho }_{s} = \mathop{\sum }\limits_{{t \in  \mathrm{G}}}f\left( t\right) {\rho }_{s}^{-1}{\rho }_{t}{\rho }_{s} = \mathop{\sum }\limits_{{t \in  \mathrm{G}}}f\left( t\right) {\rho }_{{s}^{-1}{ts}}.
\]

取 \(u = {s}^{-1}{ts}\) 后，变为:

\[
{\rho }_{s}^{-1}{\rho }_{f}{\rho }_{s} = \mathop{\sum }\limits_{{u \in  \mathrm{G}}}f\left( {{su}{s}^{-1}}\right) {\rho }_{u} = \mathop{\sum }\limits_{{u \in  \mathrm{G}}}f\left( u\right) {\rho }_{u} = {\rho }_{f}.
\]

于是我们有 \({\rho }_{f}{\rho }_{s} = {\rho }_{s}{\rho }_{f}\) 。由命题 4 的第二部分，这表明 \({\rho }_{f}\) 为一个比率为 \(\lambda\) 的相似变换。 \(\lambda\) 的迹为 \({n\lambda }\) ； \({\rho }_{f}\) 的迹为 \(\mathop{\sum }\limits_{{t \in  \mathrm{G}}}f\left( t\right) \operatorname{Tr}\left( {\rho }_{t}\right) \; = \mathop{\sum }\limits_{{t \in  \mathrm{G}}}f\left( t\right) \chi \left( t\right)\) 。因此 \(\lambda  = \left( {1/n}\right) \mathop{\sum }\limits_{{t \in  \mathrm{G}}}f\left( t\right) \chi \left( t\right)  = \left( {g/n}\right) \left( {f \mid  {\chi }^{ * }}\right)\) 。

我们现在引入类函数空间 \(H\) 在 \(G\) 上；不可约特征 \({\chi }_{1},\ldots ,{\chi }_{h}\) 属于 \(H\) 。

定理 6。字符 \({\chi }_{1},\ldots ,{\chi }_{h}\) 构成 \(\mathrm{H}\) 的一组正交归一基。

定理3表明 \({\chi }_{i}\) 在 \(\mathrm{H}\) 中构成一组正交归一系。还需证明它们生成 \(\mathrm{H}\) ，为此只需说明每个与 \({\chi }_{i}^{ * }\) 正交的 \(\mathrm{H}\) 中的元素为零。设 \(f\) 为此类元素。对 \(G\) 的任一表示 \(\rho\) ，置 \({\rho }_{f} = \mathop{\sum }\limits_{{t \in  G}}f\left( t\right) {\rho }_{t}\) 。由于 \(f\) 与 \({\chi }_{i}^{ * }\) 正交，上述命题6表明只要 \(\rho\) 不可约， \({\rho }_{f}\) 即为零；由直和分解我们得出 \({\rho }_{f}\) 始终为零。将此结论应用于正则表示R(参见2.4)并计算基向量 \({e}_{1}\) 在 \({\rho }_{f}\) 下的像，得

\[
{\rho }_{f}{e}_{1} = \mathop{\sum }\limits_{{t \in  \mathrm{G}}}f\left( t\right) {\rho }_{t}{e}_{1} = \mathop{\sum }\limits_{{t \in  \mathrm{G}}}f\left( t\right) {e}_{t}.
\]

由于 \({\rho }_{f}\) 为零，我们有 \({\rho }_{f}{e}_{1} = 0\) ，且上述公式表明对任意 \(t \in  \mathrm{G}\) 都有 \(f\left( t\right)  = 0\) ；因此 \(f = 0\) ，证明完毕。

回想两元素 \(t\) 与 \({t}^{\prime }\) 属于 \(\mathrm{G}\) 若存在 \(s \in  \mathrm{G}\) 使得 \({t}^{\prime } = {st}{s}^{-1}\) ，则称其互为共轭；这是一种等价关系，将 \(\mathbf{G}\) 划分为若干类(亦称共轭类)。

定理7. \(\mathbf{G}\) 的不既约表示(到同构)数等于 \(\mathbf{G}\) 的共轭类数。

设 \({C}_{1},\ldots ,{C}_{k}\) 为G的不重复类。说函数 \(f\) 在G上是类函数等价于说它在每个 \({\mathbf{C}}_{1},\ldots ,{\mathbf{C}}_{k}\) 上为常值；因此它由在 \({\mathbf{C}}_{i}\) 上的取值 \({\lambda }_{i}\) 决定，且这些取值可任意选择。于是类函数空间H的维数等于 \(k\) 。另一方面，按定理6，该维数等于 \(G\) 的不既约表示(到同构)的数。由此得证。

以下是定理6的另一个推论:

命题7. 设 \(s \in  \mathrm{G}\) ，并设 \(c\left( s\right)\) 为 \(s\) 的共轭类中元素的个数。

(a) 我们有 \(\mathop{\sum }\limits_{{i = 1}}^{{i = h}}{\chi }_{i}{\left( s\right) }^{ * }{\chi }_{i}\left( s\right)  = g/c\left( s\right)\) 。

(b) 若 \(t \in  \mathrm{G}\) 不与 \(s\) 共轭，则有 \(\mathop{\sum }\limits_{{i = 1}}^{{i = h}}{\chi }_{i}{\left( s\right) }^{ * }{\chi }_{i}\left( t\right)  = 0\) 。(对 \(s = 1\) 而言，此得出命题5的系论2。)

令 \({f}_{s}\) 为在 \(s\) 的类上取1、其余处取0的函数。由于它是类函数，按定理6可写成

\[
{f}_{s} = \mathop{\sum }\limits_{{i = 1}}^{{i = h}}{\lambda }_{i}{\chi }_{i},\;\text{ with }{\lambda }_{i} = \left( {{f}_{s} \mid  {\chi }_{i}}\right)  = \frac{c\left( s\right) }{g}{\chi }_{i}{\left( s\right) }^{ * }.
\]

于是对每个 \(t \in  \mathrm{G}\) ，我们有

\[
{f}_{s}\left( t\right)  = \frac{c\left( s\right) }{g}\mathop{\sum }\limits_{{i = 1}}^{{i = h}}{\chi }_{i}{\left( s\right) }^{ * }{\chi }_{i}\left( t\right) .
\]

若 \(t = s\) ，则得出(a)；若 \(t\) 不与 \(s\) 共轭，则得出(b)。

例。取 \(\mathrm{G}\) 为三字母置换群。我们有 \(g = 6\) ，且有三个共轭类:恒等元 1、三个互换和两个循环置换。令 \(t\) 为一个互换， \(c\) 为一个循环置换。我们有 \({t}^{2} = 1,{c}^{3} = 1,{tc} = {c}^{2}t\) ；因此仅有两个一次度的特征:单位特征 \({\chi }_{1}\) 和给出置换符号的特征 \({\chi }_{2}\) 。定理 7 表明还存在另一个不可约特征 \(\theta\) ；若 \(n\) 为其次数，必有 \(1 + 1 + {n}^{2} = 6\) ，因此 \(n = 2\) 。由 \({\chi }_{1} + {\chi }_{2} \; + {2\theta }\) 是 \(G\) 的正则表示的特征(参见命题 5)可推得 \(\theta\) 的取值。由此得到 \(G\) 的特征表:

\begin{center}
\adjustbox{max width=\textwidth}{
\begin{tabular}{|c|c|c|c|}
\hline
\phantom{X} & 1 & \(t\) & C \\
\cline{1-4}
\({x}_{1}\) & 1 & 1 & 1 \\
\cline{1-4}
\({x}_{2}\) & 1 & -1 & 1 \\
\cline{1-4}
\(\theta\) & 2 & 0 & -1 \\
\cline{1-4}
\hline
\end{tabular}
}
\end{center}

我们通过让 \(\mathbf{G}\) 对满足方程 \(x + y \; + z = 0\) 的 \({\mathbf{C}}^{3}\) 的坐标进行置换来得到一个特征为 \(\theta\) 的不可约表示(参见例 2.6c))。

\section*{2.6 表示的规范分解}

设 \(\rho  : G \rightarrow  {GL}\left( V\right)\) 为群 \(G\) 的线性表示。我们将定义 \(V\) 的一个比不可约表示分解“更粗”的直和分解，但该分解的优点是唯一。其构造如下:

设 \({\chi }_{1},\ldots ,{\chi }_{h}\) 为 \(\mathrm{G}\) 的不可约表示 \({\mathrm{W}}_{1},\ldots ,{\mathrm{W}}_{h}\) 的不同特征标， \({n}_{1},\ldots ,{n}_{h}\) 为它们的次数。设 \(\mathrm{V} = {\mathrm{U}}_{1} \oplus  \cdots \; \oplus  {\mathrm{U}}_{m}\) 是 \(\mathrm{V}\) 分解为不可约表示的直和。对于 \(i = 1,\ldots ,h\) ，用 \({\mathrm{V}}_{i}\) 表示那些与 \({\mathrm{W}}_{i}\) 同构的 \({\mathrm{U}}_{1},\ldots ,{\mathrm{U}}_{m}\) 的直和。显然，我们有:

\[
\mathrm{V} = {\mathrm{V}}_{1} \oplus  \cdots  \oplus  {\mathrm{V}}_{h}
\]

(换言之，我们已将 V 分解为不可约表示的直和，并把同构的表示合并在一起。)

这就是我们所说的规范分解。其性质如下:

\section*{定理 8}

(i) 分解 \(\mathrm{V} = {\mathrm{V}}_{1} \oplus  \cdots  \oplus  {\mathrm{V}}_{h}\) 不依赖于最初选定的将 \(\mathrm{V}\) 分解成不可约表示的方式。

(ii) 与该分解相关的将 \(\mathbf{V}\) 投影到 \({\mathbf{V}}_{i}\) 的投影算子 \({p}_{i}\) 由下式给出:

\[
{p}_{i} = \frac{{n}_{i}}{g}\mathop{\sum }\limits_{{t \in  \mathrm{G}}}{\chi }_{i}{\left( t\right) }^{ * }{\rho }_{t}.
\]

我们证明 (ii)。断言 (i) 将随之成立，因为投影 \({p}_{i}\) 决定了 \({\mathrm{V}}_{i}\) 。记

\[
{q}_{i} = \frac{{n}_{i}}{g}\mathop{\sum }\limits_{{t \in  \mathrm{G}}}{\chi }_{i}{\left( t\right) }^{ * }{\rho }_{t}.
\]

命题 6 表明，将 \({q}_{i}\) 限制到具有特征 \(\chi\) 且次数为 \(n\) 的不可约表示 \(W\) 上是比率为 \(\left( {{n}_{i}/n}\right) \left( {{\chi }_{i} \mid  \chi }\right)\) 的同比变换；因此当 \(\chi  \neq  {\chi }_{i}\) 时它为 0，当 \(\chi  = {\chi }_{i}\) 时它为 1。换言之， \({q}_{i}\) 在与 \({\mathbf{W}}_{i}\) 同构的不可约表示上是恒等，而在其他表示上为零。根据对 \({\mathbf{V}}_{i}\) 的定义，得出 \({q}_{i}\) 在 \({\mathbf{V}}_{i}\) 上是恒等，在 \({\mathbf{V}}_{j}\) 上为 0(对于 \(j \neq  i\) )。若将元素 \(\chi  \in  \mathbf{V}\) 分解为其分量 \({x}_{i} \in  {\mathrm{V}}_{i}\) :

\[
x = {x}_{1} + \cdots  + {x}_{h}
\]

于是我们有 \({q}_{i}\left( x\right)  = {q}_{i}\left( {x}_{1}\right)  + \cdots  + {q}_{i}\left( {x}_{h}\right)  = {x}_{i}\) 。这意味着 \({q}_{i}\) 等于将 \(\mathrm{V}\) 投影到 \({\mathrm{V}}_{i}\) 的投影 \({p}_{i}\) 。

因此，一个表示的分解 \(\mathrm{V}\) 可以分两步完成。先确定典范分解 \({\mathrm{V}}_{1} \oplus  \cdots  \oplus  {\mathrm{V}}_{n}\) ；这可以用给出投影的公式 \({p}_{i}\) 轻易完成。接着，如有需要，选择将 \({V}_{i}\) 分解为若干个各同构于 \({\mathrm{W}}_{i}\) 的不可约表示的直和:

\[
{\mathbf{V}}_{i} = {\mathbf{W}}_{i} \oplus  \cdots  \oplus  {\mathbf{W}}_{i}.
\]

最后这一步分解通常有无穷多种方式可选(参见 2.7 节及下文例 2.8)；其随意性如同在向量空间中选取基。

例。取 \(G\) 为两个元素的群 \(\{ 1,s\}\) 且具有 \({s}^{2} = 1\) 。该群有两个度为 \(1,{\mathrm{\;W}}^{ + }\) 与 \({\mathrm{W}}^{ - }\) 的不可约表示，分别对应于 \({\rho }_{s} =  + 1\) 与 \({\rho }_{s} =  - 1\) 。表示 \(V\) 的典范分解为 \(V = {V}^{ + } \oplus  {V}^{ - }\) ，其中 \({V}^{ + }\) (分别 \({V}^{ - }\) )由对称(分别反对称)元素 \(x \in  V\) 组成，即满足 \({\rho }_{s}x = x\) (分别 \({\rho }_{s}x =  - x\) )。相应的投影为:

\[
{p}^{ + }x = \frac{1}{2}\left( {x + {\rho }_{s}x}\right) ,\;{p}^{ - }x = \frac{1}{2}\left( {x - {\rho }_{s}x}\right) .
\]

将 \({\mathrm{V}}^{ + }\) 和 \({\mathrm{V}}^{ - }\) 分解为不可约分量，仅仅意味着把这些空间分解成若干直和的直线。

\section*{习题}

2.8. 设 \({\mathrm{H}}_{i}\) 为线性映射空间 \(h : {\mathrm{W}}_{i} \rightarrow  \mathrm{V}\) ，满足对所有 \(s \in  \mathbf{G}\) 有 \({\rho }_{s}h = h{\rho }_{s}\) 。每个 \(h \in  {\mathrm{H}}_{i}\) 将 \({\mathrm{W}}_{i}\) 映到 \({\mathrm{V}}_{i}\) 。

(a) 证明 \({\mathrm{H}}_{i}\) 的维数等于 \({\mathrm{W}}_{i}\) 在 \(\mathrm{V}\) 中出现的次数，即 \(\dim {\mathrm{V}}_{i}/\dim {\mathrm{W}}_{i}\) [化约到 \(\mathrm{V} = {\mathrm{W}}_{i}\) 的情况并使用 Schur 引理]。

(b) 令 \(\mathrm{G}\) 通过在 \({\mathrm{H}}_{i}\) 上的平凡表示与在 \({\mathrm{W}}_{i}\) 上的给定表示的张量积作用于 \({\mathrm{H}}_{i} \otimes  {\mathrm{W}}_{i}\) 。证明映射

\[
\mathrm{F} : {\mathrm{H}}_{i} \otimes  {\mathrm{W}}_{i} \rightarrow  {\mathrm{V}}_{i}
\]

由公式定义

\[
\mathrm{F}\left( {\sum {h}_{\alpha } \cdot  {w}_{\alpha }}\right)  = \sum {h}_{\alpha }\left( {w}_{\alpha }\right)
\]

是将 \({\mathrm{H}}_{i} \otimes  {\mathrm{W}}_{i}\) 同构到 \({\mathrm{V}}_{i}\) 的同构映射。[方法同上。]

(c) 设 \(\left( {{h}_{1},\ldots ,{h}_{k}}\right)\) 为 \({\mathrm{H}}_{i}\) 的一组基，并构成 \(k\) 个 \({\mathrm{W}}_{i}\) 的直和记为 \({\mathrm{W}}_{i} \oplus  \cdots  \oplus  {\mathrm{W}}_{i}\) 。由 \(\left( {{h}_{1},\ldots ,{h}_{k}}\right)\) 明显定义了从 \({\mathrm{W}}_{i} \oplus  \cdots  \oplus  {\mathrm{W}}_{i}\) 到 \({\mathrm{V}}_{i}\) 的线性映射 \(h\) ；证明它是表示的同构，且每一同构都可由此得到[可应用 \(\left( b\right)\) 或直接论证]。特别地，要将 \({\mathrm{V}}_{i}\) 分解为与 \({\mathrm{W}}_{i}\) 同构的表示的直和，相当于为 \({\mathrm{H}}_{i}\) 选一组基。

\section*{2.7 表示的显式分解}

保持前一节的记号，并令

\[
\mathrm{V} = {\mathrm{V}}_{1} \oplus  \cdots  \oplus  {\mathrm{V}}_{h}
\]

为给定表示的典型分解。我们已见如何通过相应的投影(定理8)确定第 \(i\) 分量 \({V}_{i}\) 。现给出一个将 \({V}_{i}\) 显式构造成与 \({W}_{i}\) 同构的子表示直和的方法。设 \({\mathrm{W}}_{i}\) 在基 \(\left( {{e}_{1},\ldots ,{e}_{n}}\right)\) 下以矩阵形式给出为 \(\left( {{r}_{\alpha \beta }\left( s\right) }\right)\) ；我们有 \({\chi }_{i}\left( s\right)  = \mathop{\sum }\limits_{\alpha }{r}_{\alpha \alpha }\left( s\right)\) 与 \(n = {n}_{i} = \dim {\mathrm{W}}_{i}\) 。对每一对从1到 \(n\) 的整数 \(\alpha ,\beta\) ，令 \({p}_{\alpha \beta }\) 表示由下式在 \(\mathrm{V}\) 上定义的线性映射

(*)
\[
{p}_{\alpha \beta } = \frac{n}{g}\mathop{\sum }\limits_{{t \in  \mathrm{G}}}{r}_{\beta \alpha }\left( {t}^{-1}\right) {\rho }_{t}.
\]

\section*{命题 8}

(a) 映射 \({p}_{\alpha \alpha }\) 是一个投影；它在 \({\mathbf{V}}_{j},j \neq  i\) 上为零。其像 \({\mathbf{V}}_{i,\alpha }\) 包含于 \({\mathbf{V}}_{i}\) ，且对于所有 \(1 \leq  \alpha  \leq  n\) ， \({\mathbf{V}}_{i}\) 为这些 \({\mathbf{V}}_{i,\alpha }\) 的直和。我们有 \({p}_{i} = \mathop{\sum }\limits_{\alpha }{p}_{\alpha \alpha }\) 。

(b) 线性映射 \({p}_{\alpha \beta }\) 在 \({\mathrm{V}}_{j},j \neq  i\) 上为零，且对于所有 \(\gamma  \neq  \beta\) ，在那些 \({\mathrm{V}}_{i,\gamma }\) 上也为零；它把 \({\mathrm{V}}_{i,\beta }\) 同构到 \({\mathrm{V}}_{i,\alpha }\) 。

(c) 设 \({x}_{1}\) 为 \({\mathrm{V}}_{i,1}\) 中的元素 \(\neq  0\) ，并令 \({x}_{\alpha } = {p}_{\alpha 1}\left( {x}_{1}\right)  \in  {\mathrm{V}}_{i,\alpha }\) 。那些 \({x}_{\alpha }\) 线性无关并生成在 \(\mathrm{G}\) 下不变且维数为 \(n\) 的向量子空间 \(\mathrm{W}\left( {x}_{1}\right)\) 。对于每个 \(s \in  \mathrm{G}\) ，我们有

\[
{\rho }_{s}\left( {x}_{\alpha }\right)  = \mathop{\sum }\limits_{\beta }{r}_{\beta \alpha }\left( s\right) {x}_{\beta }
\]

(尤其， \(\mathrm{W}\left( {x}_{1}\right)\) 与 \({\mathrm{W}}_{i}\) 同构)。

(d) 如果 \(\left( {{x}_{1}^{\left( 1\right) },\ldots ,{x}_{1}^{\left( m\right) }}\right)\) 是 \({\mathbf{V}}_{i,1}\) 的一个基，则表示 \({\mathbf{V}}_{i}\) 是由在 \(c)\) 中定义的子表示 \(\mathrm{W}\left( {x}_{1}^{\left( 1\right) }\right) ,\ldots ,\mathrm{W}\left( {x}_{1}^{\left( m\right) }\right)\) 的直和。

(因此选择 \({V}_{i,1}\) 的一组基给出 \({V}_{i}\) 分解为若干与 \({\mathrm{W}}_{i}\) 同构的表示的直和。)

我们首先注意，上述公式 \(\left( *\right)\) 允许我们在任意 \(\mathbf{G}\) 的表示中定义 \({p}_{\alpha \beta }\) ，特别是在不可约表示 \({\mathbf{W}}_{j}\) 中。对于 \({\mathbf{W}}_{i}\) ，我们有

\[
{p}_{\alpha \beta }\left( {e}_{\gamma }\right)  = \frac{n}{g}\mathop{\sum }\limits_{{t \in  \mathrm{G}}}{r}_{\beta \alpha }\left( {t}^{-1}\right) {\rho }_{t}\left( {e}_{\gamma }\right)  = \frac{n}{g}\mathop{\sum }\limits_{\delta }\mathop{\sum }\limits_{{t \in  \mathrm{G}}}{r}_{\beta \alpha }\left( {t}^{-1}\right) {r}_{\delta \gamma }\left( t\right) {e}_{\delta }.
\]

由命题4 的推论3 我们得到

\[
{p}_{\alpha \beta }\left( {e}_{\gamma }\right)  = \left\{  \begin{array}{ll} {e}_{\alpha } & \text{ if }\gamma  = \beta \\  0 & \text{ otherwise }. \end{array}\right.
\]

由此可得 \(\mathop{\sum }\limits_{\alpha }{p}_{\alpha \alpha }\) 是 \({\mathbf{W}}_{i}\) 的恒等映射，以及如下公式

\[
{p}_{\alpha \beta } \circ  {p}_{\gamma \delta } = \left\{  \begin{array}{ll} {p}_{\alpha \delta } & \text{ if }\beta  = \gamma \\  0 & \text{ otherwise } \end{array}\right.
\]

\[
{\rho }_{s} \circ  {p}_{\alpha \gamma } = \mathop{\sum }\limits_{\beta }{r}_{\beta \alpha }\left( s\right) {p}_{\beta \gamma }
\]

对于 带有 \(j \neq  i\) 的 \({\mathrm{W}}_{j}\) ，我们使用命题4 的推论2 并用相同论证证明所有 \({p}_{\alpha \beta }\) 都为零。

完成此项后，我们将 V 分解为与 \({W}_{i}\) 同构的子表示的直和，并对这些表示中的每一个应用上述结论。断言 (a) 与 (b) 随之成立；此外，上述公式在 V 中仍然成立。在 (c) 的假设下，我们有

\[
{\rho }_{s}\left( {x}_{\alpha }\right)  = {\rho }_{s} \circ  {p}_{\alpha 1}\left( {x}_{1}\right)  = \mathop{\sum }\limits_{\beta }{r}_{\beta \alpha }\left( s\right) {p}_{\beta 1}\left( {x}_{1}\right)  = \mathop{\sum }\limits_{\beta }{r}_{\beta \alpha }\left( s\right) {x}_{\beta },
\]

这证明了 (c)。最后 (d) 从 (a)、(b) 和 (c) 推出。

\section*{练习}

2.9. 设 \({\mathrm{H}}_{i}\) 为线性映射空间 \(h : {\mathrm{W}}_{i} \rightarrow  \mathrm{V}\) ，使得 \(h \circ  {\rho }_{s} = {\rho }_{s} \circ  h\) ，参见练习 2.8。证明映射 \(h \mapsto  h\left( {e}_{\alpha }\right)\) 是从 \({\mathrm{H}}_{i}\) 到 \({\mathrm{V}}_{i,\alpha }\) 的同构。

2.10. 设 \(x \in  {\mathrm{V}}_{i}\) ，并设 \(\mathrm{V}\left( x\right)\) 是包含 \(x\) 的 \(\mathrm{V}\) 的最小子表示。设 \({x}_{1}^{\alpha }\) 是 \(x\) 在 \({p}_{1\alpha }\) 下的像；证明 \(\mathrm{V}\left( x\right)\) 是表示 \(\mathrm{W}\left( {x}_{1}^{\alpha }\right) ,\alpha  = 1,\ldots ,n\) 的和。由此推断 \(\mathrm{V}\left( x\right)\) 是至多 \(n\) 个同构于 \({\mathrm{W}}_{i}\) 的子表示的直和。

\section*{第3章 子群、直积、诱导表示}

下述所有群都假定为有限群。

\section*{3.1 阿贝尔子群}

设 \(\mathrm{G}\) 为一个群。若对所有 \(s,t \in  \mathrm{G}\) 都有 \({st} = {ts}\) ，则称 \(\mathrm{G}\) 为阿贝尔(或交换)群。这等价于说 \(\mathrm{G}\) 的每个共轭类只包含一个元素，也等价于每个关于 \(\mathrm{G}\) 的函数都是类函数。此类群的线性表示格外简单:

定理9. 下列性质等价:

(i) \(\mathrm{G}\) 是阿贝尔群。

(ii) \(\mathrm{G}\) 的所有不可约表示的维数均为 1 。

设 \(g\) 为 \(\mathrm{G}\) 的阶，且 \(\left( {{n}_{1},\ldots ,{n}_{h}}\right)\) 为 \(\mathrm{G}\) 不同不可约表示的维数；我们知道，参见第 2 章， \(h\) 是 \(\mathbf{G}\) 的类数，并且 \(g = {n}_{1}^{2} + \cdots  + {n}_{h}^{2}\) 。因此， \(g\) 当且仅当所有 \({n}_{i}\) 都等于 1 时等于 \(h\) ，这证明了定理。

推论。设 \(\mathrm{A}\) 为 \(\mathrm{G}\) 的一个阿贝尔子群，设其阶为 a， \(g\) 为 \(\mathrm{G}\) 的阶。 \(\mathrm{G}\) 的每个不可约表示的维数为 \(\leq  g/a\) 。

(商 \(g/a\) 是 \(\mathbf{G}\) 中 \(\mathbf{A}\) 的指标。)

设 \(\rho  : \mathbf{G} \rightarrow  \mathbf{{GL}}\left( \mathbf{V}\right)\) 为 \(\mathbf{G}\) 的一个不可约表示。通过限制到子群 \(A\) ，它在 A 上给出一个表示 \({\rho }_{A} : A \rightarrow  {GL}\left( V\right)\) 。设 \(W \subset  V\) 为 \({\rho }_{A}\) 的一个不可约子表示；由定理 9，我们有 \(\dim \left( W\right)  = 1\) 。令 \({V}^{\prime }\) 为由 \(V\) 中的像 \({\rho }_{s}W\) (当 \(W,s\) 在 G 中取遍时)生成的向量子空间。显然 \({V}^{\prime }\) 在 \(G\) 下不变；由于 \(\rho\) 是不可约的，我们因此有 \({\mathrm{V}}^{\prime } = \mathrm{V}\) 。但对于 \(s \in  \mathrm{G}\) 和 \(t \in  \mathrm{A}\) 我们有

\[
{\rho }_{st}\mathbf{W} = {\rho }_{s}{\rho }_{t}\mathbf{W} = {\rho }_{s}\mathbf{W}.
\]

由此可得不同 \({\rho }_{s}\mathrm{\;W}\) 的数量至多等于 \(g/a\) ，因此得到所需的不等式 \(\dim \left( \mathbf{V}\right)  \leq  g/a\) ，因为 \(\mathbf{V}\) 是这些 \({\rho }_{s}\mathbf{W}\) 的和。

例:二面体群含有一个指数为 2 的循环子群。因此其不可约表示的维数为 1 或 2；我们将在后文(5.3)确定它们。

练习

3.1. 直接用舒尔引理证明:阿贝尔群的每个不可约表示(无论是否有限)维数均为 1。

3.2. 设 \(\rho\) 为 \(G\) 的一个不可约表示，维数为 \(n\) ，特征为 \(\chi\) ；设 C 为 \(G\) 的中心(即满足 \({st} = {ts}\) 对所有 \(t \in  G\) 成立的那些 \(s \in  G\) 的集合)，并设其阶为 \(c\) 。

(a) 证明对每个 \(s \in  \mathbf{C}\) ， \({\rho }_{s}\) 是一个齐性变换。[使用舒尔引理。] 从中推得对所有 \(s \in  \mathbf{C}\) 有 \(\left| {\chi \left( s\right) }\right|  = n\) 。

(b) 证明不等式 \({n}^{2} \leq  g/c\) 。[结合 (a) 使用公式 \(\mathop{\sum }\limits_{{s \in  G}}{\left| \chi \left( s\right) \right| }^{2} = g\) 。]

(c) 证明:若 \(\rho\) 具有效忠性(即对 \(s \neq  1\) 满足 \({\rho }_{s} \neq  1\) )，则群 \(\mathbf{C}\) 为循环群。

3.3. 设 \(G\) 为阶为 \(g\) 的阿贝尔群，且令 \(\widehat{G}\) 为 \(G\) 的不可约特征的集合。若 \({\chi }_{1},{\chi }_{2}\) 属于 \(\widehat{G}\) ，则其乘积 \({\chi }_{1}{\chi }_{2}\) 亦属此集合。证明这使得 \(\widehat{G}\) 成为阶为 \(g\) 的阿贝尔群；群 \(\widehat{G}\) 称为群 G 的对偶。对任意 \(x \in  G\) ，映射 \(\chi  \mapsto  \chi \left( x\right)\) 是 \(\widehat{G}\) 的一个不可约特征，因而为 \(\widehat{G}\) 的一个元素。证明由此得到的从 G 到 G 的映射是单同态；比较两群的阶，得出该映射为同构。

\section*{3.2 两群的乘积}

设 \({G}_{1}\) 和 \({G}_{2}\) 为两群，令 \({G}_{1} \times  {G}_{2}\) 为它们的乘积，即由所有对 \(\left( {{s}_{1},{s}_{2}}\right)\) 组成的集合，其中 \({s}_{1} \in  {\mathrm{G}}_{1}\) 与 \({s}_{2} \in  {\mathrm{G}}_{2}\) 分别属于各自的群。

定义如下

\[
\left( {{s}_{1},{s}_{2}}\right)  \cdot  \left( {{t}_{1},{t}_{2}}\right)  = \left( {{s}_{1}{s}_{2},{t}_{1}{t}_{2}}\right) ,
\]

我们在 \({G}_{1} \times  {G}_{2}\) 上定义群结构；赋予此结构后， \({\mathrm{G}}_{1} \times  {\mathrm{G}}_{2}\) 称为 \({\mathrm{G}}_{1}\) 与 \({\mathrm{G}}_{2}\) 的群乘积。若 \({\mathrm{G}}_{1}\) 的阶为 \({g}_{1}\) ，且 \({\mathrm{G}}_{2}\) 的阶为 \({g}_{2},{\mathrm{G}}_{1} \times  {\mathrm{G}}_{2}\) ，则 \({G}_{1} \times  {G}_{2}\) 的阶为 \(g = {g}_{1}{g}_{2}\) 。群 \({\mathrm{G}}_{1}\) 可被识别为 \({G}_{1} \times  {G}_{2}\) 的子群，包含形如 \(\left( {{s}_{1},1}\right)\) 的元素，其中 \({s}_{1}\) 在 \({G}_{1}\) 上变动；类似地， \({G}_{2}\) 可识别为 \({G}_{1} \times  {G}_{2}\) 的一个子群。按此识别，每个来自 \({G}_{1}\) 的元素与每个来自 \({G}_{2}\) 的元素互相交换。

反之，设 \(G\) 为包含子群 \({G}_{1}\) 与 \({G}_{2}\) 的群，且满足下列两条条件:

(i) 每个 \(s \in  \mathrm{G}\) 可唯一写成 \(s = {s}_{1}{s}_{2}\) 的形式，其中 \({s}_{1} \in  {\mathrm{G}}_{1}\) 与 \({s}_{2} \in  {\mathrm{G}}_{2}\) 如上。

(ii) 对于任意 \({s}_{1} \in  {\mathrm{G}}_{1}\) 与 \({s}_{2} \in  {\mathrm{G}}_{2}\) ，有 \({s}_{1}{s}_{2} = {s}_{2}{s}_{1}\) 。

则两个元素 \(s = {s}_{1}{s}_{2},t = {t}_{1}{t}_{2}\) 的乘积可以写成

\[
{st} = {s}_{1}{s}_{2}{t}_{1}{t}_{2} = \left( {{s}_{1}{t}_{1}}\right) \left( {{s}_{2}{t}_{2}}\right) .
\]

由此可知，若令 \(\left( {{s}_{1},{s}_{2}}\right)  \in  {\mathrm{G}}_{1} \times  {\mathrm{G}}_{2}\) 对应于 \(\mathrm{G}\) 的元素 \({s}_{1}{s}_{2}\) ，则得到从 \({\mathrm{G}}_{1} \times  {\mathrm{G}}_{2}\) 到 \(\mathrm{G}\) 的同构。在此情况下，我们也说 \(G\) 是其子群 \({G}_{1}\) 和 \({G}_{2}\) 的乘积(或直积)，并将其识别为 \({G}_{1} \times  {G}_{2}\) 。

现在令 \({\rho }^{1} : {G}_{1} \rightarrow  \operatorname{GL}\left( {V}_{1}\right)\) 和 \({\rho }^{2} : {G}_{2} \rightarrow  \operatorname{GL}\left( {V}_{2}\right)\) 分别为群 \({G}_{1}\) 和 \({G}_{2}\) 的线性表示。类似于 1.5 的方法，我们通过如下方式定义群 \({\mathrm{G}}_{1} \times  {\mathrm{G}}_{2}\) 到 \({\mathrm{V}}_{1} \otimes  {\mathrm{V}}_{2}\) 的线性表示 \({\rho }^{1} \otimes  {\rho }^{2}\) :

\[
\left( {{\rho }^{1} \otimes  {\rho }^{2}}\right) \left( {{s}_{1},{s}_{2}}\right)  = {\rho }^{1}\left( {s}_{1}\right)  \otimes  {\rho }^{2}\left( {s}_{2}\right) .
\]

该表示称为表示 \({\rho }^{1}\) 与 \({\rho }^{2}\) 的张量积。若 \({\chi }_{i}\) 是 \({\rho }_{i}\left( {i = 1,2}\right)\) 的迹，则 \({\rho }^{1} \otimes  {\rho }^{2}\) 的迹 \(\chi\) 由下式给出:

\[
\chi \left( {{s}_{1},{s}_{2}}\right)  = {\chi }_{1}\left( {s}_{1}\right)  \cdot  {\chi }_{2}\left( {s}_{2}\right) .
\]

当 \({G}_{1}\) 与 \({G}_{2}\) 都等于同一群 \(G\) 时，上述定义的表示 \({\rho }^{1} \otimes  {\rho }^{2}\) 是对 \(\mathbf{G} \times  \overline{\mathbf{G}}\) 的表示。将其限制到 \(\mathbf{G} \times  \mathbf{G}\) 的对角子群(由 \(\left( {s,s}\right)\) 构成，其中 \(s\) 在 \(\mathbf{G}\) 上取遍)时，得到 1.5 中记作 \({\rho }^{1} \otimes  {\rho }^{2}\) 的 \(G\) 的表示；尽管记号相同，但需区分这两种表示。

\section*{定理 10}

(i) 若 \({\rho }^{1}\) 与 \({\rho }^{2}\) 不可约，则 \({\rho }^{1} \otimes  {\rho }^{2}\) 是 \({\mathrm{G}}_{1} \times  {\mathrm{G}}_{2}\) 的不可约表示。

(ii) 每个 \({\mathrm{G}}_{1} \times  {\mathrm{G}}_{2}\) 的不可约表示同构于某个表示 \({\rho }^{1} \otimes  {\rho }^{2}\) ，其中 \({\rho }^{i}\) 是 \({\mathrm{G}}_{i} \; \left( {i = 1,2}\right)\) 的不可约表示。

若 \({\rho }^{1}\) 与 \({\rho }^{2}\) 不可约，则有(参见 2.3):

\[
\frac{1}{{g}_{1}}\mathop{\sum }\limits_{{s}_{1}}{\left| {\chi }_{1}\left( {s}_{1}\right) \right| }^{2} = 1,\;\frac{1}{{g}_{2}}\mathop{\sum }\limits_{{s}_{2}}{\left| {\chi }_{2}\left( {s}_{2}\right) \right| }^{2} = 1.
\]

通过乘法，可得:

\[
\frac{1}{g}\mathop{\sum }\limits_{{{s}_{1},{s}_{2}}}{\left| \chi \left( {s}_{1},{s}_{2}\right) \right| }^{2} = 1
\]

由此可知 \({\rho }^{1} \otimes  {\rho }^{2}\) 是不可约的(定理 5)。为证明 (ii)，只需说明每个在 \({G}_{1} \times  {G}_{2}\) 上与形如 \({\chi }_{1}\left( {s}_{1}\right) {\chi }_{2}\left( {s}_{2}\right)\) 的迹正交的类函数 \(f\) 必为零。于是假设有:

\[
\mathop{\sum }\limits_{{{s}_{1},{s}_{2}}}f\left( {{s}_{1},{s}_{2}}\right) {\chi }_{1}{\left( {s}_{1}\right) }^{ * }{\chi }_{2}{\left( {s}_{2}\right) }^{ * } = 0.
\]

固定 \({\chi }_{2}\) 并置入 \(g\left( {s}_{1}\right)  = \mathop{\sum }\limits_{{s}_{2}}f\left( {{s}_{1},{s}_{2}}\right) {\chi }_{2}{\left( {s}_{2}\right) }^{ * }\) ，我们得到:

\[
\mathop{\sum }\limits_{{s}_{1}}g\left( {s}_{1}\right) {\chi }_{1}{\left( {s}_{1}\right) }^{ * } = 0\text{ for all }{\chi }_{1}.
\]

由于 \(g\) 是类函数，这意味着 \(g = 0\) ，并且对每个 \({\chi }_{2}\) 情况相同，因此用同样的论证我们得出 \(f\left( {{s}_{1},{s}_{2}}\right)  = 0\) 。

[也可以通过计算表示的维数平方之和 \({\rho }^{1} \otimes  {\rho }^{2}\) ，并应用 2.4 来证明 (ii)。]

上述定理将对 \({G}_{1} \times  {G}_{2}\) 的表示的研究完全归结为对 \({G}_{1}\) 的表示和对 \({G}_{2}\) 的表示。

\section*{3.3 诱导表示}

\section*{子群的左陪集}

回忆以下定义:设 \(\mathrm{H}\) 为群 G 的一个子群。对于 \(s \in  \mathrm{G}\) ，我们用 \(s\mathrm{H}\) 表示由 \({st}\) 与 \(t \in  \mathrm{H}\) 的乘积构成的集合，并称 \(s\mathrm{H}\) 为包含 \(s\) 的 \(\mathrm{H}\) 的左陪集。若两个属于 \(\mathrm{G}\) 的元素 \(s,{s}^{\prime }\) 属于同一左陪集，即当且仅当 \({s}^{-1}{s}^{\prime }\) 属于 \(\mathrm{H}\) ，则称它们模 \(\mathrm{H}\) 同余；我们记作 \({s}^{\prime } \equiv  s\left( {\;\operatorname{mod}\;\mathrm{H}}\right)\) 。左陪集的集合记为 \(\mathrm{G}/\mathrm{H}\) ；它是 \(\mathrm{G}\) 的一个划分。若 \(\mathrm{G}\) 有 \(g\) 个元素且 \(\mathrm{H}\) 有 \(h\) 个元素，则 \(\mathrm{G}/\mathrm{H}\) 有 \(g/h\) 个元素；整数 \(g/h\) 是 \(\mathrm{H}\) 在 \(\mathrm{G}\) 中的指标，记作 (G:H)。

若我们从每个 \(H\) 的左陪集中选取一个元素，就得到 \(G\) 的一个子集 \(R\) ，称为 \(G/H\) 的一组代表元；每个在 \(G\) 中的 \(s\) 都可以唯一地写成 \(s = {rt}\) ，其中 \(r \in  \mathrm{R}\) 与 \(t \in  \mathrm{H}\) 。

\section*{诱导表示的定义}

令 \(\rho  : \mathrm{G} \rightarrow  \mathbf{{GL}}\left( \mathrm{V}\right)\) 为 \(\mathrm{G}\) 的线性表示， \({\rho }_{\mathrm{H}}\) 为其在 \(\mathrm{H}\) 上的限制。令 \(\mathrm{W}\) 为 \({\rho }_{\mathrm{H}}\) 的子表示，即在 \(\mathrm{V}\) 中对 \({\rho }_{t},t \in  \mathrm{H}\) 闭合的向量子空间。记由此在 \(\mathrm{H}\) 上于 \(\mathrm{W}\) 中得到的表示为 \(\theta  : \mathrm{H} \rightarrow  \mathbf{{GL}}\left( \mathrm{W}\right)\) 。令 \(s \in  \mathrm{G}\) ；向量空间 \({\rho }_{s}\mathrm{\;W}\) 仅依赖于左陪集 \(s\mathrm{H}\) 的所属 \(s\) ；事实上，若把 \(s\) 换成 \({st}\) ，且 \(t \in  \mathrm{H}\) ，则有 \({\rho }_{st}\mathrm{\;W} = {\rho }_{s}{\rho }_{t}\mathrm{\;W} = {\rho }_{s}\mathrm{\;W}\) ，因为 \({\rho }_{t}\mathrm{\;W} = \mathrm{W}\) 。若 \(\sigma\) 是 \(\mathrm{H}\) 的一个左陪集，则我们可将子空间 \({\mathrm{W}}_{\sigma }\) 定为任一 \(s \in  \sigma\) 的 \({\rho }_{s}\mathrm{\;W}\) 。显然 \({\mathrm{W}}_{\sigma }\) 在 \({\rho }_{s},s \in  \mathrm{G}\) 下相互置换。它们的和 \(\mathop{\sum }\limits_{{\sigma  \in  \mathrm{G}/\mathrm{H}}}{\mathrm{W}}_{\sigma }\) 因此是 \(\mathrm{V}\) 的一个子表示。

定义:若 \(\mathrm{V}\) 等于各个 \({\mathrm{W}}_{\sigma } \; \left( {\sigma  \in  \mathrm{G}/\mathrm{H}}\right)\) 的和且该和为直和(即若 \(\mathrm{V} = {\bigoplus }_{\sigma  \in  \mathrm{G}/\mathrm{H}}{\mathrm{W}}_{\sigma }\) )，则称 \(\rho\) 在 \(\mathrm{G}\) 中的表示由 \(\mathrm{H}\) 在 \(\mathrm{W}\) 中的表示 \(\theta\) 所诱导。

我们可用若干方式重述此条件:

(i) 每个 \(x \in  \mathrm{V}\) 都可唯一写成 \(\mathop{\sum }\limits_{{\sigma  \in  \mathrm{G}/\mathrm{H}}}{x}_{\sigma }\) ，其中 \({x}_{\sigma } \in  {\mathrm{W}}_{\sigma }\) 用于

(ii) 若 \(R\) 是 \(G/H\) 的一组代表元，则向量空间 \(V\) 是各个 \({\rho }_{r}W\) 的直和，且 \(r \in  R\) 。

特别地，我们有 \(\dim \left( V\right)  = \mathop{\sum }\limits_{{r \in  R}}\dim \left( {{\rho }_{r}W}\right)  = \left( {G : H}\right)  \cdot  \dim \left( W\right)\) 。

例 1. 令 \(V\) 为 \(G\) 的正则表示；空间 \(V\) 有一基 \({\left( {e}_{t}\right) }_{t \in  \mathrm{G}}\) ，使得 \({\rho }_{s}{e}_{t} = {e}_{st}\) 对于 \(s \in  \mathrm{G},t \in  \mathrm{G}\) 成立。设 \(\mathrm{W}\) 为 \(\mathrm{V}\) 的子空间，基为 \({\left( {e}_{t}\right) }_{t \in  \mathrm{H}}\) 。在 \(\mathrm{W}\) 上的表示 \(\theta\) 是 \(\mathrm{H}\) 的正则表示，且显然 \(\rho\) 由 \(\theta\) 诱导。

2. 取 \(\mathrm{V}\) 为一个向量空间，它有一个以 \(\mathrm{G}/\mathrm{H}\) 的元素 \(\sigma\) 为指标的基 \(\left( {e}_{\sigma }\right)\) ，并通过 \({\rho }_{s}{e}_{\sigma } = {e}_{s\sigma }\) (其中 \(s \in  \mathrm{G}\) 且 \(\sigma  \in  \mathrm{G}/\mathrm{H}\) )来定义 \(\mathrm{G}\) 在 \(\mathrm{V}\) 中的一个表示 \(\rho\) (这个公式是有意义的，因为，如果 \(\sigma\) 是 \(\mathrm{H}\) 的一个左陪集，那么 \({s\sigma }\) 也是)。这样我们就得到了 \(\mathrm{G}\) 的一个表示，它是 \(\mathrm{G}\) 的与 \(\mathrm{G}/\mathrm{H}\) 相关联的置换表示[参见 1.2，例 (c)]。对应于陪集 \(\mathrm{H}\) 的向量 \({e}_{\mathrm{H}}\) 在 \(\mathrm{H}\) 下是不变的； \(\mathrm{H}\) 在子空间 \(\mathbf{C}{e}_{\mathrm{H}}\) 中的表示因此是 \(\mathrm{H}\) 的单位表示，并且显然这个表示诱导出了 \(\mathbf{G}\) 在 \(\mathbf{V}\) 中的表示。

3. 若 \({\rho }_{1}\) 由 \({\theta }_{1}\) 诱导，且 \({\rho }_{2}\) 由 \({\theta }_{2}\) 诱导，则 \({\rho }_{1} \oplus  {\rho }_{2}\) 由 \({\theta }_{1} \oplus  {\theta }_{2}\) 诱导。

4. 若 \(\left( {V,\rho }\right)\) 由 \(\left( {W,\theta }\right)\) 诱导，且 \({W}_{1}\) 为 \(W\) 的不变子空间，则 \(V\) 的子空间 \({V}_{1} = \mathop{\sum }\limits_{{r \in  R}}{\rho }_{r}{W}_{1}\) 在 \(G\) 下不变，且 \(G\) 在 \({V}_{1}\) 上的表示由 \(H\) 在 \({W}_{1}\) 上的表示诱导而来。

5. 若 \(\rho\) 由 \(\theta\) 诱导，若 \({\rho }^{\prime }\) 是 \(G\) 的一个表示，且 \({\rho }_{H}^{\prime }\) 是 \({\rho }^{\prime }\) 限制到 \(\mathrm{H}\) ，则 \(\rho  \otimes  {\rho }^{\prime }\) 由 \(\theta  \otimes  {\rho }_{\mathrm{H}}^{\prime }\) 诱导。

诱导表示的存在性与唯一性

引理 1. 设 \(\left( {\mathbf{V},\rho }\right)\) 由 \(\left( {\mathbf{W},\theta }\right)\) 诱导。令 \({\rho }^{\prime } : \mathbf{G} \rightarrow  \mathbf{{GL}}\left( {\mathbf{V}}^{\prime }\right)\) 是 \(\mathrm{G}\) 的线性表示，且令 \(f : \mathrm{W} \rightarrow  {\mathrm{V}}^{\prime }\) 为一线性映射，使得对所有 \(t \in  \mathrm{H}\) 与 \(w \in  \mathrm{W}\) 都有 \(f\left( {{\theta }_{t}w}\right)  = {\rho }_{t}^{\prime }f\left( w\right)\) 。则存在唯一的线性映射 \(\mathrm{F} : \mathrm{V} \rightarrow  {\mathrm{V}}^{\prime }\) ，它扩张 \(f\) 且对所有 \(s \in  \mathbf{G}\) 满足 \(\mathrm{F} \circ  {\rho }_{s} = {\rho }_{s}^{\prime } \circ  \mathrm{F}\) 。

如果 \(\mathrm{F}\) 满足这些条件，且如果 \(x \in  {\rho }_{s}\mathrm{\;W}\) ，我们有 \({\rho }_{s}^{-1}x \in  \mathrm{W}\) ；因此

\[
\mathrm{F}\left( x\right)  = \mathrm{F}\left( {{\rho }_{s}{\rho }_{s}^{-1}x}\right)  = {\rho }_{s}^{\prime }\mathrm{F}\left( {{\rho }_{s}^{-1}x}\right)  = {\rho }_{s}^{\prime }f\left( {{\rho }_{s}^{-1}x}\right) .
\]

该公式在 \({\rho }_{s}\mathrm{\;W}\) 上确定了 \(\mathrm{F}\) ，因此在 \(\mathrm{V}\) 上也是如此，因为 \(\mathrm{V}\) 是这些 \({\rho }_{s}\mathrm{\;W}\) 的和。这证明了 \(\mathrm{F}\) 的唯一性。

现在令 \(x \in  {\mathrm{W}}_{\sigma }\) ，并选择 \(s \in  \sigma\) ；我们按上述公式 \(F\left( x\right)  = {\rho }_{s}^{\prime }f\left( {{\rho }_{s}^{-1}x}\right)\) 定义 \(\mathrm{F}\left( x\right)\) 。此定义不依赖于在 \(\sigma\) 中对 \(s\) 的选择；确实，若我们以满足 \(t \in  \mathrm{H}\) 的 \({st}\) 代替 \(s\) ，则有

\[
{\rho }_{st}^{\prime }f\left( {{\rho }_{st}^{-1}x}\right)  = {\rho }_{s}^{\prime }{\rho }_{t}^{\prime }f\left( {{\theta }_{t}^{-1}{\rho }_{s}^{-1}x}\right)  = {\rho }_{s}^{\prime }\left( {{\theta }_{t}{\theta }_{t}^{-1}{\rho }_{s}^{-1}x}\right)  = {\rho }_{s}^{\prime }f\left( {{\rho }_{s}^{-1}x}\right) .
\]

由于 \(\mathbf{V}\) 是这些 \({\mathbf{W}}_{\sigma }\) 的直和，存在唯一的线性映射

F: \(V \rightarrow  {V}^{\prime }\) ，其扩展了先前在各 \({W}_{\sigma }\) 上定义的部分映射。容易检查对所有 \(s \in  G\) 有 \(F \circ  {\rho }_{s} = {\rho }_{s}^{\prime } \circ  F\) 。

定理 11. 设 \(\left( {\mathbf{W},\theta }\right)\) 为 \(\mathbf{H}\) 的线性表示。存在由 \(\left( {\mathbf{W},\theta }\right)\) 诱导的 \(\mathbf{G}\) 的线性表示 \(\left( {\mathbf{V},\rho }\right)\) ，且该表示在同构意义下唯一。

我们先证明诱导表示 \(\rho\) 的存在。参照上面的例 3，我们可以假定 \(\theta\) 是不可约的。在此情形下， \(\theta\) 同构于 H 的正则表示的一个子表示，该正则表示可诱导为 G 的正则表示(参见例 1)。应用例 4，我们得出 \(\theta\) 本身可以被诱导。

剩下要证明的是 \(\rho\) 在同构意义下的唯一性。令 \(\left( {V,\rho }\right)\) 和 \(\left( {{V}^{\prime },{\rho }^{\prime }}\right)\) 为由 \(\left( {W,\theta }\right)\) 诱导的两个表示。将引理 1 应用于 \(W\) 注入 \({V}^{\prime }\) 的映射，我们得到存在一个在线性映射 \(F : V \rightarrow  {V}^{\prime }\) ，它在 \(W\) 上为恒等，并对所有 \(s \in  G\) 满足 \(F \circ  {\rho }_{s} = {\rho }_{s}^{\prime } \circ  F\) 。因而 \(F\) 的像包含所有 \({\rho }_{s}^{\prime }W\) ，从而等于 \({V}^{\prime }\) 。由于 \({V}^{\prime }\) 与 \(V\) 维数相同 (G: H) \(\cdot  \dim \left( W\right)\) ，我们得到 \(F\) 是同构，从而证得该定理。(有关定理 11 的更自然的证明，见 7.1。)

\section*{诱导表示的特征}

假设 \(\left( {\mathrm{V},\rho }\right)\) 由 \(\left( {\mathrm{W},\theta }\right)\) 导出，令 \({\chi }_{\rho }\) 和 \({\chi }_{\theta }\) 分别为 \(G\) 和 \(H\) 的相应特征。既然 \(\left( {W,\theta }\right)\) 可将 \(\left( {V,\rho }\right)\) 唯一确定到同构，我们应能由 \({\chi }_{\theta }\) 计算出 \({\chi }_{\rho }\) 。下述定理说明了方法:

定理 12. 设 \(h\) 为 \(\mathrm{H}\) 的阶，令 \(\mathrm{R}\) 为 \(\mathrm{G}/\mathrm{H}\) 的一组代表元。对每个 \(u \in  \mathrm{G}\) ，有

\[
{\chi }_{\rho }\left( u\right)  = \mathop{\sum }\limits_{\substack{{r \in  \mathrm{R}} \\  {{r}^{-1}{ur} \in  \mathrm{H}} }}{\chi }_{\theta }\left( {{r}^{-1}{ur}}\right)  = \frac{1}{h}\mathop{\sum }\limits_{\substack{{s \in  \mathrm{G}} \\  {{s}^{-1}{us} \in  \mathrm{H}} }}{\chi }_{\theta }\left( {{s}^{-1}{us}}\right) .
\]

(特别地， \({\chi }_{\rho }\left( u\right)\) 是 \({\chi }_{\theta }\) 在 \(\mathrm{H}\) 与 G 中 \(u\) 的共轭类交集上的值的线性组合。)

空间 \(\mathrm{V}\) 是各 \({\rho }_{r}\mathrm{\;W},r \in  \mathrm{R}\) 的直和。此外， \({\rho }_{u}\) 在这些 \({\rho }_{r}\mathrm{\;W}\) 之间作置换。更精确地说，若将 \({ur}\) 写成 \({r}_{u}t\) 的形式，且 \({r}_{u} \in  \mathrm{R}\) 与 \(t \in  \mathrm{H}\) 为……，则可见 \({\rho }_{u}\) 将 \({\rho }_{r}\mathrm{\;W}\) 送入 \({\rho }_{{r}_{u}}\mathrm{\;W}\) 。为确定 \({\chi }_{o}\left( u\right)  = {\operatorname{Tr}}_{\mathrm{V}}\left( {\rho }_{u}\right)\) ，我们可用一个由各 \({\rho }_{r}\mathrm{\;W}\) 的基组成的基。使得 \({r}_{u} \neq  r\) 的指数 \(r\) 给出零对角项；其余的给出 \({\rho }_{u}\) 在相应 \({\rho }_{r}\mathrm{\;W}\) 上的迹。由此得到:

\[
{\chi }_{\rho }\left( u\right)  = \mathop{\sum }\limits_{{r \in  {\mathrm{R}}_{u}}}{\operatorname{Tr}}_{{\rho }_{r}\mathrm{\;W}}\left( {\rho }_{u,r}\right)
\]

其中 \({\mathbf{R}}_{u}\) 表示那些满足 \({r}_{u} = r\) 并且 \({\rho }_{u,r}\) 为 \({\rho }_{u}\) 限制到 \({\rho }_{r}\mathrm{\;W}\) 的 \(r \in  \mathbf{R}\) 的集合。注意，当且仅当 ur 可写成 \({rt}\) 且具有 \(t \in  \mathrm{H}\) ，即当 \({r}^{-1}{ur}\) 属于 \(\mathrm{H}\) 时， \(r\) 属于 \({\mathrm{R}}_{u}\) 。

剩下要计算的是 \({\operatorname{Tr}}_{{\rho }_{r}\mathrm{\;W}}\left( {\rho }_{u,r}\right)\) ，针对 \(r \in  {\mathbf{R}}_{u}\) 。为此，注意 \({\rho }_{r}\) 给出一个将 \(\mathrm{W}\) 同构到 \({\rho }_{r}\mathrm{\;W}\) 的同构，并且我们有

\[
{\rho }_{r} \circ  {\theta }_{t} = {\rho }_{u,r} \circ  {\rho }_{r},\;\text{ with }t = {r}^{-1}{ur} \in  \mathrm{H}.
\]

因此 \({\rho }_{u,r}\) 的迹等于 \({\theta }_{t}\) 的迹，即等于 \({\chi }_{\theta }\left( t\right)  = {\chi }_{\theta }\left( {{r}^{-1}{ur}}\right)\) 。我们确实得到:

\[
{\chi }_{\rho }\left( u\right)  = \mathop{\sum }\limits_{{r \in  {\mathrm{R}}_{u}}}{\chi }_{\theta }\left( {{r}^{-1}{ur}}\right)
\]

给出 \({\chi }_{\rho }\left( u\right)\) 的第二个公式可由第一个推得，因为位于左陪集 \(r\mathrm{H}\left( {r \in  {\mathrm{R}}_{u}}\right)\) 中的所有元素 \(s\) 属于 \(\mathrm{G}\) 满足 \({\chi }_{\theta }\left( {{s}^{-1}{us}}\right) \; = {\chi }_{\theta }\left( {{r}^{-1}{ur}}\right)\) 。

读者将在第二部分发现诱导表示的其他性质。值得注意的是:

(i) Frobenius 互反公式

\[
{\left( {f}_{\mathrm{H}} \mid  {\chi }_{\theta }\right) }_{\mathrm{H}} = {\left( f \mid  {\chi }_{\rho }\right) }_{\mathrm{G}}
\]

其中 \(f\) 是 \(G\) 的类函数， \({f}_{H}\) 是它在 \(H\) 上的限制，标量积分别在 \(\mathrm{H}\) 和 \(\mathrm{G}\) 上计算。

(ii) Mackey 判据，说明何时诱导表示是不可约的。

(iii) Artin 定理(分别 Brauer 定理)，说明群 G 的每个特征都是从循环子群诱导的表示的特征(分别从“基本”子群诱导的表示的特征)的有理(分别整数)系数线性组合。

\section*{习题}

3.4. 证明每个 \(G\) 的不可约表示包含在由 \(H\) 的一个不可约表示诱导的表示中。[利用不可约表示包含于正则表示的事实。] 从中得到定理 9 的推论的另一个证明。

3.5. 设 \(\left( {\mathrm{W},\theta }\right)\) 为 \(\mathrm{H}\) 的一个线性表示。令 \(\mathrm{V}\) 为由满足 \(f\left( {tu}\right)  = {\theta }_{t}f\left( u\right)\) (对所有 \(u \in  \mathrm{G},t \in  \mathrm{H}\) 成立)的函数 \(f : \mathrm{G} \rightarrow  \mathrm{W}\) 组成的向量空间。令 \(\rho\) 为在 \(\mathrm{V}\) 上由 \(\mathrm{G}\) 作用的表示，其定义为对所有 \(s,u \in  \mathrm{G}\) ， \(\left( {{\rho }_{s}f}\right) \left( u\right)  = f\left( {us}\right)\) 。对 \(w \in  \mathbf{W}\) ，令 \({f}_{w} \in  \mathbf{V}\) 由对 \(t \in  \mathbf{H}\) 满足 \({f}_{w}\left( t\right)  = {\theta }_{t}w\) 且对 \(s \notin  \mathrm{H}\) 满足 \({f}_{w}\left( s\right)  = 0\) 来定义。证明 \(w \mapsto  {f}_{w}\) 是把 \(\mathbf{W}\) 同构到由在 H 外取零的函数组成的子空间 \({\mathbf{W}}_{0}\) 的同构。证明若以此方式将 W 与 \({\mathbf{W}}_{0}\) 识别，则表示 \(\left( {\mathbf{V},\rho }\right)\) 是由表示 \(\left( {\mathbf{W},\theta }\right)\) 诱导而来。

3.6. 假设 \(G\) 是两个子群 \(H\) 与 \(K\) 的直积(参见 3.2)。令 \(\rho\) 为由 \(H\) 的表示 \(\theta\) 诱导得到的 \(G\) 的表示。证明 \(\rho\) 同构于 \(\theta  \otimes  {r}_{\mathrm{K}}\) ，其中 \({r}_{\mathrm{K}}\) 表示 \(\mathrm{K}\) 的正则表示。

\section*{第四章 紧群}

本章旨在说明先前结果如何推广到任意紧群(不必为有限群)；证明详见书目中引用的 [1]、[4]、[6]。

除例 5.2、5.5 和 5.6 外，下列结论在后文中不会被使用。

\section*{4.1 紧群}

拓扑群 G 是带有拓扑结构的群，使得乘法 \(s \cdot  t\) 和求逆 \({s}^{-1}\) 都是连续的。若其拓扑是紧空间的拓扑，即满足 Borel–Lebesgue 定理，则称该群为紧群。例如，欧氏平面(或三维等)中围绕一点的旋转群具有自然的拓扑，使其成为紧群；其闭子群也同样是紧群。

作为非紧群的例子，可举平移群 \(x \mapsto  x + a\) 和保持二次型 \({x}^{2} + {y}^{2} + {z}^{2} - {t}^{2}\) 的线性变换群(“洛伦兹群”)。这些群的线性表示与紧群情形下具有完全不同的性质。

\section*{4.2 紧群上的不变测度}

在研究有限群 \(G\) (阶为 \(g\) )的线性表示时，我们大量使用了对 \(G\) 的平均操作，即对定义在 \(G\) 上的函数 \(f\) 赋予元素 \(\left( {1/g}\right) \mathop{\sum }\limits_{{t \in  G}}f\left( t\right)\) ( \(f\) 的取值可以是复数或更一般的向量空间元素)。对紧群也存在类似的操作；当然，不是有限求和，而是对某个测度 \({dt}\) 的积分 \({\int }_{\mathrm{G}}f\left( t\right) {dt}\) 。

更确切地，可以证明存在且唯一有一测度 dt 支持在 G 上并满足下列两条性质:

(i) \({\int }_{\mathrm{G}}f\left( t\right) {dt} = {\int }_{\mathrm{G}}f\left( {ts}\right) {dt}\) 对每个连续函数 \(f\) 与每个 \(s \in  \mathrm{G}\) 保持不变(dt 在右平移下的不变性)。

(ii) \({\int }_{\mathrm{G}}{dt} = 1\) (总质量 \({dt}\) 等于 1 \()\) )。

并且可证明 \({dt}\) 在左平移下也不变，即:

(i')
\[
{\int }_{\mathrm{G}}f\left( t\right) {dt} = {\int }_{\mathrm{G}}f\left( {st}\right) {dt}
\]

该测度 \({dt}\) 称为群 G 的不变测度(或 Haar 测度)。给出两个例子(参见第 5 章):

(1) 若 \(\mathbf{G}\) 是有限群，阶为 \(g\) ，则测度 \({dt}\) 可通过给每个元素 \(t \in  \mathbf{G}\) 赋予质量 \(1/g\) 来构造。

(2) 若 \(G\) 是平面内的旋转群 \({C}_{\infty }\) ，并把元素 \(t \in  \mathrm{G}\) 表示为 \(t = {e}^{i\alpha }\) ( \(\alpha\) 按 \({2\pi }\) 取模)，则不变测度为 \(\left( {1/{2\pi }}\right) {d\alpha }\) ；因子 \(1/{2\pi }\) 用以满足条件 (ii)。

\section*{4.3 紧群的线性表示}

设 \(\mathbf{G}\) 为紧群， \(\mathbf{V}\) 为复数域上的有限维向量空间。 \(\mathbf{G}\) 在 \(\mathbf{V}\) 上的线性表示是一个到 GL(V) 的同态 \(\rho  : G \rightarrow\) 且连续；此条件等价于说 \({\rho }_{s}x\) 是关于两个变量 \(s \in  \mathrm{G},x \in  \mathrm{V}\) 的连续函数。对希尔伯特空间中 \(\mathrm{G}\) 的线性表示亦可类似定义；并且可以证明此类表示同构于有限维幺正表示的(希尔伯特)直和，因此可将注意力限制在后者。

有限群表示的大多数性质可推广到紧群的表示；只需将“ \(\left( {1/g}\right) \mathop{\sum }\limits_{{t \in  \mathrm{G}}}f\left( t\right)\) ”改为“ \({\int }_{\mathrm{G}}f\left( t\right) {dt}\) ”。例如，两个函数 \(\phi\) 与 \(\psi\) 的标量积 \(\left( {\phi  \mid  \psi }\right)\) 为

\[
\left( {\phi  \mid  \psi }\right)  = {\int }_{\mathrm{G}}\phi \left( t\right) \psi {\left( t\right) }^{ * }{dt}
\]

更精确地说:

(a) 定理1、2、3、4和5及其证明可不变地推广。命题1、2、3和4亦然。

(b) 在2.4节中，有必要将规则表示 \(\mathbf{R}\) 定义为在 \(\mathbf{G}\) 上的平方可和函数构成的希尔伯特空间，群作用为 \(\left( {{\rho }_{s}f}\right) \left( t\right)  = f\left( {{s}^{-1}t}\right)\) 。若 \(\mathbf{G}\) 非有限，则此表示为无限维，不能再谈其特征，故命题5不再有意义。然而，每个不可约表示仍包含在 \(\mathrm{R}\) 中，重数等于其度数。

(c) 命题6与定理6可不变地推广(在定理6中，把 H 取为 G 上平方可和函数的希尔伯特空间)。

(d) 当 \(\mathbf{G}\) 非有限时，定理7 仍然成立(但无趣):共轭类无限多，不可约表示也无限多。

(e) 定理8与命题8及其证明可不变地推广。正交分解(定理8)中的投影 \({p}_{i}\) 由下列公式给出

\[
{p}_{i}x = {n}_{i}{\int }_{\mathrm{G}}{\chi }_{i}{\left( t\right) }^{ * }{\rho }_{t}{xdt}
\]

(f) 定理9和10及其证明可不变地推广。关于定理10，注意乘积 \({\mathrm{G}}_{1} \times  {\mathrm{G}}_{2}\) 的不变测度是群 \({\mathrm{G}}_{1}\) 与 \({\mathrm{G}}_{2}\) 的不变测度的乘积 \(d{s}_{1}d{s}_{2}\) 。

(g) 只要 \(\mathrm{H}\) 为 \(\mathrm{G}\) 的闭子群且指标有限，由3.3节定义的由 \(\mathrm{H}\) 的表示诱导出的 \(\mathrm{G}\) 的表示，以及定理11和12，仍然成立。当 H 的指标为无限时，由 \(\left( {\mathbf{W},\theta }\right)\) 诱导的表示定义为在 \(G\) 上取值于 \(W\) 的平方可和函数构成的希尔伯特空间 \(f\) ，满足对每个 \(t \in  \mathrm{H}\) 有 \(f\left( {tu}\right)  = {\theta }_{t}f\left( u\right)\) ，且 \(\mathrm{G}\) 以 \({\rho }_{s}f\left( u\right)  = f\left( {us}\right)\) 的方式作用于该空间，参见练习3.5。

\section*{第5章 例子}

\section*{5.1 循环群 \({\mathrm{C}}_{n}\)}

这是一个阶为 \(n\) 的群，由某元素 \(r\) 的幂 \(1,r,\ldots ,{r}^{n - 1}\) 组成，满足 \({r}^{n} = 1\) 。它可视为绕某轴按角度 \({2k\pi }/n\) 的旋转群。它是一个阿贝尔群。

根据定理 9， \({\mathrm{C}}_{n}\) 的不可约表示都是 1 维。这样的表示将一个复数 \(\chi \left( r\right)  = w\) 关联给 \(r\) ，并将数 \(\chi \left( {r}^{k}\right)  = {w}^{k}\) 关联给 \({r}^{k}\) ；由于 \({r}^{n} = 1\) ，我们有 \({w}^{n} = 1\) ，即 \(w = {e}^{{2\pi ih}/n}\) ，伴随 \(h = 0,1,\ldots ,n - 1\) 。由此我们得到 \(n\) 个 1 维不可约表示，其特征 \({\chi }_{0},{\chi }_{1},\ldots ,{\chi }_{n - 1}\) 如下所示

\[
{\chi }_{h}\left( {r}^{k}\right)  = {e}^{{2\pi ihk}/n}.
\]

我们有 \({\chi }_{h} \cdot  {\chi }_{{h}^{\prime }} = {\chi }_{h + {h}^{\prime }}\) ，约定当 \(h + {h}^{\prime } \geq  n\) 时 \({\chi }_{h + {h}^{\prime }} = {\chi }_{h + {h}^{\prime } - n}\) (换言之， \({\chi }_{h}\) 的指标 \(h\) 取模 \(n\) )。

例如，对于 \(n = 3\) ，特征表如下:

\begin{center}
\adjustbox{max width=\textwidth}{
\begin{tabular}{|c|c|c|c|}
\hline
\phantom{X} & 1 & \(r\) & \({r}^{2}\) \\
\cline{1-4}
\({x}_{0}\) & 1 & 1 & 1 \\
\cline{1-4}
\({x}_{1}\) & 1 & W & \({w}^{2}\) \\
\cline{1-4}
\({x}_{2}\) & 1 & \({w}^{2}\) & w \\
\cline{1-4}
\hline
\end{tabular}
}
\end{center}

哪里

\[
w = {e}^{{2\pi i}/3} =  - \frac{1}{2} + i\frac{\sqrt{3}}{2}.
\]

我们有

\[
{\chi }_{0} \cdot  {\chi }_{i} = {\chi }_{i},{\chi }_{1} \cdot  {\chi }_{1} = {\chi }_{2},{\chi }_{2} \cdot  {\chi }_{2} = {\chi }_{1}\text{ and }{\chi }_{1} \cdot  {\chi }_{2} = {\chi }_{0}.
\]

\section*{5.2 群 \({\mathrm{C}}_{\infty }\)}

这是平面旋转群。如果我们用 \({r}_{\alpha }\) 表示绕角度 \(\alpha\) 的旋转(以 \({2\pi }\) 为模确定)，则 \({\mathrm{C}}_{\infty }\) 上的不变测度为 \(\left( {1/{2\pi }}\right) {d\alpha }\) (参见 4.2)。

\({\mathrm{C}}_{\infty }\) 的不可约表示为 1 维。它们由下式给出:

\[
{\chi }_{n}\left( {r}_{\alpha }\right)  = {e}^{in\alpha }\;\left( {n\text{ an arbitrary integer }}\right) .
\]

正交关系在此处给出众所周知的公式:

\[
\frac{1}{2\pi }{\int }_{0}^{2\pi }{e}^{-{in\alpha }} \cdot  {e}^{im\alpha }{d\alpha } = {\delta }_{nm},
\]

且定理 6 给出了周期函数作为傅里叶级数的展开。

\section*{5.3 二面体群 \({\mathrm{D}}_{n}\)}

这是保持有 \(n\) 个顶点的正多边形的不变的平面旋转和反射群。它包含 \(n\) 个旋转，构成与 \({\mathrm{C}}_{n}\) 同构的子群，以及 \(n\) 个反射。其阶为 \({2n}\) 。若我们将绕角度 \({2\pi }/n\) 的旋转记为 \(r\) ，且 \(s\) 是任一反射，则有:

\[
{r}^{n} = 1,\;{s}^{2} = 1,\;\operatorname{srs} = {r}^{-1}.
\]

\({D}_{n}\) 的每个元素可唯一地写成两种形式之一:若属于 \({\mathrm{C}}_{n}\) ，则写为 \({r}^{k}\) ，其中 \(0 \leq  k \leq  n - 1\) ；若不属于 \({\mathrm{C}}_{n}\) ，则写为 \(s{r}^{k}\) ，其中 \(0 \leq  k \; \leq  n - 1\) 。注意关系 \({srs} = {r}^{-1}\) 蕴含 \(s{r}^{k}s = {r}^{-k}\) ，从而得出 \({\left( s{r}^{k}\right) }^{2} = 1\) 。

将 \({\mathrm{D}}_{n}\) 实现为三维空间的刚体运动群

有几种这样的实现:

(a) 常见的实现(传统记为 \({D}_{n}\) ，参见 Eyring [5])。取绕 Oz 轴的旋转为旋转，取通过平面 Oxy 的 \(n\) 条直线的反射为反射，这些直线成角为 \(\pi /n\) 的倍数。

(b) 通过群 \({\mathrm{C}}_{nv}\) 的实现(Eyring [5] 的记法):不取关于 Oxy 直线的反射，而取关于包含 Oz 轴的平面的反射。

(c) 群 \({\mathrm{D}}_{2n}\) 也可以实现为群 \({\mathrm{D}}_{nd}\) (Eyring [5] 的记法)。

群 \({\mathrm{D}}_{n}\) 的不可约表示(n 为偶数 \(\geq  2\) )

首先，有4个一维表示，通过将 \(\pm  1\) 以所有可能方式对应到 \(r\) 和 \(s\) 而得到。它们的特征 \({\psi }_{1},{\psi }_{2},{\psi }_{3},{\psi }_{4}\) 由下表给出:

\begin{center}
\adjustbox{max width=\textwidth}{
\begin{tabular}{|c|c|c|}
\hline
\phantom{X} & \({r}^{k}\) & \(S{r}^{k}\) \\
\cline{1-3}
\({\psi }_{1}\) & 1 & 1 \\
\cline{1-3}
\({\psi }_{2}\) & 1 & -1 \\
\cline{1-3}
\({\psi }_{3}\) & \({\left( -1\right) }^{k}\) & \({\left( -1\right) }^{k}\) \\
\cline{1-3}
\({\psi }_{4}\) & \({\left( -1\right) }^{k}\) & \({\left( -1\right) }^{k + 1}\) \\
\cline{1-3}
\hline
\end{tabular}
}
\end{center}

接下来我们考虑 2 维表示。设定 \(w = {e}^{{2\pi i}/n}\) ，并令 \(h\) 为任意整数。我们通过如下给出一个 \({\rho }^{h}\) 对 \({\mathrm{D}}_{n}\) 的表示:

\[
{\rho }^{h}\left( {r}^{k}\right)  = \left( \begin{matrix} {w}^{hk} & 0 \\  0 & {w}^{-{hk}} \end{matrix}\right) ,\;{\rho }^{h}\left( {s{r}^{k}}\right)  = \left( \begin{matrix} 0 & {w}^{-{hk}} \\  {w}^{hk} & 0 \end{matrix}\right) .
\]

直接计算表明这确实是一个表示。此表示由 \({C}_{n}\) 的表示(在 3.3 节意义下)诱导而来，其特征是 \({\chi }_{h}\) (5.1)。它仅依赖于 \(n\) 模 \(h\) 的剩余类；此外 \({\rho }^{h}\) 与 \({\rho }^{n - h}\) 同构。因此我们可假定 \(0 \leq  h \; \leq  n/2\) 。极端情形 \(h = 0\) 与 \(h = n/2\) 无趣:对应的表示是可约的，特征分别为 \({\psi }_{1} + {\psi }_{2}\) 与 \({\psi }_{3} + {\psi }_{4}\) 。另一方面，对于 \(0 < h < n/2\) ，表示 \({\rho }^{h}\) 是不可约的:由于 \({w}^{h} \neq  {w}^{-h}\) ，在 \({\rho }^{h}\left( r\right)\) 下不变的唯一直线是坐标轴，而这些坐标轴在 \({\rho }^{h}\left( s\right)\) 下并不不变。相同的论证表明这些表示两两不相同构。相应的特征 \({\chi }^{h}\) 如下所示:

\[
{\chi }_{h}\left( {r}^{k}\right)  = {w}^{hk} + {w}^{-{hk}} = 2\cos \frac{2\pi hk}{n}
\]

\[
{\chi }_{h}\left( {s{r}^{k}}\right)  = 0.
\]

上述构造的 1 维与 2 维不可约表示(按同构算)是 \({\mathrm{D}}_{n}\) 的全部不可约表示。确实，它们维数平方之和等于 \(4 \times  1 + \left( {\left( {n/2}\right)  - 1}\right)  \times  4 = {2n}\) ，而 \(4 \times  1 + \left( {\left( {n/2}\right)  - 1}\right)  \times  4 = {2n}\) 是 \({\mathbf{D}}_{n}\) 的阶。

例。群 \({D}_{6}\) 有 4 个 1 维表示，其特征为 \({\psi }_{1},{\psi }_{2},{\psi }_{3},{\psi }_{4}\) ，并有 2 个 2 维不可约表示，其特征为 \({\chi }_{1}\) 与 \({\chi }_{2}\) 。

群 \({\mathrm{D}}_{n}\left( {n\text{ odd }}\right)\) 的不可约表示

仅有两个 1 维表示，它们的特征 \({\psi }_{1}\) 与 \({\psi }_{2}\) 列在下表:

\begin{center}
\adjustbox{max width=\textwidth}{
\begin{tabular}{|c|c|c|}
\hline
\phantom{X} & \({r}^{k}\) & \({\mathit{{sr}}}^{k}\) \\
\cline{1-3}
\({\psi }_{1}\) & 1 & 1 \\
\cline{1-3}
\({\psi }_{2}\) & 1 & -1 \\
\cline{1-3}
\hline
\end{tabular}
}
\end{center}

度为2的表示 \({\rho }^{h}\) 由与 \(n\) 为偶数时相同的公式定义。与 \(0 < h < n/2\) 对应的那些表示是不可约且两两非同构的(注意，由于 \(n\) 为奇数，条件 \(h < n/2\) 也可写作 \(h \leq  \left( {n - 1}\right) /2\) )。给出它们特征的公式相同。

这些表示是唯一的。事实上，它们的次数平方和等于 \(2 \times  1 + \frac{1}{2}\left( {n - 1}\right)  \times  4 = {2n}\) ，而这正是 \({D}_{n}\) 的阶。

\section*{练习}

5.1. 证明在 \({D}_{n},n\) 为偶数(分别为奇数)时，反射构成两类共轭元(分别一类)，且 \({C}_{n}\) 的元素构成 \(\left( {n/2}\right)  + 1\) 类(分别 \(\left( {n + 1}\right) /2\) 类)。由此得到 \({D}_{n}\) 的共轭类数，并检验其与不可约特征数相符。

5.2. 证明 \({\chi }_{h} \cdot  {\chi }_{{h}^{\prime }} = {\chi }_{h + {h}^{\prime }} + {\chi }_{h - {h}^{\prime }}\) 。特别地，我们有

\[
{\chi }_{n} \cdot  {\chi }_{h} = {\chi }_{2h} + {\chi }_{0} = {\chi }_{2h} + {\psi }_{1} + {\psi }_{2}.
\]

证明 \({\psi }_{2}\) 是 \({\rho }^{h}\) 交错平方的特征，且 \({\chi }_{2h} + {\psi }_{1}\) 是其对称平方的特征(参见1.5和命题3)。

5.3. 证明将 \({D}_{n}\) 常作为在 \({R}^{3}\) 中的刚性运动群来实现(Eyring [5])是可约的且其迹为 \({\chi }_{1} + {\psi }_{2}\) ，并证明将 \({D}_{n}\) 实现为 \({C}_{nv}\) (同上)其迹为 \({\chi }_{1} + {\psi }_{1}\) 。

\section*{5.4 群 \({\mathrm{D}}_{nh}\)}

该群为乘积 \({\mathrm{D}}_{n} \times  \mathrm{I}\) ，其中 \(\mathrm{I}\) 是阶为2的一群，由满足 \({\iota }^{2} = 1\) 的元素 \(\{ 1,\iota \}\) 组成。它的阶为 \({4n}\) 。若 \({\mathrm{D}}_{n}\) 按通常方式作为三维空间的旋转与反射群实现(参见5.3,(a))，则 \({\mathrm{D}}_{nh}\) 可实现为由 \({\mathrm{D}}_{n}\) 与穿过原点的反射 \(\iota\) 生成的群。

根据定理10， \({\mathrm{D}}_{nh}\) 的不可约表示是 \({\mathrm{D}}_{n}\) 与I的表示的张量积。群I只有两个不可约表示，均为1维。它们的特征 \(g\) 和 \(u\) 由下表给出:

\begin{center}
\adjustbox{max width=\textwidth}{
\begin{tabular}{|c|c|c|}
\hline
\phantom{X} & 1 & \(l\) \\
\cline{1-3}
\(g\) & 1 & 1 \\
\cline{1-3}
\(u\) & 1 & -1 \\
\cline{1-3}
\hline
\end{tabular}
}
\end{center}

因此， \({\mathrm{D}}_{nh}\) 的不可约表示数是 \({\mathrm{D}}_{n}\) 的两倍。更确切地说，每个不可约特征 \(\chi\) 在 \({\mathrm{D}}_{n}\) 上都会在 \({\mathrm{D}}_{nh}\) 上定义两个不可约特征 \({\chi }_{g}\) 和 \({\chi }_{u}\) ，定义如下:

\[
\left( {x \in  {\mathrm{D}}_{n}}\right)
\]

\begin{center}
\adjustbox{max width=\textwidth}{
\begin{tabular}{|c|c|c|}
\hline
\phantom{X} & \(x\) & 文字 \\
\cline{1-3}
\({\chi }_{g}\) & \(\chi \left( x\right)\) & \(\chi \left( x\right)\) \\
\cline{1-3}
\({\chi }_{u}\) & \(\chi \left( x\right)\) & \(- \chi \left( x\right)\) \\
\cline{1-3}
\hline
\end{tabular}
}
\end{center}

例如， \({\mathrm{D}}_{n}\) 的字符 \({\chi }_{1}\) 产生字符 \({\chi }_{1g}\) 和 \({\chi }_{1u}\) :

\begin{center}
\adjustbox{max width=\textwidth}{
\begin{tabular}{|c|c|c|c|c|}
\hline
\phantom{X} & \({r}^{k}\) & \(s{r}^{k}\) & \(t{r}^{k}\) & \({\iota s}{r}^{k}\) \\
\cline{1-5}
\({\chi }_{1g}\) & \(2\cos {2\pi k}/n\) & 0 & \(2\cos {2\pi k}/n\) & 0 \\
\cline{1-5}
\({\chi }_{1u}\) & 2 cos \({2\pi k}/n\) & 0 & \(- 2\cos {2\pi k}/n\) & 0 \\
\cline{1-5}
\hline
\end{tabular}
}
\end{center}

同样的结论也适用于 \({\mathrm{D}}_{n}\) 的其他元素。

\section*{5.5 群 \({\mathrm{D}}_{\infty }\)}

这是保持原点不变的平面旋转与反射的群。它包含旋转群 \({\mathrm{C}}_{\infty }\)  \({r}_{\alpha }\) ；若 \(s\) 是任意一个反射，则有如下关系:

\[
{s}^{2} = 1,\;s{r}_{\alpha }s = {r}_{-\alpha }.
\]

\({\mathrm{D}}_{\infty }\) 的每个元素可唯一表示为两种形式之一:若其属于 \({C}_{\infty }\) ，则为 \({r}_{\alpha }\) 的形式；若不属于 \({C}_{\infty }\) ，则为 \(s{r}_{\alpha }\) 的形式；作为拓扑空间， \({\mathrm{D}}_{\infty }\) 由两个不相交的圆组成。 \({\mathrm{D}}_{\infty }\) 的不变测度是测度 \({d\alpha }/{4\pi }\) 。更精确地，函数 \(f\) 的平均值 \({\int }_{\mathrm{G}}f\left( t\right) {dt}\) 由下式给出

\[
{\int }_{\mathrm{G}}f\left( t\right) {dt} = \frac{1}{4\pi }{\int }_{0}^{2\pi }f\left( {r}_{\alpha }\right) {d\alpha } + \frac{1}{4\pi }{\int }_{0}^{2\pi }f\left( {s{r}_{\alpha }}\right) {d\alpha }.
\]

特别地，2.6 中的投影 \({p}_{i}\) 为:

\[
{p}_{i}x = \frac{{n}_{i}}{4\pi }{\int }_{0}^{2\pi }{\chi }_{i}{\left( {r}_{\alpha }\right) }^{ * }{\rho }_{{r}_{\alpha }}\left( x\right) {d\alpha } + \frac{{n}_{i}}{4\pi }{\int }_{0}^{2\pi }{\chi }_{i}{\left( s{r}_{\alpha }\right) }^{ * }{\rho }_{s{r}_{\alpha }}\left( x\right) {d\alpha }.
\]

作为三维空间中刚体运动群的 \({\mathrm{D}}_{\alpha }\) 的实现

共有两种:

(a) 常见的实现(在 Eyring [5] 中记作 \({\mathrm{D}}_{\infty }\) )。旋转以 Oz 为轴，反射关于通过 O 的 Oxy 平面上的直线。

(b) 通过群 \({\mathrm{C}}_{\infty v}\) 的实现(采用 Eyring [5] 的记法):反射改为关于通过 Ox 的平面，而不是 Oxy 的直线。

\section*{群 \({\mathrm{D}}_{\infty }\) 的不可约表示}

它们的构造类似于 \({\mathrm{D}}_{n}\) 的表示。首先有两个一维表示，其特征 \({\psi }_{1}\) 与 \({\psi }_{2}\) 由下表给出:

\begin{center}
\adjustbox{max width=\textwidth}{
\begin{tabular}{|c|c|c|}
\hline
\phantom{X} & \({r}_{a}\) & \(s{r}_{a}\) \\
\cline{1-3}
\({\psi }_{1}\) & 1 & 1 \\
\cline{1-3}
\({\psi }_{2}\) & 1 & -1 \\
\cline{1-3}
\hline
\end{tabular}
}
\end{center}

有一列不可约表示 \({\rho }^{h}\) ，其次数为 2 \(\left( {h = 1,2,\ldots }\right)\) ，由以下公式给出:

\[
{\rho }^{h}\left( {r}_{\alpha }\right)  = \left( \begin{matrix} {e}^{ih\alpha } & 0 \\  0 & {e}^{-{ih\alpha }} \end{matrix}\right) ,\;{\rho }^{h}\left( {s{r}_{\alpha }}\right)  = \left( \begin{matrix} 0 & {e}^{-{ih\alpha }} \\  {e}^{ih\alpha } & 0 \end{matrix}\right) .
\]

它们的特征 \({\chi }_{1},{\chi }_{2},\ldots\) 在下列取值:

\[
{\chi }_{h}\left( {r}_{\alpha }\right)  = 2\cos \left( {h\alpha }\right) ,\;{\chi }_{h}\left( {s{r}_{\alpha }}\right)  = 0.
\]

可以证明这些就是 \({\mathrm{D}}_{\infty }\) 的所有不可约表示(模同构)。

\section*{5.6 群 \({\mathrm{D}}_{\infty h}\)}

该群是乘积 \({\mathrm{D}}_{\infty } \times  \mathrm{I}\) ；也可表示为由 \({\mathrm{D}}_{\infty }\) 及经原点的反射 \(\iota\) 所生成的群。其元素可唯一写成以下四种形式之一:

\[
{r}_{\alpha },\;s{r}_{\alpha },\;\iota {r}_{\alpha },\;{\iota s}{r}_{\alpha }.
\]

作为拓扑空间，它是四个不相交圆的并。 \({\mathrm{D}}_{\infty h}\) 的不变测度是 \(\left( {1/{8\pi }}\right) {d\alpha }\) 。如上所述，这意味着函数 \(f\) 在 \({\mathrm{D}}_{\infty h}\) 上的平均值 \({\int }_{\mathrm{G}}f\left( t\right) {dt}\) 由下式给出:

\[
{\int }_{\mathrm{G}}f\left( t\right) {dt} = \frac{1}{8\pi }{\int }_{0}^{2\pi }f\left( {r}_{\alpha }\right) {d\alpha } + \frac{1}{8\pi }{\int }_{0}^{2\pi }f\left( {s{r}_{\alpha }}\right) {d\alpha } + \frac{1}{8\pi }{\int }_{0}^{2\pi }f\left( {u}_{\alpha }\right) {d\alpha }
\]

\[
+ \frac{1}{8\pi }{\int }_{0}^{2\pi }f\left( {{\iota s}{r}_{\alpha }}\right) {d\alpha }
\]

留给读者推导 2.6 中投影 \({p}_{i}\) 的显式表达式。

与 \({\mathrm{D}}_{nh}\) 的情形类似， \({\mathrm{D}}_{\infty h}\) 的不可约表示由 \({\mathrm{D}}_{\infty }\) 成对出现: \({\mathrm{D}}_{\infty }\) 的每个特征 \(\chi\) 会在 \({\mathrm{D}}_{\infty h}\) 上生成两个特征 \({\chi }_{g}\) 和 \({\chi }_{u}\) 。

因此，例如， \({\mathrm{D}}_{\infty }\) 的特征 \({\chi }_{3}\) 给出:

\begin{center}
\adjustbox{max width=\textwidth}{
\begin{tabular}{|c|c|c|c|c|}
\hline
\phantom{X} & \({r}_{a}\) & \(S{r}_{\alpha }\) & \(\iota {r}_{a}\) & \({LS}{T}_{\alpha }\) \\
\cline{1-5}
\({\chi }_{3g}\) & \(2\cos {3\alpha }\) & 0 & \(2\cos {3\alpha }\) & 0 \\
\cline{1-5}
\({\chi }_{3u}\) & 2 cos \({3\alpha }\) & 0 & \(- 2\cos {3\alpha }\) & 0 \\
\cline{1-5}
\hline
\end{tabular}
}
\end{center}

\section*{5.7 交替群 \({\mathfrak{A}}_{4}\)}

这是在集合 \(\{ a,b,c,d\}\) 上所有偶置换组成的群，含 4 个元素；它同构于在 \({\mathbf{R}}^{3}\) 中保持以原点为重心的正四面体不变的旋转群。它有 12 个元素:

恒等元 1；

3 个阶为 \(2,x = \left( {ab}\right) \left( {cd}\right) ,y = \left( {ac}\right) \left( {bd}\right) ,z = \left( {ad}\right) \left( {bc}\right)\) 的元素，对应于通过连接两条对边中点的直线的四面体反射；

8 个三阶元素: (abc), (acb), ..., (bcd)，它们对应于绕连结顶点与对面重心的直线对 \(\pm  {120}^{ \circ  }\) 的旋转。

我们用 (abc) 表示循环置换 \(a \mapsto  b,b \mapsto  c,c \mapsto  a,d \mapsto  d\) ；同样， \(\left( {ab}\right) \left( {cd}\right)\) 表示置换 \(a \mapsto  b,b \mapsto  a,c \mapsto  d,d \mapsto  c\) ，它是换位 \(\left( {ab}\right)\) 和 \(\left( {cd}\right)\) 的乘积。

设 \(t = \left( {abc}\right) ,\mathrm{K} = \left\{  {1,t,{t}^{2}}\right\}\) 与 \(\mathrm{H} = \{ 1,x,y,z\}\) 。我们有

\[
{tx}{t}^{-1} = z,\;{tz}{t}^{-1} = y,\;{ty}{t}^{-1} = x;
\]

此外 \(\mathrm{H}\) 与 \(\mathrm{K}\) 是 \({\mathfrak{A}}_{4},\mathrm{H}\) 的子群，且 \({\mathfrak{A}}_{4},\mathrm{H}\) 为正规子群，并且 \(\mathrm{H} \cap  \mathrm{K} = \{ 1\}\) 。易见 \({\mathfrak{A}}_{4}\) 的每个元素都能唯一地写成乘积 \(h \cdot  k\) ，其中 \(h \in  \mathrm{H}\) 与 \(k \in  \mathrm{K}\) 。

也有人说 \({\mathfrak{A}}_{4}\) 是由正规子群 \(\mathrm{H}\) 的半直积由 \(\mathrm{K}\) ；注意这不是直积，因为 K 的元素与 \(\mathrm{H}\) 的元素不交换。

在 \({\mathfrak{A}}_{4} : \{ 1\} ,\{ x,y,z\} ,\{ t,{tx},{ty},{tz}\}\) 中有 4 个共轭类，且 \(\left\{  {{t}^{2},{t}^{2}x}\right.\) 、 \(\left. {{t}^{2}y,{t}^{2}z}\right\}\) ，因此有 4 个不可约特征。存在三个次数为 1 的特征，对应于群 \(\mathrm{K}\) 的三个特征 \({\chi }_{0},{\chi }_{1}\) 与 \({\chi }_{2}\) (参见 5.1)，通过对 \({\mathfrak{A}}_{4}\) 设定 \({\chi }_{i}\left( {h \cdot  k}\right)  = {\chi }_{i}\left( k\right)\) 对 \(h \in  \mathrm{H}\) 与 \(k \in  \mathrm{K}\) 的取值而扩张得到。最后一个特征 \(\psi\) 可由引理 2 至命题 5 推定；它被发现是 \({\mathfrak{A}}_{4}\) 在 \({\mathbf{R}}^{3}\) 中自然表示的特征(通过线性扩张到 \({\mathbf{C}}^{3}\) )。因此我们得到 \({\mathfrak{A}}_{4}\) 的下列特征表:

\begin{center}
\adjustbox{max width=\textwidth}{
\begin{tabular}{|c|c|c|c|c|}
\hline
\phantom{X} & 1 & \(x\) & \(t\) & \({t}^{2}\) \\
\cline{1-5}
\({\chi }_{0}\) & 1 & 1 & 1 & 1 \\
\cline{1-5}
\({\chi }_{1}\) & 1 & 1 & w & \({w}^{2}\) \\
\cline{1-5}
\({\chi }_{2}\) & 1 & 1 & \({w}^{2}\) & w \\
\cline{1-5}
\(\psi\) & 3 & -1 & 0 & 0 \\
\cline{1-5}
\hline
\end{tabular}
}
\end{center}

with

\[
w = {e}^{{2\pi i}/3} =  - \frac{1}{2} + i\frac{\sqrt{3}}{2}.
\]

Exercise

5.4. 设 \(\theta \left( 1\right)  = \theta \left( x\right)  = 1\) 与 \(\theta \left( y\right)  = \theta \left( z\right)  =  - 1\) ；这是 \(H\) 的一维表示。由 \(\theta\) 诱导的 \({\mathfrak{A}}_{4}\) 的表示(参见 3.3)是三维的；证明它是不可约的，且其迹为 \(\psi\) 。

\section*{5.8 对称群 \({\mathfrak{S}}_{4}\)}

这是对 \(\{ a,b,c,d\}\) 的所有置换构成的群；它同构于固定正四面体的所有刚性运动构成的群。它有 24 个元素，分为 5 个共轭类:

恒等元 1；

6 个换位: (ab), (ac), (ad), (bc), (bd), (cd)；

\({\mathfrak{A}}_{4} : x = \left( {ab}\right) \left( {cd}\right) ,y = \left( {ac}\right) \left( {bd}\right) ,z = \left( {ad}\right) \left( {bc}\right)\) 中 3 个阶为 2 的元素；

8 个阶为 3 的元素: \(\left( {abc}\right) ,\ldots ,\left( {bcd}\right)\) ；

6 个阶为 4 的元素:(abcd), (abdc), (acbd), (acdb), (adbc), (adcb)。

设 \(\mathrm{H} = \{ 1,x,y,z\}\) 且令 \(\mathrm{L}\) 为固定 \(d\) 的置换群。如前一节所示， \({\mathfrak{S}}_{4}\) 为由正规子群 \(H\) 的半直积 \(L\) 。每个 \(L\) 的表示 \(\rho\) 按公式 \(\rho \left( {h \cdot  l}\right)  = \rho \left( l\right)\) 对 \(h \in  \mathrm{H},l \in  \mathrm{L}\) 进行扩展以成为 \({\mathfrak{S}}_{4}\) 的表示。由此得到 \({\mathfrak{S}}_{4}\) 的三个不可约表示(参见 2.5)，维数为 1、1 与 2。另一方面， \({\mathfrak{S}}_{4}\) 在 \({\mathbf{C}}^{3}\) 中的自然表示是不可约的(因其限制到 \({\mathfrak{A}}_{4}\) 是不可约)，其与 \({\mathfrak{S}}_{4}\) 的非平凡一维表示的张量积也同样不可约。由此得出 \({\mathfrak{S}}_{4}\) 的下列迹表:

\begin{center}
\adjustbox{max width=\textwidth}{
\begin{tabular}{|c|c|c|c|c|c|}
\hline
\phantom{X} & 1 & (ab) & (ab)(cd) & (abc) & (abcd) \\
\cline{1-6}
\({\chi }_{0}\) & 1 & 1 & 1 & 1 & 1 \\
\cline{1-6}
\(\varepsilon\) & 1 & -1 & 1 & 1 & -1 \\
\cline{1-6}
\(\theta\) & 2 & 0 & 2 & -1 & 0 \\
\cline{1-6}
\(\psi\) & 3 & 1 & -1 & 0 & -1 \\
\cline{1-6}
εψ & 3 & -1 & -1 & 0 & 1 \\
\cline{1-6}
\hline
\end{tabular}
}
\end{center}

注意 \({\mathfrak{S}}_{4}\) 的字符值是整数；这是对称群表示的一般性质(参见 13.1)。

\section*{5.9 立方体的群}

在 \({\mathbf{R}}^{3}\) 中考虑顶点为 \(\left( {x,y,z}\right)\) ， \(x =  \pm  1,y =  \pm  1\) 和 \(z =  \pm  1\) 的立方体 \(\mathbf{C}\) 。设 \(G\) 为将 \({\mathbf{R}}^{3}\) 自同构并稳定立方体 C(即置换其八个顶点)的一组同构；这个群 \(G\) 可以通过几种方式描述:

(i) 群 \(G\) 包含置换 \(\{ x,y,z\}\) 的群 \({\mathfrak{S}}_{3}\) 以及由下列变换组成的阶为 8 的群 \(M\) 。

\[
\left( {x,y,z}\right)  \mapsto  \left( {\pm x, \pm  y, \pm  z}\right) .
\]

很容易检查 \(\mathrm{G}\) 是 \({\mathfrak{S}}_{3}\) 的半直积，正规子群为 M；其阶为 \(6 \cdot  8 = {48}\) 。

(ii) 记 \(\left( {x,y,z}\right)  \mapsto  \left( {-x, - y, - z}\right)\) 为关于原点的反射 \(\iota\) 。设 \(T\) 为顶点为 \(\left( {1,1,1}\right) ,\left( {1, - 1, - 1}\right)\) 、 \(\left( {-1,1, - 1}\right) ,\left( {-1, - 1,1}\right)\) 的四面体，且设 \({T}^{\prime } = {\iota T}\) ； \(C\) 的每个顶点都是 \(T\) 或 \({T}^{\prime }\) 的一个顶点。记 \(S\left( T\right)\) 为将 \({R}^{3}\) 自同构并稳定 \(\mathrm{T}\) 的同构群；对于 \(s \in  \mathrm{S}\left( \mathrm{T}\right)\) 我们有 \(s{\mathrm{\;T}}^{\prime } = {s\iota }\mathrm{T} = {\iota s}\mathrm{\;T} = {\mathrm{T}}^{\prime }\) ，这表明 \(s\) 稳定 \(C\) 的顶点集，从而属于 \(G\) 。因此 \(S\left( T\right)  \subset  G\) ，并且我们立刻看到 \(G\) 是 \(S\left( T\right)\) 与群 \(I = \{ 1,\iota \}\) 的直积。由于 \(S\left( T\right)  = {\mathfrak{S}}_{4}\) ， \(G\) 的不可约特征由 \({\mathfrak{S}}_{4}\) 的那些成对得到，正如 \({D}_{nh}\) 的由 \({D}_{n}\) 的得到那样。因此有 4 个 1 维不可约特征、2 个 2 维和 4 个 3 维；其精确描述留给读者。

\section*{练习}

5.4 从分解 \(G = {\mathfrak{S}}_{4} \times  I\) 和 \({\mathfrak{S}}_{4} = {\mathfrak{S}}_{3} \cdot  L\) 恢复半直积分解 \(G = {\mathfrak{S}}_{3} \cdot  M\) (参见 5.8)。

5.5 设 \({G}_{ + }\) 为 \(G\) 的行列式为 1 的子群(立方体的旋转群)。证明:若将 \(G\) 分解为 \(S\left( T\right)  \times  I\) ，投影 \(G \rightarrow  S\left( T\right)\) 对 \({G}_{ + }\) 给出到 \(S\left( T\right)  = {\mathfrak{S}}_{4}\) 的同构。

参考书目:第一部分

若干书籍论述了有限群的表示论。首先提及一本经典著作:

[1] H. Weyl. The Theory of Groups and Quantum Mechanics. Dover Publ., 1931.

另见:

[2] M. Hall. The Theory of Groups. Macmillan, New York, 1959.

[3] G. G. Hall. Applied Group Theory. Mathematical Physics Series, Longmans, 1967.

关于诱导表示及其应用的讨论见:

[4] A. J. Coleman. Induced and subduced representations. Group Theory and Its Applications, edited by M. Loebl. Academic Press, New York, 1968.

关于 \({\mathbf{R}}^{3}\) 中刚体运动群的标准符号和特征标表，见:

[5] H. Eyring, J. Walter, and G. Kimball. Quantum Chemistry. John Wiley and Sons, New York, 1944.

(表见附录 VII，第 376–388 页。)

关于紧群，见 [1]、[4]，以及:

[6] L. Loomis. An Introduction to Abstract Harmonic Analysis. Van Nostrand, New York, 1953.

对特征标理论历史感兴趣的读者可查阅 Frobenius 的原始论文:

[7] F. G. Frobenius. Gesammelte Abhandlungen, Bd. III. Springer-Verlag, 1969.

\section*{II 零特征下的表示论}

除非另有明确说明，所有群均假定为有限群，所有向量空间(或所有模)均假定为有限维(或有限生成)。

在第6到11章(6.1除外)中，底域为复数域 \(\mathbf{C}\) 。

\section*{第6章 群代数}

\section*{6.1 表示与模}

设 \(\mathrm{G}\) 为一个有限阶为 \(g\) 的群， \(\mathrm{K}\) 为一交换环。我们用 \(\mathrm{K}\left\lbrack  \mathrm{G}\right\rbrack\) 表示群代数 \(\mathrm{G}\) 在 \(\mathrm{K}\) 上的代数；该代数有一组以 \(\mathrm{G}\) 元素为指标的基，大多时候我们将这组基同 \(\mathrm{G}\) 视为相同。任意元素 \(f\) ∈ \(\mathrm{K}\left\lbrack  \mathrm{G}\right\rbrack\) 可唯一写成以下形式

\[
f = \mathop{\sum }\limits_{{s \in  \mathrm{G}}}{a}_{s}s,\;\text{ with }{a}_{s} \in  \mathrm{K},
\]

并且 \(\mathrm{K}\left\lbrack  \mathrm{G}\right\rbrack\) 中的乘法扩展自 \(\mathrm{G}\) 中的乘法。

设 \(\mathrm{V}\) 为一个 \(\mathrm{K}\) -模， \(\rho  : \mathrm{G} \rightarrow  \mathrm{{GL}}\left( \mathrm{V}\right)\) 是在 \(\mathbf{V}\) 上的 \(\mathbf{G}\) 的线性表示。对任意 \(s \in  \mathbf{G}\) 和 \(x \in  \mathbf{V}\) ，设 \({sx} = {\rho }_{s}x\) ；由线性可定义 \({fx}\) ，对任意 \(f \in  \mathrm{K}\left\lbrack  \mathrm{G}\right\rbrack\) 和 \(x \in  \mathrm{V}\) 亦然。因此 \(\mathrm{V}\) 具有左 K[G]-模结构；反过来，此类结构决定了 G 在 V 上的一个线性表示。下文中我们不加区分地使用“线性表示”或“模”这一术语。

命题9. 若 \(\mathrm{K}\) 为特征零的域，则代数 \(\mathrm{K}\left\lbrack  \mathrm{G}\right\rbrack\) 为半单代数。

(关于半单代数的基本事实，参见例如 Bourbaki [8] 或 Lang [10]。)

说 \(\mathrm{K}\left\lbrack  \mathrm{G}\right\rbrack\) 是半单代数等价于说每个 K[G]-模 \(\mathbf{V}\) 是半单的，即每个子模 \({\mathbf{V}}^{\prime }\) 在作为 \(\mathrm{K}\left\lbrack  \mathrm{G}\right\rbrack\) -模的 \(\mathrm{V}\) 中都是一个直和因子。此点可用与1.3中相同的平均法证明:先选取一个从 \(\mathbf{V}\) 到 \({\mathbf{V}}^{\prime }\) 的 K-线性投影 \(p\) ，然后取其在 \(\mathbf{G}\) 作用下的变换的平均 \({p}^{0} = \left( {1/g}\right) \mathop{\sum }\limits_{{s \in  \mathrm{G}}}{sp}{s}^{-1}\) 。

由此得到的投影 \({p}^{0}\) 是 \(\mathrm{K}\left\lbrack  \mathrm{G}\right\rbrack\) -线性的，这意味着 \({\mathbf{V}}^{\prime }\) 作为 \(\mathbf{K}\left\lbrack  \mathbf{G}\right\rbrack\) -模是 \(\mathbf{V}\) 的一个直和因子。

推论。代数 \(\mathbf{K}\left\lbrack  \mathbf{G}\right\rbrack\) 是若干矩阵代数在以 \(\mathbf{K}\) 为中心的有限度斜域上的直积。

这是半单代数结构定理的一个结果(同上出处)。

\section*{习题}

6.1. 设 \(\mathrm{K}\) 为特征为 \(p > 0\) 的域。证明以下两条性质等价:

(i) \(\mathrm{K}\left\lbrack  \mathrm{G}\right\rbrack\) 是半单的。

(ii) \(p\) 不整除 \(g\) 的阶 \(\mathbf{G}\) 。

(事实上 (ii) \(\Rightarrow\) (i) 如上所证。为证明逆命题，说明若 \(p\) 整除 \(g\) ，则由在 \(\mathrm{K}\left\lbrack  \mathrm{G}\right\rbrack\) 中满足 \(\sum {a}_{s}s\) 的所有 \(\sum {a}_{s} = 0\) 所成的理想作为 \(\mathrm{K}\left\lbrack  \mathrm{G}\right\rbrack\) 的模不是直和因子。)

\section*{6.2 C[G] 的分解}

从此我们取 \(\mathrm{K} = \mathrm{C}\) (尽管任何特征零的代数闭域同样适用)，因此每个在 \(\mathbf{C}\) 上的有限度斜域都等于 \(\mathbf{C}\) 。命题 9 的推论表明 \(\mathbf{C}\left\lbrack  \mathbf{G}\right\rbrack\) 是若干矩阵代数 \({\mathbf{M}}_{{n}_{i}}\left( \mathbf{C}\right)\) 的直积。更准确地，设 \({\rho }_{i} : \mathbf{G} \rightarrow  \mathbf{{GL}}\left( {\mathrm{W}}_{i}\right)\) 、 \(1 \leq  i \leq  h\) 为 \(\mathbf{G}\) 的不同不可约表示(至同构)，并设 \({n}_{i} = \dim \left( {\mathbf{W}}_{i}\right)\) ，使得端同态环 \(\operatorname{End}\left( {\mathbf{W}}_{i}\right)\) 与 \({\mathbf{W}}_{i}\) 同构为 \({\mathbf{M}}_{{n}_{i}}\left( \mathbf{C}\right)\) 。映射 \({\rho }_{i} : \mathbf{G} \rightarrow  \mathbf{{GL}}\left( {\mathbf{W}}_{i}\right)\) 线性扩展得代数同态 \({\widetilde{\rho }}_{i} : \mathbf{C}\left\lbrack  \mathrm{G}\right\rbrack   \rightarrow  \operatorname{End}\left( {\mathrm{W}}_{i}\right)\) ；族 \(\left( {\widetilde{\rho }}_{i}\right)\) 给出一个同态

\[
\widetilde{\rho } : \mathbf{C}\left\lbrack  \mathrm{G}\right\rbrack   \rightarrow  \mathop{\prod }\limits_{{i = 1}}^{{i = h}}\operatorname{End}\left( {\mathrm{\;W}}_{i}\right)  \simeq  \mathop{\prod }\limits_{{i = 1}}^{{i = h}}{\mathbf{M}}_{{n}_{i}}\left( \mathbf{C}\right) .
\]

命题 10. 上述定义的同态 \(\widetilde{\rho }\) 是一个同构。

这是半单代数的一般性质。在本例中，可按如下方式验证:首先， \(\widetilde{\rho }\) 是满的。否则将存在一个在 \(\prod {\mathbf{M}}_{{n}_{i}}\left( \mathbf{C}\right)\) 上非零的线性型，使其在 \(\widetilde{\rho }\) 的像上取零；这会在表示的系数 \({\rho }_{i}\) 上给出非平凡的关系，而根据 2.2 的正交公式这是不可能的。另一方面， \(\mathbf{C}\left\lbrack  \mathrm{G}\right\rbrack\) 和 \(\prod {\mathbf{M}}_{{n}_{i}}\left( \mathbf{C}\right)\) 的维数均为 \(g = \sum {n}_{i}^{2}\) ，参见 2.4；因此既然 \(\widetilde{\rho }\) 是满的，它必是双射的。

可以描述作为 \(\widetilde{\rho }\) 逆映射的那个同构:

命题11(傅里叶反演公式)。设 \({\left( {u}_{i}\right) }_{1 \leq  i \leq  h}\) 为 \(\prod \operatorname{End}\left( {\mathrm{W}}_{i}\right)\) 中的一个元素，令 \(u = \mathop{\sum }\limits_{{s \in  \mathrm{G}}}u\left( s\right) s\) 为 \(\mathbf{C}\left\lbrack  \mathrm{G}\right\rbrack\) 中的元素，且对所有 \(i\) 有 \(\widetilde{\rho }\left( u\right)  = {u}_{i}\) 。第s项系数 \(u\left( s\right)\) 的 \(u\) 由下式给出

\[
u\left( s\right)  = \frac{1}{g}\mathop{\sum }\limits_{{i = 1}}^{{i = h}}{n}_{i}{\operatorname{Tr}}_{{\mathrm{W}}_{i}}\left( {{\rho }_{i}\left( {s}^{-1}\right) {u}_{i}}\right) ,\;\text{ where }{n}_{i} = \dim \left( {\mathrm{W}}_{i}\right) .
\]

由线性性，只需在 \(u\) 等于 \(t\) 的 \(G\) 元素时检验该公式。我们有

\[
u\left( s\right)  = {\delta }_{st}\;\text{ and }\;{\operatorname{Tr}}_{{\mathrm{W}}_{i}}\left( {{\rho }_{i}\left( {s}^{-1}\right) {u}_{i}}\right)  = {\chi }_{i}\left( {{s}^{-1}t}\right) ,
\]

其中 \({\chi }_{i}\) 为与 \({\mathrm{W}}_{i}\) 对应的 \(\mathrm{G}\) 的不可约特征。因此还需证明

\[
{\delta }_{st} = \frac{1}{g}\mathop{\sum }\limits_{{i = 1}}^{{i = h}}{n}_{i}{\chi }_{i}\left( {{s}^{-1}t}\right)
\]

这是命题2.4中命题5的推论1和2的一个直接结果。

练习

6.2(普朗歇雷尔公式)。设 \(u = \sum u\left( s\right) s\) 和 \(v = \sum v\left( s\right) s\) 为 \(\mathbf{C}\left\lbrack  \mathrm{G}\right\rbrack\) 中的两个元素，置 \(\langle u,v\rangle  = \left( {1/g}\right) \mathop{\sum }\limits_{{s \in  \mathrm{G}}}u\left( {s}^{-1}\right) v\left( s\right)\) 。证明下列公式

\[
\langle u,v\rangle  = \mathop{\sum }\limits_{{i = 1}}^{{i = h}}{n}_{i}{\operatorname{Tr}}_{{\mathbf{W}}_{i}}\left( {{\widetilde{\rho }}_{i}\left( {uv}\right) }\right) .
\]

[约化到 \(u\) 和 \(v\) 属于G的情形。]

6.3 设 \(U\) 为包含G的乘法群 \(C\left\lbrack  G\right\rbrack\) 的有限子群。令 \(u = \sum u\left( s\right) s\) 和 \({u}^{\prime } = \sum {u}^{\prime }\left( s\right) s\) 为 \(U\) 中的两个元素，满足 \(u \cdot  {u}^{\prime } = 1\) ；令 \({u}_{i}\) (分别 \({u}_{i}^{\prime }\) )为 \({\widetilde{\rho }}_{i}\) 作用下 \(u\) (分别 \({u}^{\prime }\) )在End \(\left( {W}_{i}\right)\) 中的像。

(a) 证明 \({\rho }_{i}\left( {s}^{-1}\right) {u}_{i} = {\widetilde{\rho }}_{i}\left( {{s}^{-1}u}\right)\) 的特征值都是单位根。由此得出，对于任意 \(s \in  \mathbf{G}\) 与任意 \(i\) ，有

\[
{\operatorname{Tr}}_{{\mathrm{W}}_{i}}{\left( {\rho }_{i}\left( {s}^{-1}\right) {u}_{i}\right) }^{ * } = {\operatorname{Tr}}_{{\mathrm{W}}_{i}}\left( {{u}_{i}^{\prime }{\rho }_{i}\left( s\right) }\right)  = {\operatorname{Tr}}_{{\mathrm{W}}_{i}}\left( {{\rho }_{i}\left( s\right) {u}_{i}^{\prime }}\right) ,
\]

从而，应用命题 \({11},u{\left( s\right) }^{ * } = {u}^{\prime }\left( {s}^{-1}\right)\) 。

(b) 证明 \(\mathop{\sum }\limits_{{s \in  \mathrm{G}}}{\left| u\left( s\right) \right| }^{2} = 1\) 。[用(a)。]

(c) 设 \(U\) 包含于 \(\mathbf{Z}\left\lbrack  G\right\rbrack\) ，使得 \(u\left( s\right)\) 为整数。证明除了一个等于 \(\pm  1\) 的以外，其余所有 \(u\left( s\right)\) 均为零。由此得出 \(U\) 包含于由形如 \(\pm  t\) 的元素构成的群 \(\pm  G\) ，其中 \(t \in  \mathbf{G}\) 。

(d) 设 \(G\) 为阿贝尔群。证明乘法群 \(\mathbf{Z}\left\lbrack  G\right\rbrack\) 中所有有限阶元素均包含于 \(\pm  G\) (Higman定理)。

\section*{6.3 C[G]的中心}

这是 \(\mathbf{C}\left\lbrack  \mathbf{G}\right\rbrack\) 的元素集合，这些元素与 \(\mathbf{C}\left\lbrack  G\right\rbrack\) 中的所有元素都可交换(或等价地，与 \(\mathbf{G}\) 的所有元素可交换)。

对于 \(c\) 为 \(\mathbf{G}\) 的一个共轭类，设为 \({e}_{c} = \mathop{\sum }\limits_{{s \in  c}}s\) 。立刻可检验出 \({e}_{c}\) 构成了 \(\mathbf{C}\left\lbrack  \mathbf{G}\right\rbrack\) 中心的基；因此后者的维数为 \(h\) ，其中 \(h\) 是 \(\mathbf{G}\) 的类数，参见 2.5。设

\[
{\rho }_{i} : \mathbf{G} \rightarrow  \mathbf{{GL}}\left( {\mathbf{W}}_{i}\right)
\]

为 \(G\) 的一个不可约表示，其特征为 \({\chi }_{i}\) 、次数为 \({n}_{i}\) ，并令 \({\widetilde{\rho }}_{i} : \mathbf{C}\left\lbrack  \mathrm{G}\right\rbrack   \rightarrow  \operatorname{End}\left( {\mathrm{W}}_{i}\right)\) 为对应的代数同态(参见 6.2)。

命题 12. 同态 \({\widetilde{\rho }}_{i}\) 将 \(\mathbf{C}\left\lbrack  \mathrm{G}\right\rbrack\) 的中心映入 \({\mathrm{W}}_{i}\) 的同质伸缩映射集合，并给出一个代数同态

\[
{\omega }_{i} : \text{ Cent. }\mathbf{C}\left\lbrack  \mathbf{G}\right\rbrack   \rightarrow  \mathbf{C}\text{ . }
\]

若 \(u = \sum u\left( s\right) s\) 是 Cent. \(\mathbf{C}\left\lbrack  \mathbf{G}\right\rbrack\) 的一个元素，则有

\[
{\omega }_{i}\left( u\right)  = \frac{1}{{n}_{i}}{\operatorname{Tr}}_{{\mathrm{W}}_{i}}\left( {\widetilde{\rho }\left( u\right) }\right)  = \frac{1}{{n}_{i}}\mathop{\sum }\limits_{{s \in  \mathrm{G}}}u\left( s\right) {\chi }_{i}\left( s\right) .
\]

这只是 2.5 中命题 6 的一种重述。

命题 13. 族 \({\left( {\omega }_{i}\right) }_{1 \leq  i \leq  h}\) 给出 Cent. C[G] 到代数 \({\mathbf{C}}^{h} = \mathbf{C} \times  \cdots  \times  \mathbf{C}\) 的一个同构。

若我们将 \(\mathbf{C}\left\lbrack  \mathrm{G}\right\rbrack\) 识别为各个 \(\operatorname{End}\left( {\mathrm{W}}_{i}\right)\) 的直积，则 \(\mathbf{C}\left\lbrack  \mathrm{G}\right\rbrack\) 的中心成为各个 \(\operatorname{End}\left( {\mathrm{W}}_{i}\right)\) 中心的直积。但 End \(\left( {\mathrm{W}}_{i}\right)\) 的中心由同质伸缩构成。由此得到 Cent. C[G] 到 \(\mathbf{C} \times  \cdots  \times  \mathbf{C}\) 的一个同构，并且显然它就是命题 13 中的那个同构。

练习

6.4. 设

\[
{p}_{i} = \frac{{n}_{i}}{g}\mathop{\sum }\limits_{{s \in  \mathrm{G}}}{\chi }_{i}\left( {s}^{-1}\right) s.
\]

证明 \({p}_{i}\left( {1 \leq  i \leq  h}\right)\) 构成 Cent. C[G] 的一组基，并且 \({p}_{i}^{2} = {p}_{i}\) 、 \({p}_{i}{p}_{j} = 0\) 对于 \(i \neq  j\) ，以及 \({p}_{1} + \cdots  + {p}_{h} = 1\) 。由此得到 2.6 定理 8 的另一证法。证明 \({\omega }_{i}\left( {p}_{j}\right)  = {\delta }_{ij}\) 。

6.5. 证明每个从 Cent. C[G] 到 \(\mathbf{C}\) 的同态都等于某个 \({\omega }_{i}\) 。

\section*{6.4 整数的基本性质}

设 \(\mathrm{R}\) 为交换环，令 \(x \in  \mathrm{R}\) 。若存在整数 \(n \geq  1\) 及 \(\mathbf{Z}\) 中的元素 \({a}_{1},\ldots ,{a}_{n}\) 使得，则称 \(x\) 关于 \(\mathbf{Z}\) 是整性的

\[
{x}^{n} + {a}_{1}{x}^{n - 1} + \cdots  + {a}_{n} = 0.
\]

关于 \(\mathbf{Z}\) 为整性的复数称为代数整数。每一个单位根都是代数整数。若 \(x \in  \mathbf{Q}\) 是代数整数，则有 \(x \in  \mathbf{Z}\) ；否则我们可以把 \(x\) 写成 \(p/q\) 的形式，且 \(p,q \in  \mathbf{Z}\) 、 \(q \geq  2\) 、 \(p,q\) 两两互素。方程 \(\left( *\right)\) 将导致

\[
{p}^{n} + {a}_{1}q{p}^{n - 1} + \cdots  + {a}_{n}{q}^{n} = 0,
\]

因此 \({p}^{n} \equiv  0\left( {{\;\operatorname{mod}\;.}q}\right)\) 与 \(p\) 和 \(q\) 互素的事实矛盾。

命题14. 设 \(x\) 为交换环 R 的一元。下列性质等价:

(i) \(x\) 关于 \(\mathbf{Z}\) 是整的。

(ii) 由 \(x\) 生成的 \(\mathbf{R}\) 的子环 \(\mathbf{Z}\left\lbrack  x\right\rbrack\) 作为一个 \(\mathbf{Z}\) - 模是有限生成的。

(iii) 存在包含 \(\mathbf{Z}\left\lbrack  x\right\rbrack\) 的有限生成子 Z-模在 \(\mathrm{R}\) 中。

(ii) 与 (iii) 的等价性来自于:有限生成的 \(\mathbf{Z}\) -模的子模是有限生成的，因 \(\mathbf{Z}\) 为诺特环。另一方面，若 \(x\) 满足某个方程

\[
{x}^{n} + {a}_{1}{x}^{n - 1} + \cdots  + {a}_{n} = 0,\;\text{ with }{a}_{i} \in  \mathbf{Z},
\]

由此由 \(1,x,\ldots ,{x}^{n - 1}\) 在 \(R\) 中生成的子 Z-模在乘以 \(x\) 下稳定，因而与 \(\mathbf{Z}\left\lbrack  x\right\rbrack\) 重合，这证明了 (i) \(\Rightarrow\) (ii)。反之，设 (ii) 成立，记由 \(1,x,\ldots ,{x}^{n - 1}\) 在 R 中生成的子 Z-模为 \({\mathbf{R}}_{n}\) 。这些 \({\mathbf{R}}_{n}\) 构成递增序列，其并集是 \(\mathbf{Z}\left\lbrack  x\right\rbrack\) ；既然 \(\mathbf{Z}\left\lbrack  x\right\rbrack\) 是有限生成的，必有 \({\mathbf{R}}_{n} = \mathbf{Z}\left\lbrack  x\right\rbrack\) 对足够大的 \(n\) 成立。这表明 \({x}^{n}\) 是以整数系数对 \(1,x,\ldots ,{x}^{n - 1}\) 的线性组合，遂得 (i)。

推论1. 若 \(\mathrm{R}\) 是有限生成的 \(\mathbf{Z}\) -模，则 \(\mathrm{R}\) 中的每一元对 \(\mathbf{Z}\) 都是整的。

此即由 (iii) \(\Rightarrow\) (i) 的推论。

推论2. 在 \(\mathbf{R}\) 中关于 \(\mathbf{Z}\) 整的元素构成 \(\mathbf{R}\) 的一个子环。

设 \(x,y \in  \mathbf{R}\) ; 若 \(x,y\) 对 \(\mathbf{Z}\) 是整扩张，则环 \(\mathbf{Z}\left\lbrack  x\right\rbrack\) 与 \(\mathbf{Z}\left\lbrack  y\right\rbrack\) 在 \(\mathbf{Z}\) 上为有限生成。其张量积 \(\mathbf{Z}\left\lbrack  x\right\rbrack   \otimes  \mathbf{Z}\left\lbrack  y\right\rbrack\) 及其在 \(\mathbf{R}\) 中的像 \(\mathbf{Z}\left\lbrack  {x,y}\right\rbrack\) 也同样成立。于是所有 \(\mathbf{Z}\left\lbrack  {x,y}\right\rbrack\) 的元素都对 \(\mathbf{Z}\) 是整的。

注。在前述定义与结论中，可将 \(\mathbf{Z}\) 替换为任意交换诺特环；对于 (i) \(\Leftrightarrow\) (ii) 甚至无需假设该环为诺特环。

6.5 表征的整性性质。应用

命题15. 设 \(\chi\) 为有限群 G 的表示 \(\rho\) 的特征。则对每个 \(s \in  \mathrm{G}\) ， \(\chi \left( s\right)\) 为代数整数。

确然 \(\chi \left( s\right)\) 是 \(\rho \left( s\right)\) 的迹，故为 \(\rho \left( s\right)\) 的特征值之和，而这些特征值是单位根。

命题16. 设 \(u = \sum u\left( s\right) s\) 为 Cent. C[G] 中的一元，使得 \(u\left( s\right)\) 为代数整数。则 \(u\) 对 \(\mathbf{Z}\) 是整的。

(此命题有意义，因为 Cent. C[G] 是交换环。)

设 \({c}_{i}\left( {1 \leq  i \leq  h}\right)\) 为 \(\mathbf{G}\) 的共轭类并置 \({e}_{i} = \mathop{\sum }\limits_{{s \in  {c}_{i}}}s\) ，参见6.3。对于 \({s}_{i} \in  {c}_{i}\) 我们可将 \(u\) 写成 \(u = \mathop{\sum }\limits_{{i = 1}}^{{i = h}}u\left( {s}_{i}\right) {e}_{i}\) 形式。依据命题14的推论2，仅需证明 \({e}_{i}\) 对 \(\mathbf{Z}\) 是整的。但显然每个乘积 \({e}_{i}{e}_{j}\) 可用整数系数线性组合表示为若干 \({e}_{k}\) 。子群 \(\mathrm{R} = \mathbf{Z}{e}_{1} \oplus  \cdots  \oplus  \mathbf{Z}{e}_{h}\) 在 Cent. C[G] 中因此为子环；因其在 \(\mathbf{Z}\) 上有限生成，其每一元素对 \(\mathbf{Z}\) 都是整的(命题14的推论1)。由此得证。

推论1. 设 \(\rho\) 为群 \(\mathrm{G}\) 的一不可约表示，次数为 \(n\) ，特征为 \(\chi\) 。若 \(u\) 如上，则数 \(\left( {1/n}\right) \mathop{\sum }\limits_{{s \in  \mathrm{G}}}u\left( s\right) \chi \left( s\right)\) 为代数整数。

确实，这个数是同态下 \(u\) 的像

\[
\omega  : \text{ Cent. }\mathbf{C}\left\lbrack  \mathbf{G}\right\rbrack   \rightarrow  \mathbf{C}
\]

与 \(\rho\) 相关(参见命题 12)。由于 \(u\) 在 \(\mathbf{Z}\) 上是整的，其在 \(\omega\) 下的像也同样是整的。

推论 2. \(\mathbf{G}\) 的不可约表示的维数整除 \(\mathbf{G}\) 的阶。

令 \(g\) 为 \(\mathbf{G}\) 的阶。我们对元素 \(u \; = \mathop{\sum }\limits_{{s \in  \mathrm{G}}}\chi \left( {s}^{-1}\right) s\) 应用推论 1，这是合法的，因为 \(\chi\) 是类函数且 \(\chi \left( s\right)\) 是代数整数(命题 15)；于是得到该数

\[
\frac{1}{n}\mathop{\sum }\limits_{{s \in  \mathrm{G}}}\chi \left( {s}^{-1}\right) \chi \left( s\right)  = \frac{g}{n}\langle \chi ,\chi \rangle  = \frac{g}{n}
\]

是一个代数整数。由于该数是有理的，故它属于 \(\mathbf{Z}\) ，即 \(n\) 整除 \(g\) 。

推论 2 可稍作加强(参见 8.1，命题 24 的推论)。以下是朝此方向的第一条结果:

命题 17. 令 \(\mathbf{C}\) 为 \(\mathbf{G}\) 的中心。 \(\mathrm{G}\) 的不可约表示的维数整除 ( \(\mathrm{G} : \mathrm{C}\) )。

令 \(g\) 为 \(\mathrm{G}\) 的阶， \(c\) 为 \(\mathrm{C}\) 的阶，且令 \(\rho  : \mathrm{G} \rightarrow  \mathbf{{GL}}\left( \mathbf{W}\right)\) 为一个度为 \(n\) 的 \(\mathbf{G}\) 的不可约表示。如果 \(s \in  \mathbf{C},\rho \left( s\right)\) 与所有 \(\rho \left( t\right) ,t \in  \mathbf{G}\) 交换；由舒尔引理， \(\rho \left( s\right)\) 为一比例算子。若记作 \(\lambda \left( s\right)\) ，则映射 \(\lambda  : s \mapsto  \lambda \left( s\right)\) 是将 \(\mathbf{C}\) 同态到 \({\mathbf{C}}^{ * }\) 的同态。令 \(m\) 为一个整数 \(\geq  0\) ，并构造张量积

\[
{\rho }^{m} : {\mathrm{G}}^{m} \rightarrow  \mathrm{{GL}}\left( {\mathrm{W} \otimes  \cdots  \otimes  \mathrm{W}}\right)
\]

由 \(m\) 份表示 \(\rho\) 构成；这是群 \({\mathrm{G}}^{m} = \mathrm{G} \times  \cdots  \times  \mathrm{G}\) 的一個不可约表示，参见 3.2 定理 10。映射 \({\rho }^{m}\) 将 \({\mathbf{C}}^{m}\) 中的元素 \(\left( {{s}_{1},\ldots ,{s}_{m}}\right)\) 的像送为比率为 \(\lambda \left( {{s}_{1}\cdots {s}_{m}}\right)\) 的相似变换。子群 \(\mathrm{H}\) (位于 \({\mathrm{C}}^{m}\) )由那些在 \(\mathrm{W} \otimes  \cdots  \otimes  \mathrm{W}\) 上被 \({s}_{1}\cdots {s}_{m} = 1\) 平凡作用的 \(\left( {{s}_{1},\ldots ,{s}_{m}}\right)\) 构成，因此取商后得到 \({\mathrm{G}}^{m}/\mathrm{H}\) 的一个不可约表示。由命题 16 的推论 2 可得，该表示的维数 \({n}^{m}\) 整除 \({\mathbf{G}}^{m}/\mathrm{H}\) 的阶 \({g}^{m}/{c}^{m - 1}\) 。因此对所有 \(m\) 有 \({\left( g/cn\right) }^{m} \in  {c}^{-1}\mathbf{Z}\) ，这意味着 \(\left( {g/{cn}}\right)\) 为整数(例如参见命题 14)。

(此证明归功于 J. Tate。)

\section*{练习}

6.6. 证明环 \(\mathbf{Z}{e}_{1} \oplus  \cdots  \oplus  \mathbf{Z}{e}_{h}\) 是 \(\mathbf{Z}\left\lbrack  \mathrm{G}\right\rbrack\) 的中心。

6.7. 设 \(\rho\) 为群 \(G\) 的一個不可約表示，次数為 \(n\) ，特徵為 \(\chi\) 。若 \(s \in  \mathbf{G}\) ，則證明 \(\left| {\chi \left( s\right) }\right|  \leq  n\) ，且當且僅當 \(\rho \left( s\right)\) 為齊次算子時取等號〔注意 \(\rho \left( s\right)\) 是 \(n\) 個單位根之和〕。推得 \(\rho \left( s\right)  = 1 \Leftrightarrow  \chi \left( s\right)  = n.\)

6.8. 设 \({\lambda }_{1},\ldots ,{\lambda }_{n}\) 为单位根，且设 \(a = \frac{1}{n}\sum {\lambda }_{i}\) 。证明若 \(a\) 是代数整数，则要么 \(a = 0\) ，要么 \({\lambda }_{1} = \cdots  = {\lambda }_{n} = a\) 。［设 \(\mathbf{A}\) 为 \(a\) 在 \(\mathbf{Q}\) 上的共轭的乘积；证明 \(\left| A\right|  \leq  1\) 。］

6.9. 设 \(\rho\) 是 \(\mathbf{G}\) 的一个不可约表示，次数为 \(n\) ，特征为 \(\chi\) 。设 \(s \in  \mathbf{G}\) 和 \(c\left( s\right)\) 为与 \(s\) 共轭类中元素的个数。证明 \(\left( {c\left( s\right) /n}\right) \chi \left( s\right)\) 是代数整数(取 \(u\) 为 \(s\) 共轭元的和，应用命题16的推论1)。证明若 \(c\left( s\right)\) 与 \(n\) 互素且若 \(\chi \left( s\right)  \neq  0\) ，则 \(\rho \left( s\right)\) 是一個相似变换(注意 \(\left( {1/n}\right) \chi \left( s\right)\) 是代数整数，并应用习题6.8)。

6.10. 设 \(s \in  \mathrm{G},s \neq  1\) 。假定包含 \(s\) 的共轭类的元素个数 \(c\left( s\right)\) 是素数 \(p\) 的幂。证明存在一个不可约特征 \(\chi\) ，不等于恒等特征，使得 \(\chi \left( s\right)  \neq  0\) 与 \(\chi \left( 1\right)  \neq  0\left( {\;\operatorname{mod}\;.p}\right)\) 。［利用公式 \(1 + \mathop{\sum }\limits_{{\chi  \neq  1}}\chi \left( 1\right) \chi \left( s\right)  = 0\) ，参见命题5的推论2，说明若不存在这样的特征 \(\chi\) 则数 \(1/p\) 将是代数整数。］设 \(\rho\) 为特征为 \(\chi\) 的表示，证明 \(\rho \left( s\right)\) 是一個相似变换(用习题6.9)。由此得出，如果 \(\mathrm{N}\) 是 \(\rho\) 的核，则有 \(\mathrm{N} \neq  \mathrm{G}\) ，且 \(s\) 在 \(\mathrm{G}/\mathrm{N}\) 中的像属于 \(\mathrm{G}/\mathrm{N}\) 的中心。

\section*{第7章 诱导表示；Mackey 判别法}

\section*{7.1 诱导}

设 H 是群 G 的子群，R 是 H 的左陪集代表系。设 V 为 C[G]-模，W 为 V 的一个 C[H]-子模。回顾(参见3.3)，若存在 \(\mathrm{W}\) 使得 V 被 \(\mathrm{W}\) 诱导，即有 \(\mathrm{V} = { \oplus  }_{s \in  \mathrm{R}}s\mathrm{\;W}\) ，即 \(\mathrm{V}\) 是像 \(s\mathrm{\;W},s \in  \mathrm{R}\) 的直和(此条件与 \(R\) 的选择无关)。此性质可改写为:

设

\[
{\mathrm{W}}^{\prime } = \mathbf{C}\left\lbrack  \mathrm{G}\right\rbrack  { \otimes  }_{\mathbf{C}\left\lbrack  \mathrm{H}\right\rbrack  }\mathrm{W}
\]

为通过从 \(\mathbf{C}\left\lbrack  \mathrm{H}\right\rbrack\) 到 C[G] 的标量扩张由 \(\mathbf{W}\) 得到的 \(\mathbf{C}\left\lbrack  \mathrm{G}\right\rbrack\) -模。注入映射 \(\mathrm{W} \rightarrow  \mathrm{V}\) 按线性性扩张为一个 C[G]-同态 \(i : {\mathrm{W}}^{\prime } \rightarrow  \mathrm{V}\) 。

命题18. 为了使 \(\mathrm{V}\) 被 \(\mathrm{W}\) 诱导，必要且充分的条件是同态

\[
i : \mathbf{C}\left\lbrack  \mathrm{G}\right\rbrack  { \otimes  }_{\mathbf{C}\left\lbrack  \mathrm{H}\right\rbrack  }\mathrm{W} \rightarrow  \mathrm{V}
\]

为同构。

这是因为作为右 C[H]-模考察时，R 的元素构成 C[G] 的一组基。

\section*{备注}

(1) 由 W 诱导的表示的这一刻画表明诱导表示显然存在且唯一(参见 3.3，定理 11)。

下文中，由 W 诱导的 G 的表示记为 \({\operatorname{Ind}}_{H}^{G}\left( W\right)\) ，若无歧义可简记为 \(\operatorname{Ind}\left( W\right)\) 。

(2) 若 \(V\) 由 \(W\) 诱导，且 \(E\) 为一个 \(C\left\lbrack  G\right\rbrack\) -模，则存在一个自然同构

\[
{\operatorname{Hom}}^{\mathrm{H}}\left( {\mathrm{W},\mathrm{E}}\right)  \cong  {\operatorname{Hom}}^{\mathrm{G}}\left( {\mathrm{V},\mathrm{E}}\right) ,
\]

其中 \({\operatorname{Hom}}^{G}\left( {V,E}\right)\) 表示从 \(C\left\lbrack  G\right\rbrack\) 到 E 的 \(V\) -同态的向量空间， \({\operatorname{Hom}}^{H}\left( {W,E}\right)\) 作类似定义。这源于张量积的一个基本性质(另见 3.3，引理 1)。

(3) 诱导是可传递的:若 \(\mathrm{G}\) 是群 \(\mathrm{K}\) 的子群，则有

\[
{\operatorname{Ind}}_{\mathrm{G}}^{\mathrm{K}}\left( {{\operatorname{Ind}}_{\mathrm{H}}^{\mathrm{G}}\left( \mathrm{W}\right) }\right)  \cong  {\operatorname{Ind}}_{\mathrm{H}}^{\mathrm{K}}\left( \mathrm{W}\right) .
\]

这可以直接看出，或利用张量积的结合性来证明。

命题 19. 设 \(\mathrm{V}\) 为一个 \(\mathrm{C}\left\lbrack  \mathrm{G}\right\rbrack\) -模，且可写为向量子空间的直和 \(\mathrm{V} = { \oplus  }_{i \in  \mathrm{I}}{\mathrm{W}}_{i}\) ，这些子空间被 \(\mathbf{G}\) 所可传递地置换。设 \({i}_{0} \in  \mathbf{I},\mathbf{W} = {\mathbf{W}}_{{i}_{0}}\) 并令 \(\mathrm{H}\) 为 \(\mathrm{G}\) 中稳定 \(\mathrm{W}\) 的稳定子(即所有满足 \(s\mathrm{\;W} = \mathrm{W}\) 的 \(s \in  \mathrm{G}\) 的集合)。则 \(\mathrm{W}\) 在子群 \(\mathrm{H}\) 下不变，且 \(\mathrm{C}\left\lbrack  \mathrm{G}\right\rbrack\) -模 \(\mathrm{V}\) 由 \(\mathbf{C}\left\lbrack  \mathrm{H}\right\rbrack\) -模 \(\mathbf{W}\) 诱导。

此处显然。

备注。为将命题 19 应用于不可约表示 \(V =  \oplus  {W}_{i}\) 的 \(G\) ，只需检验各 \({W}_{i}\) 在 \(G\) 下彼此置换；传递性条件是自动的，因为 \(G\) 在各 \({W}_{i}\) 的集合上的每个轨道都定义了 \(V\) 的子表示。

例子。当 \({W}_{i}\) 维数为 1 时，表示 \(V\) 称为单项表示。

\section*{7.2 诱导表示的特征；互反公式}

保持前述记号。若 \(f\) 为 \(\mathrm{H}\) 上的类函数，考虑在 \(\mathrm{G}\) 上由下式定义的函数 \({f}^{\prime }\) :

\[
{f}^{\prime }\left( s\right)  = \frac{1}{h}\mathop{\sum }\limits_{\substack{{t \in  \mathrm{G}} \\  {{t}^{-1}{st} \in  \mathrm{H}} }}f\left( {{t}^{-1}{st}}\right) \;\text{ where }h = \operatorname{Card}\left( \mathrm{H}\right) .
\]

我们称 \({f}^{\prime }\) 由 \(f\) 诱导，并记作 \({\operatorname{Ind}}_{\mathrm{H}}^{\mathrm{G}}\left( f\right)\) 或 \(\operatorname{Ind}\left( f\right)\) 。

\section*{命题 20。}

(i) 函数 \(\operatorname{Ind}\left( f\right)\) 是 \(\mathrm{G}\) 上的类函数。

(ii) 若 \(f\) 是 \(\mathrm{H}\) 的表示 \(\mathrm{W}\) 的特征，则 \(\operatorname{Ind}\left( f\right)\) 是诱导表示 \(\operatorname{Ind}\left( \mathrm{W}\right)\) 的特征，作用于 \(\mathrm{G}\) 。

断言 (ii) 已在 (3.3, th. 12) 中证明。断言 (i) 可由直接计算得到，或由 (ii) 及每个类函数可表示为若干特征的线性组合这一观察推得。

回想，对于 \({\varphi }_{1}\) 和 \({\varphi }_{2}\) 两个在 \(G\) 上的类函数，我们定义

\[
\left\langle  {{\varphi }_{1},{\varphi }_{2}}\right\rangle   = \frac{1}{g}\mathop{\sum }\limits_{{s \in  \mathrm{G}}}{\varphi }_{1}\left( {s}^{-1}\right) {\varphi }_{2}\left( s\right) ,\;\text{ where }g = \operatorname{Card}\left( \mathrm{G}\right) ,
\]

参见 2.2；当欲更明确指出群 G 时，我们写作 \({\left\langle  {\varphi }_{1},{\varphi }_{2}\right\rangle  }_{\mathrm{G}}\) 而不是 \(\left\langle  {{\varphi }_{1},{\varphi }_{2}}\right\rangle\) 。

此外，若 \({V}_{1}\) 和 \({V}_{2}\) 是两个 \(C\left\lbrack  G\right\rbrack\) -模，我们定义

\[
{\left\langle  {V}_{1},{V}_{2}\right\rangle  }_{G} = \dim  \cdot  {\operatorname{Hom}}^{G}\left( {{V}_{1},{V}_{2}}\right) .
\]

引理 2. 若 \({\varphi }_{1}\) 和 \({\varphi }_{2}\) 分别是 \({V}_{1}\) 和 \({V}_{2}\) 的特征，则有

\[
{\left\langle  {\varphi }_{1},{\varphi }_{2}\right\rangle  }_{\mathrm{G}} = {\left\langle  {V}_{1},{V}_{2}\right\rangle  }_{\mathrm{G}}.
\]

将 \({V}_{1}\) 与 \({V}_{2}\) 分解为直和后，我们可假设它们不可约，此时该引理由特征的正交关系 (2.3, th. 3) 得出。

若 \(\varphi\) (或记作 V)是 G 上的函数(或 G 的表示)，我们用 Res \(\varphi\) (或 Res V)表示其对子群 H 的限制。

定理 13(Frobenius 对偶性)。若 \(\psi\) 是 \(\mathrm{H}\) 上的类函数且 \({\varphi a}\) 是 \(\mathrm{G}\) 上的类函数，则有

\[
\langle \psi ,\operatorname{Res}\varphi {\rangle }_{\mathrm{H}} = \langle \operatorname{Ind}\psi ,\varphi {\rangle }_{\mathrm{G}}.
\]

由于每个类函数可表示为特征的线性组合，我们可假设 \(\psi\) 是一个 \(\mathbf{C}\left\lbrack  \mathrm{H}\right\rbrack\) -模 \(\mathbf{W}\) 的特征且 \(\varphi\) 是 C[G]-模 E 的特征。根据引理 2，只需证明

(*)
\[
\langle \mathrm{W},\operatorname{Res}\mathrm{E}{\rangle }_{\mathrm{H}} = \langle \operatorname{Ind}\mathrm{W},\mathrm{E}{\rangle }_{\mathrm{G}},
\]

即，

\[
\dim {\mathrm{{Hom}}}^{\mathrm{H}}\left( {\mathrm{W},\operatorname{Res}\mathrm{E}}\right)  = \dim {\mathrm{{Hom}}}^{\mathrm{G}}\left( {\operatorname{Ind}\mathrm{W},\mathrm{E}}\right) ,
\]

此项由 7.1 的备注 2(或等价的 3.3 的引理 1)可得。当然也可通过直接计算证明定理 13。

\section*{备注}

(1) 定理 13 表明映射 Res 与 Ind 互为伴随。

(2) 除了双线性形式 \(\langle \alpha ,\beta \rangle\) ，我们也可使用 2.3 中定义的标量积 \(\left( {\alpha  \mid  \beta }\right)\) 。我们有相同的公式:

\[
{\left( \psi  \mid  \operatorname{Res}\varphi \right) }_{\mathrm{H}} = {\left( \operatorname{Ind}\psi  \mid  \varphi \right) }_{\mathrm{G}}.
\]

(3) 我们还提到下列有用的公式

\[
\operatorname{Ind}\left( {\psi  \cdot  \operatorname{Res}\varphi }\right)  = \left( {\operatorname{Ind}\psi }\right)  \cdot  \varphi .
\]

可由简单计算验算，或从公式 Ind \(\left( \mathrm{W}\right)  \otimes  \mathrm{E} \cong\) Ind \(\left( {\mathrm{W} \otimes  \text{ Res E }}\right)\) 推出，参见 3.3，例 5。

命题21. 设 \(\mathrm{W}\) 为 \(\mathrm{H}\) 的一個不可約表示， \(\mathrm{E}\) 為 \(\mathrm{G}\) 的一個不可約表示。則 \(\mathrm{W}\) 在 Res E 中出現的次數等於 E 在 Ind W 中出現的次數

此結果可由定理13 得出，將其應用於 \(\mathrm{W}\) 上的字符 \(\psi\) 與 \(E\) 上的字符 \(\varphi\) (亦可應用公式(*))。

\section*{練習}

7.1.(誘導表示概念的推廣。)設 \(\alpha  : \mathrm{H} \rightarrow  \mathrm{G}\) 為一群的同構(未必單射)，並令 \(\widetilde{\alpha } : \mathbf{C}\left\lbrack  \mathrm{H}\right\rbrack \; \rightarrow  \mathbf{C}\left\lbrack  \mathrm{G}\right\rbrack\) 為相應的代數同態。若 \(\mathrm{E}\) 是一個 \(\mathbf{C}\left\lbrack  \mathrm{G}\right\rbrack\) -模，我們用 \({\operatorname{Res}}_{\alpha }\mathrm{E}\) 表示通過 \(\widetilde{\alpha }\) 從 \(\mathrm{E}\) 得到的 \(\mathrm{C}\left\lbrack  \mathrm{H}\right\rbrack\) -模；若 \(\varphi\) 為 \(E\) 的字符，則 \({\operatorname{Res}}_{\alpha }E\) 的字符為 \({\operatorname{Res}}_{\alpha }\varphi  = \varphi  \circ  \alpha\) 。若 \(W\) 是一個 \(C\left\lbrack  H\right\rbrack\) -模，我們用 \({\operatorname{Ind}}_{\alpha }\mathbf{W}\) 表示 \(\mathbf{C}\left\lbrack  \mathrm{G}\right\rbrack\) -模 \(\mathbf{C}\left\lbrack  \mathrm{G}\right\rbrack  { \otimes  }_{\mathbf{C}\left\lbrack  \mathrm{H}\right\rbrack  }\mathbf{W}\) ，若 \(\psi\) 為 \(\mathbf{W}\) 的字符，則用 \({\operatorname{Ind}}_{\alpha }\psi\) 表示 \({\operatorname{Ind}}_{\alpha }\mathbf{W}\) 的字符。

(a) 證明我們仍然有互反公式

\[
{\left\langle  \psi ,{\operatorname{Res}}_{\alpha }\varphi \right\rangle  }_{\mathrm{H}} = {\left\langle  {\operatorname{Ind}}_{\alpha }\psi ,\varphi \right\rangle  }_{\mathrm{G}}.
\]

(b) 假定 \(\alpha\) 為滿射並將 \(G\) 識別為 \(H\) 被 \(\alpha\) 的核 \(\mathrm{N}\) 所商的商群。證明 \({\operatorname{Ind}}_{\alpha }\mathrm{W}\) 同構於令 \(G = H/N\) 在 \(W\) 中作用於由 \(\mathrm{N}\) 不變元素構成的子空間而得到的模。由此推導出該公式

\[
\left( {{\operatorname{Ind}}_{\alpha }\psi }\right) \left( s\right)  = \frac{1}{n}\mathop{\sum }\limits_{{\alpha \left( t\right)  = s}}\psi \left( t\right) \;\text{ where }n = \operatorname{Card}\left( \mathrm{\;N}\right)
\]

7.2. 設 \(\mathrm{H}\) 為 \(\mathrm{G}\) 的一個子群，且令 \(\chi\) 為與 \(\mathrm{G}/\mathrm{H}\) 相關的置換表示的字符(參見1.2)。證明 \(\chi  = {\operatorname{Ind}}_{\mathrm{H}}^{\mathrm{G}}\left( 1\right)\) ，且 \(\psi  = \chi  - 1\) 為 \(G\) 一個表示的字符；確定在何種條件下該表示為不可約[使用練習2.6，或應用互反公式]。

7.3. 设 \(\mathrm{H}\) 为 \(\mathrm{G}\) 的子群。假设对每个 \(t \notin  \mathrm{H}\) 都有 \(\mathbf{H} \cap  t\mathbf{H}{t}^{-1} = \{ 1\}\) ，在这种情形下 \(\mathbf{H}\) 称为 \(\mathbf{G}\) 的 Frobenius 子群。记 \(\mathbf{N}\) 为 \(\mathbf{G}\) 中不与 \(\mathbf{H}\) 的任何元素共轭的元素集。

(a) 设 \(g = \operatorname{Card}\left( \mathrm{G}\right)\) 并设 \(h = \operatorname{Card}\left( \mathrm{H}\right)\) 。证明 \(\mathrm{N}\) 的元素个数为 \(\left( {g/h}\right)  - 1\) 。

(b) 设 \(f\) 为 \(\mathrm{H}\) 上的类函数。证明存在唯一的类函数 \(\widetilde{f}\) 在 \(\mathrm{G}\) 上延拓 \(f\) 且在 \(\mathrm{N}\) 上取值为 \(f\left( 1\right)\) 。

(c) 证明 \(\widetilde{f} = {\operatorname{Ind}}_{\mathrm{H}}^{\mathbf{G}}f - f\left( 1\right) \psi\) ，其中 \(\psi\) 是 G 的特征 \({\operatorname{Ind}}_{\mathrm{H}}^{\mathbf{G}}\left( 1\right)  - 1\) ，参见习题 7.2。

(d) 证明 \({\left\langle  {f}_{1},{f}_{2}\right\rangle  }_{\mathrm{H}} = {\left\langle  {\widetilde{f}}_{1},{\widetilde{f}}_{2}\right\rangle  }_{\mathrm{G}}\) 。

(e) 取 \(f\) 为 H 的一个不可约特征。利用 (c) 和 (d) 证明 \(\langle \widetilde{f},\widetilde{f}{\rangle }_{\mathbf{G}} = 1,\widetilde{f}\left( 1\right)  \geq  0\) ，且 \(\widetilde{f}\) 是 G 的不可约特征的整系数线性组合。由此得出 \(\widetilde{f}\) 为 G 的不可约特征。若 \(\rho\) 是对应的 G 表示，证明对每个 \(s \in  \mathrm{N}\) 有 \(\rho \left( s\right)  = 1\) [用习题 6.7]。

(f) 证明每个 \(\mathrm{H}\) 的线性表示都可扩展为一个其核含有 \(\mathrm{N}\) 的 \(\mathrm{G}\) 的线性表示。由此得出 \(\mathrm{N} \cup  \{ 1\}\) 是 \(G\) 的正规子群，且 \(G\) 是 \(H\) 与 \(N \cup  \{ 1\}\) 的半直积(Frobenius 定理)。

(g) 反过来，设 G 是 H 与正规子群 A 的半直积。证明 \(H\) 当且仅当对每个 \(s \in  H - \{ 1\}\) 和每个 \(t \in  A - \{ 1\}\) 都有 \({st}{s}^{-1} \neq  t\) (即 \(H\) 在 \(A - \{ 1\} )\) 上自由作用)时为 \(G\) 的 Frobenius 子群。(若 \(H \neq  \{ 1\}\) ，此性质由 Thompson 定理蕴含 \(A\) 为幂零群。)

\section*{7.3 限制到子群}

设 \(\mathrm{H}\) 与 \(\mathrm{K}\) 为 \(\mathrm{G}\) 的两个子群， \(\rho  : \mathrm{H} \rightarrow  \mathrm{{GL}}\left( \mathrm{W}\right)\) 为 \(\mathrm{H}\) 的一个线性表示， \(\mathrm{V} = {\operatorname{Ind}}_{\mathrm{H}}^{\mathrm{G}}\left( \mathrm{V}\right)\) 为对应的诱导表示(由 \(\mathrm{G}\) 得到)。我们将确定将 \(\mathrm{V}\) 限制到 \(\mathrm{K}\) 的表示 \({\operatorname{Res}}_{K}\mathrm{V}\) 。

首先选择一组代表 \(S\) 来表示 \(\left( {H,K}\right)\) 的双陪集的 \(G\) ；这意味着 \(G\) 是不重叠的并且由 \({KsH}\) (对 \(s \in  S\) )组成(我们也可以写成 \(s \in  \mathrm{K} \smallsetminus  \mathrm{G}/\mathrm{H}\) )。对于 \(s \in  \mathrm{S}\) ，令 \({\mathrm{H}}_{s} = s\mathrm{H}{s}^{-1} \cap  \mathrm{K}\) ，它是 \(K\) 的一个子群。如果我们设

\[
{\rho }^{s}\left( x\right)  = \rho \left( {{s}^{-1}{xs}}\right) ,\;\text{ for }x \in  {\mathrm{H}}_{s},
\]

我们得到一个同态 \({\rho }_{s} : {\mathrm{H}}_{s} \rightarrow  \mathbf{{GL}}\left( \mathrm{W}\right)\) ，从而得到 \({\mathrm{H}}_{s}\) 的一个线性表示，记为 \({\mathrm{W}}_{s}\) 。由于 \({\mathrm{H}}_{s}\) 是 \(\mathrm{K}\) 的子群，诱导表示 \({\operatorname{Ind}}_{{\mathrm{H}}_{s}}^{\mathrm{K}}\left( {\mathrm{W}}_{s}\right)\) 是有定义的。

命题 22. 表示 \({\operatorname{Res}}_{\mathrm{K}}{\operatorname{Ind}}_{\mathrm{H}}^{\mathrm{G}}\left( \mathrm{W}\right)\) 同构于各 \({\operatorname{Ind}}_{{\mathrm{H}}_{s}}^{\mathrm{K}}\left( {\mathrm{W}}_{s}\right)\) (对所有 \(s \in  \mathrm{S} \simeq  \mathrm{K} \smallsetminus  \mathrm{G}/\mathrm{H}\) )的表示的直和。

我们知道 \(\mathbf{V}\) 是各个像 \(x\mathbf{W}\) 的直和， 对于 \(x \in  \mathbf{G}/\mathbf{H}\) 。 令 \(s \in  S\) 并设 \(V\left( s\right)\) 为由 \(V\) 的各个像 \({xW}\) (对于 \(x \in  \mathrm{K}s\mathrm{H}\) )生成的子空间；空间 \(\mathrm{V}\) 是这些 \(\mathrm{V}\left( s\right)\) 的直和，并且显然 \(\mathrm{V}\left( s\right)\) 在 \(\mathrm{K}\) 下不变。 还须证明 \(\mathrm{V}\left( s\right)\) 与 \({\operatorname{Ind}}_{{H}_{r}}^{K}\left( {W}_{s}\right)\) 在 \(\mathrm{K}\) 下同构。 但是由 \(K\) 构成的那些满足 \(x\) 使得 \(x\left( {s\mathrm{\;W}}\right)  = s\mathrm{\;W}\) 的子群显然等于 \({\mathrm{H}}_{s}\) ，且 \(\mathrm{V}\left( s\right)\) 是各个像 \(x\left( {s\mathrm{\;W}}\right) ,x \in  \mathrm{K}/{\mathrm{H}}_{s}\) 的直和。因此 \(\mathrm{V}\left( s\right)  = {\operatorname{Ind}}_{{\mathrm{H}}_{s}}^{\mathrm{K}}\left( {s\mathrm{\;W}}\right)\) 。 现在还需核查 \(s\mathrm{\;W}\) 与 \({\mathrm{W}}_{s}\) 在 \({\mathrm{H}}_{s}\) 下同构，这很直接:同构由 \(s : {\mathrm{W}}_{s} \rightarrow  s\mathrm{\;W}\) 给出。

备注。由于 \(\mathrm{V}\left( s\right)\) 仅取决于 \(s\) 在 \(\mathrm{K} \smallsetminus  \mathrm{G}/\mathrm{H}\) 中的像，我们也看到表示 \({\operatorname{Ind}}_{{H}_{ * }}^{K}\left( {W}_{s}\right)\) (至同构)仅取决于 \(s\) 的双陪集。

\section*{7.4 马基的不约性判据}

我们将前述结果应用于情形 \(\mathrm{K} = \mathrm{H}\) 。 对于 \(s \in  \mathrm{G}\) ，我们仍将 \(\mathrm{H}\) 的子群 \(s\mathrm{H}{s}^{-1} \cap  \mathrm{H}\) 记作 \({\mathrm{H}}_{s}\) ；表示 \(\rho\) 对 \(\mathrm{H}\) 的限制在 \({\mathrm{H}}_{s}\) 上定义了表示 \({\operatorname{Res}}_{s}\left( \rho \right)\) ，这一点不应与 7.3 中定义的表示 \({\rho }^{s}\) 混淆。

命题 23. 为使诱导表示 \(\mathrm{V} = {\operatorname{Ind}}_{\mathrm{H}}^{\mathrm{G}}\mathrm{W}\) 不可约，当且仅当下列两条同时成立:

(a) \(\mathrm{W}\) 不可约。

(b) 对每个 \(s \in  \mathrm{G} - \mathrm{H}\) ，群 \({\mathrm{H}}_{s}\) 上的两表示 \({\rho }^{s}\) 与 \({\operatorname{Res}}_{s}\left( \rho \right)\) 互不相容。

(若群 \(K\) 的两表示 \({V}_{1}\) 与 \({V}_{2}\) 没有共同的不可约成分，即若 \({\left\langle  {V}_{1},{V}_{2}\right\rangle  }_{K} = 0\) ，则称它们互不相容。)

为了使 \(V\) 不可约，必要且充分的条件是 \(\langle \mathbf{V},\mathbf{V}{\rangle }_{\mathbf{G}} = 1\) 。但是，根据 Frobenius 对偶律，我们有:

\[
\langle \mathbf{V},\mathbf{V}{\rangle }_{\mathbf{G}} = \langle \mathbf{W},{\operatorname{Res}}_{\mathbf{H}}\mathbf{V}{\rangle }_{\mathbf{H}}.
\]

然而，由 7.3 我们有:

\[
{\operatorname{Res}}_{\mathrm{H}}\mathrm{V} = \mathop{\oplus }\limits_{{s \in  \mathrm{H} \smallsetminus  \mathrm{G}/\mathrm{H}}}{\operatorname{Ind}}_{{\mathrm{H}}_{s}}^{\mathrm{H}}\left( {\rho }^{s}\right) .
\]

再次应用 Frobenius 公式，我们得到:

\[
\langle \mathrm{V},\mathrm{V}{\rangle }_{\mathrm{G}} = \mathop{\sum }\limits_{{s \in  \mathrm{H} \smallsetminus  \mathrm{G}/\mathrm{H}}}{d}_{s},\;\text{ with }{d}_{s} = {\left\langle  {\operatorname{Res}}_{s}\left( \rho \right) ,{\rho }^{s}\right\rangle  }_{{\mathrm{H}}_{s}}.
\]

对于 \(s = 1\) 我们有 \({d}_{s} = \langle \rho ,\rho \rangle  \geq  1\) 。为了使 \(\langle \mathbf{V},\mathbf{V}{\rangle }_{\mathbf{G}} = 1\) ，因此对 \(s \neq  1\) 而言，必要且充分的条件是 \({d}_{1} = 1\) 和 \({d}_{s} = 0\) ；这正是条件 (a) 和 (b)。

推论。设 \(\mathrm{H}\) 在 \(\mathrm{G}\) 中为正规子群。为了使 \({\operatorname{Ind}}_{\mathrm{H}}^{\mathrm{G}}\left( \rho \right)\) 不可约，必要且充分的条件是 \(\rho\) 不可约且不与其任何共轭 \({\rho }^{s}\) 对于 \(s \notin  \mathrm{H}\) 同构。

事实上，我们于是有 \({\mathrm{H}}_{s} = \mathrm{H}\) 和 \({\operatorname{Res}}_{s}\left( \rho \right)  = \rho\) 。

练习

7.4. 设 \(k\) 为有限域，设 \(G = {\mathbf{{SL}}}_{2}\left( k\right)\) 并令 \(H\) 为 \(G\) 的子群，包含那些矩阵 \(\left( \begin{array}{ll} a & b \\  c & d \end{array}\right)\) 使得 \(c = 0\) 。令 \(\omega\) 为从 \({k}^{ * }\) 到 \({\mathbf{C}}^{ * }\) 的同态，令 \({\chi }_{\omega }\) 为由下式定义的 \(\mathbf{H}\) 的一次特征:

\[
{\chi }_{\omega }\left( \begin{array}{ll} a & b \\  c & d \end{array}\right)  = \omega \left( a\right) .
\]

证明:若 \({\omega }^{2} \neq  1\) 则由 \({\chi }_{\omega }\) 诱导的 \(G\) 表示是不可约的。

\section*{第8章 诱导表示的例子}

8.1 正规子群；不可约表示次数的应用

命题 24。设 A 为群 G 的正规子群， \(\rho  : \mathrm{G} \rightarrow  \mathrm{{GL}}\left( \mathrm{V}\right)\) 为 \(\mathrm{G}\) 的一个不可约表示。则:

(a) 或者存在一个子群 \(\mathrm{H}\) ，它包含 A 且不等于 \(\mathrm{G}\) ，以及 \(\mathrm{H}\) 的一个不可约表示 \(\sigma\) ，使得 \(\rho\) 由 \(\sigma\) 诱导而来

(b) 或者 \(\rho\) 限制到 \(\mathrm{A}\) 是同型类的。

(当一个表示是同型类的时，意指它是若干同构不可约表示的直和。)

令 \(\mathbf{V} =  \oplus  {\mathbf{V}}_{i}\) 为表示 \(\rho\) (限制到 A)分解为同型类表示直和的典范分解(参见 2.6)。对于 \(s \in  \mathrm{G}\) ，通过“结构搬运”可见 \(\rho \left( s\right)\) 置换这些 \({\mathbf{V}}_{i}\) ；由于 \(\mathbf{V}\) 是不可约的，它对它们做传递置换。取其中一个 \({V}_{{i}_{0}}\) ；若 \({V}_{{i}_{0}}\) 等于 \(V\) ，则为情形 (b)。否则，令 \(H\) 为 \(G\) 的子群，由那些满足 \(s \in  G\) 的 \(\rho \left( s\right) {V}_{{i}_{0}} = {V}_{{i}_{0}}\) 构成。我们有 \(A \subset  H\) ， \(\mathrm{H} \neq  \mathrm{G}\) ，且 \(\rho\) 由 \(\mathrm{H}\) 在 \({\mathrm{V}}_{{i}_{0}}\) 中的自然表示 \(\sigma\) 诱导得到，即情形 (a)。

注。若 \(\mathbf{A}\) 是阿贝尔群，则 (b) 等价于对每个 \(a \in  \mathbf{A}\) ， \(\rho \left( \mathbf{A}\right)\) 为一个齐次缩放变换(homothety)。

推论。若 \(\mathbf{A}\) 为 \(\mathbf{G}\) 的阿贝尔正规子群，则 \(\mathbf{G}\) 的每个不可约表示 \(\rho\) 的次数整除 \(\mathbf{A}\) 在 \(\mathbf{G}\) 中的指数 (G: A)。

证明按 |G| 作归纳。在前一命题的情形 (a) 中，归纳假设表明 \(\sigma\) 的次数整除 (H: A)，将此关系乘以 (G: H) 得到 \(\rho\) 的次数整除 (G: A)。在情形 (b) 中，取 \({G}^{\prime } = \rho \left( G\right)\) 与 \({A}^{\prime } = \rho \left( A\right)\) ；由于典范映射 \(G/A \rightarrow  {G}^{\prime }/{A}^{\prime }\) 是满的， \(\left( {{G}^{\prime } : {A}^{\prime }}\right)\) 整除 \(\left( {G : A}\right)\) 。我们前述的备注现表明 \({A}^{\prime }\) 的元素是齐次缩放变换，从而包含于 \({G}^{\prime }\) 的中心。由 6.5 的命题 17，可得 \(\rho\) 的次数整除 \(\left( {{G}^{\prime } : {A}^{\prime }}\right)\) ，从而亦整除 (G: A)。

注。若 A 是 G 的一个阿贝尔子群(不必正规)，则一般不再成立 \(\deg \left( \rho \right)\) 整除 (G: A)，但我们仍有 \(\deg \left( \rho \right)  \leq\) (G: A)，参见 3.1，定理 9 的推论。

\section*{8.2 由阿贝尔群作半直积}

设 \(A\) 和 \(H\) 是群 \(G\) 的两个子群，且 \(A\) 为正规子群。作出以下假设:

(i) A 是阿贝尔群。

(ii) \(\mathrm{G}\) 是 \(\mathrm{H}\) 被 \(\mathrm{A}\) 半直积。

[回忆 (ii) 意味着 \(G = A \cdot  H\) 与 \(A \cap  H = \{ 1\}\) ，换言之， \(G\) 的每个元素可唯一表示为乘积 \({ah}\) ，其中 \(a \in  \mathbf{A}\) 与 \(h \in  \mathbf{H}\) 。]

我们将证明 \(\mathbf{G}\) 的不可约表示可由 \(H\) 的某些子群的表示构造(这是 Wigner 与 Mackey 的“小群”方法)。

由于 \(\mathbf{A}\) 是阿贝尔的，其不可约特征为一维并形成一群 \(X = \operatorname{Hom}\left( {A,{C}^{ * }}\right)\) 。群 \(G\) 通过下列方式作用于 \(X\) :

\[
\left( {s\chi }\right) \left( a\right)  = \chi \left( {{s}^{-1}{as}}\right) \;\text{ for }s \in  \mathrm{G},\chi  \in  \mathrm{X},a \in  \mathrm{A}.
\]

设 \({\left( {\chi }_{i}\right) }_{i \in  \mathrm{X}/\mathrm{H}}\) 为 \(\mathrm{H}\) 在 \(\mathrm{X}\) 中轨道的一组代表。对每个 \(i \in  \mathrm{X}/\mathrm{H}\) ，令 \({\mathrm{H}}_{i}\) 为 \(\mathrm{H}\) 的子群，包含那些满足 \(h{\chi }_{i} = {\chi }_{i}\) 的元素 \(h\) ，并令 \({\mathrm{G}}_{i} = \mathrm{A} \cdot  {\mathrm{H}}_{i}\) 为 G 中相应的子群。将函数 \({\chi }_{i}\) 通过规定扩展到 \({\mathrm{G}}_{i}\) ，如下

\[
{\chi }_{i}\left( {ah}\right)  = {\chi }_{i}\left( a\right) \;\text{ for }a \in  \mathrm{A},h \in  {\mathrm{H}}_{i}.
\]

利用 \(h{\chi }_{i} = {\chi }_{i}\) 对所有 \(h \in  {\mathrm{H}}_{i}\) 都成立这一事实，我们看到 \({\chi }_{i}\) 是 \({\mathrm{G}}_{i}\) 的一维特征。现在令 \(\rho\) 为 \({\mathrm{H}}_{i}\) 的一个不可约表示；通过与典型投影 \({\mathrm{G}}_{i} \rightarrow  {\mathrm{H}}_{i}\) 复合，我们得到 \({\mathrm{G}}_{i}\) 的一个不可约表示 \(\widetilde{\rho }\) 。最后，通过对 \({\chi }_{i}\) 与 \(\widetilde{\rho }\) 张量积，我们得到 \({\mathrm{G}}_{i}\) 的一个不可约表示 \({\chi }_{i} \otimes  \widetilde{\rho }\) ；令 \({\theta }_{i,p}\) 为对应的对 \(\mathrm{G}\) 的诱导表示。

\section*{命题 25}

(a) \({\theta }_{i,\rho }\) 是不可约的。

(b) 若 \({\theta }_{i,\rho }\) 与 \({\theta }_{{i}^{\prime },{\rho }^{\prime }}\) 同构，则 \(i = {i}^{\prime }\) 与 \(\rho\) 同构于 \({\rho }^{\prime }\) 。

(c) 每个 \(\mathrm{G}\) 的不可约表示都同构于这些 \({\theta }_{i,\rho }\) 中的某一个。

(因此我们得到了 G 的所有不可约表示。)

我们用 Mackey 判据 (7.4，命题 23) 证明 (a):设 \(s \notin  {\mathrm{G}}_{i} = \mathrm{A} \cdot  {\mathrm{H}}_{i}\) ，且设 \({\mathrm{K}}_{s} = {\mathrm{G}}_{i} \cap  s{\mathrm{G}}_{i}{s}^{-1}\) 。需证明将表示 \({\chi }_{i} \otimes  \widetilde{\rho }\) 的 \({\mathrm{G}}_{i}\) 与由 \(x \mapsto  x\) 和 \(x \mapsto  {s}^{-1}{xs}\) 定义的两个注入 \({\mathrm{K}}_{s} \rightarrow  {\mathrm{G}}_{i}\) 复合后，会得到 \({\mathrm{K}}_{s}\) 的两个互不相交的表示。为此，只需检验这些表示限制到 \({\mathrm{K}}_{s}\) 的子群 \(A\) 时是互不相交的。但第一个限制为 \({\chi }_{i}\) 的若干倍，第二个为 \(s{\chi }_{i}\) 的若干倍；由于 \(s \notin  A \cdot  {\mathrm{H}}_{i}\) 我们有 \(s{\chi }_{i} \neq  {\chi }_{i}\) ，因此所述两个表示确实互不相交。

下面证明 (b)。首先， \({\theta }_{i,\rho }\) 限制到 \(\mathbf{A}\) 时只涉及属于 \({\chi }_{i}\) 的轨道 \(\mathrm{H}{\chi }_{i}\) 中的特征 \(\chi\) 。这表明 \({\theta }_{i,\rho }\) 决定 \(i\) 。接着，设 \(\mathbf{W}\) 为 \({\theta }_{i,\rho }\) 的表示空间，令 \({\mathbf{W}}_{i}\) 为对应于 \({\chi }_{i}\) 的子空间 [即满足 对所有 \(a \in  A\) 都有 \({\theta }_{i,\rho }\left( a\right) x = {\chi }_{i}\left( a\right) x\) 的 \(x \in  \mathbf{W}\) 的集合]。子空间 \({W}_{i}\) 在 \({H}_{i}\) 下不变，且立刻可检验 \({H}_{i}\) 在 \({W}_{i}\) 中的表示与 \(\rho\) 同构；因此 \({\theta }_{i,\rho }\) 决定 \(\rho\) 。

最后，设 \(\sigma  : \mathrm{G} \rightarrow  \mathbf{{GL}}\left( \mathrm{W}\right)\) 为 \(\mathrm{G}\) 的一个不可约表示。设 \(W = { \oplus  }_{x \in  X}W\) 为 \({\operatorname{Res}}_{A}W\) 的典范分解。至少有一个 \({W}_{\chi }\) 非零；若 \(s \in  G,\sigma \left( s\right)\) 将 \({W}_{\chi }\) 变为 \({W}_{s\left( \chi \right) }\) 。群 \({H}_{i}\) 将 \({\mathbf{W}}_{{\chi }_{i}}\) 映入自身；取 \({\mathbf{W}}_{i}\) 为 \({\mathbf{W}}_{{\chi }_{i}}\) 的一个不可约子- \(\mathbf{C}\left\lbrack  {\mathrm{H}}_{i}\right\rbrack\) -模，令 \(\rho\) 为对应的 \({\mathrm{H}}_{i}\) 的表示。显然 \({\mathrm{G}}_{i} = \mathrm{A} \cdot  {\mathrm{H}}_{i}\) 的该表示与 \({\chi }_{i} \otimes  \widetilde{\rho }\) 同构。于是 \(\sigma\) 限制到 \({\mathrm{G}}_{i}\) 至少包含一次 \({\chi }_{i} \otimes  \widetilde{\rho }\) 。由命题21，得 \(\sigma\) 在诱导表示 \({\theta }_{i,\rho }\) 中至少出现一次；因 \({\theta }_{i,\rho }\) 不可约，此即表明 \(\sigma\) 与 \({\theta }_{i,\rho }\) 同构，从而证明了 (c)。

\section*{练习}

8.1. 设 \(a,h,{h}_{i}\) 分别为 \(\mathrm{A},\mathrm{H},{\mathrm{H}}_{i}\) 的阶。证明 \(a = \sum \left( {h/{h}_{i}}\right)\) 。证明:对固定的 \(i\) ，表示 \({\theta }_{i,\rho }\) 的维数平方和为 \({h}^{2}/{h}_{i}\) 。由此从另一角度推出 (c) 的结论。

8.2. 用命题25 重新计算群 \({D}_{n}\) 、 \({\mathfrak{A}}_{4}\) 和 \({\mathfrak{S}}_{4}\) 的不可约表示(参见第5章)。

\section*{8.3 一些有限群类的回顾}

关于本节及下节结果的更多细节，请参见 Bourbaki，Alg. I，第7节。

可解群。若存在一列序列使得 \(\mathbf{G}\) 可解，则称 \(\mathbf{G}\) 为可解群

\[
\{ 1\}  = {G}_{0} \subset  {G}_{1} \subset  \cdots  \subset  {G}_{n} = G
\]

对 \(G\) 的子群列，其中 \({G}_{i - 1}\) 在 \({G}_{i}\) 中为正规子群且 \({G}_{i}/{G}_{i - 1}\) 为阿贝尔群，

则称 \(1 \leq  i \leq  n\) 可解。(等价定义:G 由平凡群 \{1\} 经有限次以阿贝尔核的扩张得到。)

超可解群。与上述相同，但要求所有 \({\mathrm{G}}_{i}\) 在 \(\mathrm{G}\) 中正规且 \({\mathrm{G}}_{i}/{\mathrm{G}}_{i - 1}\) 为循环群。

幂零群。如上，只是要求 \({\mathrm{G}}_{i}/{\mathrm{G}}_{i - 1}\) 在 \(\mathrm{G}/{\mathrm{G}}_{i - 1}\) 的中心以便于 \(1 \leq  i \leq  n\) 。 (等价定义: \(\mathrm{G}\) 由群 \(\{ 1\}\) 经有限次中心扩张得到。)

显然超可解的 \(\Rightarrow\) 是可解的。另一方面，立刻可检验每个超可解群的中心扩张仍为超可解；因此幂零的 \(\Rightarrow\) 是超可解的。

\(p\) -群。如果 \(p\) 是素数，阶为 \(p\) 的幂的群称为 \(p\) -群。

定理 14. 每个 p-群都是幂零的(因此超可解)。

鉴于前述，只需证明每个非平凡的 \(p\) -群 \(\mathbf{G}\) 的中心非平凡。下面的引理可得此结论:

引理 3. 设 \(\mathrm{G}\) 为在有限集合 \(\mathrm{X}\) 上作用的 p-群，令 \({\mathrm{X}}^{\mathrm{G}}\) 为被 \(\mathrm{G}\) 固定的 \(\mathrm{X}\) 的元素集合。我们有

\[
\operatorname{Card}\left( \mathrm{X}\right)  \equiv  \operatorname{Card}\left( {\mathrm{X}}^{\mathrm{G}}\right) \;\left( {\;\operatorname{mod}\;p}\right) .
\]

确实， \(X - {X}^{G}\) 是 \(G\) 的若干非平凡轨道的并，每个轨道的基数是 \(p\) 的某次幂 \({p}^{\alpha }\) ，且 \(\alpha  \geq  1\) ；因此 \(\operatorname{Card}\left( {X - {X}^{G}}\right)\) 可被 \(p\) 整除。

现在把该引理应用到 \(X = G\) 的情形，且 \(G\) 通过内自同构作用。集合 \({X}^{G}\) 恰为 \(G\) 的中心 \(C\) 。因此

\[
\operatorname{Card}\left( \mathrm{C}\right)  \equiv  \operatorname{Card}\left( \mathrm{G}\right)  \equiv  0\left( {\;\operatorname{mod}\;.p}\right) ,
\]

从而 \(C \neq  \{ 1\}\) ，这就证明了该定理。

我们记录引理 3 的另一应用，这将在第三部分中使用:

命题 26. 设 \(\mathrm{V}\) 是域 \(k\) (特征为 \(p\) )上的向量空间 \(\neq  0\) ，且 \(\rho  : \mathrm{G} \rightarrow  \mathbf{{GL}}\left( \mathrm{V}\right)\) 是 \(p\) -群 \(\mathrm{G}\) 在 \(\mathrm{V}\) 上的线性表示。则存在被所有 \(\rho \left( s\right) ,s \in  \mathrm{G}\) 固定的非零元素。

取 \(\mathrm{V}\) 中的非零元素 \(x\) ，令由 \(\rho \left( s\right) x,s \in  \mathrm{G}\) 生成的 \(\mathrm{V}\) 的子群为 \(\mathrm{X}\) 。对 \(\mathrm{X}\) 应用引理 3，注意到 \(\mathrm{X}\) 有限且阶为 \(p\) 的幂。因此 \({\mathrm{X}}^{\mathrm{G}} \neq  \{ 0\}\) ，这证明了该命题。

推论. 在特征为 \(p\) 的情形下，p-群的唯一不可约表示是平凡表示。

\section*{习题}

8.3. 证明二面体群 \({D}_{n}\) 是超可解的，并且当且仅当 \(n\) 是 2 的幂时它是幂零的。

8.4. 证明交错群 \({\mathfrak{A}}_{4}\) 是可解的，但非超可解。对群 \({\mathfrak{S}}_{4}\) 也问同样的问题。

8.5. 证明可解群(分别超可解、幂零群)的每个子群和每个商群仍然是可解(分别超可解、幂零)。

8.6. 设 \(p\) 与 \(q\) 为互异素数，且令 \(\mathrm{G}\) 为阶为 \({p}^{a}{q}^{b}\) 的群，其中 \(a\) 与 \(b\) 为整数 \(> 0\) 。

(i) 假设 \(\mathrm{G}\) 的中心是 \(\{ 1\}\) 。对 \(s \in  \mathbf{G}\) ，记 \(c\left( s\right)\) 为 \(s\) 的共轭类中的元素个数。证明存在 \(s \neq  1\) 使得 \(c\left( s\right)  \text{ ≢ } 0\left( {{\;\operatorname{mod}\;.}q}\right)\) 。(否则 \(G - \{ 1\}\) 的元素个数将被 \(q\) 整除。) 对这样的 \(s,c\left( s\right)\) 是 \(p\) 的幂；由此推得存在一个既不等于 \(\{ 1\}\) 也不等于 \(G\) 的正规子群 [参见练习 6.10]。

(ii) 证明 \(\mathbf{G}\) 是可解的(Burnside 定理)。[对 \(\mathbf{G}\) 的阶做归纳，并区分两种情况，取决于 \(\mathbf{G}\) 的中心是否等于 \{1\}。]

(iii) 用例子说明 \(\mathbf{G}\) 不必定是超可解的(参见练习 8.4)。

(iv) 给出一个非可解群的例子，其阶仅被三个素数 \(\left\lbrack  {{\mathfrak{S}}_{5},{\mathfrak{S}}_{6},{\mathbf{{GL}}}_{2}\left( {\mathbf{F}}_{7}\right) \text{ will do }}\right\rbrack\) 整除。

\section*{8.4 Sylow 定理}

设 \(p\) 为素数，且令 \(\mathbf{G}\) 为阶为 \(g = {p}^{n}m\) 的群，其中 \(m\) 与 \(p\) 互素。 \(\mathbf{G}\) 中阶为 \({p}^{n}\) 的子群称为 G 的 Sylow \(p\) -子群。

\section*{定理 15}

(a) 存在 Sylow p-子群。

(b) 它们在内自同构下是共轭的。

(c) 每个 \(\mathbf{G}\) -子群包含于某个 Sylow \(\mathbf{G}\) -子群中。

为证明 (a) 我们对 \(\mathbf{G}\) 的阶作归纳。我们可假定 \(n \geq  1\) ，即 Card (G) \(\equiv  0\left( {{\;\operatorname{mod}\;.}p}\right)\) 。设 \(\mathrm{C}\) 为 \(\mathrm{G}\) 的中心。若 Card (C) 的阶可被 \(p\) 整除，则一个基本论证表明 \(\mathrm{C}\) 含有一循环子群 D，其阶为 \(p\) 。由归纳假设，G/D 存在一个 Sylow \(p\) -子群，而该子群在 \(G\) 中的逆像就是 G 的一个 Sylow \(p\) -子群。若 Card (C) \(\text{ ≢ } 0\left( {{\;\operatorname{mod}\;.}p}\right)\) ，令 G 通过内部自同构作用于 G - C；这将把 \(G - C\) 划分为若干轨道(共轭类)。由于 \(\operatorname{Card}\left( {G - C}\right)  \text{ ≢ } 0\left( {{\;\operatorname{mod}\;.}p}\right)\) ，其中一个轨道的基数与 \(p\) 互素。由此可得存在一个子群 \(H\) (不同于 \(G\) )，使得 (G: H) \(\text{ ≢ } 0\left( {{\;\operatorname{mod}\;.}p}\right)\) 。因此 \(\mathrm{H}\) 的阶可被 \({p}^{n}\) 整除，归纳假设说明 \(\mathrm{H}\) 含有阶为 \({p}^{n}\) 的子群。

现在令 P 为 G 的一个 Sylow \(p\) -子群，Q 为 G 的一个 \(p\) -子群。 \(p\) -群 \(\mathrm{Q}\) 通过左平移作用于 \(\mathrm{X} = \mathrm{G}/\mathrm{P}\) 。由 8.3 引理 3 我们有

\[
\operatorname{Card}\left( {\mathrm{X}}^{\mathrm{Q}}\right)  \equiv  \operatorname{Card}\left( \mathrm{X}\right)  \text{ ≢ } 0\left( {\;\operatorname{mod}\;.p}\right) ,
\]

因此 \({X}^{Q} \neq  \varnothing\) 。于是存在元素 \(x \in  G\) 使得 \({QxP} \; = x\mathrm{P}\) ，因此 \(\mathrm{Q} \subset  x\mathrm{P}{x}^{-1}\) ，这证明了 (c)。如果另外 \(\operatorname{Card}\left( \mathrm{Q}\right)  = {p}^{n}\) ，则群 \(\mathbf{Q}\) 和 \(x\mathbf{P}{x}^{-1}\) 阶相同，且 \(\mathbf{Q} = x\mathbf{P}{x}^{-1}\) ，这证明了 (b)。

练习

8.7. 设 \(\mathrm{H}\) 为群 \(\mathrm{G}\) 的正规子群，且令 \({\mathrm{P}}_{\mathrm{H}}\) 为 \(\mathrm{G}/\mathrm{H}\) 的一个 Sylow \(p\) -子群。

(a) 证明存在一个 Sylow \(p\) -子群 \(\mathrm{P}\) 于 \(\mathrm{G}\) ，其在 \(\mathrm{G}/\mathrm{H}\) 中的像为 \({P}_{H}\) [利用 Sylow 子群的共轭性]。

(b) 若 \(H\) 为 \(p\) -群或 \(H\) 位于 \(G\) 的中心，则证明 \(P\) 唯一 [化约到 \(H\) 的阶与 \(p\) 互素的情形，并利用任一从 \({P}_{H}\) 到 \(H\) 的同态都是平凡的这一事实]。

8.8. 设 \(G\) 为幂零群。证明对每一素数 \(p,G\) ， \(G\) 包含唯一的 Sylow \(p\) -子群，且该子群为正规子群 [对 G 的阶进行归纳，并将归纳假设应用于 G 除以其中心的商，参见练习 8.7(b)]。由此得出 \(\mathbf{G}\) 是若干 \(p\) -群的直积。

8.9. 设 \(\mathbf{G} = {\mathbf{{GL}}}_{n}\left( k\right)\) ，其中 \(k\) 为特征为 \(p\) 的有限域。证明由对角线上全为 1 的上三角矩阵组成的 \(\mathbf{G}\) 的子群是 \(\mathbf{G}\) 的一个 Sylow \(p\) -子群。

8.5 超可解群的线性表示

引理 4. 设 \(\mathrm{G}\) 为非阿贝尔的超可解群。则存在一个正规阿贝尔子群，且不包含于 \(\mathrm{G}\) 的中心。

令 \(\mathrm{C}\) 为 \(\mathrm{G}\) 的中心。商群 \(\mathrm{H} = \mathrm{G}/\mathrm{C}\) 是超可解的，因此有一条组成列，其第一个非平凡项 \({\mathrm{H}}_{1}\) 是 \(\mathrm{H}\) 的循环正规子群。 \({\mathrm{H}}_{1}\) 在 \(\mathrm{G}\) 中的逆像具有所需性质。

定理 16. 设 \(\mathrm{G}\) 为超可解群。则 \(\mathrm{G}\) 的每一个不可约表示都由 \(\mathrm{G}\) 的某个子群的 1 维表示诱导而来(即是单项表示)。

我们用群阶的归纳法证明该定理。因此我们仅需考虑那些忠实的不可约表示 \(\rho\) ，即满足 \(\operatorname{Ker}\left( \rho \right)  = \{ 1\}\) 的表示。如果 \(\mathbf{G}\) 是阿贝尔的，则这样的 \(\rho\) 为1维，无需证明。设 \(\mathbf{G}\) 非阿贝尔，令 \(\mathbf{A}\) 为 \(G\) 的一个正规阿贝尔子群且不包含在 \(G\) 的中心中(参见引理4)。由于 \(\rho\) 是忠实的，这意味着 \(\rho \left( A\right)\) 不包含在 \(\rho \left( \mathrm{G}\right)\) 的中心中；因此存在 \(a \in  \mathrm{A}\) 使得 \(\rho \left( a\right)\) 不是一个齐次相似变换。故 \(\rho\) 限制到 \(\mathrm{A}\) 并非同型块。由命题24，这意味着 \(\rho\) 由 \(G\) 的一个真子群 \(H\) 上的不可约表示诱导而来。对 \(H\) 应用归纳即可得出定理。

\section*{练习}

8.10. 将定理16推广到那些由一个超可解群与一个阿贝尔正规子群半直积的群〔用命题25化简到超可解情形〕。

8.11. 设 \(\mathbf{H}\) 为基于 \(\mathbf{R}\) 的四元数域，基为 \(\{ 1,i,j,k\}\) 且满足

\[
{i}^{2} = {j}^{2} = {k}^{2} =  - 1,\;{ij} =  - {ji} = k,\;{jk} =  - {kj} = i,
\]

\[
{ki} =  - {ik} = j.
\]

设 \(E\) 为 \({H}^{ * }\) 的一个子群，包含八个元素 \(\pm  1, \pm  i, \pm  j, \pm  k\) (四元数群)，令 \(G\) 为 \(E\) 与十六个元素 \(\left( {\pm 1 \pm  i \pm  j \pm  k}\right) /2\) 的并。证明 \(G\) 是 \({H}^{ * }\) 的一个可解子群，且为一个三阶循环群与正规子群E的半直积。利用同构 \(\mathbf{H}{ \otimes  }_{\mathbf{R}}\mathbf{C} = {\mathbf{M}}_{2}\left( \mathbf{C}\right)\) 来构造G的一个2维不可约表示。证明该表示非单项(注意 \(\mathbf{G}\) 没有指数2的子群)。【群 \(\mathbf{G}\) 为Hurwitz“整四元数”环的可逆元群；在特征2下亦为椭圆曲线 \({y}^{2} - y = {x}^{3}\) 的自同构群。它同构于 \({\mathbf{{SL}}}_{2}\left( {\mathbf{F}}_{3}\right)\) 。】

8.12. 设 \(G\) 为一个 \(p\) -群。证明对每个 \(G\) 的不可约特征 \(\chi\) ，有 \(\sum {\chi }^{\prime }{\left( 1\right) }^{2} \equiv  0\left( {{\;\operatorname{mod}\;.}\chi {\left( 1\right) }^{2}}\right)\) ，求和在所有满足 \({\chi }^{\prime }\left( 1\right)  < \chi \left( 1\right)\) 的不可约特征 \({\chi }^{\prime }\) 之上。〔利用 \(\chi \left( 1\right)\) 为 \(p\) 的幂这一事实，并对命题5应用推论2(a)。〕

第9章 阿廷定理

\section*{9.1 环 R(G)}

令 \(G\) 为有限群， \({\chi }_{1},\ldots ,{\chi }_{h}\) 为其互不相同的不可约特征。定义在 \(G\) 上的类函数当且仅当它是各个 \({\chi }_{i}\) 的非负整数系数线性组合时为特征。记 \({R}^{ + }\left( G\right)\) 为这些函数的集合，记 \(R\left( G\right)\) 为由 \({R}^{ + }\left( G\right)\) 生成的群，即两个特征之差的集合。我们有

\[
\mathrm{R}\left( \mathrm{G}\right)  = \mathbf{Z}{\chi }_{1} \oplus  \cdots  \oplus  \mathbf{Z}{\chi }_{h}.
\]

\(R\left( G\right)\) 的一个元素被称为虚特征标。由于两个特征标的乘积是一个特征标，所以 \(R\left( G\right)\) 是 \(G\) 上取值为复数的类函数环 \({F}_{C}\left( G\right)\) 的一个子环。由于 \({\chi }_{i}\) 构成了 \({F}_{C}\left( G\right)\) 在 \(\mathbf{C}\) 上的一组基，我们可知 \(\mathbf{C} \otimes  \mathrm{R}\left( \mathrm{G}\right)\) 可与 \({\mathrm{F}}_{\mathbf{C}}\left( \mathrm{G}\right)\) 等同。

我们也可以将 \(\mathrm{R}\left( \mathrm{G}\right)\) 视为有限生成的 C[G]-模范畴的格罗滕迪克群；这将在第三部分中使用。

若 \(\mathrm{H}\) 为 \(\mathrm{G}\) 的子群，则限制运算给出一个环同态 \(\mathrm{R}\left( \mathrm{G}\right)  \rightarrow  \mathrm{R}\left( \mathrm{H}\right)\) ，记作 \({\operatorname{Res}}_{\mathrm{H}}^{\mathrm{G}}\) 或 Res。

类似地，诱导运算 (7.2) 给出一个阿贝尔群同态 \(R\left( H\right)  \rightarrow  R\left( G\right)\) ，记作 \({\operatorname{Ind}}_{H}^{G}\) 或 Ind。Ind 与 Res 在双线性形式 \(\langle \varphi ,\psi {\rangle }_{\mathrm{H}}\) 和 \(\langle \varphi ,\psi {\rangle }_{\mathrm{G}}\) 下互为伴随，参见定理 13。此外，公式

\[
\operatorname{Ind}\left( {\varphi  \cdot  \operatorname{Res}\left( \psi \right) }\right)  = \operatorname{Ind}\left( \varphi \right)  \cdot  \psi
\]

表明 Ind 的像: \(R\left( H\right)  \rightarrow  R\left( G\right)\) 是环 \(R\left( G\right)\) 的一个理想。

若 \(\mathbf{A}\) 是交换环，Res 和 Ind 同态可按线性扩展为 A-线性映射:

\[
\mathrm{A} \otimes  \text{ Res: }\mathrm{A} \otimes  \mathrm{R}\left( \mathrm{G}\right)  \rightarrow  \mathrm{A} \otimes  \mathrm{R}\left( \mathrm{H}\right)
\]

\[
\mathrm{A} \otimes  \text{ Ind: }\mathrm{A} \otimes  \mathrm{R}\left( \mathrm{H}\right)  \rightarrow  \mathrm{A} \otimes  \mathrm{R}\left( \mathrm{G}\right)
\]

练习

9.1. 设 \(\varphi\) 为 G 上的实值类函数。假设 \(\langle \varphi ,1\rangle  = 0\) 且对每个 \(\varphi \left( s\right)  \leq  0\) 有 \(s \neq  1\) 。证明对每个特征 \(\chi\) ， \(\langle \varphi ,\chi \rangle\) 的实部为 \(\geq  0\) ［利用事实:对所有 \(\varphi \left( {s}^{-1}\right) \chi \left( s\right)\) ， \(\varphi \left( {s}^{-1}\right) \chi \left( 1\right)\) 的实部不小于 \(s\rbrack\) 的实部。得出结论:若 \(\varphi\) 属于 \(\mathbf{R}\left( \mathbf{G}\right) ,\varphi\) ，则为特征。］

9.2. 设 \(\chi  \in  \mathbf{R}\left( G\right)\) 。证明 \(\chi\) 当且仅当 \(\langle \chi ,\chi \rangle  = 1\) 与 \(\chi \left( 1\right)  \geq  0\) 为不可约特征。

9.3. 若 \(f\) 是定义在 \(G\) 上的函数，且 \(k\) 为一个整数，记 \({\Psi }^{k}\left( f\right)\) 为函数 \(s \mapsto  f\left( {s}^{k}\right)\) 。

(a) 设 \(\rho\) 是带特征 \(\chi\) 的 \(\mathbf{G}\) 的一个表示。对每个整数 \(k \geq  0\) ，记 \({\chi }_{\sigma }^{k}\) (分别 \({\chi }_{\lambda }^{k}\) )为 \(k\) 次对称幂(分别 \(k\) 次外幂) \(\rho\) 的特征(在 \(k = 2\) 情形参见 2.1)。设

\[
{\sigma }_{\mathrm{T}}\left( \chi \right)  = \mathop{\sum }\limits_{{k = 0}}^{\infty }{\chi }_{\sigma }^{k}{\mathrm{\;T}}^{k}\;\text{ and }\;{\lambda }_{\mathrm{T}}\left( \chi \right)  = \mathop{\sum }\limits_{{k = 0}}^{\infty }{\chi }_{\lambda }^{k}{\mathrm{\;T}}^{k},
\]

其中 \(\mathrm{T}\) 是一个不定元。证明:对 \(s \in  \mathrm{G}\) 有

\[
{\sigma }_{\mathrm{T}}\left( \chi \right) \left( s\right)  = 1/\det \left( {1 - \rho \left( s\right) \mathrm{T}}\right) \;\text{ and }\;{\lambda }_{\mathrm{T}}\left( \chi \right) \left( s\right)  = \det \left( {1 + \rho \left( s\right) \mathrm{T}}\right) .
\]

推出下列公式

\[
{\sigma }_{\mathrm{T}}\left( \chi \right)  = \exp \left\{  {\mathop{\sum }\limits_{{k = 1}}^{\infty }{\Psi }^{k}\left( \chi \right) {\mathrm{T}}^{k}/k}\right\}
\]

\[
{\lambda }_{\mathrm{T}}\left( \chi \right)  = \exp \left\{  {\mathop{\sum }\limits_{{k = 1}}^{\infty }{\left( -1\right) }^{k - 1}{\Psi }^{k}\left( \chi \right) {\mathrm{T}}^{k}/k}\right\}
\]

和

\[
n{\chi }_{\sigma }^{n} = \mathop{\sum }\limits_{{k = 1}}^{n}{\Psi }^{k}\left( \chi \right) {\chi }_{\sigma }^{n - k},\;n{\chi }_{\lambda }^{n} = \mathop{\sum }\limits_{{k = 1}}^{n}{\left( -1\right) }^{k - 1}{\Psi }^{k}\left( \chi \right) {\chi }_{\lambda }^{n - k},
\]

这推广了 2.1 中的那些公式。

(b) 从 a) 得出 \(\mathbf{R}\left( \mathbf{G}\right)\) 在算子 \({\Psi }^{k},k \in  \mathbf{Z}\) 下是稳定的。

9.4. 设 \(n\) 为与 \(\mathbf{G}\) 的阶互素的整数。

(a) 设 \(\chi\) 为 \(\mathbf{G}\) 的一个不可约特征。证明 \({\Psi }^{n}\left( \chi \right)\) 是 \(\mathbf{G}\) 的一个不可约特征 [利用前两题]。

(b) 将 \(x \mapsto  {x}^{n}\) 按线性性扩展为作用在向量空间 C[G] 上的一个自同态 \({\psi }_{n}\) 。证明 \({\psi }_{n}\) 限制到 Cent. C[G] 后是代数 Cent. C[G] 的一个自同构。

\section*{9.2 阿廷定理的陈述}

内容如下:

定理 17。设 \(\mathrm{X}\) 为有限群 G 的一族子群。令 Ind: \({ \oplus  }_{\mathrm{H} \in  \mathrm{X}}\mathrm{R}\left( \mathrm{H}\right)  \rightarrow  \mathrm{R}\left( \mathrm{G}\right)\) 为由该族 \({\operatorname{Ind}}_{\mathbf{H}}^{\mathbf{G}},\mathbf{H} \in  \mathbf{X}\) 定义的同态。则下列性质等价:

(i) \(\mathrm{G}\) 是属于 \(\mathrm{X}\) 的子群的共轭类的并。

(ii) Ind: \({\bigoplus }_{\mathrm{H} \in  \mathrm{X}}\mathrm{R}\left( \mathrm{H}\right)  \rightarrow  \mathrm{R}\left( \mathrm{G}\right)\) 的余核是有限的。

由于 \(\mathrm{R}\left( \mathrm{G}\right)\) 作为群是有限生成的，我们可以将 (ii) 重新表述为:

(ii’) 对于每个 \(\mathrm{G}\) 的特征 \(\chi\) ，存在虚特征 \({\chi }_{\mathrm{H}} \in  \mathrm{R}\left( \mathrm{H}\right) ,\mathrm{H} \in  \mathrm{X}\) 以及一个整数 \(d \geq  1\) ，使得

\[
{d\chi } = \mathop{\sum }\limits_{{\mathrm{H} \in  \mathrm{X}}}{\operatorname{Ind}}_{\mathrm{H}}^{\mathrm{G}}\left( {\chi }_{\mathrm{H}}\right)
\]

注意，G 的循环子群族满足 (i)。因此:

推论。 \(\mathrm{G}\) 的每个特征都是由由循环子群的特征诱导而来的特征以有理系数的线性组合。

在下一节我们将看到，上述命题在把“有理”替换为“整数”并把“循环”替换为“初等”后仍然成立。

\section*{练习}

9.5. 取交替群 \({\mathfrak{A}}_{4}\) 为 \(G\) ，取 \(X\) 为 G 的循环子群族。令 \(\left\{  {{\chi }_{0},{\chi }_{1},{\chi }_{2},\psi }\right\}\) 为 G 的不同不可约特征(参见 5.7)。证明 \(\oplus  {\mathrm{R}}^{ + }\left( \mathrm{H}\right)\) 在 Ind 下的像由以下五个特征生成:

\[
{\chi }_{0} + {\chi }_{1} + {\chi }_{2} + \psi ,\;{2\psi },\;{\chi }_{0} + \psi ,\;{\chi }_{1} + \psi ,\;{\chi }_{2} + \psi .
\]

得出结论:元素 \(\chi\) 属于 \(R\left( G\right)\) 的像 Ind 的充要条件是 \(\chi \left( 1\right)  \equiv  0\) (模2)。证明没有任何字符 \({\chi }_{0},{\chi }_{1},{\chi }_{2}\) 是从循环子群诱导的字符的带正有理系数线性组合。

\section*{9.3 第一种证明}

首先我们证明 (ii) \(\Rightarrow\) (i)。令 \(\mathrm{S}\) 为属于 \(\mathrm{X}\) 的子群 \(\mathrm{H}\) 共轭类的并。任意形如 \(\sum {\operatorname{Ind}}_{\mathrm{H}}^{\overline{\mathrm{G}}}\left( {f}_{\mathrm{H}}\right)\) 的函数，带有 \({f}_{\mathrm{H}} \in  \mathrm{R}\left( \mathrm{H}\right)\) ，在 S 外消失。若满足 (ii)，则每个在 \(\mathbf{G}\) 上的类函数在 \(\mathbf{S}\) 外消失，由此说明 \(\mathbf{S} = \mathbf{G}\) 。因此 (i) 成立。

反过来，假设 (i) 成立。为证明 (ii)，只需证明 Q 线性映射

\[
\mathbf{Q} \otimes  \text{ Ind: }\underset{\mathrm{H} \in  \mathrm{X}}{ \oplus  }\mathbf{Q} \otimes  \mathrm{R}\left( \mathrm{H}\right)  \rightarrow  \mathbf{Q} \otimes  \mathrm{R}\left( \mathrm{G}\right) .
\]

是满的，这也等价于 C 线性映射

\[
\mathrm{C} \otimes  \text{ Ind: }\underset{\mathrm{H} \in  \mathrm{X}}{ \oplus  }\mathrm{C} \otimes  \mathrm{R}\left( \mathrm{H}\right)  \rightarrow  \mathrm{C} \otimes  \mathrm{R}\left( \mathrm{G}\right)
\]

由对偶性这等价于伴随映射的单射性

\[
\mathrm{C} \otimes  \mathrm{{Res}} : \mathrm{C} \otimes  \mathrm{R}\left( \mathrm{G}\right)  \rightarrow  \underset{\mathrm{H} \in  \mathrm{X}}{ \oplus  }\mathrm{C} \otimes  \mathrm{R}\left( \mathrm{H}\right) .
\]

但这个单射性是显然的:它相当于说，如果一个关于 \(G\) 的类函数在每个循环子群上限制为 0，则它为零。由此得证。

\section*{练习题}

我们假设族 \(X\) 在共轭与取子群下保持不变，且 \(G\) 是属于 X 的子群的并。(例:G 的循环子群族。)

9.6. 记由同态的核为 \(\mathrm{N}\)

\[
\mathbf{Q} \otimes  \text{ Ind: }\underset{\mathrm{H} \in  \mathrm{X}}{ \oplus  }\mathbf{Q} \otimes  \mathrm{R}\left( \mathrm{H}\right)  \rightarrow  \mathbf{Q} \otimes  \mathrm{R}\left( \mathrm{G}\right) .
\]

(a) 令 \(\mathrm{H},{\mathrm{H}}^{\prime } \in  \mathrm{X}\) ，伴随 \({\mathrm{H}}^{\prime } \subset  \mathrm{H}\) ，令 \({\chi }^{\prime } \in  \mathrm{R}\left( {\mathrm{H}}^{\prime }\right)\) 与 \(\chi  = {\operatorname{Ind}}_{{\mathrm{H}}^{\prime }}^{\mathrm{H}}\left( \chi \right) \; \in  R\left( H\right)\) 。证明 \(\chi  - {\chi }^{\prime }\) 属于 \(N\) 。

(b) 设 \(\mathrm{H} \in  \mathrm{X}\) 与 \(s \in  \mathrm{G}\) 。 置 \({}^{s}\mathrm{H} = s\mathrm{H}{s}^{-1}\) 。 设 \(\chi  \in  \mathrm{R}\left( \mathrm{H}\right)\) 并令 \({}^{s}\chi\) 为由 \({}^{s}\chi \left( {{sh}{s}^{-1}}\right)  = \chi \left( h\right)\) 在 \(\mathrm{R}\left( {{}^{s}\mathrm{H}}\right)\) 中为 \(h \in  \mathrm{H}\) 定义的元素。 证明 \(\chi  - {}^{s}\chi\) 属于 \(\mathrm{N}\) 。

(c) 证明 \(\mathrm{N}\) 由上述类型 (a) 与 (b) 的元素在 \(\mathbf{Q}\) 上生成。[把标量扩张到 \(\mathbf{C}\) 并利用对偶性。可归结为证明:若对每个 \(\mathrm{H} \in  \mathrm{X}\) 给定一个类函数 \({f}_{\mathrm{H}}\) 在 \(\mathrm{H}\) 上，且这些 \({f}_{\mathrm{H}}\) 满足与上述 (a) 与 (b) 类似的限制与共轭条件，则存在一个在 \(G\) 上的类函数 \(f\) ，使得对每个 H 有 \({\operatorname{Res}}_{\mathbf{H}}^{G}f = {f}_{\mathbf{H}}\) 。]*

9.7. 证明 \(\mathbf{Q} \otimes  \mathrm{R}\left( \mathrm{G}\right)\) 有如下形式的生成元与关系的表示:

生成元:符号 \(\left( {\mathrm{H},\chi }\right)\) ，带有 \(\mathrm{H} \in  \mathrm{X}\) 与 \(\chi  \in  \mathbf{Q} \otimes  \mathrm{R}\left( \mathrm{H}\right)\) 。 关系:

\HRule

* 本习题给出以诱导特征 \(\left( {\mathbf{H},\chi }\right)\) 表示的 \(\mathbf{Q} \otimes  \mathbf{R}\left( \mathbf{G}\right)\) 的“表示”。为应用于 L-函数理论，若能为 R(G) 自身(不与 Q 张量)给出此类表示将非常可取。当 G 可解时，Langlands–Deligne 已完成此事(Lecture Notes in Math. 349, p. 517, th. 4)。

\HRule

(i) \(\left( {\mathbf{H},{\lambda \chi } + {\lambda }^{\prime }{\chi }^{\prime }}\right)  = \lambda \left( {\mathbf{H},\chi }\right)  + {\lambda }^{\prime }\left( {\mathbf{H},{\chi }^{\prime }}\right)\) 对于 \(\lambda ,{\lambda }^{\prime } \in  \mathbf{Q}\) ，并且 \(x,{x}^{\prime } \in  \mathbf{Q} \; \otimes  \mathbf{R}\left( \mathbf{H}\right)\) 。

(ii) 对于 \({\mathrm{H}}^{\prime } \subset  \mathrm{H},{\chi }^{\prime } \in  \mathrm{R}\left( {\mathrm{H}}^{\prime }\right)\) ，与 \(\chi  = {\operatorname{Ind}}_{{\mathrm{H}}^{\prime }}^{\mathrm{H}}\left( {\chi }^{\prime }\right)\) ，有 \(\left( {\mathrm{H},\chi }\right)  = \left( {{\mathrm{H}}^{\prime },{\chi }^{\prime }}\right)\) 。

(iii) 对于 \(\mathrm{H} \in  \mathrm{X},s \in  \mathrm{G},\chi  \in  \mathrm{R}\left( \mathrm{H}\right)\) ，我们有 \(\left( {\mathrm{H},\chi }\right)  = \left( {{}^{s}\mathrm{H},{}^{s}\chi }\right)\) ，参见习题 9.6(b) 的记法。

[参见习题 9.6].

\section*{9.4 (i) \(\Rightarrow\) 的第二种证明 (ii)}

首先令 \(\mathbf{A}\) 为循环群，令 \(a\) 为其阶。用以下公式在 A 上定义函数 \({\theta }_{\mathbf{A}}\) :

\[
{\theta }_{\mathrm{A}}\left( x\right)  = \left\{  \begin{array}{ll} a & \text{ if }x\text{ generates }\mathrm{A} \\  0 & \text{ otherwise } \end{array}\right.
\]

命题 27. 若 \(G\) 为阶数为 \(g\) 的有限群，则

\[
g = \mathop{\sum }\limits_{{\mathrm{A} \subset  \mathrm{G}}}{\operatorname{Ind}}_{\mathrm{A}}^{\mathrm{G}}\left( {\theta }_{\mathrm{A}}\right)
\]

其中 A 遍历 \(\mathbf{G}\) 的所有循环子群。

(在此公式中，字母 \(g\) 表示恒等于 \(g\) 的常函数。)

取 \({\theta }_{\mathrm{A}}^{\prime } = {\operatorname{Ind}}_{\mathrm{A}}^{\mathrm{G}}\left( {\theta }_{\mathrm{A}}\right)\) 。对于 \(x \in  \mathrm{G}\) 我们有

\[
{\theta }_{\mathrm{A}}^{\prime }\left( x\right)  = \frac{1}{a}\mathop{\sum }\limits_{\substack{{y \in  \mathrm{G}} \\  {{yx}{y}^{-1} \in  \mathrm{A}} }}{\theta }_{\mathrm{A}}\left( {{yx}{y}^{-1}}\right)
\]

\[
= \frac{1}{a}\mathop{\sum }\limits_{\substack{{y \in  \mathrm{G}} \\  {{yx}{y}^{-1}\text{ gen. }\mathrm{A}} }}a = \mathop{\sum }\limits_{\substack{{y \in  \mathrm{G}} \\  {{yx}{y}^{-1}\text{ gen. }\mathrm{A}} }}1.
\]

然而，每个 \(y \in  \mathrm{G},{yx}{y}^{-1}\) 都生成 \(\mathrm{G}\) 的唯一一个循环子群。因此我们有:

\[
\mathop{\sum }\limits_{{\mathrm{A} \subset  \mathrm{G}}}{\theta }_{\mathrm{A}}^{\prime }\left( x\right)  = \mathop{\sum }\limits_{{y \in  \mathrm{G}}}1 = g.
\]

命题 28. 若 \(\mathrm{A}\) 为循环群，则 \({\theta }_{\mathrm{A}} \in  \mathrm{R}\left( \mathrm{A}\right)\) 。

证明按 \(\mathrm{A}\) 的阶 \(a\) 归纳， \(a = 1\) 的情形显然。由命题 27 我们有

\[
a = \mathop{\sum }\limits_{{\mathrm{B} \subset  \mathrm{A}}}{\operatorname{Ind}}_{\mathrm{B}}^{\mathrm{A}}\left( {\theta }_{\mathrm{B}}\right)  = {\theta }_{\mathrm{A}} + \mathop{\sum }\limits_{{\mathrm{B} \neq  \mathrm{A}}}{\operatorname{Ind}}_{\mathrm{B}}^{\mathrm{A}}\left( {\theta }_{\mathrm{B}}\right) .
\]

归纳假设给出对于 \(B \neq  A\) 的 \({\theta }_{B} \in  R\left( B\right)\) ，因此 \({\operatorname{Ind}}_{B}^{A}\left( {\theta }_{B}\right)\) 属于 \(R\left( A\right)\) ；另一方面显然 \(a \in  R\left( A\right)\) ，故得出 \({\theta }_{\mathrm{A}}\) 属于 \(\mathrm{R}\left( \mathrm{A}\right)\) 。

应用于 (i) \(\Rightarrow\) (ii) 的证明

首先注意，如果 \({A}^{\prime }\) 包含在 \(A\) 的共轭子中，则 \({\operatorname{Ind}}_{{A}^{\prime }}^{G}\) 的像包含在 \({\operatorname{Ind}}_{A}^{G}\) 的像中。因此我们可以假定 \(X\) 为 G 的所有循环子群的族。命题 27 和 28 随即表明

\[
g = \mathop{\sum }\limits_{{\mathrm{A} \in  \mathrm{X}}}{\operatorname{Ind}}_{\mathrm{A}}^{\mathrm{G}}\left( {\theta }_{\mathrm{A}}\right) ,\;\text{ with }{\theta }_{\mathrm{A}} \in  \mathrm{R}\left( \mathrm{A}\right) .
\]

因此元素 \(g\) 属于 Ind 的像。由于该像是 \(R\left( G\right)\) 的一个理想，参见 9.1，它包含所有形如 \({g\chi }\) 的元素，且 \(\chi  \in  R\left( G\right)\) ，这证明了 (ii')(甚至更强，因为我们有显式的分母，即 \(\mathrm{G}\) 的阶)。

\section*{练习}

9.8. 若 \(\mathrm{A}\) 为阶为 \(a\) 的循环群，则令 \({\lambda }_{\mathrm{A}} = \varphi \left( a\right) {r}_{\mathrm{A}} - {\theta }_{\mathrm{A}}\) ，其中 \(\varphi \left( a\right)\) 为 \(\mathrm{A}\) 的生成元数， \({r}_{a}\) 为正规表示的特征。证明 \({\lambda }_{\mathbf{A}}\) 是与单位特征正交的 \(\mathbf{A}\) 的一个特征 [应用练习 9.1]。证明当 A 遍历一个阶为 \(g\) 的群 G 的所有循环子群时，有

(*)
\[
\mathop{\sum }\limits_{{A \subset  G}}{\operatorname{Ind}}_{A}^{G}\left( {\lambda }_{A}\right)  = g\left( {{r}_{G} - 1}\right)
\]

其中 \({r}_{\mathrm{G}}\) 为 \(\mathrm{G}\) 的正规表示的特征 [用命题 27]。

[应用(Aramata–Brauer):设 \(\mathrm{F}\) 为数域 \(E\) 的有限扩张，且令其狄利克雷 η 函数的商为 \(\Phi \left( s\right)  = {\zeta }_{\mathrm{F}}\left( s\right) /{\zeta }_{\mathrm{E}}\left( s\right)\) 。已知 \(\Phi\) 在整个复平面上为亚纯函数。现在假设 L/K 为伽罗瓦扩张，伽罗瓦群为 G。那么上面的公式 (*) 蕴含恒等式

\[
\Phi {\left( s\right) }^{g} = \mathop{\prod }\limits_{\mathrm{A}}{\mathrm{L}}_{\mathrm{F}/{\mathrm{F}}_{\mathrm{A}}}\left( {s,{\lambda }_{\mathrm{A}}}\right) ,
\]

其中 \({F}_{A}\) 表示与循环子群 \(A\) 对应的子域 \(F\) 。函数 \({L}_{F/{F}_{A}}\left( {s,{\lambda }_{A}}\right)\) 是“阿贝尔” L-函数，因此全纯。因此我们看到 \(\Phi\) 本身是全纯的，即 \({\zeta }_{\mathrm{E}}\) 整除 \({\zeta }_{\mathrm{F}}\) ；尚不清楚该结论是否对非伽罗瓦扩张仍然成立(这将从 Artin 的猜想推论出)。]

第十章 布劳尔定理 在第 10.1 至 10.4 节中字母 \(p\) 表示一个素数。

10.1 p-正则元； \(p\) -元子群

设 \(x\) 为有限群 G 的一个元素。若 \(x\) 的阶为 \(p\) 的幂，则称 \(x\) 为 p-元(或称为 \(p\) -单根)；若其阶与 \(p\) 互素，则称 \(x\) 为 \({p}^{\prime }\) -元(或称为 \(p\) -正则)。

每个 \(x \in  \mathbf{G}\) 都可唯一写成 \(x = {x}_{u}{x}_{r}\) 的形式，其中 \({x}_{u}\) 为 \(p\) -单根， \({x}_{r}\) 为 \(p\) -正则，且 \({x}_{u}\) 与 \({x}_{r}\) 互相交换；此外， \({x}_{u}\) 与 \({x}_{r}\) 都是 \(x\) 的幂。可通过将由 \(x\) 生成的循环子群分解为其 \(p\) -分量与其 \({p}^{\prime }\) -分量的直积来看到这一点。元素 \({x}_{u}\) (分别 \({x}_{r}\) )称为 \(x\) 的 \(p\) -分量(分别 \({p}^{\prime }\) -分量)。

群 \(\mathrm{H}\) 称为 \(p\) -元子群，如果它是一个阶与 \(p\) 互素的循环群 \(\mathbf{C}\) 与一个 \(p\) -群 \(\mathbf{P}\) 的直积。这样的群是幂零群，其分解 \(\mathrm{C} \times  \mathrm{P}\) 是唯一的: \(\mathrm{C}\) 是 \(\mathrm{H}\) 的 \({p}^{\prime }\) -元元的集合，而 \(\mathrm{P}\) 是 \(p\) -元元的集合。

设 \(x\) 为有限群 \(\mathrm{G}\) 的一个 \({p}^{\prime }\) -元素，令 \(\mathrm{C}\) 为由 \(x\) 所生成的循环子群，令 \(Z\left( x\right)\) 为 \(x\) 的中心化子(所有满足 \({sx} = {xs}\) 的 \(s \in  \mathrm{G}\) 的集合)。若 \(\mathrm{P}\) 为 \(\mathrm{Z}\left( x\right)\) 的一个 Sylow \(p\) -子群，则群 \(\mathrm{H} = \mathrm{C} \cdot  \mathrm{P}\) 是 \(\mathrm{G}\) 中与 \(x\) 关联的一个 \(p\) -元素子群；在 \(\mathbf{Z}\left( x\right)\) 中该子群共轭意义下唯一。

\section*{练习}

10.1. 设 \(\mathrm{H} = \mathrm{C} \cdot  \mathrm{P}\) 为有限群 \(\mathrm{G}\) 的一个 \(p\) -元素子群，且令 C 的生成元为 \(x\) 。证明 \(\mathrm{H}\) 包含于一个与 \(x\) 关联的 \(p\) -元素子群 \({\mathrm{H}}^{\prime }\) 中。

10.2. 设 \(\mathbf{G} = {\mathbf{{GL}}}_{n}\left( k\right)\) ，其中 \(k\) 是特征为 \(p\) 的有限域。证明元素 \(x \in  \mathbf{G}\) 当且仅当其所有特征值均等于 1(即若 \(1 - x\) 为幂零)时为 \(p\) -元素；当且仅当它是半单的(即在 \(k\) 的某个有限扩张中可对角化)时为 \({p}^{\prime }\) -元素。

\section*{10.2 来自 \(p\) -元素子群的诱导特征}

本节及接下来的两节的目的是证明以下结果:

定理 18. 设 \(\mathrm{G}\) 为有限群，令 \({\mathrm{V}}_{p}\) 为由从 \(\mathrm{G}\) 的 p-元素子群的特征诱导得到的特征生成的子群。则 \({\mathrm{V}}_{p}\) 在 \(\mathrm{R}\left( \mathrm{G}\right)\) 中的指数为有限且与 p 互素。

令 \(\mathrm{X}\left( p\right)\) 为 \(\mathrm{G}\) 的 \(p\) -元素子群族。群 \({\mathrm{V}}_{p}\) 是由诱导同态 \(\mathrm{X}\left( p\right)\) 定义的同态的像。

\[
\text{ Ind: }\mathop{\oplus }\limits_{{\mathrm{H} \in  \mathrm{X}\left( p\right) }}\mathrm{R}\left( \mathrm{H}\right)  \rightarrow  \mathrm{R}\left( \mathrm{G}\right)
\]

由诱导同态 \({\operatorname{Ind}}_{H}^{G},H \in  X\left( p\right)\) 定义。于是 \({V}_{p}\) 是 \(R\left( G\right)\) 的一个理想，要证明该定理，只需证明存在一个与 \(p\) 互素的整数 \(m\) ，使得 \(m \in  {\mathbf{V}}_{p}\) 。事实上，我们证明了以下更精确的结果:

定理 18’. 设 \(g = {p}^{n}l\) 为 \(\mathrm{G}\) 的阶，且有 \(\left( {p,l}\right)  = 1\) 。则 \(l \in  {\mathrm{V}}_{p}\) 。

证明(归功于 Roquette 和 Brauer-Tate [12])使用了由 \(g\) 次单位根生成的子环 A。该环作为 Z-模是自由且有限生成的；其元素为代数整数。我们有 \(\mathbf{Q} \cap  \mathbf{A} = \mathbf{Z}\) ，因为该交集的元素既是有理数又是代数整数(参见 6.4)。商群 \(A/Z\) 是有限生成且无挠的，因此是自由的；由此(通过在 A 中提升 \(A/Z\) 的一组基)可得 A 有一组包含元素 1 的基 \(\left\{  {1,{\alpha }_{1},\ldots ,{\alpha }_{c}}\right\}\) 。

通过与 A 张量，同态 Ind 定义了一个 A 线性映射

\[
\mathrm{A} \otimes  \mathrm{{Ind}} : \underset{\mathrm{H} \in  \mathrm{X}\left( p\right) }{ \oplus  }\mathrm{A} \otimes  \mathrm{R}\left( \mathrm{H}\right)  \rightarrow  \mathrm{A} \otimes  \mathrm{R}\left( \mathrm{G}\right) .
\]

基 \(\left\{  {1,{\alpha }_{1},\ldots ,{\alpha }_{c}}\right\}\) 的存在随后蕴含下述结论:

引理 5。 \(\mathrm{A} \otimes\) Ind 的像为 \(\mathrm{A} \otimes  {\mathrm{V}}_{p}\) ；此外我们有

\[
\left( {\mathrm{A} \otimes  {\mathrm{V}}_{p}}\right)  \cap  \mathrm{R}\left( \mathrm{G}\right)  = {\mathrm{V}}_{p}\text{ . }
\]

因此，要证明常函数 \(l\) 属于 \({\mathbf{V}}_{p}\) ，只需证明 \(l\) 属于 \(A \otimes\) Ind 的像，换言之， \(l\) 可表示为 \(\mathop{\sum }\limits_{\mathrm{H}}{a}_{\mathrm{H}}{\operatorname{Ind}}_{\mathrm{H}}^{\mathrm{G}}\left( {f}_{\mathrm{H}}\right)\) 的形式，其中 \({a}_{\mathrm{H}} \in  \mathrm{A}\) 与 \({f}_{\mathrm{H}} \in  \mathrm{R}\left( \mathrm{H}\right)\) 。

\section*{备注}

(1) 环 \(\mathbf{A}\) 相对于环 \(\mathbf{Z}\) 的优点在于 \(G\) 的所有特征值都在 \(A\) 中，因为这些值是 \(g\) 次单位根的和。由此可知 \(A \otimes  R\left( G\right)\) 是取值在 A 的 \(G\) 上类函数环的一个子环。

(2) 可证明 \(A\) 即为圆乘域 \(\mathbf{Q} \cdot  \mathbf{A}\) 的代数整数集，但我们不需要此事实。

\section*{10.3 特征的构造}

引理 6。任一在 \(\mathbf{G}\) 上取值为可被 \(g\) 整除的整数的类函数，均可表示为从 \(\mathrm{G}\) 的循环子群的特征诱导来的特征的 A 线性组合。

(在此及以下所有情形中，“整数值”一词意指“取值在 Z 中”。)

设 \(f\) 为此类函数，并写成 \({g\chi }\) 的形式，其中 \(\chi\) 是取整数值的类函数。若 \(\mathbf{C}\) 为 \(\mathbf{G}\) 的循环子群，令 \({\theta }_{\mathbf{C}}\) 为 9.4 节中定义的 \(R\left( C\right)\) 中的元素。我们有

\[
g = \mathop{\sum }\limits_{\mathrm{C}}{\operatorname{Ind}}_{\mathrm{C}}^{\mathrm{G}}\left( {\theta }_{\mathrm{C}}\right) ,\;\text{ cf. prop. }{27},
\]

由此

\[
f = {g\chi } = \mathop{\sum }\limits_{\mathrm{C}}{\operatorname{Ind}}_{\mathrm{C}}^{\mathrm{G}}\left( {\theta }_{\mathrm{C}}\right) \chi  = \mathop{\sum }\limits_{\mathrm{C}}{\operatorname{Ind}}_{\mathrm{C}}^{\mathrm{G}}\left( {{\theta }_{\mathrm{C}} \cdot  {\operatorname{Res}}_{\mathrm{C}}^{\mathrm{G}}\chi }\right) .
\]

仍需证明对每个 \(\mathrm{C}\) ， \({\theta }_{\mathrm{C} \cdot  {\mathrm{{Res}}}_{\mathrm{{Cv}}}}\) 属于 \(\mathrm{A} \otimes  \mathrm{R}\left( \mathrm{C}\right)\) 。但 \({\chi }_{\mathrm{C}} = {\theta }_{\mathrm{C}} \cdot  {\operatorname{Res}}_{\mathrm{C}}\chi\) 的值可被 \(\mathrm{C}\) 的阶整除，故若 \(\psi\) 为 \(\mathrm{C}\) 的一个特征，则有 \(\left\langle  {{\chi }_{\mathrm{C}},\psi }\right\rangle   \in  \mathrm{A}\) ，这表明 \({\chi }_{\mathrm{C}}\) 可用 \(\mathrm{A}\) -线性组合的 \(\mathrm{C}\) 的特征表示，因此 \({\chi }_{\mathrm{C}} \in  \mathrm{A} \otimes  \mathrm{R}\left( \mathrm{C}\right)\) 。

引理 7. 设 \(\chi\) 为取整数值的 \(\mathrm{A} \otimes  \mathrm{R}\left( \mathrm{G}\right)\) 的一个元素，设 \(x \in  \mathrm{G}\) ，并设 \({x}_{r}\) 为 \(x\) 的 \({p}^{\prime }\) 分量(参见 10.1)。则

\[
\chi \left( x\right)  \equiv  \chi \left( {x}_{r}\right) \text{ (mod. }p\text{ ). }
\]

通过限制，我们归结到 G 为由 \(x\) 生成的循环群的情形。现在 \(\chi  = \sum {a}_{i}{\chi }_{i}\) ，其与 \({a}_{i} \in  \mathrm{A}\) 以及遍历 G 的不同的一次特征的 \({\chi }_{i}\) 。若 \(q\) 为 \(p\) 的充分大次幂，则有 \({x}^{q} = {x}_{r}^{q}\) ，从而对所有 \(i\) 有 \({\chi }_{i}{\left( x\right) }^{q} = {\chi }_{i}{\left( {x}_{r}\right) }^{q}\) 。因此

\[
\chi {\left( x\right) }^{q} = {\left( \sum {a}_{i}{\chi }_{i}\left( x\right) \right) }^{q} \equiv  \sum {a}_{i}^{q}{\chi }_{i}{\left( x\right) }^{q}
\]

\[
\equiv  \sum {a}_{i}^{q}{\chi }_{i}{\left( {x}_{r}\right) }^{q}) \equiv  \chi {\left( {x}_{r}\right) }^{q}\left( {{\;\operatorname{mod}\;.}p\text{ A }}\right) .
\]

由于 \(p\mathbf{A} \cap  \mathbf{Z} = \mathbf{Z}\) ，这意味着

\[
\chi {\left( x\right) }^{q} \equiv  \chi {\left( {x}_{r}\right) }^{q}\left( {\;\operatorname{mod}\;.p}\right)
\]

因此 \(\chi \left( x\right)  \equiv  \chi \left( {x}_{r}\right) \left( {{\;\operatorname{mod}\;.}p}\right)\) ，因为对所有 \(\lambda  \in  \mathbf{Z}\) 有 \({\lambda }^{q} = \lambda \left( {{\;\operatorname{mod}\;.}p}\right)\) 。

引理 8. 设 \(x\) 为 \({p}^{\prime }\) 的一个 \(\mathbf{G}\) 元素，且设 \(\mathbf{H}\) 为与 \(x\left( {10.1}\right)\) 相关的 \(\mathrm{G}\) 的一个 p-初等子群。则存在一个取整数值的函数 \(\psi  \in  \mathrm{A} \otimes  \mathrm{R}\left( \mathrm{H}\right)\) ，使得诱导函数 \({\psi }^{\prime } = {\operatorname{Ind}}_{\mathrm{H}}^{\mathrm{G}}\psi\) 具有下列性质:

(a) \({\psi }^{\prime }\left( x\right)  \text{ ≢ } 0\left( {\;\operatorname{mod}\;.p}\right)\) .

(b) 对于每个不与 \(x\) 共轭的 \({p}^{\prime }\) -元素的 \(\mathrm{G}\) ，均有 \({\psi }^{\prime }\left( s\right)  = 0\) 。

令 \(\mathrm{C}\) 为由 \(\mathrm{G}\) 中的 \(x\) 生成的循环子群，令 \(\mathrm{Z}\left( x\right)\) 为 \(\mathrm{G}\) 中 \(x\) 的中心化子。我们有 \(\mathrm{H} = \mathrm{C} \times  \mathrm{P}\) ，其中 \(\mathrm{P}\) 是 \(Z\left( x\right)\) 的一个 Sylow \(p\) -子群。令 \(c\) 为 \(C\) 的阶，令 \({p}^{a}\) 为 \(P\) 的阶。令 \({\psi }_{C}\) 为在 \(\mathbf{C}\) 上定义的函数

\[
{\psi }_{\mathrm{C}}\left( x\right)  = c\;\text{ and }\;{\psi }_{\mathrm{C}}\left( y\right)  = 0\;\text{ if }y \neq  x.
\]

我们有 \({\psi }_{\mathrm{C}} = \mathop{\sum }\limits_{\chi }\chi \left( {x}^{-1}\right) \chi\) ，其中 \(\chi\) 遍历 \(\mathbf{C}\) 的不可约特征的集合；因此 \({\psi }_{\mathbf{C}}\) 属于 \(\mathbf{A} \otimes  \mathbf{R}\left( \mathbf{C}\right)\) (这也可由引理 6 得出)。

令 \(\psi\) 为在 \(H = C \times  P\) 上由 \(\psi \left( {xy}\right)  = {\psi }_{C}\left( x\right)\) 为 \(x \in  \mathrm{C}\) 和 \(y \in  \mathrm{P}\) 定义的函数。这是投影 \(H \rightarrow  C\) 的逆像 \({\psi }_{\mathrm{C}}\) 。所以我们有 \(\psi  \in  A \otimes  R\left( H\right)\) 。现证明 \(\psi\) 满足该引理的条件:

如果 \(s\) 是 \(\mathbf{G}\) 的一个 \({p}^{\prime }\) -元，而 \(y \in  \mathbf{G},{ys}{y}^{-1}\) 也是一个 \({p}^{\prime }\) -元；若 \({ys}{y}^{-1}\) 属于 \(H\) 则属于 \(C\) ，并且只要 \({ys}{y}^{-1} \neq  x\) 就有 \(\psi \left( {{ys}{y}^{-1}}\right)  = 0\) 。若 \(s\) 不与 \(x\) 共轭，则得 \({\psi }^{\prime }\left( s\right)  = 0\) ，这证明了 (b)。此外:

\[
{\psi }^{\prime }\left( x\right)  = \frac{1}{c \cdot  {p}^{a}}\mathop{\sum }\limits_{{{yx}{y}^{-1} = x}}\psi \left( x\right)  = \frac{1}{{p}^{a}}\mathop{\sum }\limits_{{{yx}{y}^{-1} = x}}1 = \frac{\operatorname{Card}\left( {\mathrm{Z}\left( x\right) }\right) }{{p}^{a}}
\]

因此 \({\psi }^{\prime }\left( x\right)  \text{ ≢ } 0\left( {\;\operatorname{mod}\;.p}\right)\) ，因为 \({p}^{a} = \operatorname{Card}\left( \mathrm{P}\right)\) 是整除 \(\operatorname{Card}\left( {\mathrm{Z}\left( x\right) }\right)\) 的最大的 \(p\) 幂。

引理 9. 存在一个属于 \(\mathrm{A} \otimes  {\mathrm{V}}_{p}\) 的元素 \(\psi\) (取整数值)，使得 对每个 \(x \in  \mathrm{G}\) 都有 \(\psi \left( x\right)  \text{ ≢ } 0\left( {{\;\operatorname{mod}\;.}p}\right)\) 。

设 \({\left( {x}_{i}\right) }_{i \in  \mathrm{I}}\) 为 \(p\) -正规类的一组代表(即那些由 \({p}^{\prime }\) -元组成的类)。引理 8 给出我们在 \(\mathrm{A} \otimes  {\mathrm{V}}_{p}\) 中的一个取整数值的元素 \({\psi }_{i}\) ，使得

\[
{\psi }_{i}\left( {x}_{i}\right)  \text{ ≢ } 0\left( {\;\operatorname{mod}\;p}\right) \;\text{ and }\;{\psi }_{i}\left( {x}_{j}\right)  \equiv  0\left( {\;\operatorname{mod}\;.\;p}\right) \text{ for }j \neq  i.
\]

记 \(\psi  = \sum {\psi }_{i}\) 。显然 \(\psi\) 属于 \(\mathrm{A} \otimes  {\mathrm{V}}_{p}\) 并取整数值。对于 \(x \in  \mathrm{G}\) ， \(x\) 的 \({p}^{\prime }\) 分量与唯一的 \({x}_{i}\) 共轭。由引理 7 我们得到

\[
\psi \left( x\right)  \equiv  \psi \left( {x}_{i}\right)  \equiv  {\psi }_{i}\left( {x}_{i}\right)  \text{ ≢ } 0\left( {\;\operatorname{mod}\;.p}\right) .
\]

\section*{练习}

10.3. 将引理 6 推广到取值于 \(\mathbf{A}\) 的理想 \(g\mathbf{A}\) 的类函数。

10.4. 设 \(\mathfrak{p}\) 为 \(\mathbf{A}\) 的一个素理想，使得 \(\mathfrak{p} \cap  \mathbf{Z} = p\mathbf{Z}\) (等价于说 \(\mathbf{A}/\mathfrak{p}\) 是特征为 \(p\) 的有限域)。设 \(\chi  \in  \mathbf{A} \otimes  \mathbf{R}\left( \mathbf{G}\right)\) ，设 \(x \in  \mathbf{G}\) ，并令 \({x}_{r}\) 为 \(x\) 的 \({p}^{\prime }\) 分量。证明 \(\chi \left( x\right)  \equiv  \chi \left( {x}_{r}\right)\) (模 p)(证明同引理 7)，但我们不再总是有 \(\chi \left( x\right)  \equiv  \chi \left( {x}_{r}\right) \left( {{\;\operatorname{mod}\;.}\mathfrak{p}\mathrm{\;A}}\right) .\)

\section*{10.4 定理 18 及 \({18}^{\prime }\) 的证明}

设 \(g = {p}^{n}l\) 为 \(\mathrm{G}\) 的阶，且 \(\left( {p,l}\right)  = 1\) 。只需证明 \(l\) 属于 \(A \otimes  {V}_{p}\) ，参见 10.2。

设 \(\psi\) 为满足引理 9 条件的 \(A \otimes  {V}_{p}\) 的一个元素。 \(\psi\) 的取值为 \(\text{ ≢ } 0\left( {{\;\operatorname{mod}\;.}p}\right)\) 。设 \(\mathrm{N} = \varphi \left( {p}^{n}\right)\) 为群 \({\left( \mathbf{Z}/{p}^{n}\mathbf{Z}\right) }^{ * }\) 的阶，因此对每个与 \(p\) 互素的整数 \(\lambda\) 有 \({\lambda }^{\mathrm{N}} \equiv  1\left( {{\;\operatorname{mod}\;.}{p}^{n}}\right)\) 。因此对所有 \(x \in  \mathrm{G}\) 有 \(\psi {\left( x\right) }^{\mathrm{N}} \equiv  1\left( {{\;\operatorname{mod}\;.}{p}^{n}}\right)\) ，并且函数 \(l\left( {{\psi }^{\mathrm{N}} - 1}\right)\) 取可被 \(l{p}^{n} = g\) 整除的整数值。由引理 6，该函数是从 G 的循环子群诱导出的特征的 A-线性组合。由于每个循环群是 \(p\) -元的，我们有 \(l\left( {{\psi }^{\mathrm{N}} - 1}\right)  \in  \mathrm{A} \otimes  {\mathrm{V}}_{p}\) 。但 \(\mathrm{A} \otimes  {\mathrm{V}}_{p}\) 是 \(A \otimes  R\left( G\right)\) 的一个理想，故有 \(l{\psi }^{N} \in  A \otimes  {V}_{p}\) 。相减得出 \(l\) 属于 \(\mathrm{A} \otimes  {\mathrm{V}}_{p}\) ，从而完成证明。

\section*{10.5 布劳尔定理}

我们称 \(\mathrm{G}\) 的子群为元(elementary)，如果它对至少一个素数 \(p\) 是 \(p\) -元的。

定理 19. \(\mathrm{G}\) 的每个特征都可以表示为从元子群的特征诱导而来的特征的带整数系数的线性组合。

设 \({\mathrm{V}}_{p}\) 为定理 18 中定义的 \(\mathrm{R}\left( \mathrm{G}\right)\) 的子群。只需证明对素数 \(p\) ，由所有 \({\mathrm{V}}_{p}\) 的和 \(\mathrm{V}\) 等于 \(\mathrm{R}\left( \mathrm{G}\right)\) 。现在 \(\mathrm{V}\) 包含 \({\mathrm{V}}_{p}\) ，因此 \(\mathrm{V}\) 在 \(\mathrm{R}\left( \mathrm{G}\right)\) 中的指数整除 \({\mathrm{V}}_{p}\) 的指数，因而由定理 18 与 \(p\) 互素。既然对所有 \(p\) 都成立，则该指数为 1，从而证明定理。

定理 20. \(\mathrm{G}\) 的每个特征都可以表示为幂特征(monomial characters)的带整数系数的线性组合。

(回想:若一个特征可由某子群的 1 次特征诱导而来，则称为幂特征。)

此命题由定理 19 推出，并且每个元群的特征都是幂特征，因为此类群是幂零的(参见 8.5，定理 16)。

\section*{备注}

(1) 定理 19 和 20 中出现的线性组合其系数可为正也可为负。通常不可能将给定特征表示为仅有正系数(整数或甚至实数)的幂特征的线性组合，见下文习题 10.5。

(2) 定理 20 在表示论的许多应用中发挥关键作用:在很大程度上，它把任意特征 \(\chi\) 的问题约化到 \(\chi\) 为 1 次(因此来自循环群的特征)的情形。例如，布劳尔正是用此方法证明了阿廷 L 函数在整个复平面上解析继续为亚纯函数。我们稍后还会看到其他应用。

\section*{习题}

10.5. 设 \(\chi\) 是群 G 的一个不可约特征。

(a) 设 \(\chi\) 是若干单项特征以正实系数的线性组合。证明存在一个整数 \(m \geq  1\) 使得 \({m\chi }\) 是单项的。

(b) 取交替群为 \({\mathfrak{A}}_{5}\) 的 \(\mathbf{G}\) 。相应的置换表示是平凡表示与一个4维不可约表示的直和；令 \(\chi\) 为后者的特征。若 \({m\chi }\) 由某子群 \(\mathrm{H}\) 的1维特征所诱导，则 \(\mathrm{H}\) 的阶将等于 \({15}/m\) ，且 \(m\) 只能取 1、3、5、15。并且， \(\chi\) 限制到 \(\mathrm{H}\) 时必须包含一个1维特征且重数为 \(m\) (注意 \(\mathbf{G}\) 无阶为15的子群)。由此得出 \(\chi\) 不能表示为若干单项特征以正实系数的线性组合。

10.6.(A. Weil 建议)我们要证明:每个 \(f \in  R\left( G\right)\) ，使得 \(f\left( 1\right)  = 0\) 是由形如 \({\operatorname{Ind}}_{\mathrm{E}}^{G}\left( {\alpha  - 1}\right)\) 的元素的 \(\mathbf{Z}\) -线性组合，其中 \(\mathrm{E}\) 是 \(\mathrm{G}\) 的元群子群，且 \(\alpha\) 是1维特征。

(a) 令 \({\mathrm{R}}_{0}^{\prime }\left( \mathrm{G}\right)\) 为由 \({\operatorname{Ind}}_{\mathrm{E}}^{\mathrm{G}}\left( {\alpha  - 1}\right)\) 生成的 \(\mathrm{R}\left( \mathrm{G}\right)\) 的子群，且令 \({\mathrm{R}}^{\prime }\left( \mathrm{G}\right)  = \mathbf{Z} + {\mathrm{R}}_{0}^{\prime }\left( \mathrm{G}\right)\) 。证明:若 \(\mathrm{H}\) 是 \(\overline{\mathrm{G}},{\operatorname{Ind}}_{\mathrm{H}}^{\mathrm{G}}\) 的子群，则 \({\mathrm{R}}_{0}^{\prime }\left( \mathrm{H}\right)\) 将映入 \({\mathrm{R}}_{0}^{\prime }\left( \mathrm{G}\right)\) 。

(b) 假设 \(H\) 在 \(G\) 中正规，且 \(G/H\) 是阿贝尔群。证明 \({\operatorname{Ind}}_{H}^{G}\) 将 \({R}^{\prime }\left( H\right)\) 映到 \({R}^{\prime }\left( G\right)\) 中。[只需证明 \({\operatorname{Ind}}_{H}^{G}\left( 1\right)\) 属于 \({\mathrm{R}}^{\prime }\left( \mathrm{G}\right)\) ，这可由以下事实推出: \({\operatorname{Ind}}_{\mathrm{H}}^{\mathrm{G}}\left( 1\right)\) 是 \(\mathrm{G}\) 的 \(\left( {\mathrm{G} : \mathrm{H}}\right)\) 个次数为 1 的特征标的和，且这些特征标的核包含 \(\mathrm{H}\) 。]

(c) 设 \(\mathrm{G}\) 为素元群。令 \(\mathrm{Y}\) 为 \(\mathrm{G}\) 的极大子群集合。证明若 \(\mathrm{H} \in  \mathrm{Y}\) ，则 \(\mathrm{H}\) 在 \(\mathrm{G}\) 中为正规子群，且 \(\mathrm{G}/\mathrm{H}\) 的阶为素数〔利用 \(G\) 为幂零群这一事实〕。推得 \(R\left( G\right)\) 由 \(G\) 的一维特征构与 \({\operatorname{Ind}}_{H}^{G}\left( {R\left( H\right) }\right)\) 所生成，其中 \(H\) 遍历 \(Y\) 〔应用定理 16〕。证明 \({R}^{\prime }\left( G\right)  = R\left( G\right)\) 〔对 \(G\) 的阶用归纳法，并用 (b) 证明 \({\operatorname{Ind}}_{H}^{G}\left( {R\left( H\right) }\right)\) 包含于 \(\left. {{\mathrm{R}}^{\prime }\left( \mathrm{G}\right) }\right\rbrack\) 〕。

(d) 回到一般情形，记 \(X\) 为 \(G\) 的素元子群集合。由定理 19 我们有 \(1 = \sum {\operatorname{Ind}}_{E}^{G}\left( {f}_{E}\right)\) ，且 \({f}_{E} \in  R\left( E\right)\) 。若 \(\varphi  \in  \mathbf{R}\left( G\right)\) 则得

\[
\varphi  = \mathop{\sum }\limits_{{E \in  X}}{\operatorname{Ind}}_{E}^{G}\left( {\varphi }_{E}\right) \;\text{ where }{\varphi }_{E} = {f}_{E} \cdot  {\operatorname{Res}}_{E}^{G}\left( \varphi \right) .
\]

若 \(\varphi \left( 1\right)  = 0\) ，则由 (c) 我们有 \({\varphi }_{E} \in  {R}_{0}^{\prime }\left( E\right)\) 。由此得出 \(\varphi\) 属于 \({R}_{0}^{\prime }\left( G\right)\) ，从而 \({R}^{\prime }\left( G\right)  = R\left( G\right)\) 。

第 11 章 布劳尔定理的应用

11.1 特征值的刻画

设 \(B\) 为 \(C\) 的子环，且令 \(G\) 为有限群。

定理 21. 设 \(\varphi\) 为作用于 \(G\) 的类函数，使得对每一素元子群 \(\mathrm{H}\) of \(\mathrm{G}\) ，我们有 \({\operatorname{Res}}_{\mathrm{H}}^{\mathrm{G}}\varphi  \in  \mathrm{B} \otimes  \mathrm{R}\left( \mathrm{H}\right)\) 。则 \(\varphi  \in  \mathrm{B} \otimes  \mathrm{R}\left( \mathrm{G}\right)\) 。

令 \(\mathrm{X}\) 为 G 的所有素元子群的集合。由定理 11，我们可将常值函数 1 写成

\[
1 = \mathop{\sum }\limits_{{\mathrm{H} \in  \mathrm{X}}}{\operatorname{Ind}}_{\mathrm{H}}^{\mathrm{G}}{f}_{\mathrm{H}},\;\text{ with }{f}_{\mathrm{H}} \in  \mathrm{R}\left( \mathrm{H}\right) .
\]

两边乘以 \(\varphi\) ，得

\[
\varphi  = \mathop{\sum }\limits_{{H \in  X}}\varphi  \cdot  {\operatorname{Ind}}_{H}^{G}{f}_{H} = \mathop{\sum }\limits_{{H \in  X}}{\operatorname{Ind}}_{H}^{G}\left( {{f}_{H} \cdot  {\operatorname{Res}}_{H}^{G}\varphi }\right) .
\]

由于 \({f}_{\mathrm{H}}\) 属于 \(\mathrm{R}\left( \mathrm{H}\right)\) 且 \({\operatorname{Res}}_{\mathrm{H}}\varphi\) 属于 \(\mathrm{B} \otimes  \mathrm{R}\left( \mathrm{H}\right)\) ，二者的乘积属于 \(B \otimes  R\left( H\right)\) 。因此 \(\varphi\) 属于 \(B \otimes  R\left( G\right)\) 。

类似的论证，利用阿廷定理(第 9 章)可得:

定理21’。假设 \(\mathrm{B}\) 包含 \(\mathrm{Q}\) 。若对 \(\mathrm{G}\) 的每个循环子群 \(\mathrm{H}\) 都有 \({\operatorname{Res}}_{\mathrm{H}}^{\mathrm{G}}\varphi  \in  \mathrm{B} \otimes  \mathrm{R}\left( \mathrm{H}\right)\) ，则 \(\varphi  \in  \mathrm{B} \otimes  \mathrm{R}\left( \mathrm{G}\right)\) 。

注。定理21 可解释为一种相容性性质。假设对每个 \(H \in  X\) 给定一个元素 \({\varphi }_{H}\) 属于 \(B \otimes  R\left( H\right)\) ，并且满足下列性质:

(i) 若 \({\mathrm{H}}^{\prime } \subset  \mathrm{H}\) ，则 \({\varphi }_{{\mathrm{H}}^{\prime }} = {\operatorname{Res}}_{{\mathrm{H}}^{\prime }}^{\mathrm{H}}\left( {\varphi }_{\mathrm{H}}\right)\) 。

(ii) 若 \({\mathrm{H}}^{\prime } = s\mathrm{H}{s}^{-1}\) ，且有 \(s \in  \mathrm{G}\) ，则 \({\varphi }_{{\mathrm{H}}^{\prime }}\) 由 \({\varphi }_{\mathrm{H}}\) 经由同构 \(x \mapsto  {sx}{s}^{-1}\) 得到。

则存在唯一的元素 \(\varphi\) 属于 \(B \otimes  R\left( G\right)\) ，使得对所有 \(\mathrm{H} \in  \mathrm{X}\) 都有 \({\operatorname{Res}}_{H}^{G}\varphi \; = {\varphi }_{\mathrm{H}}\) 。

定理22。设 \(\varphi\) 是定义在 \(\mathbf{G}\) 上的类函数，使得对 \(\mathrm{G}\) 的每个基本子群 \(\mathrm{H}\) 及其每个一度特征 \(\chi\) ，数

\[
{\left\langle  \chi ,{\operatorname{Res}}_{\mathrm{H}}\varphi \right\rangle  }_{\mathrm{H}} = \frac{1}{\operatorname{Card}\left( \mathrm{H}\right) }\mathop{\sum }\limits_{{s \in  \mathrm{H}}}\chi \left( {s}^{-1}\right) \varphi \left( s\right)
\]

属于 B。则 \(\varphi\) 属于 \(\mathrm{B} \otimes  \mathrm{R}\left( \mathrm{G}\right)\) 。

令 \(\mathrm{H}\) 为 \(\mathrm{G}\) 的一个基本子群。设

\[
{\operatorname{Res}}_{\mathrm{H}}^{\mathrm{G}}\varphi  = \mathop{\sum }\limits_{\omega }{c}_{\omega }\omega ,\;\text{ where }{c}_{\omega } = {\left\langle  \omega ,{\operatorname{Res}}_{\mathrm{H}}\varphi \right\rangle  }_{\mathrm{H}},
\]

为 \({\operatorname{Res}}_{\mathrm{H}}^{\mathrm{G}}\varphi\) 按不可约特征 \(\omega\) (属于 \(\mathrm{H}\) )的分解。由定理16，每个特征 \(\omega\) 都由 \(\mathrm{H}\) 的某子群 \({\mathrm{H}}_{\omega }\) 上的一度特征 \({\chi }_{\omega }\) 诱导而来。由弗罗贝尼乌斯互反性，我们有

\[
{c}_{\omega } = {\left\langle  {\chi }_{\omega },{\operatorname{Res}}_{{\mathrm{H}}_{\omega }}^{\mathrm{G}}\varphi \right\rangle  }_{{\mathrm{H}}_{\omega }}.
\]

由于 \({\mathrm{H}}_{\omega }\) 是基本群，关于 \(\varphi\) 的假设保证 \({c}_{\omega }\) 属于 B。因此， \({\operatorname{Res}}_{\mathrm{H}}\varphi  = \sum {c}_{\omega }\omega\) 属于 \(\mathrm{B} \otimes  \mathrm{R}\left( \mathrm{H}\right)\) ，由定理21 得出结论。

推论。要使 \(\varphi\) 成为一个虚特征(即 \(\varphi  \in  \mathbf{R}\left( \mathrm{G}\right)\) )，当且仅当，对任意基本子群 \(\mathrm{H}\) 与任意同态 \(\chi  : \mathrm{H} \rightarrow  {\mathbf{C}}^{ * }\) ，都有 \({\left\langle  \chi ,{\operatorname{Res}}_{\mathrm{H}}\varphi \right\rangle  }_{\mathrm{H}} \in  \mathbf{Z}\) 。

这是特殊情形 \(B = \mathbf{Z}\) 。

令 Res 表示由限制同态 \({\operatorname{Res}}_{H}^{G}\) 定义的从 \(R\left( G\right)\) 到 \(\_  \oplus  \_ R\left( H\right)\) 的同态。

命题 29. 同态 Res: \(\mathrm{R}\left( \mathrm{G}\right)  \rightarrow  \underset{\mathrm{H} \in  \mathrm{X}}{ \oplus  }\mathrm{R}\left( \mathrm{H}\right)\) 是一个分裂单射。

(若存在 \(r : \mathrm{M} \rightarrow  \mathrm{L}\) 使得 \(r \circ  f = 1\) ，则称模同态 \(f : \mathrm{L} \rightarrow  \mathrm{M}\) 为分裂单射；这等价于 \(f\) 是单射且 \(f\left( \mathrm{\;L}\right)\) 是 \(\mathrm{M}\) 的直和因子。)

显然 Res 是单射。为证明其为分裂同态，只需证明其余核为无挠元素，因为所考虑的群是有限生成的自由 Z-模。因此我们要证明:若 \(f = {\left( {f}_{\mathrm{H}}\right) }_{\mathrm{H} \in  \mathrm{X}}\) 为 \(\oplus  \mathrm{R}\left( \mathrm{H}\right)\) 的元素，且存在非零的 \(n\) 使得 \({nf} = \operatorname{Res}\varphi\) ，并且 \(\varphi  \in  \mathrm{R}\left( \mathrm{G}\right)\) ，则 \(f \in  \operatorname{Im}\left( \operatorname{Res}\right)\) 。但这可由定理 21 应用于函数 \(\varphi /n\) 与环 \(\mathbf{Z}\) 得出。

[该论证亦可用对偶性的表述:由于所涉群为有限生成的自由 Z-模，证明 Res 为分裂同态等价于证明其转置为满射。但其转置是

\[
\text{ Ind: } \oplus  \mathrm{R}\left( \mathrm{H}\right)  \rightarrow  \mathrm{R}\left( \mathrm{G}\right) \text{ , }
\]

而这正是由布劳尔定理保证的满射。]

\section*{11.2 佛罗贝尼乌斯定理}

与第 10 章相同，我们记 A 为由 \(g\) 次单位根生成的 \(\mathbf{C}\) 的子环，其中 \(g = \operatorname{Card}\left( \mathrm{G}\right)\) 。

设 \(n\) 为整数 \(\geq  1\) ，令 \(\left( {g,n}\right)\) 为 \(g\) 与 \(n\) 的最大公约数。若 \(f\) 是定义在 \(\mathrm{G}\) 上的函数，则记 \({\Psi }^{n}f\) 为函数 \(x \mapsto  f\left( {x}^{n}\right)\) 。易验证(参见习题 9.3)算子 \({\Psi }^{n}\) 将 \(\mathrm{R}\left( \mathrm{G}\right)\) 映入自身。此外:

定理 23. 若 \(f\) 是取值于 \(\mathrm{A}\) 的 \(\mathrm{G}\) 上的类函数，则函数 \(\left( {g/\left( {g,n}\right) }\right) {\Psi }^{n}f\) 属于 \(\mathrm{A} \otimes  \mathrm{R}\left( \mathrm{G}\right)\) 。

若 \(c\) 是 \(\mathrm{G}\) 的共轭类，记 \({f}_{c}\) 为其特征函数，对 \(c\) 取值 1，在 \(\mathrm{G} - c\) 上取值 0。函数 \({\Psi }^{n}{f}_{c}\) 由下式给出:

\[
{\Psi }^{n}{f}_{c}\left( x\right)  = \left\{  \begin{array}{ll} 1 & \text{ if }{x}^{n} \in  c \\  0 & \text{ otherwise }. \end{array}\right.
\]

每个取值于 \(\mathbf{A}\) 的类函数都是那些 \({f}_{c}\) 的线性组合。因此定理 23 等价于:

定理23’。对于每个 \(\mathrm{G}\) 的共轭类 \(c\) ，函数 \(\left( {g/\left( {g,n}\right) }\right) {\Psi }^{n}{f}_{c}\) 属于 \(\mathrm{A} \otimes  \mathrm{R}\left( \mathrm{G}\right)\) 。

这可以用另一种方式表述:

定理23"。对于每个 \(\mathrm{G}\) 的共轭类 \(c\) ，以及 \(\mathrm{G}\) 的每个特征 \(\chi\) ，我们有 \(1/\left( {g,n}\right) \mathop{\sum }\limits_{{{x}^{n} \in  c}}\chi \left( x\right)  \in  \mathrm{A}\) 。

当将单位特征取为 \(\chi\) 时，得到:

推论1。满足 \({x}^{n} \in  c\) 是 \(\left( {g,n}\right)\) 的倍数的元素数量为 \(x \in  \mathrm{G}\) 。

特别地:

推论2。如果 \(n\) 整除 \(\mathrm{G}\) 的阶，则满足 \({x}^{n} = 1\) 的 \(x \in  \mathrm{G}\) 的个数是 \(n\) 的倍数。

(此处提及弗罗贝尼乌斯的一个猜想:若集合 \({\mathrm{G}}_{n}\) 为那些 \(s \in  \mathrm{G}\) 使得 \({s}^{n} = 1\) 有 \(n\) 个元素的元素集合，则 \({\mathrm{G}}_{n}\) 是 \(\mathrm{G}\) 的子群。)

定理23的证明。(R. Brauer。)由定理21，现只需证明将函数 \(\left( {g/\left( {g,n}\right) }\right) {\Psi }^{n}f\) 限制到每个 \(G\) 的初等子群 \(H\) 上属于 \(A \otimes  R\left( H\right)\) 。现在，若 \(h\) 是 \(H\) 的阶，则 \(g/\left( {g,n}\right)\) 可被 \(h/\left( {h,n}\right)\) 整除。所以只需证明

\[
\frac{h}{\left( h,n\right) }{\Psi }^{n}\left( {{\operatorname{Res}}_{\mathrm{H}}f}\right)
\]

属于 \(\mathrm{A} \otimes  \mathrm{R}\left( \mathrm{H}\right)\) ，即证明归结为初等群的情形。由于初等群是 \(p\) -群的直积，足以处理一 个 \(p\) -群的情形。利用此类群的不可约特征由一维特征所诱导这一事实，最终我们被引导去证明如下命题:

引理10。设 \(\mathrm{G}\) 为 p 群， \(c\) 为其共轭类， \(\chi\) 为 \(\mathrm{G}\) 的一维特征，且设 \({a}_{c} = \mathop{\sum }\limits_{{{x}^{n} \in  c}}\chi \left( x\right)\) 。则 \({a}_{c} \equiv  0\left( {\operatorname{mod}.\left( {g,n}\right) \mathrm{A}}\right)\) 。

首先注意，对于固定的 \(\chi\) 且变动的 \(c\) ，这些 \({a}_{c}\) 的和等于 \(\mathop{\sum }\limits_{{x \in  \mathrm{G}}}\chi \left( x\right)\) ，即当 \(\chi  = 1\) 时为 \(g\) ，否则为 0。因此

\[
\mathop{\sum }\limits_{c}{a}_{c} \equiv  0\left( {{\;\operatorname{mod}\;.}\left( {g,n}\right) }\right) .
\]

因此只需对那些不同于单位类的类 \(c\) 证明引理10。

将 \(n\) 写成 \({p}^{a}m\) 的形式，且满足 \(\left( {p,m}\right)  = 1\) 。设 \({p}^{b}\) 为 \(c\) 元素的阶，设 \(C\) 为满足 \(x \in  G\) 的 \({x}^{n} \in  c\) 的集合。由于 \({x}^{n} = {x}^{{p}^{a}m}\) 的阶为 \({p}^{b} > 1\) ，且由于 \(G\) 是一个 \(p\) -群，故 \(x\) 的阶为 \({p}^{a + b}\) 。由此可得:若 \(z\) 为整数 \(\equiv  1\left( {\;\operatorname{mod}\;.{p}^{b}}\right)\) ，则 \({\left( {x}^{z}\right) }^{n} = {x}^{n}\) ，从而 \(x \in  \mathbf{C}\) ；此外，当且仅当 \(z \equiv  1\left( {{\;\operatorname{mod}\;.}{p}^{a + b}}\right)\) 时等号 \({x}^{z} = x\) 成立。换言之，由与 \(1{\;\operatorname{mod}\;{p}^{b}}\) 同余的元素组成的子群 \(\Gamma\) 在 \({\left( \mathbf{Z}/{p}^{a + b}\mathbf{Z}\right) }^{ * }\) 上自由作用。现在集合 \(\mathbf{C}\) 在 \(\Gamma\) 的作用下被划分为若干轨道，且只需证明每个轨道上 \(\chi \left( x\right)\) 的和在环 A 中可被 \(\left( {g,n}\right)\) 整除。这样的一个轨道由若干元素 \({x}^{1 + {p}^{b}t}\) 组成，且 \(t \in  \mathbf{Z}/{p}^{a}\mathbf{Z}\) 。因此 \(\chi\) 在该轨道上取值的和等于

\[
{a}_{c}\left( x\right)  = \chi \left( x\right) \mathop{\sum }\limits_{{t{\;\operatorname{mod}\;{p}^{a}}}}{z}^{t},\;\text{ where }z = \chi \left( {x}^{{p}^{b}}\right) .
\]

但 \(\chi \left( x\right)\) 是一个 \({p}^{a + b}\) 次单位根，且 \(z\) 是一个 \({p}^{a}\) 次单位根。因此

\[
\mathop{\sum }\limits_{{t{\;\operatorname{mod}\;{p}^{a}}}}{z}^{t} = \left\{  \begin{array}{ll} {p}^{a} & \text{ if }z = 1, \\  0 & \text{ if }z \neq  1. \end{array}\right.
\]

因此 \({a}_{c}\left( x\right)\) 可被 \({p}^{a}\) 整除，更进一步可被 \(\left( {g,n}\right)\) 整除。

EXERCISE

11.1. 设 \(f\) 为定义在 \(\mathbf{G}\) 上、取值于 \(\mathbf{Q}\) 的类函数，且对于所有与 \(g\) 互素的 \(m\) 有 \(f\left( {x}^{m}\right)  = f\left( x\right)\) 。证明 \(f\) 属于 \(\mathbf{Q} \otimes  \mathrm{R}\left( \mathrm{G}\right)\) [用定理 \(2{1}^{\prime }\) 化简到循环情形]。由定理 23 推断:若此外 \(f\) 的取值属于 \(\mathbf{Z}\) ，则函数 \(\left( {g/\left( {g,n}\right) }\right) {\Psi }^{n}f\) 属于 \(\mathbf{R}\left( \mathbf{G}\right)\) 。将此结论应用于单位类的特征函数。

\section*{11.3 对布劳尔定理的一个逆定理}

字母 \(\mathbf{A}\) 与 \(g\) 在本节中与前一节含义相同。

引理11。设 \(p\) 为素数。设 \(x\) 为 \({p}^{\prime }\) -元的 \(\mathbf{G},\mathbf{C}\) 子群由 \(x\) 所生成，且令 \(\mathrm{P}\) 为中央化子 \(\mathbf{Z}\left( x\right)\) 在 \(\mathrm{G}\) 中的一个 Sylow p-子群。设 \(\mathrm{H}\) 为包含不含任何 \(\mathrm{C} \times  \mathrm{P}\) 共轭元的 \(\mathrm{G}\) 的子群，设 \(\psi\) 为定义在 \(\mathrm{H}\) 上取值于 \(\mathrm{A}\) 的类函数，且设 \({\psi }^{\prime } = {\operatorname{Ind}}_{\mathrm{H}}^{\mathrm{G}}\psi\) 。则 \({\psi }^{\prime }\left( x\right)  \equiv  0\left( {{\;\operatorname{mod}\;.}p\mathrm{\;A}}\right) .\)

令 \(\mathrm{S}\left( x\right)\) 为 \(x\) 的所有共轭元的集合。则

\[
{\psi }^{\prime }\left( x\right)  = \frac{\operatorname{Card}\mathrm{Z}\left( x\right) }{\operatorname{Card}\mathrm{H}}\mathop{\sum }\limits_{{y \in  \mathrm{S}\left( x\right)  \cap  \mathrm{H}}}\psi \left( y\right) .
\]

令 \({\left( {Y}_{i}\right) }_{i \in  I}\) 为包含于 \(S\left( x\right)  \cap  H\) 中的不同 \(H\) -共轭类，并在每个 \({Y}_{i}\) 中选取一个元素 \({y}_{i}\) 。元素 \({y}_{i}\) 在 \(H\) 中的共轭数等于 \(\operatorname{Card}{Y}_{i}\) ，亦等于 \(\left( {H : H \cap  Z\left( {y}_{i}\right) }\right)\) 。因此

\[
{\psi }^{\prime }\left( x\right)  = \frac{\operatorname{Card}Z\left( x\right) }{\operatorname{Card}H}\mathop{\sum }\limits_{{i \in  I}}\operatorname{Card}{Y}_{i} \cdot  \psi \left( {y}_{i}\right) ,
\]

\[
= \mathop{\sum }\limits_{{i \in  \mathrm{I}}}{n}_{i}\psi \left( {y}_{i}\right) ,\;\text{ with }{n}_{i} = \frac{\operatorname{Card}\mathrm{Z}\left( {y}_{i}\right) }{\operatorname{Card}\left( {\mathrm{H} \cap  \mathrm{Z}\left( {y}_{i}\right) }\right) }.
\]

假定对某个 \(i \in  \mathrm{I}\) 我们有 \({n}_{i} \text{ ≢ } 0\left( {\;\operatorname{mod}\;.p}\right)\) 。则 \(\operatorname{Card}Z\left( {y}_{i}\right)\) 与 \(\operatorname{Card}\left( {\mathrm{H} \cap  \mathrm{Z}\left( {y}_{i}\right) }\right)\) 可被同一 \(p\) 的幂整除；故 \(\mathrm{H} \cap  \mathrm{Z}\left( {y}_{i}\right)\) 的一个 Sylow \(p\) -子群 \({\mathrm{P}}_{i}\) 也是 \(\mathrm{Z}\left( {y}_{i}\right)\) 的一个 Sylow \(p\) -子群。若 \({\mathrm{C}}_{i}\) 为由 \({y}_{i}\) 所生成的循环群，则 \({\mathrm{C}}_{i} \times  {\mathrm{P}}_{i}\) 包含于 \(\mathrm{H}\) ，并且是群 G 中与 \({y}_{i}\) 相关的一个 \(p\) -元子群。由于 \({y}_{i}\) 与 \(x\) 在 \(\mathrm{G}\) 中共轭，群 \({\mathrm{C}}_{i} \times  {\mathrm{P}}_{i}\) 与 \(\mathrm{C} \times  \mathrm{P}\) 共轭。这与对 \(\mathrm{H}\) 的假设矛盾。于是对所有 \(i\) 有 \({n}_{i} \equiv  0\left( {\;\operatorname{mod}\;.p}\right)\) ，从而

\[
{\psi }^{\prime }\left( x\right)  \equiv  0\left( {{\;\operatorname{mod}\;.}p\mathrm{\;A}}\right) .
\]

定理23(J. Green)。设 \({\left( {\mathrm{H}}_{i}\right) }_{i \in  \mathrm{I}}\) 为 \(\mathrm{G}\) 的一族子群，满足 \(\mathrm{R}\left( \mathrm{G}\right)  = \mathop{\sum }\limits_{{i \in  \mathrm{I}}}\operatorname{Ind}\mathrm{R}\left( {\mathrm{H}}_{i}\right)\) 。则每个 \(\mathrm{G}\) 的元子群都包含于某个 \({\mathrm{H}}_{i}\) 的共轭子中。

设 \(\mathrm{C} \times  \mathrm{P}\) 为 \(p\) -初等子群。我们可假定该子群是极大的，因而与 \(\mathrm{G}\) 的一个 \({p}^{\prime }\) -素元 \(x\) 相关联。若 \(\mathrm{C} \times  \mathrm{P}\) 不包含于任何 \({\mathrm{H}}_{i}\) 的共轭子中，前述引理将对所有 \(\chi  \in  \sum\) Ind \(\mathrm{R}\left( {\mathrm{H}}_{i}\right)\) 证明 \(\chi \left( x\right)  \equiv  0\left( {{\;\operatorname{mod.}\;}p\mathrm{\;A}}\right)\) ，尤其当 \(\chi\) 为 \(\mathbf{G}\) 的单位特征时，这将导致矛盾。

换言之，初等子群的族是使布劳尔定理成立的“最小”族。

\section*{11.4 \(A \otimes  R\left( G\right)\) 的谱}

回想若 \(\mathbf{C}\) 为交换环，则其谱 Spec(C) 为 C 的素理想集，参见 Bourbaki, Alg. Comm., Ch. II。

我们要确定环 \(A \otimes  R\left( G\right)\) 的谱。(我们也可描述 \(R\left( G\right)\) 的谱，但那更为复杂。)

设 \({Cl}\left( G\right)\) 为 \(G\) 的共轭类集。环 \({A}^{{Cl}\left( G\right) }\) 可被识别为取值于 \(A\) 的关于 \(G\) 的类函数环；若 \(f\) 属于此环，且 \(c\) 是一共轭类，则令 \(f\left( c\right)\) 表示 \(f\) 在任一属于 \(c\) 的元上的值。注入映射 \(A \rightarrow  A \otimes  R\left( G\right)  \rightarrow  {A}^{{Cl}\left( G\right) }\) 给出映射

\[
\operatorname{Spec}\left( {\mathrm{A}}^{\mathrm{{Cl}}\left( \mathrm{G}\right) }\right)  \rightarrow  \operatorname{Spec}\left( {\mathrm{A} \otimes  \mathrm{R}\left( \mathrm{G}\right) }\right)  \rightarrow  \operatorname{Spec}\left( \mathrm{A}\right) .
\]

这些映射是满射；例如，这可由 \({\mathrm{A}}^{\mathrm{{Cl}}\left( \mathrm{G}\right) }\) 在 \(\mathrm{A}\) (甚至在 \(\mathbf{Z}\) )上整性导出，参见 Bourbaki, Alg. Comm., Ch. IV, §2。

另一方面，我们知道 Spec(A) 由零理想和 A 的极大理想组成。此外，若 M 在 A 中为极大理想，则商域 A/M 为有限域；其特征称为 M 的剩余特征。

\({A}^{{Cl}\left( G\right) }\) 的谱可被识别为 \({Cl}\left( G\right)  \times\) Spec \(\left( A\right)\) :对每个 \(c \in  {Cl}\left( G\right)\) 与每个 \(M \in  \operatorname{Spec}\left( A\right)\) ，我们关联素理想 \({M}_{c}\) ，其由所有满足 \(f\left( c\right)  \in  M\) 的 \(f \in  {A}^{{Cl}\left( G\right) }\) 构成。 \({M}_{c}\) 在 Spec \(\left( {A \otimes  R\left( G\right) }\right)\) 中的像是素理想 \({P}_{M,c} = {M}_{c} \cap  \left( {A \otimes  R\left( G\right) }\right)\) 。

命题 30. 如果

(i) 对每个类 \(c \in  \mathrm{{Cl}}\left( \mathrm{G}\right)\) 我们关联 \({\mathrm{P}}_{0,c}\) ,

(ii) 对每个 p-正则类 \(c\) 以及每个剩余特征为 \(p\) 的极大理想 \(\mathbf{M}\) of \(\mathbf{A}\) 我们关联 \({\mathrm{P}}_{\mathbf{M},c}\) ,

则我们恰好获得每个 \(\mathrm{A} \otimes  \mathrm{R}\left( \mathrm{G}\right)\) 的素理想一次且仅一次。

(若共轭类由 \({p}^{\prime }\) 元素组成，则称该共轭类为 \(p\) -正则，参见 10.1。)

由于 \(\operatorname{Spec}\left( {A}^{\mathrm{{Cl}}\left( G\right) }\right)  \rightarrow  \operatorname{Spec}\left( {A \otimes  R\left( G\right) }\right)\) 是满射(参见上文)，每个 \(A \otimes  R\left( G\right)\) 的素理想 \(\mathfrak{p}\) 都为某个 \({P}_{M,c}\) 的形式；由于 \(\mathfrak{p} \cap  A\) 是 \(M\) ，我们看到 \(\mathfrak{p}\) 决定了 \(M\) ，余下的只是确定哪些类对 \({c}_{1}\) 和 \({c}_{2}\) 满足 \({\mathrm{P}}_{\mathrm{M},{c}_{1}} = {\mathrm{P}}_{\mathrm{M},{c}_{2}}\) 。因此该命题由下列命题推出:

命题 \({30}^{\prime }\) 。

(i) 若 \(\mathrm{M} = 0,{\mathrm{P}}_{0,{c}_{1}} = {\mathrm{P}}_{0,{c}_{2}}\) 等价于 \({c}_{1} = {c}_{2}\) 。

(ii) 设 \(\mathrm{M} \neq  0\) 的剩余特征为 \(p\) 。令 \({c}_{1}^{\prime }\) (或 \({c}_{2}^{\prime }\) )为由 \({c}_{1}\) (或 \({c}_{2}\) ) 的元素的 \(p\) -素部分组成的类。则 \({\mathrm{P}}_{\mathrm{M},{c}_{2}} = {\mathrm{P}}_{\mathrm{M},{c}_{2}}\) 等价于 \({c}_{1}^{\prime } = {c}_{2}^{\prime }\) 。

\begin{center}
\includegraphics[max width=0.6\textwidth]{images/bo_d6kn08k601uc73e1e9r0_94_487_300_662_446_0.jpg}
\end{center}
\hspace*{3em} 

为证明 (i) 我们须表明，若 \({c}_{1} \neq  {c}_{2}\) ，则存在元素 \(f \in  \mathrm{A} \otimes  \mathrm{R}\left( \mathrm{G}\right)\) 使得 \(f\left( {c}_{1}\right)  \neq  0\) 与 \(f\left( {c}_{2}\right)  = 0\) ，这显而易见(取 \(f\) 为在 \({c}_{1}\) 上恒等于 \(g\) 而在其它处为 0 的函数)。

若 \(M\) 的特征为 \(p\) ，一个类似于引理 7 证明的简单论证表明 \({\mathrm{P}}_{\mathbf{M},{c}_{1}} = {\mathrm{P}}_{\mathbf{M},{c}_{1}^{\prime }}\) (参见习题 10.4)。另一方面，引理 8 显示若 \({c}_{1}^{\prime } \neq  {c}_{2}^{\prime }\) 则 \({\mathbf{P}}_{\mathbf{M},{c}_{1}^{\prime }} \neq  {\mathbf{P}}_{\mathbf{M},{c}_{2}^{\prime }}\) 。由此得出 (ii)。

\section*{备注}

(1) 设 I 是 \(A \otimes  R\left( G\right)\) 的一个理想。要证明 I 等于 \(A \otimes  R\left( G\right)\) ，只需证明 I 不包含于任何素理想 \({P}_{M,c}\) ；这是 Brauer 定理证明中采用的方法(见下文练习 11.7)。

(2) 我们可以将 Spec \(\left( {A \otimes  R\left( G\right) }\right)\) 在图上表示为若干“直线”的并集 \({\mathrm{D}}_{c}\) ，对应于各个类 \(c\) ，每一条这些直线代表 Spec(A)。这些直线以如下方式“相交”: \({\mathrm{D}}_{{c}_{1}}\) 和 \({\mathrm{D}}_{c}\) 在某极大理想 \(\mathbf{M}\) (属于 \(\mathbf{A}\) )之上的一个共同点处相交，且该极大理想的剩余类特征为 \(p\) 当且仅当 \({p}^{\prime }\) -分量的 \({c}_{1}\) 和 \({c}_{2}\) 相等。

命题 31. 在 Zariski 拓扑下 \(\operatorname{Spec}\left( {\mathrm{A} \otimes  \mathrm{R}\left( \mathrm{G}\right) }\right)\) 是连通的。

(若 \(\mathrm{C}\) 是交换环，则 \(\operatorname{Spec}\left( \mathrm{C}\right)\) 的子集 \(\mathrm{F}\) 在 Zariski 拓扑下为闭集当且仅当存在 \(\mathrm{H} \subset  \mathrm{C}\) 使得 \(\mathfrak{p} \in  \mathrm{F} \Leftrightarrow  \mathfrak{p} \supset  \mathrm{H}\) 。)

设 \(x\) 是 \(\mathrm{G}\) 的一个元素，阶为 \({p}_{1}^{{n}_{1}} \cdot  {p}_{2}^{{n}_{2}}\cdots {p}_{k}^{{n}_{k}};x\) 可分解为积 \(x = {x}_{{p}_{1}} \cdot  {x}_{{p}_{2}}\cdots {x}_{{p}_{k}}\) ，其中 \({x}_{{p}_{i}}\) 的阶为 \({p}_{i}^{{n}_{i}}\) 。与 \(x\) 和 \({x}_{{p}_{2}}\cdots {x}_{{p}_{k}}\) 相关的类具有相同的 \({p}_{1}\) -整常分量。因此， \(\operatorname{Spec}\left( {A \otimes  R\left( G\right) }\right)\) 的相应“直线”相交；此外，每条直线同构于 Spec(A)，因此是连通的。逐步推进直至单位元，可见 Spec(A ⊗ R(G)) 是连通的。

推论。Spec R(G) 是连通的。

确实，这是 \(\operatorname{Spec}\left( {A \otimes  R\left( G\right) }\right)\) 在一个连续映射下的像。

例。取 \(G\) 为对称群 \({\sigma }_{3}\) 。共有三类: \(1,{c}_{2}\) (由阶为 2 的元素组成)，以及 \({c}_{3}\) (阶为 3 的元素)。在 \(A\) 中存在唯一一个剩余特征为 2 的素理想 \({\mathfrak{p}}_{2}\) ，对特征 3 亦然。 \(A \otimes  R\left( G\right)\) 的谱由三条如下面所示相交的“直线”组成:

\begin{center}
\includegraphics[max width=0.6\textwidth]{images/bo_d6kn08k601uc73e1e9r0_95_433_705_725_418_0.jpg}
\end{center}
\hspace*{3em} 

注。该节的结果已由 G. Segal 推广到紧李群(Publ. Math. I.H.E.S., 34, 1968)。

练习

11.2. 证明 \({\mathrm{P}}_{\mathrm{M},c}\) 的剩余域是 \(\mathrm{A}/\mathrm{M}\) 。

11.3. 如果 \(B\) 是一个 A-代数，按照命题 30 和 \({30}^{\prime }\) 的证明，用 \(\operatorname{Spec}\left( B\right)\) 表示出 \(\operatorname{Spec}\left( {B \otimes  R\left( G\right) }\right)\) 。

11.4. 设 \(\mathrm{K}\) 为 \(\mathrm{A}\) 的分式域， \(\Gamma\) 为 \(\mathrm{K}/\mathrm{Q}\) 的伽罗瓦群。我们知道 \(\Gamma\) 同构于 \({\left( \mathbf{Z}/g\mathbf{Z}\right) }^{ * }\) 。令 \(\Gamma\) 通过其在 A 上的作用作用于 \(\mathrm{A} \otimes  \mathrm{R}\left( \mathrm{G}\right)\) ，并确定它在 \(\operatorname{Spec}\left( {\mathrm{A} \otimes  \mathrm{R}\left( \mathrm{G}\right) }\right)\) 上相应的作用。通过观察 \(R\left( G\right)\) 是由 \(\Gamma\) 不动的 \(A \otimes  R\left( G\right)\) 的子环，得到 Spec \(\left( {R\left( G\right) }\right)\) 。

11.5. 当 \(G\) 为交换群时确定 Spec \(\left( {A\left\lbrack  G\right\rbrack  }\right)\) (注意可将 \(A\left\lbrack  G\right\rbrack\) 识别为 \(A \otimes  R\left( \widehat{G}\right)\) ，其中 \(\widehat{G}\) 为 \(G\) 的对偶，参见习题 3.3)。

11.6. 令 \(B\) 为 \({A}^{{Cl}\left( G\right) }\) 的子环，包含那些函数 \(f\) ，满足:对任一带有剩余特征 \(p\) 的极大理想 \(M\) ⊂ \(A\) ，以及任一具有 \(p\) -正则分量 \({c}^{\prime }\) 的类 \(c\) ，我们都有 \(f\left( c\right)  \equiv  f\left( {c}^{\prime }\right) \left( {{\;\operatorname{mod}\;.}\mathrm{M}}\right)\) 。证明 \(A \otimes  R\left( G\right)  \subset  B\) ，并证明这两个环具有相同的谱；给出一个它们不同的例子。

11.7. 令 \(H\) 为 \(G\) 的一个子群，且令 \({I}_{H}\) 为 \(A \otimes  R\left( G\right)\) 中由 \(A \otimes  {\operatorname{Ind}}_{\mathrm{H}}^{G}\) 映入的理想。

(a) 令 \(c\) 为 \(G\) 的一个类。证明当且仅当 \(\mathbf{H} \cap  c = \varnothing\) 时， \({I}_{H}\) 包含于 \({P}_{0,c}\) 。

(b) 令 \(c\) 为一个 \(p\) -正则类，且令 \(\mathbf{M}\) 为包含 \(p\) 的 \(\mathbf{A}\) 中的 \(p\) -正则理想。证明当且仅当 \(\mathrm{H}\) 不含与 \(c\) 中元素相关的任何 \(p\) -初等子群时， \({\mathrm{I}}_{\mathrm{H}}\) 包含于 \({\mathrm{P}}_{\mathrm{M},c}\) 。

(c) 从 (b) 得出定理 18 和 23 的另一种证明。

\section*{第 12 章 有理性问题}

到目前为止我们只研究了在复数域 \(\mathbf{C}\) 上定义的表示。事实上，前面各节的所有证明在特征零的代闭域上仍然成立，例如 \(\mathbf{Q}\) 的一个代数闭包。现在我们来考察当域非代闭时会发生什么。

\section*{12.1 环 \({R}_{K}\left( G\right)\) 和 \({\bar{R}}_{K}\left( G\right)\)}

在本节中， \(\mathrm{K}\) 表示特征零的域， \(\mathrm{C}\) 表示 K 的代数闭包。如果 \(V\) 是一个 K-向量空间，我们令 \({V}_{C}\) 表示通过将标量从 K 扩张到 C 而从 V 得到的 C-向量空间 C \({ \otimes  }_{\mathrm{K}}\) V。如果 G 是一个有限群，则在域 \(\mathrm{K}\) 上的每个线性表示 \(\rho  : \mathrm{G} \rightarrow  \mathbf{{GL}}\left( \mathrm{V}\right)\) 都定义了一个表示

\[
{\rho }_{\mathrm{C}} : \mathrm{G} \rightarrow  \mathbf{{GL}}\left( \mathrm{V}\right)  \rightarrow  \mathbf{{GL}}\left( {\mathrm{V}}_{\mathrm{C}}\right)
\]

在域 \(C\) 上。用“模”的术语(参见 6.1)，我们有

\[
{\mathrm{V}}_{\mathrm{C}} = \mathrm{C}\left\lbrack  \mathrm{G}\right\rbrack  { \otimes  }_{\mathrm{K}\left\lbrack  \mathrm{G}\right\rbrack  }\mathrm{V}.
\]

\(\rho\) 的表象 \({\chi }_{\rho } = \operatorname{Tr}\left( \rho \right)\) 与 \({\rho }_{\mathrm{C}}\) 的相同；它是定义在 \(\mathbf{G}\) 上、取值于 \(\mathbf{K}\) 的类函数。

我们用 \({R}_{K}\left( G\right)\) 表示由在 \(G\) 上、以 \(K\) 为系数的表示的特征生成的群；它是第 9、10、11 章中研究的环 \(R\left( G\right)  = {R}_{C}\left( G\right)\) 的一个子环。

我们也可以将 \({\mathrm{R}}_{\mathrm{K}}\left( \mathrm{G}\right)\) 定义为有限型 K[G]-模范畴的格罗滕迪克群，参见第三部分，第14章。

命题 32. 设 \(\left( {{\mathrm{V}}_{i},{\rho }_{i}}\right)\) 为在 \(\mathrm{G}\) 上以 \(\mathrm{K}\) 为系数的不重复(按同构)不可约线性表示，且令 \({\chi }_{i}\) 为相应的特征。则

(a) 这些 \({\chi }_{i}\) 构成 \({\mathrm{R}}_{\mathrm{K}}\left( \mathrm{G}\right)\) 的一组基。

(b) 这些 \({\chi }_{i}\) 两两正交。

[和往常一样，此处所指的正交是相对于双线性型 \(\langle \varphi ,\chi \rangle  = \left( {1/g}\right) \mathop{\sum }\limits_{{s \in  \mathrm{G}}}\varphi \left( {s}^{-1}\right) \chi \left( s\right)\) 的正交。]

显然这些 \({\chi }_{i}\) 生出 \({\mathrm{R}}_{\mathrm{K}}\left( \mathrm{G}\right)\) 。另一方面，如果 \(i \neq  j\) 我们有 \({\operatorname{Hom}}^{G}\left( {{\mathrm{\;V}}_{i},{\mathrm{\;V}}_{j}}\right)  = 0\) 。但在一般情况下，若 \(\mathrm{V}\) 和 \(\mathrm{W}\) 的特征分别为 \({\chi }_{\mathrm{V}}\) 和 \({\chi }_{\mathrm{W}}\) ，我们有

\[
{\dim }_{\mathrm{K}}{\operatorname{Hom}}^{\mathrm{G}}\left( {\mathrm{V},\mathrm{W}}\right)  = {\dim }_{\mathrm{C}}{\operatorname{Hom}}^{\mathrm{G}}\left( {{\mathrm{V}}_{\mathrm{C}},{\mathrm{W}}_{\mathrm{C}}}\right)  = \left\langle  {{\chi }_{\mathbf{V}},{\chi }_{\mathbf{W}}}\right\rangle  ,
\]

参见 7.2，引理 2。由此可得若 \(\left\langle  {{\chi }_{i},{\chi }_{j}}\right\rangle   = 0\) 则 \(i \neq  j\) ，并且 \(\left\langle  {{\chi }_{i},{\chi }_{i}}\right\rangle \; = \dim {\operatorname{End}}^{G}\left( {V}_{i}\right)\) 是一个整数 \(\geq  1\) (当且仅当 \({V}_{C}\) 是不可约的，即当 V 在绝对意义上不可约时等于 1，参见 Bourbaki [8], §13, no. 4)。特别地，这些 \({\chi }_{i}\) 是线性无关的。

若将 \(\mathrm{G}\) 在 \(\mathrm{C}\) 上的线性表示与某个形如 \({\rho }_{\mathrm{C}}\) 的表示同构，并且 \(\rho\) 是 \(\mathrm{G}\) 在 \(\mathrm{K}\) 上的线性表示，则称该表示可在 \(\mathrm{K}\) 上实现(或在 \(\mathrm{K}\) 上有理)；这相当于说可由系数属于 \(\mathrm{K}\) 的矩阵来实现。

命题33. 要使 \(\mathrm{G}\) 在 \(\mathrm{C}\) 上的线性表示可在 \(\mathrm{K}\) 上实现，当且仅当其特征属于 \({\mathrm{R}}_{\mathrm{K}}\left( \mathrm{G}\right)\) 。

该条件显然是必要的。反过来，设给定表示的特征为 \(\chi\) 且该条件成立。根据命题32，我们有 \(\chi  = \sum {n}_{i}{\chi }_{i}\) ，其中 \({n}_{i} \in  \mathbf{Z}\) ，由此得到:

\[
\left\langle  {\chi ,{\chi }_{i}}\right\rangle   = {n}_{i}\left\langle  {{\chi }_{i},{\chi }_{i}}\right\rangle  \;\text{ for all }i.
\]

由于 \(\chi\) 是 \(\mathrm{G}\) 在 \(\mathrm{C}\) 上某表示的特征，标量乘积 \(\left\langle  {\chi ,{\chi }_{i}}\right\rangle\) 等于 \(\geq  0\) 。因此 \({n}_{i}\) 为正，给定表示可实现为若干 \({V}_{i}\) 的直和，每个重复 \({n}_{i}\) 次。

同样的论证表明，该实现在 K-同构意义下是唯一的。

除环 \({R}_{K}\left( G\right)\) 外，我们还考虑子环 \({\bar{R}}_{K}\left( G\right)\) ，由那些取值在 \(K\) 的 \(R\left( G\right)\) 元素组成。显然， \({\mathbf{R}}_{\mathbf{K}}\left( \mathbf{G}\right)  \subset  {\overline{\mathbf{R}}}_{\mathbf{K}}\left( \mathbf{G}\right)\) 。此外:

命题34. 群 \({\mathrm{R}}_{\mathrm{K}}\left( \mathrm{G}\right)\) 在 \({\overline{\mathrm{R}}}_{\mathrm{K}}\left( \mathrm{G}\right)\) 中的指数有限。

首先注意，每个在 C 上不可约的 G 表示都能在某个由相应矩阵表示系数生成的 \(K\) 的有限扩张上实现。因此存在 K 的有限扩张 L，使得 \({R}_{L}\left( G\right)  = R\left( G\right)\) 。设该扩张的次数为 \(d = \left\lbrack  {L : K}\right\rbrack\) ；由下述引理可得命题成立:

引理12. 我们有 \(d \cdot  {\overline{\mathbf{R}}}_{\mathbf{K}}\left( \mathbf{G}\right)  \subset  {\mathbf{R}}_{\mathbf{K}}\left( \mathbf{G}\right)\) 。

首先，让 \(\mathbf{V}\) 成为在 \(\mathbf{L}\) 上的 \(\mathbf{G}\) 的线性表示，特征为 \(\chi\) ；通过限制标量，我们可以将 \(\mathbf{V}\) 视为一个 K-向量空间(维数增大了 \(d\) 倍)，甚至作为在 \(\mathbf{K}\) 上的 \(\mathbf{G}\) 的线性表示。我们立刻看到该表示的特征等于 \({\operatorname{Tr}}_{\mathrm{L}/\mathrm{K}}\left( \chi \right)\) ，其中 \({\operatorname{Tr}}_{\mathrm{L}/\mathrm{K}}\) 表示与扩张 \(\mathrm{L}/\mathrm{K}\) 相关的迹。由线性性可得，对于 \({\mathrm{R}}_{\mathrm{L}}\left( \mathrm{G}\right)\) 的每个元素 \(\chi\) ，都有 \({\operatorname{Tr}}_{\mathrm{L}/\mathrm{K}}\left( \chi \right)  \in  {\mathrm{R}}_{\mathrm{K}}\left( \mathrm{G}\right)\) 。

特别地，取 \(\chi  \in  {\overline{\mathbf{R}}}_{\mathbf{K}}\left( \mathrm{G}\right)\) ，即假设 \(\chi\) 的取值属于 \(\mathrm{K}\) 。于是 \({\operatorname{Tr}}_{\mathrm{L}/\mathrm{K}}\left( \chi \right)  = d \cdot  \chi\) ；因此 \(d \cdot  \chi  \in  {\mathrm{R}}_{\mathrm{K}}\left( \mathrm{G}\right)\) ，证毕。

\section*{12.2 Schur 指数}

前一节的结果也可以通过半单代数理论得到，甚至可以加以细化。我们简要勾勒如下:

代数 \(K\left\lbrack  G\right\rbrack\) 是若干简单代数 \({A}_{i}\) 的直积，对应于在 \(\mathbf{K}\) 上的不同不可约表示 \({\mathbf{V}}_{i}\) (对 \(\mathbf{G}\) )。若 \({\mathbf{D}}_{i} = {\operatorname{Hom}}^{\mathbf{G}}\left( {{\mathbf{V}}_{i}\text{ , }}\right. \; {\mathbf{V}}_{i}\) ) 是 \(\operatorname{End}\left( {\mathbf{V}}_{i}\right)\) 中与 \(\mathbf{G}\) 交换的代数，那么 \({\mathbf{D}}_{i}\) 是一个域(通常是非交换的)，且 \({\mathrm{A}}_{i}\) 可被识别为 \({\operatorname{End}}_{{\mathrm{D}}_{i}}\left( {\mathrm{V}}_{i}\right)\) ，即 \({D}_{i}\) -向量空间 \({V}_{i}\) 的端om代数。若 \(\left\lbrack  {{V}_{i} : {D}_{i}}\right\rbrack   = {n}_{i}\) ，则 \({A}_{i} \cong  {\mathbf{M}}_{n}\left( {D}_{i}^{ \circ  }\right)\) ，其中 \({D}_{i}^{ \circ  }\) 是 \({D}_{i}\) 的对合环。此外， \({D}_{i}\) 在其中心 \({K}_{i}\) 上的次数为完全平方数，记作 \({m}_{i}^{2}\) ；整数 \({m}_{i}\) 被称为表示 \({\mathrm{V}}_{i}\) (或分量 \({\mathrm{A}}_{i}\) )的 Schur 指数。

设 \(s \in  \mathrm{G}\) ，并令 \({\rho }_{i}\left( s\right)\) 为对应的 \({\mathrm{V}}_{i}\) 的端om。我们须考虑 \({\rho }_{i}\left( s\right)\) 的三类“迹”:

(a) 作为 K-端同态的迹；这是记作 \({\chi }_{i}\left( s\right)\) 的 K 中的元素；

(b) 作为 \({\mathrm{K}}_{i}\) -端同态的迹；这是 \({\mathrm{K}}_{i}\) 中的一个元素，我们记作 \({\varphi }_{i}\left( s\right)\) ；

(c) 在单代数 \({\mathrm{A}}_{i}\) 中的约化迹(参见例如 [8], no.12.3)；这是 \({\mathrm{K}}_{i}\) 中的一个元素，我们记作 \({\psi }_{i}\left( s\right)\) 。

各种迹由下列公式联系起来

\[
{\chi }_{i}\left( s\right)  = {\operatorname{Tr}}_{{\mathrm{K}}_{i}/\mathrm{K}}\left( {{\varphi }_{i}\left( s\right) }\right) \;\text{ and }\;{\varphi }_{i}\left( s\right)  = {m}_{i}{\psi }_{i}\left( s\right) .
\]

现在令 \({\sum }_{i}\) 为域 \({\mathbf{K}}_{i}\) 到代数闭域 C 的 K-同构的集合。若 \(\sigma  \in  {\sum }_{i}\) ，通过 \(\sigma\) 将标量从 \(\mathbf{K}\) 扩张到 \(\mathbf{C}\) 使得 \({\mathbf{D}}_{i}\) 变成矩阵代数 \({\mathbf{M}}_{{m}_{i}}\left( \mathbf{C}\right)\) ，并且 \({\mathbf{A}}_{i}\) 变为 \({\mathbf{M}}_{{n}_{i}{m}_{i}}\left( \mathbf{C}\right)\) 。复合 \(\mathbf{G} \rightarrow  {\mathbf{A}}_{i} \rightarrow  {\mathbf{M}}_{{n}_{i}{m}_{i}}\left( \mathbf{C}\right)\) ，我们得到在 \(\mathbf{C}\) 上的不可约表示，来自 \(\mathbf{G}\) ，度为 \({n}_{i}{m}_{i}\) ，其特征为 \({\psi }_{i,\sigma } = \sigma \left( {\psi }_{i}\right)\) 。对固定的 \(i\) ，这些特征 \({\psi }_{i,\sigma }\) 是共轭的: \(\mathbf{C}\) 对 \(\mathbf{K}\) 的伽罗瓦群对它们作传递置换。此外，G 在 \(\mathbf{C}\) 上的每个不可约特征都等于这 \({\psi }_{i,\sigma }\) 中的一个。我们有

\[
{\chi }_{i} = {\mathrm{{Tr}}}_{{\mathrm{K}}_{i}/\mathrm{K}}\left( {\varphi }_{i}\right)  = \mathop{\sum }\limits_{{\sigma  \in  {\sum }_{i}}}\sigma \left( {\varphi }_{i}\right)  = {m}_{i}\mathop{\sum }\limits_{{\sigma  \in  {\sum }_{i}}}{\psi }_{i,\sigma },
\]

由此给出 \({\chi }_{i}\) 在 C 上作为不可约特征和的分解。

现在令 \(\chi  = \mathop{\sum }\limits_{{i,\sigma }}{d}_{i,\sigma }{\psi }_{i,\sigma }\) 为 \(\mathrm{R}\left( \mathrm{G}\right)\) 中的一个元素，其中 \({d}_{i,\sigma }\) 为整数。为了使 \(\chi\) 取值于 \(\mathrm{K}\) ，必要且充分的条件是它在 \(\mathbf{C}\) 对 \(\mathbf{K}\) 的伽罗瓦群下不变，即这些 \({d}_{i,\sigma }\) 仅依赖于 \(i\) 。若确实如此，并令 \({d}_{i}\) 表示它们的公共值，我们有

\[
\chi  = \mathop{\sum }\limits_{i}{d}_{i}{\psi }_{i} = \mathop{\sum }\limits_{i}{d}_{i}{\chi }_{i}/{m}_{i}.
\]

因此我们有下述命题，它细化了命题 34:

命题 35. 这些特征 \({\psi }_{i} = {\chi }_{i}/{m}_{i}\) 构成 \({\overline{\mathbf{R}}}_{\mathbf{K}}\left( \mathrm{G}\right)\) 的一组基。

称 \(\mathrm{K}\left\lbrack  \mathrm{G}\right\rbrack\) 为拟分裂的，当且仅当 \({\mathrm{D}}_{i}\) 是交换的，或者等价地，当 Schur 指数 \({m}_{i}\) 都等于 1。于是命题 35 蕴含:

推论。为了使 \({\mathrm{R}}_{\mathrm{K}}\left( \mathrm{G}\right)  = {\overline{\mathbf{R}}}_{\mathrm{K}}\left( \mathrm{G}\right)\) ，必要且充分的条件是 K[G] 为拟分裂。

特别地，在下列每种情形中我们都有 \({R}_{K}\left( G\right)  = {\bar{R}}_{K}\left( G\right)\) :

(i) \(\mathrm{G}\) 是阿贝尔群(因为此时 \(\mathrm{K}\left\lbrack  \mathrm{G}\right\rbrack\) 与 \({\mathrm{D}}_{i}\) 是可交换的)。

(ii) 所有 \(\mathrm{K}\) 的有限扩张的布劳尔群均平凡。

\section*{习题}

12.1. 证明第5章所考虑的有限群的所有舒尔指标均等于1。

12.2. 取 \(G\) 为交错群 \({\mathfrak{A}}_{4}\) ，参见5.7。证明 Q[G] 分解为简单因子具有如下形式

\[
\mathbf{Q}\left\lbrack  \mathbf{G}\right\rbrack   = \mathbf{Q} \times  \mathbf{Q}\left( \omega \right)  \times  {\mathbf{M}}_{3}\left( \mathbf{Q}\right) ,
\]

其中 \(\mathbf{Q}\left( \omega \right)\) 是通过向 \(\mathbf{Q}\) 添加入一个三次单位根 \(\omega\) 而得到的 \(\mathbf{Q}\) 的二次扩张。

12.3. 取 \(\mathbf{G}\) 为四元数群 \(\{  \pm  1, \pm  i, \pm  j, \pm  k\}\) 。群 \(\mathbf{G}\) 有4个值域在 \(\{  \pm  1\}\) 的一维特征。另一方面，将 \(\mathbf{G}\) 自然嵌入到以 \(\mathbf{Q}\) 为基域的四元数除环 \({\mathbf{H}}_{\mathbf{O}}\) 中定义了一个满同态 \(\mathbf{Q}\left\lbrack  G\right\rbrack   \rightarrow  {\mathbf{H}}_{\mathbf{O}}\) 。证明 Q[G] 分解为简单分量为

\[
\mathbf{Q}\left\lbrack  \mathrm{G}\right\rbrack   = \mathbf{Q} \times  \mathbf{Q} \times  \mathbf{Q} \times  \mathbf{Q} \times  {\mathbf{H}}_{\mathbf{Q}}.
\]

最后一分量的舒尔指标等于2。对应的特征 \(\psi\) 由下式给出

\[
\psi \left( 1\right)  = 2,\;\psi \left( {-1}\right)  =  - 2,\;\psi \left( s\right)  = 0\;\text{ for }s \neq   \pm  1.
\]

因此 \(\mathbf{K}\left\lbrack  \mathbf{G}\right\rbrack\) 在若且唯若 \(\mathbf{K} \otimes  {\mathbf{H}}_{\mathbf{O}}\) 与 \({\mathbf{M}}_{2}\left( \mathbf{K}\right)\) 同构；证明这等价于在 \(\mathbf{K}\) 中 -1 可表示为两个平方和。

12.4. 证明舒尔指标 \({m}_{i}\) 整除中心 \(\mathbf{G}\) 的指数 \(a\) 。[注意带有特征 \({\psi }_{i,\sigma }\) 的不可约表示的度为 \({n}_{i}{m}_{i}\) ，并应用命题17。] 推出 \(a \cdot  {\overline{\mathbf{R}}}_{\mathbf{K}}\left( \mathbf{G}\right)\) 包含于 \({\mathbf{R}}_{\mathbf{K}}\left( \mathbf{G}\right)\) 中。

12.5. 设 \(\mathrm{L}\) 是 \(\mathrm{K}\) 的有限扩张。证明:若 \(\mathrm{L}\left\lbrack  \mathrm{G}\right\rbrack\) 为准分裂，则 \(\left\lbrack  {\mathrm{L} : \mathrm{K}}\right\rbrack\) 可被每一个舒尔指标 \({m}_{i}\) 整除。

\section*{12.3 在循环单位根域上的可实现性}

我们沿用前几节的记法，并以 \(m\) 表示 \(G\) 元素阶的最小公倍数；它是 \(g\) 的一个因子。

定理24(布劳尔)。若 \(\mathrm{K}\) 含有第 \(m\) 次单位根，则 \({\mathrm{R}}_{\mathrm{K}}\left( \mathrm{G}\right) \; = \mathrm{R}\left( \mathrm{G}\right)\) 。

结合命题33，可得:

推论。每一个 \(G\) 的线性表示都可以在 \(K\) 上实现。

(该结论曾由 Schur 猜想。)

设 \(\chi  \in  \mathrm{R}\left( \mathrm{G}\right)\) 。根据 10.5 的定理 20，我们可以把 \(\chi\) 写成如下形式

\[
\chi  = \sum {n}_{i}{\operatorname{Ind}}_{{\mathrm{H}}_{i}}^{\mathrm{G}}\left( {\varphi }_{i}\right) ,\;\left( {{n}_{i} \in  \mathbf{Z}}\right)
\]

其中 \({\varphi }_{i}\) 是子群 \({\mathrm{H}}_{i}\) 的一维字符。 \({\varphi }_{i}\) 的取值是 \(m\) 次单位根；它们属于 \(\mathbf{K}\) 。因此 \({\varphi }_{i} \in  {\mathbf{R}}_{K}\left( {\mathrm{H}}_{i}\right)\) 。但是，若 \(\mathrm{H}\) 是 \(\mathrm{G}\) 的子群，显然 \({\operatorname{Ind}}_{\mathrm{H}}^{\mathrm{G}}\) 将 \({\mathrm{R}}_{\mathrm{K}}\left( \mathrm{H}\right)\) 映入 \({\mathrm{R}}_{K}\left( \mathrm{G}\right)\) 。因此对所有 \(i\) 有 \({\operatorname{Ind}}_{{\mathrm{H}}_{i}}^{\mathbf{G}}\left( {\varphi }_{i}\right)  \in  {\mathrm{R}}_{\mathrm{K}}\left( \mathbf{G}\right)\) ，由此证明了该定理。

\section*{练习}

12.6. 证明 \(G\) 的 Schur 指数(在任意域上)整除欧拉函数 \(\varphi \left( m\right)\) [使用练习 12.5].

\section*{12.4 \({R}_{K}\left( G\right)\) 的秩}

我们现在回到特征为零的任意域 \(K\) 的情形。我们将确定 \({R}_{K}\left( G\right)\) 的秩，或等价地，确定 \(\mathrm{G}\) 在 \(\mathrm{K}\) 上不可约表示的个数。

选择一个整数 \(m\) ，它是 G 中各元素阶的倍数(例如，它们的最小公倍数或 \(\mathbf{G}\) 的阶 \(g\) )，并令 L 为通过在 \(\mathrm{K}\) 中添入第 \(m\) 个单位根得到的域。我们知道(参见例如 Bourbaki，Alg. V，§11)扩张 L/K 是 Galois 扩张，其 Galois 群 \(\mathrm{{Gal}}\left( {\mathrm{L}/\mathrm{K}}\right)\) 是可逆元乘法群 \({\left( \mathbf{Z}/m\mathbf{Z}\right) }^{ * }\) 的一个子群， \(\mathbf{Z}/m\mathbf{Z}\) 。更准确地说，若 \(\sigma  \in  \mathrm{{Gal}}\left( {\mathrm{L}/\mathrm{K}}\right)\) ，则存在唯一的元素 \(t \in  {\left( \mathbf{Z}/m\mathbf{Z}\right) }^{ * }\) 使得

\[
\sigma \left( \omega \right)  = {\omega }^{t}\;\text{ if }{\omega }^{m} = 1.
\]

我们以 \({\Gamma }_{\mathrm{K}}\) 表示 \(\mathrm{{Gal}}\left( {\mathrm{L}/\mathrm{K}}\right)\) 在 \({\left( \mathbf{Z}/m\mathbf{Z}\right) }^{ * }\) 中的像，若 \(t \in  {\Gamma }_{\mathrm{K}}\) ，则令 \({\sigma }_{t}\) 表示对应的 \(\mathrm{{Gal}}\left( {\mathrm{L}/\mathrm{K}}\right)\) 的元。前一节讨论的是 \({\Gamma }_{\mathrm{K}} = \{ 1\}\) 的情形。

设 \(s \in  \mathbf{G}\) ，并令 \(n\) 为整数。则 \({s}^{n}\) 在 \(\mathbf{G}\) 中的元只取决于 \(n\) 模 \(s\) 的同类，因此更不用说模 \(m\) ；特别地，对每个 \(t \in  {\Gamma }_{\mathrm{K}}\) 定义了 \({s}^{t}\) 。群 \({\Gamma }_{\mathrm{K}}\) 在 \(G\) 的底层集上作为置换群作用。若存在 \(t \in  {\Gamma }_{\mathrm{K}}\) 使得 \({s}^{\prime }\) 与 \({s}^{t}\) 被 \(G\) 的某一元共轭，则称 \(G\) 中的两个元 \(s,{s}^{\prime }\) 为 \({\Gamma }_{\mathrm{K}}\) -共轭。由此定义的关系是等价关系且不依赖于 \(m\) 的选择；其类称为 G 的 \({\Gamma }_{\mathrm{K}}\) -类(或 K-类)。

定理25. 一个取值于 \(\mathrm{L}\) 的类函数 \(f\) 在 \(\mathrm{G}\) 上属于 \(\mathrm{K}{ \otimes  }_{\mathbf{Z}}\mathrm{R}\left( \mathrm{G}\right)\) 的充要条件是

(*)
\[
{\sigma }_{t}\left( {f\left( s\right) }\right)  = f\left( {s}^{t}\right) \text{ for all }s \in  \mathrm{G}\text{ and all }t \in  {\Gamma }_{\mathrm{K}}.
\]

(换言之，我们对于所有 \(t \in  {\Gamma }_{\mathrm{K}}\) 必有 \({\sigma }_{t}\left( f\right)  = {\Psi }^{t}\left( f\right)\) ，参见11.2。)

设 \(\rho\) 为 \(G\) 的一表示，其特征是 \(\chi\) 。对 \(s \in  G\) ，算子 \(\rho \left( s\right)\) 的特征值 \({\omega }_{i}\) 是 \(m\) 次单位根，而 \(\rho \left( {s}^{t}\right)\) 的特征值是 \({\omega }_{i}^{t}\) 。因此我们有

\[
{\sigma }_{t}\left( {\chi \left( s\right) }\right)  = \sigma \left( {\sum {\omega }_{i}}\right)  = \sum {\omega }_{i}^{t} = \chi \left( {s}^{t}\right) ,
\]

这表明 \(\chi\) 满足条件(*)。由线性性，同样适用于 \(\mathrm{K} \otimes  \mathrm{R}\left( \mathrm{G}\right)\) 的所有元素。

反之，设 \(f\) 是在 \(\mathrm{G}\) 上满足条件(*) 的类函数。则

\[
f = \sum {c}_{\chi }\chi ,\;\text{ with }{c}_{\chi } = \langle f,\chi \rangle ,
\]

其中 \(\chi\) 遍历 G 的不可约特征的集合。我们须证明 \({c}_{\chi }\) 属于 \(\mathrm{K}\) ，根据伽罗瓦理论，这等价于证明它们在 \({\sigma }_{t},t \in  {\Gamma }_{\mathrm{K}}\) 下不变。但若 \(\varphi\) 与 \(\chi\) 是 \(\mathbf{G}\) 上的两个类函数，则有

\[
\left\langle  {{\Psi }^{t}\varphi ,{\Psi }^{t}\chi }\right\rangle   = \langle \varphi ,\chi \rangle ,
\]

这点可容易验证。故

\[
{c}_{\chi } = \langle f,\chi \rangle  = \left\langle  {{\Psi }^{t}f,{\Psi }^{t}\chi }\right\rangle   = \left\langle  {{\sigma }_{t}\left( f\right) ,{\sigma }_{t}\left( \chi \right) }\right\rangle   = {\sigma }_{t}\left( {\langle f,\chi \rangle }\right)  = {\sigma }_{t}\left( {c}_{\chi }\right) ,
\]

从而完成证明。

推论1. 一个取值于 \(\mathbf{K}\) 的类函数 \(f\) 在 \(\mathbf{G}\) 上属于 \(\mathbf{K} \otimes  {\mathbf{R}}_{\mathbf{K}}\left( \mathbf{G}\right)\) 的充要条件是它在 \(\mathrm{G}\) 的 \({\Gamma }_{\mathbf{K}}\) -类上为常值。

如果 \(f \in  \mathrm{K} \otimes  {\mathrm{R}}_{\mathrm{K}}\left( \mathrm{G}\right)\) ，则对所有 \(s \in  \mathrm{G}\) 有 \(f\left( s\right)  \in  \mathrm{K}\) ，并且公式 \(\left( *\right)\) 表明对所有 \(t \in  {\Gamma }_{\mathbf{K}}\) 有 \(f\left( s\right)  = f\left( {s}^{t}\right)\) 。因此 \(f\) 在 \(\mathbf{G}\) 的 \({\Gamma }_{\mathbf{K}}\) 类上是常数。

反之，假设 \(f\) 的取值属于 \(\mathrm{K}\) ，且在 G 的 \({\Gamma }_{\mathrm{K}}\) 类上保持常数。则条件 (*) 满足，我们可以写成

\[
f = \sum \langle f,\chi \rangle \chi ,\;\text{ with }\langle f,\chi \rangle  \in  \mathrm{K}
\]

如上。此外， \(f\) 在 \({\sigma }_{t},t \in  {\Gamma }_{\mathbf{K}}\) 下不变，这表明 \(\langle f,\chi \rangle  = \left\langle  {f,{\sigma }_{t}\left( \chi \right) }\right\rangle\) ，因此两个共轭字符 \(\chi\) 与 \({\sigma }_{t}\left( \chi \right)\) 的系数相同。将同一共轭类的字符归并，我们可以把 \(f\) 写成形如 \({\operatorname{Tr}}_{\mathrm{L}/\mathrm{K}}\left( \chi \right)\) 的字符的线性组合。由于后者属于 \({R}_{K}\left( G\right)\) ，参见 12.1，这证明了推论。

[或者:令 \({\Gamma }_{\mathrm{K}}\) 通过 \(\mathrm{K} \otimes  \mathrm{R}\left( \mathrm{G}\right)\) 在 \(f \mapsto  {\sigma }_{t}\left( f\right)  = {\Psi }^{t}\left( f\right)\) 上作用，并注意不动点集为 \(\mathrm{K} \otimes  {\mathrm{R}}_{\mathrm{K}}\left( \mathrm{G}\right)\) 。]

推论 2. 令 \({\chi }_{i}\) 为在 \(\mathbf{K}\) 上不同不可约表示的各个字符，则这些 \({\chi }_{i}\) 构成在 \(\mathbf{G}\) 上对 \({\Gamma }_{\mathrm{K}}\) 类保持常数的函数空间的一组基，它们的个数等于 \({\Gamma }_{\mathrm{K}}\) 类的数目。

此结论来自推论 1。

注。在推论 1 中，我们可以用 \({\bar{R}}_{K}\left( G\right)\) 代替 \({R}_{K}\left( G\right)\) 。确实，命题 34 表明

\[
\mathbf{Q} \otimes  {\mathrm{R}}_{\mathrm{K}}\left( \mathrm{G}\right)  = \mathbf{Q} \otimes  {\overline{\mathrm{R}}}_{\mathrm{K}}\left( \mathrm{G}\right) \text{ , whence }\mathrm{K} \otimes  {\mathrm{R}}_{\mathrm{K}}\left( \mathrm{G}\right)  = \mathrm{K} \otimes  {\overline{\mathrm{R}}}_{\mathrm{K}}\left( \mathrm{G}\right) \text{ . }
\]

12.5 阿廷定理的推广

如果 \(\mathrm{H}\) 是 \(\mathrm{G}\) 的子群，显然有

\[
{\operatorname{Res}}_{\mathrm{H}} : \mathrm{R}\left( \mathrm{G}\right)  \rightarrow  \mathrm{R}\left( \mathrm{H}\right) \;\text{ and }\;{\operatorname{Ind}}_{\mathrm{H}} : \mathrm{R}\left( \mathrm{H}\right)  \rightarrow  \mathrm{R}\left( \mathrm{G}\right)
\]

将 \({R}_{K}\left( G\right)\) 映入 \({R}_{K}\left( H\right)\) 且将 \({R}_{K}\left( H\right)\) 映入 \({R}_{K}\left( G\right)\) 。因此我们可以问当用 \({\mathrm{R}}_{\mathrm{K}}\) 代替 \(\mathrm{R}\) 时，阿廷和布劳尔的定理是否仍然成立。在阿廷定理的情形，答案是肯定的:

定理 26. 设 T 为 G 的循环子群的集合。则由映射

\[
\mathbf{Q} \otimes  \text{ Ind: }\underset{\mathrm{H} \in  \mathrm{T}}{ \oplus  }\mathbf{Q} \otimes  {\mathrm{R}}_{\mathrm{K}}\left( \mathrm{H}\right)  \rightarrow  \mathbf{Q} \otimes  {\mathrm{R}}_{\mathrm{K}}\left( \mathrm{G}\right)
\]

由映射 \(\mathbf{Q} \otimes  {\operatorname{Ind}}_{\mathbf{H}}^{\mathbf{G}},\mathrm{H} \in  \mathrm{T}\) 定义，是满射。

第 9 章给出的两种证明不需修改即可适用。第一种是对偶性论证；须证明映射

\[
\mathbf{Q} \otimes  \text{ Res: }\underset{\mathrm{H} \in  \mathrm{T}}{ \oplus  }\mathbf{Q} \otimes  {\mathrm{R}}_{\mathrm{K}}\left( \mathrm{G}\right)  \rightarrow  \underset{\mathrm{H} \in  \mathrm{T}}{ \oplus  }\mathbf{Q} \otimes  {\mathrm{R}}_{\mathrm{K}}\left( \mathrm{H}\right)
\]

是单射，这显然。

第二种证明由使用公式

\[
g = \mathop{\sum }\limits_{{\mathrm{H} \in  \mathrm{T}}}{\operatorname{Ind}}_{\mathrm{H}}^{\mathrm{G}}\left( {\theta }_{\mathrm{H}}\right) ,\;\text{ cf. prop. }{27}\left( {9.4}\right) ,
\]

并证明 \({\theta }_{\mathrm{H}}\) 属于 \({\mathrm{R}}_{\mathrm{K}}\left( \mathrm{H}\right)\) 。后者可以通过对 \(\mathrm{H}\) 的阶数归纳来验证，或者观察到 \({\mathbf{\theta }}_{\mathrm{H}}\) 取整数值因此属于 \({\overline{\mathbf{R}}}_{\mathbf{K}}\left( \mathrm{H}\right)\) ；由于 \(\mathrm{H}\) 为阿贝尔群，我们有 \({\mathrm{R}}_{\mathbf{K}}\left( \mathrm{H}\right)  = {\overline{\mathbf{R}}}_{\mathbf{K}}\left( \mathrm{H}\right)\) 。现在上述恒等式表明常值函数 1 属于 \(\mathbf{Q} \otimes\) Ind 的像。由于该像是一个理想，它必定是整个环 \(\mathbf{Q} \otimes  {\mathrm{R}}_{\mathrm{K}}\left( \mathrm{G}\right)\) 。

\section*{12.6 布劳尔定理的推广}

我们保留前几节的记号。不难看出，如果 \(\mathrm{X}\) 是 \(\mathrm{G}\) 的初等子群族，则映射

\[
\text{ Ind: } \oplus  {\mathrm{R}}_{\mathrm{K}}\left( \mathrm{H}\right)  \rightarrow  {\mathrm{R}}_{\mathrm{K}}\left( \mathrm{G}\right)
\]

一般情况下并不总是满射(例: \(G = {\mathfrak{S}}_{3},K = \mathbf{R}\) )。有必要用稍大的子群族 \({X}_{K}\) 来替代 \(X\) ，即“ \({\Gamma }_{K}\) -初等”子群族:

设 \(p\) 是素数。子群 \(H\) 属于 \(G\) 被称为 \({\Gamma }_{K} - p\) -元的，如果它是 \(p\) -群 \(P\) 与一个阶数与 \(p\) 互素的循环群 C 的半直积，且满足*:

\(\left( { * }_{\mathrm{K}}\right)\) 对于每个 \(y \in  \mathrm{P}\) ，存在 \(t \in  {\Gamma }_{\mathrm{K}}\) ，使得对于每个 \(x \in  \mathbf{C}\) 都有 \({yx}{y}^{-1} = {x}^{t}\) 。(当 \({\Gamma }_{\mathrm{K}} = \{ 1\}\) 时，此条件仅意味着 \(\mathrm{C}\) 和 \(\mathrm{P}\) 可交换，因此 \(\mathbf{H} = \mathbf{C} \times  \mathbf{P}\) 是一个 \(p\) -初等群。)若 \(\mathbf{G}\) 的一个子群对于至少一个素数 \(p\) 是 \({\Gamma }_{\mathrm{K}} - p\) -初等的，则称该子群为 \({\Gamma }_{\mathrm{K}}\) -初等子群。

\HRule

\begin{itemize}\relax
\item\relax 子群 \(C\) 不应与 12.1 中选定的 \(K\) 的代数闭包混淆；后者不会在本节出现。
\end{itemize}

\HRule

记 \({\mathrm{X}}_{\mathrm{K}}\) (分别记为 \({\mathrm{X}}_{\mathrm{K}}\left( p\right)\) )为 G 的 \({\Gamma }_{\mathrm{K}}\) -初等(分别 \({\Gamma }_{\mathrm{K}} - p\) -初等)子群族。则我们有 th.19 的下述类似结论:

定理 27. 映射 Ind: \({\bigoplus }_{\mathrm{H} \in  {\mathrm{X}}_{\mathrm{K}}}{\mathrm{R}}_{\mathrm{K}}\left( \mathrm{H}\right)  \rightarrow  {\mathrm{R}}_{\mathrm{K}}\left( \mathrm{G}\right)\) 是满射。

如同 10.5，我们从一个更精确的关于固定素数 \(p\) 的结果得到定理 27:

定理 28。设 \(g = {p}^{n}l\) 为 \(\mathrm{G}\) 的阶，其中 \(\left( {p,l}\right)  = 1\) 。常值函数 \(l\) 属于映射的像 \({\mathrm{V}}_{\mathrm{K},p}\)

\[
\text{ Ind: }\underset{\mathrm{H} \in  {\mathrm{X}}_{\mathrm{K}}\left( p\right) }{ \oplus  }{\mathrm{R}}_{\mathrm{K}}\left( \mathrm{H}\right)  \rightarrow  {\mathrm{R}}_{\mathrm{K}}\left( \mathrm{G}\right) \text{ . }
\]

特别地， \({\mathrm{V}}_{\mathrm{K},p}\) 在 \({\mathrm{R}}_{\mathrm{K}}\left( \mathrm{G}\right)\) 中的指数是有限的并且与 \(p\) 互素。

本定理的证明与定理 \({18}^{\prime }\) 的证明完全类似(当 \(\mathrm{K}\) 代数封闭时本定理即归结为该定理)。我们将在下一节给出证明，目前只指出两个推论:

命题 36。设 \(\varphi\) 为 G 上的类函数。要使 \(\varphi\) 属于 \({\mathrm{R}}_{\mathrm{K}}\left( \mathrm{G}\right)\) ，当且仅当对 \(\mathrm{G}\) 的每个 \({\Gamma }_{\mathrm{K}}\) -元子群 \(\mathrm{H}\) ，都有 \({\operatorname{Res}}_{\mathrm{H}}^{\mathrm{G}}\varphi  \in  {\mathrm{R}}_{\mathrm{K}}\left( \mathrm{H}\right)\) 。

利用定理 27，我们有一个恒等式

\[
1 = \mathop{\sum }\limits_{{\mathrm{H} \in  {\mathrm{X}}_{\mathrm{K}}}}{\operatorname{Ind}}_{\mathrm{H}}^{\mathrm{G}}{f}_{\mathrm{H}},\;\text{ with }{f}_{\mathrm{H}} \in  {\mathrm{R}}_{\mathrm{K}}\left( \mathrm{H}\right) .
\]

两边乘以 \(\varphi\) ，得到

\[
\varphi  = \mathop{\sum }\limits_{{H \in  {X}_{K}}}\varphi  \cdot  {\mathrm{{Ind}}}_{H}^{G}{f}_{H} = \mathop{\sum }\limits_{{H \in  {X}_{K}}}{\mathrm{{Ind}}}_{H}^{G}\left( {{f}_{H} \cdot  {\mathrm{{Res}}}_{H}^{G}\varphi }\right) .
\]

因此，若对所有 \(\mathrm{H} \in  {\mathrm{X}}_{\mathrm{K}}\) 有 \({\operatorname{Res}}_{\mathrm{H}}\varphi  \in  {\mathrm{R}}_{\mathrm{K}}\left( \mathrm{H}\right)\) ，则我们有 \(\varphi  \in  {\mathrm{R}}_{\mathrm{K}}\left( \mathrm{G}\right)\) ；反之亦然。

命题 37。若各代数 \(\mathrm{K}\left\lbrack  \mathrm{H}\right\rbrack  ,\mathrm{H} \in  {\mathrm{X}}_{\mathrm{K}}\) 均为拟分裂(参见 12.2)，则 \(\mathrm{K}\left\lbrack  \mathrm{G}\right\rbrack\) 也有同样性质。

设 \(\varphi  \in  {\overline{\mathrm{R}}}_{\mathrm{K}}\left( \mathrm{G}\right)\) 。对于 \(\mathrm{H} \in  {\mathrm{X}}_{\mathrm{K}}\) ，我们有 \({\operatorname{Res}}_{\mathrm{H}}^{\mathrm{G}}\varphi  \in  {\overline{\mathrm{R}}}_{\mathrm{K}}\left( \mathrm{H}\right)\) ，且由于 \(\mathrm{K}\left\lbrack  \mathrm{H}\right\rbrack\) 为拟分裂(参见命题 35 的推论)， \({\overline{\mathrm{R}}}_{\mathrm{K}}\left( \mathrm{H}\right)\) 等于 \({\mathrm{R}}_{\mathrm{K}}\left( \mathrm{H}\right)\) 。前一命题遂表明 \(\varphi\) 属于 \({\mathrm{R}}_{\mathrm{K}}\left( \mathrm{G}\right)\) 。因此 \({\overline{\mathrm{R}}}_{\mathrm{K}}\left( \mathrm{G}\right)  = {\mathrm{R}}_{\mathrm{K}}\left( \mathrm{G}\right)\) ，且 \(\mathrm{K}\left\lbrack  \mathrm{G}\right\rbrack\) 为拟分裂。

练习

12.7. 证明诱导映射 Ind: \(\underset{\mathrm{H} \in  {\mathrm{X}}_{\mathrm{K}}}{ \oplus  }{\overline{\mathrm{R}}}_{\mathrm{K}}\left( \mathrm{H}\right)  \rightarrow  {\overline{\mathrm{R}}}_{\mathrm{K}}\left( \mathrm{G}\right)\) 是满射。[使用命题 36 的证明。]

\section*{12.7 定理 28 的证明}

我们记 \(\mathbf{A}\) 为由 \(m\) 次单位根生成的 \(\mathbf{L}\) 的子环。

引理 13. 若 \(l\) 属于 \(\mathrm{A} \otimes  {\mathrm{V}}_{\mathrm{K},p}\) ，则 \(l \in  {\mathrm{V}}_{\mathrm{K},p}\) 。

此证法与 10.2 中用于引理 5 的论证相同。

引理 14. 包含 \(p\) 的 \(\mathrm{A}\) 的素理想 \({\mathfrak{p}}_{1},\ldots ,{\mathfrak{p}}_{h}\) 只有有限个。商 \(A/{\mathfrak{p}}_{i}\) 是特征为 \(p\) 的有限域，且存在整数 \(\mathrm{N}\) 使得 \(p\mathrm{\;A} \supset  {\left( {\mathfrak{p}}_{1} \cap  \cdots  \cap  {\mathfrak{p}}_{h}\right) }^{\mathrm{N}}\) 。

\({\mathfrak{p}}_{i}\) 对应于 \(\mathrm{A}/p\mathrm{\;A}\) 的素理想，而 \(\mathrm{A}/p\mathrm{\;A}\) 是特征为 \(p\) 的有限环。前两项断言由此而来。第三项则由 \(\left( {{\mathfrak{p}}_{1} \cap  \cdots  \cap  {\mathfrak{p}}_{h}}\right) /p\mathrm{\;A}\) 为 Artinian 环 \(\mathrm{A}/p\mathrm{\;A}\) 的根得出，因此是幂零的。

引理 15. 设 \(f\) 为定义在 \(\mathrm{G}\) 上、在 \({\Gamma }_{\mathrm{K}}\) -类上常值且取值于 \(g\mathrm{\;A}\) 的函数。则 \(f\) 可写成以下形式

\[
f = \sum {\operatorname{Ind}}_{\mathrm{C}}^{\mathrm{G}}\left( {\varphi }_{\mathrm{C}}\right) ,\;\text{ with }{\varphi }_{\mathrm{C}} \in  \mathrm{A} \otimes  {\mathrm{R}}_{\mathrm{K}}\left( \mathrm{C}\right) ,
\]

其中 \(\mathrm{C}\) 在所有 \(\mathrm{G}\) 的循环子群上取遍。

令 \(\varphi  = f/g\) 。按引理 6 的记法，我们有

\[
f = \sum {\operatorname{Ind}}_{\mathrm{C}}^{\mathrm{G}}\left( {{\theta }_{\mathrm{C}} \cdot  {\operatorname{Res}}_{\mathrm{C}}^{\mathrm{G}}\varphi }\right) ,
\]

余下只需证明对任意 \(\mathrm{C}\) 有 \({\varphi }_{\mathrm{C}} = {\theta }_{\mathrm{C}} \cdot  {\operatorname{Res}}_{\mathrm{C}}^{\mathrm{G}}\varphi\) 属于 \(\mathrm{A} \otimes  {\mathrm{R}}_{\mathrm{K}}\left( \mathrm{C}\right)\) 。但 \({\varphi }_{\mathrm{C}}\) 的值可被 \(\mathrm{C}\) 的阶整除；因此若 \(\chi\) 为 \(\mathrm{C}\) 的一次特征，则有 \(\left\langle  {{\varphi }_{\mathrm{C}},\chi }\right\rangle   \in  \mathrm{A}\) 。此外， \(f\) 在 \({\Gamma }_{\mathrm{K}}\) -类上常值这一事实意味着

\[
\left\langle  {{\varphi }_{\mathrm{C}},\chi }\right\rangle   = \left\langle  {{\Psi }^{t}{\varphi }_{\mathrm{C}},{\Psi }^{t}\chi }\right\rangle   = \left\langle  {{\varphi }_{\mathrm{C}},{\Psi }^{t}\chi }\right\rangle  ,\;\text{ if }t \in  {\Gamma }_{\mathrm{K}}.
\]

故在 \({\varphi }_{\mathrm{C}}\) 中共轭于 \(\mathrm{K}\) 的特征的系数相等，因此我们可以将 \({\varphi }_{\mathrm{C}}\) 表为 \(\mathrm{K}\) 上若干特征 \(\chi\) 的迹的 A 线性组合；于是 \({\varphi }_{\mathrm{C}} \in  \mathrm{A} \otimes  {\mathrm{R}}_{\mathrm{K}}\left( \mathrm{C}\right)\) 。

引理 16. 设 \(x,y \in  \mathrm{G}\) 为若干元素，其 \({p}^{\prime }\) 分量两两为 \({\Gamma }_{\mathrm{K}}\) -共轭。

若 \(f \in  \mathrm{A} \otimes  {\mathrm{R}}_{\mathrm{K}}\left( \mathrm{G}\right)\) ，则

\[
f\left( x\right)  \equiv  f\left( y\right) \;\left( {\;\operatorname{mod}\;.{\mathfrak{p}}_{i}}\right) \;\text{ for }i = 1,\ldots ,h.
\]

我们知道 \(f\) 在 \({\Gamma }_{\mathrm{K}}\) -类上是常数(定理25的推论1)。因此我们可以假设 \(x\) 是 \({p}^{\prime }\) -分量的 \(y\) ，在这种情况下与引理7的证明中相同的论证适用。

引理17。设 \(x\) 是 \({p}^{\prime }\) 的 \(\mathbf{G}\) 元素，令 \(\mathbf{C}\) 为由 \(x\) 生成的循环子群，令 \(\mathrm{N}\left( x\right)\) 为存在 \(t \in  {\Gamma }_{\mathbf{K}}\) 使得 \({yx}{y}^{-1} = {x}^{t}\) 的那些 \(y \in  \mathbf{G}\) 的集合，且令 \(\mathrm{P}\) 为 \(\mathrm{N}\left( x\right)\) 的一个Sylow子群。则:

(a) \(\mathrm{H} = \mathrm{C} \cdot  \mathrm{P}\) 是 \({\Gamma }_{\mathrm{K}}\) -p-初等子群的 \(\mathrm{G}\) 。

(b) \(\mathrm{C}\) 在 \(\mathrm{K}\) 上的每个线性表示都可扩展到 \(\mathrm{H}\) 。

(c) 映射 Res: \({\mathrm{R}}_{\mathrm{K}}\left( \mathrm{H}\right)  \rightarrow  {\mathrm{R}}_{\mathrm{K}}\left( \mathrm{G}\right)\) 是满射。

断言(a) 显然。为证明(b)，仅需考虑在 K 上的不可约表示。这类表示可通过选择一个同态 \(\chi  : \mathrm{C} \rightarrow  {\mathrm{L}}^{ * }\) ，取作为向量空间由 \(\chi \left( \mathrm{C}\right)\) 生成的子域 \({\mathrm{K}}_{\chi }\) ，并按下列公式定义 \(\rho  : \mathrm{C} \rightarrow  \mathbf{{GL}}\left( {\mathrm{K}}_{\chi }\right)\) 来获得

\[
\rho \left( s\right) \omega  = \chi \left( s\right) \omega \;\text{ if }s \in  \mathrm{G}\text{ and }\omega  \in  \mathrm{K}.
\]

群 \({\Gamma }_{\mathrm{K}} = \mathrm{{Gal}}\left( {\mathrm{L}/\mathrm{K}}\right)\) 在 \(\mathrm{K}\) -线性地作用于 \({\mathrm{K}}_{x}\) 。若 \(y \in  \mathrm{P}\) ，令 \(t \in  {\Gamma }_{\mathrm{K}}\) 为满足 \({yx}{y}^{-1} = {x}^{t}\) 的元素，并将 \(\rho \left( y\right)\) 定义为 \({\sigma }_{t}\) 限制到 \({\mathrm{K}}_{x}\) 。检验得出 \(\rho \left( y\right)\) 不依赖于所选的 \(t\) ，并且

\[
\rho \left( y\right) \rho \left( x\right) \rho {\left( y\right) }^{-1} = \rho \left( {x}^{t}\right) .
\]

由此可得上面定义的 \(C\) 与 \(P\) 的同态可扩展为 \(H\) 到 \({GL}\left( {K}_{\chi }\right)\) 的同态，从而证明了(b)。断言(c) 从(b) 推出。

在10.3中我们有 \({\Gamma }_{K} = \{ 1\}\) ，因此 \(H = C \times  P\) ，故上述引理是平凡的。

引理18。保留引理17的记号。则存在

\[
\psi  \in  \mathrm{A} \otimes  {\mathrm{R}}_{\mathrm{K}}\left( \mathrm{H}\right)
\]

使得诱导函数 \({\psi }^{\prime } = {\operatorname{Ind}}_{\mathrm{H}}^{\mathrm{G}}\psi\) 具有如下性质:

(i) 对于 \(i = 1,\ldots ,h\) 有 \({\psi }^{\prime }\left( x\right)  \text{ ≢ } 0\left( {\;\operatorname{mod}\;.{\mathfrak{p}}_{i}}\right)\) 。

(ii) \({\psi }^{\prime }\left( s\right)  = 0\) 对于每个不与 \(x\)  \({\Gamma }_{\mathbf{K}}\) -共轭的 \(\mathbf{G}\) 的 \({p}^{\prime }\) -元 \(s\) 。

设 \(c\) 为 \(\mathbf{C}\) 的阶，且令 \({\psi }_{\mathbf{C}}\) 为在 \(y\) 具有形式 \({x}^{t}\) 时由 \({\psi }_{\mathrm{C}}\left( y\right)  = c\) 定义在 \(\mathbf{C}\) 上的函数，在此情况下为 \(t \in  {\Gamma }_{\mathrm{K}}\) ，否则为 \({\psi }_{\mathrm{C}}\left( y\right)  = 0\) 。然后 \({\psi }_{\mathrm{C}} \in  \mathrm{A} \otimes  {\mathrm{R}}_{K}\left( \mathrm{C}\right)\) :例如这可从对 \(\mathrm{C}\) 应用引理15 得出。由引理17，存在 \(\psi  \in  \mathrm{A} \otimes  {\mathrm{R}}_{\mathrm{K}}\left( \mathrm{H}\right)\) 使得 \({\operatorname{Res}}_{\mathrm{C}}^{\mathrm{H}}\psi  = {\psi }_{\mathrm{C}}\) 。我们证明 \(\psi\) 可行:

若 \(s\) 是 \(\mathrm{G}\) 的一个 \({p}^{\prime }\) -元，且若 \(y \in  \mathrm{G},{ys}{y}^{-1}\) 是一个 \({p}^{\prime }\) -元。若 \({ys}{y}^{-1} \in  \mathrm{H}\) ，则 \({yss}{y}^{-1} \in  \mathrm{C}\) 且当 \({ys}{y}^{-1}\) 不是形如 \({x}^{t}\) 时 \(\psi \left( {{ys}{y}^{-1}}\right)\) 为零，针对 \(t \in  {\Gamma }_{\mathrm{K}}\) 。因此若 \(s\) 不与 \(x\)  \({\Gamma }_{\mathrm{K}}\) -共轭则可得出 \({\psi }^{\prime }\left( s\right)  = 0\) ，这证明了 (ii)。其余部分，令 \(Z\) 为 \({x}^{t},t \in  {\Gamma }_{\mathbf{K}}\) 的集合。然后

\[
{\psi }^{\prime }\left( x\right)  = \frac{1}{\operatorname{Card}\left( \mathrm{H}\right) }\mathop{\sum }\limits_{{{yx}{y}^{-1} \in  \mathrm{Z}}}c = \frac{\operatorname{Card}\left( {\mathrm{N}\left( x\right) }\right) }{\operatorname{Card}\left( \mathrm{P}\right) },
\]

并且由于 \(\mathrm{P}\) 是 \(\mathrm{N}\left( x\right)\) 的一个 Sylow \(p\) -子群(参见引理17)，我们看到 \({\psi }^{\prime }\left( x\right)\) 是与 \(p\) 互素的整数，因而得出 (i)。

引理19。存在 \(\varphi  \in  \mathrm{A} \otimes  {\mathrm{V}}_{\mathrm{K},p}\) 使得对每个 \(x \in  \mathrm{G}\) 和每个 \(i = 1,\ldots ,h\) 都有 \(\varphi \left( x\right)  \equiv  0\left( {\;\operatorname{mod}\;.{\mathfrak{p}}_{i}}\right)\) 。

令 \({\left( {x}_{\lambda }\right) }_{\lambda  \in  \Lambda }\) 为 \({p}^{\prime }\) -常规 \({\Gamma }_{\mathrm{K}}\) -类的一组代表元，即那些由 \({p}^{\prime }\) -元组成的类。对于每个 \(\lambda  \in  \Lambda\) ，前一引理使我们得以构造 \({\varphi }_{\lambda } \in  \mathrm{A} \otimes  {\mathrm{V}}_{\mathrm{K},p}\) 使得

\[
{\varphi }_{\lambda }\left( {x}_{\lambda }\right)  \equiv  0\left( {\;\operatorname{mod}\;{\mathfrak{p}}_{i}}\right) \;\text{ and }\;{\varphi }_{\lambda }\left( {x}_{\lambda }\right)  = 0\;\text{ if }\lambda  \neq  \mu .
\]

先置 \(\varphi  = \mathop{\sum }\limits_{\lambda }{\varphi }_{\lambda }\) 。然后 \(\varphi\) 属于 \(\mathrm{A} \otimes  {\mathrm{V}}_{\mathrm{K},p}\) ，并且对 G 中的每个 \({p}^{\prime }\) -元 \(x\) 我们都有 \(\varphi \left( x\right)  \text{ ≢ } 0\) (模 \({\mathfrak{p}}_{i}\) )。引理16表明对 \(\mathrm{G}\) 中的每个 \(x\) 也成立。

\section*{定理28证明的完成}

设 \(\varphi  \in  \mathrm{A} \otimes  {\mathrm{V}}_{\mathrm{K},p}\) 满足引理19的条件。对每个 \(x \in  \mathrm{G}\) 与每个 \(i\) ， \(\varphi \left( x\right)\) 模 \({\mathfrak{p}}_{i}\) 的类属于域 \(\mathrm{A}/{\mathfrak{p}}_{i}\) 的乘法群。由于域 \(\mathrm{A}/{\mathfrak{p}}_{i}\) 是有限的(引理14)，存在一个 \(\mathrm{M} \geq  1\) 使得对所有 \(i\) 与所有 \(x \in  \mathrm{G}\) 都有 \({\varphi }^{\mathrm{M}}\left( x\right)  \equiv  1\left( {\;\operatorname{mod}\;.{\mathfrak{p}}_{i}}\right)\) 。然后由引理14 我们有 \({\varphi }^{\mathrm{{MN}}}\left( x\right)  \equiv  1\left( {{\;\operatorname{mod}\;.}p\mathrm{\;A}}\right)\) ，并将 \({\varphi }^{\mathrm{{MN}}}\) 提到 \(n\) 次方，得到 \(\psi  \in  \mathrm{A} \otimes  {\mathrm{V}}_{\mathrm{K},p}\) ，使得

\[
\psi \left( x\right)  \equiv  1\left( {\;\operatorname{mod}\;{p}^{n}\mathbf{A}}\right) \;\text{ for all }x \in  \mathbf{G}.
\]

函数 \(l\left( {\psi  - 1}\right)\) 因此取值于 \({p}^{n}/\mathrm{A} = g\mathrm{\;A}\) 。根据引理15，我们有 \(l\left( {\psi  - 1}\right)  \in  \mathrm{A} \otimes  {\mathrm{V}}_{\mathrm{K},p}\) 。减去后，我们得到 \(l\) 属于 \(\mathrm{A} \otimes  {\mathrm{V}}_{\mathrm{K},p}\) ，现在由引理13定理即得。

\section*{练习}

12.7. 求环 \(A \otimes  R\left( G\right)\) 的谱。(结果与11.4相同，只是共轭类被 \({\Gamma }_{\mathrm{K}}\) -类取代。)

\section*{第13章 有理性问题:例子}

我们沿用第12章的记号。

\section*{13.1 域 \(\mathbf{Q}\)}

设 \(\mathrm{G}\) 为阶为 \(g\) 的有限群，且令 \(m\) 为包含 \(\mathbf{G}\) 中所有元素阶的公倍数。取基域 \(\mathbf{K}\) 为有理数域 \(\mathbf{Q}\) ，并令 \(\mathbf{Q}\left( m\right)\) 为通过将第 \(m\) 个单位根加入 \(\mathbf{Q}\) 而得到的域。 \(\mathbf{Q}\left( m\right)\) 对 \(\mathbf{Q}\) 的伽罗瓦群即12.4中记为 \({\Gamma }_{\mathbf{O}}\) 的群；它是群 \({\left( \mathbf{Z}/m\mathbf{Z}\right) }^{ * }\) 的子群。事实上:

定理(高斯)。我们有 \({\Gamma }_{\mathbf{Q}} = {\left( \mathbf{Z}/m\mathbf{Z}\right) }^{ * }\) 。

(这相当于说 \(m\) 阶的cyclotomic多项式 \({\Phi }_{m}\) 在Q上不可约。)

我们假定这一经典结果；有关证明，参见例如Lang [10], 第204页。

推论。若且唯若它们所生成的循环子群共轭， \(\mathrm{G}\) 的两个元素才是 \({\Gamma }_{\mathbf{Q}}\) -共轭的。

应用12.4的结果，我们得到:

定理29。令 \(f\) 是在 \(\mathbf{G}\) 上的取值于 \(\mathbf{Q}\left( m\right)\) 的类函数。

(a) 为了使 \(f\) 属于 \(\mathbf{Q} \otimes  \mathrm{R}\left( \mathrm{G}\right)\) ，当且仅当对每个与 \(m\) 互素的 \(t\) 有 \({\sigma }_{t}\left( f\right)  = {\Psi }^{t}\left( f\right)\) 。

(b) 为了使 \(f\) 属于 \(\mathbf{Q} \otimes  {\mathrm{R}}_{\mathbf{O}}\left( \mathrm{G}\right)\) ，当且仅当 \(f\) 的值在 \(\mathbf{Q}\) 之中，且对每个与 \(m\) 互素的 \(t\) 有 \({\Psi }^{t}\left( f\right)  = f\) (即若 \(x\) 与 \(y\) 生成同一 \(\mathrm{G}\) 的子群，则必须有 \(f\left( x\right)  = f\left( y\right)\) )。

(回想 \({\sigma }_{t}\) 是 \(\mathbf{Q}\left( m\right)\) 的自同构，其将一个 \(m\) 次单位根送为其 \(t\) 次幂，且 \({\Psi }^{t}\left( f\right)\) 是函数 \(x \mapsto  f\left( {x}^{t}\right)\) 。)

推论 1. 在 \(\mathbf{Q}\) 上 \(\mathbf{G}\) 的不可约表示的同构类数等于 \(\mathrm{G}\) 的循环子群的共轭类数。

此由定理 25 的推论 2 推出。

推论 2. 下列性质等价:

(i) \(\mathrm{G}\) 的每个特征值均取于 \(\mathbf{Q}\) 。

(i') \(\mathrm{G}\) 的每个特征值均取于 \(\mathbf{Z}\) 。

(ii) 生成同一子群的 \(\mathrm{G}\) 中两个元素是共轭的。

(i) 与 (i') 的等价性来自特征值是代数整数，因此一旦属于 \(\mathbf{Q}\) 就属于 \(\mathbf{Z}\) 。(i) 与 (ii) 的等价性由定理 29 得出。

\section*{例子}

(1) 对称群 \({\mathfrak{S}}_{n}\) 满足 (ii)，因此满足 (i)。此外，可以证明 \({\mathfrak{S}}_{n}\) 的每个表示都可在 \(\mathbf{Q}\) 上实现，即 \(\mathrm{R}\left( {\mathfrak{S}}_{n}\right) \; = {\mathrm{R}}_{\mathbf{O}}\left( {\mathfrak{S}}_{n}\right)\) 。

(2) 四元数群 \(G = \{  \pm  1, \pm  i, \pm  j, \pm  k\}\) 满足推论的条件。故 \({\bar{R}}_{O}\left( G\right)  = R\left( G\right)\) ；群 \({R}_{O}\left( G\right)\) 是 \(R\left( G\right)\) 的指数为 2 的子群，参见例 12.3。

若 \(H\) 是 \(G\) 的子群，记 \({l}_{H}\) 为 \(H\) 的恒等特征，记 \({l}_{H}^{G}\) 为由 \({l}_{H}\) 诱导到 \(G\) 的特征(即在 \(\mathrm{G}/\mathrm{H}\) 上的置换表示的特征，参见 3.3，例 2)。

定理 30. \({\mathrm{R}}_{\mathrm{O}}\left( \mathrm{G}\right)\) 的每个元素都可以表示为以 \(\mathbf{Q}\) 为系数的特征 \({1}_{\mathrm{C}}^{\mathrm{G}}\) 的线性组合，其中 \(\mathrm{C}\) 在 \(\mathrm{G}\) 的循环子群集上变化。

这等价于说 \(\mathbf{Q} \otimes  {\mathrm{R}}_{\mathbf{O}}\left( \mathrm{G}\right)\) 由 \({\mathrm{l}}_{\mathrm{C}}^{\mathrm{G}}\) 生成。由于 \(\mathbf{Q} \otimes  {\mathbf{R}}_{\mathbf{Q}}\left( \mathbf{G}\right)\) 配备了非退化的双线性型

\[
\left( {\varphi ,\psi }\right)  \mapsto  \langle \varphi ,\psi \rangle ,
\]

我们同样可以证明，若 \({\mathrm{R}}_{\mathrm{O}}\left( \mathrm{G}\right)\) 中的元素 \(\theta\) 与所有 \({1}_{\mathrm{C}}^{\mathrm{G}}\) 正交，则该元素为零。然而，我们有

\[
\left\langle  {\theta ,{\mathrm{l}}_{\mathrm{C}}^{\mathrm{G}}}\right\rangle   = \left\langle  {{\operatorname{Res}}_{\mathrm{C}}^{\mathrm{G}}\theta ,{\mathrm{l}}_{\mathrm{C}}}\right\rangle   = \frac{1}{c}\mathop{\sum }\limits_{{s \in  \mathrm{C}}}\theta \left( s\right) ,\;\text{ where }c = \mathrm{{Card}}\mathrm{C}.
\]

因此定理 30 等价于下列命题:

定理 \({30}^{\prime }\) 。若对每个 \(\mathrm{G}\) 的循环子群 \(\mathrm{C}\) 都有 \(\mathop{\sum }\limits_{{s \in  \theta }}\theta \left( s\right)  = 0\) ，且 \(\theta  \in  {\mathrm{R}}_{\mathbf{Q}}\left( \mathrm{G}\right)\) 满足该性质，则 \(\widetilde{\theta } = 0\) 。

我们按 Card(G) 做归纳证明此结果。令 \(s \in  \mathrm{G}\) ，并令 \(\mathbf{C}\left( s\right)\) 为由 \(s\) 生成的 \(\mathbf{G}\) 的循环子群。取 \(x \in  \mathbf{C}\left( s\right)\) ；若 \(x\) 生成了 \(C\left( s\right)\) ，因为 \(x\) 与 \(s\) 是 \({\Gamma }_{0}\) -共轭的，便有 \(\theta \left( x\right)  = \theta \left( s\right)\) ；若 \(x\) 仅生成 \(\mathbf{C}\left( s\right)\) 的真子群，则对该子群限制后的 \(\theta\) 应用归纳假设可得 \(\theta \left( x\right)  = 0\) 。由此我们得到

\[
\mathop{\sum }\limits_{{x \in  \mathrm{C}\left( s\right) }}\theta \left( x\right)  = a \cdot  \theta \left( s\right)
\]

其中 \(a\) 是 \(\mathrm{C}\left( s\right)\) 的生成元个数。但根据假设我们有

\[
\mathop{\sum }\limits_{{x \in  \mathrm{C}\left( s\right) }}\theta \left( x\right)  = 0,
\]

因此 \(\theta \left( s\right)  = 0\) 。

推论。设 \(\mathbf{V}\) 与 \({\mathbf{V}}^{\prime }\) 为在域 \(\mathbf{Q}\) 上的两个线性表示。要使 \(\mathbf{V}\) 与 \({\mathbf{V}}^{\prime }\) 同构，当且仅当对 \(\mathbf{G}\) 的每个循环子群 \(\mathbf{C}\) ，都有

\[
\dim {\mathbf{V}}^{\mathbf{C}} = \dim {\mathbf{V}}^{\prime \mathbf{C}}
\]

其中 \({\mathrm{V}}^{\mathrm{C}}\) (分别 \({\mathrm{V}}^{\prime \mathrm{C}}\) )表示 \(\mathrm{V}\) (分别 \({\mathrm{V}}^{\prime }\) )中在 \(\mathbf{C}\) 下不变的子空间。

必要性显然。要证明该条件充分，令 \(\chi\) 与 \({\chi }^{\prime }\) 分别为 \(\mathrm{V}\) 与 \({\mathrm{V}}^{\prime }\) 的特征。我们有:

\[
\dim {\mathbf{V}}^{\mathbf{C}} = {\left\langle  {\operatorname{Res}}_{\mathbf{C}}^{\mathbf{G}}\chi ,{\mathbf{l}}_{\mathbf{C}}\right\rangle  }_{\mathbf{C}}
\]

因此对每个 \(\mathrm{C}\) 有 \(\left\langle  {{\operatorname{Res}}_{\mathrm{C}}^{\mathrm{G}}\left( {\chi  - {\chi }^{\prime }}\right) ,{1}_{\mathrm{C}}}\right\rangle   = 0\) ，由定理 30 得到 \(\chi  - {\chi }^{\prime } = 0\) 。于是 \(\mathrm{V}\) 与 \({\mathrm{V}}^{\prime }\) 同构。

\section*{注记}

(1) 一般并非每个 \({R}_{0}\left( G\right)\) 元都是若干特征 \({1}_{\mathrm{H}}^{\mathrm{G}}\) 的整系数线性组合，即便 \(\mathrm{H}\) 遍及 \(\mathrm{G}\) 的所有子群(参见例 13.4)。

(2) 定理 30 蕴含以下结果:设 \(\mathrm{F}/\mathrm{E}\) 为数域的有限伽罗瓦扩张，且设 \(\chi\) 为可在 \(\mathbf{Q}\) 上实现的 \(\mathrm{{Gal}}\left( {\mathrm{F}/\mathrm{E}}\right)\) 的线性表示的特征标。那么，我们可以将相对于 \(\chi\) 的阿廷 L - 函数写成与 \({Gal}\left( {F/E}\right)\) 的循环子群 \(C\) 对应的 \(F\) 的子域 \({F}_{C}\) 的zeta 函数的分数次幂的乘积。

\section*{练习}

13.1. 设 \(\mathrm{G}\) 为阶为 \(n\) 的循环群。对每个 \(n\) 的约数 \(d\) ，记 \({\mathrm{G}}_{d}\) 为 \(\mathrm{G}\) 中指数为 \(d\) 的子群。

(a) 证明 \(\mathbf{G}\) 在 \(\mathbf{Q}\) 上有一个不可约表示，至同构唯一，其核等于 \({\mathrm{G}}_{d}\) 。记其特征为 \({\chi }_{d}\) ；则 \({\chi }_{d}\left( 1\right)  = \varphi \left( d\right)\) 。这些 \({\chi }_{d}\) 构成 \({\mathrm{R}}_{\mathbf{Q}}\left( \mathrm{G}\right)\) 的正交基。

(b) 给出从 \(\mathbf{Q}\left\lbrack  \mathrm{G}\right\rbrack\) 到 \(\mathop{\prod }\limits_{{d \mid  n}}\mathbf{Q}\left( d\right)\) 的同构映射。

(c) 设 \({\psi }_{d} = {1}_{{G}_{d}}^{G}\) 。证明 \({\psi }_{d} = \mathop{\sum }\limits_{{{d}^{\prime } \mid  d}}{\chi }_{{d}^{\prime }}\) 且 \({\chi }_{d} = \mathop{\sum }\limits_{{{d}^{\prime } \mid  d}}\mu \left( {d/{d}^{\prime }}\right) {\psi }_{{d}^{\prime }}\) ，其中 \(\mu\) 表示 Möbius 函数。由此推得这些 \({\psi }_{d}\) 构成 \({\mathbf{R}}_{\mathbf{Q}}\left( \mathbf{G}\right)\) 的一组基。

13.2. 通过用定理 26 将问题化约到循环情形，然后应用练习 13.1，证明定理 30。

13.3. 设 \(\rho\) 为 \(\mathbf{G}\) 在 \(\mathbf{Q}\) 上的一个不可约表示，令 \(\mathbf{A} = {\mathbf{M}}_{n}\left( \mathbf{D}\right)\) 为 \(\mathbf{Q}\left\lbrack  \mathbf{G}\right\rbrack\) 的相应简单直和成分(D 为一域，不一定交换)，并令 \(\chi\) 为 \(\rho\) 的特征。假设 \(\rho\) 可信(即 ker \(\rho  = 1\) )且 \(G\) 的每个子群都是正规子群。令 \(H\) 为 G 的一个子群。证明在 \(\mathrm{G}/\mathrm{H}\) 上的置换表示当 \(\mathrm{H} = \{ 1\}\) 时包含表示 \({\rho n}\) 次，当 \(\mathrm{H} \neq  \{ 1\}\) 时包含 0 次。由此得出:如果 \(n \geq  2\) ，则 \(\chi\) 不包含在由特征 \({l}_{H}^{G}\) 所生成的 \({R}_{Q}\left( G\right)\) 的子群中。

13.4. 设 \(E\) 为四元数群， \(C\) 为 3 阶循环群，且设 \(\mathbf{G} = \mathbf{E} \times  \mathbf{C}\) 。若 \({\mathbf{H}}_{\mathbf{O}}\) 表示通常的四元数域(在 \(\mathbf{Q}\) 上)，证明 E 与 \(\mathrm{C}\) 可嵌入乘法群 \({\mathbf{H}}_{\mathbf{O}}^{ * }\) 。这给出了一个由右乘(分别由左乘)作用于向量空间 \({\mathbf{H}}_{\mathbf{O}}\) 的 \(E\) (分别 \(C\) )作用。从中得到一个在 \(\mathbf{Q}\) 上次数为 4 的不可约表示 \(\rho\) (为 \(\mathbf{G}\) 的表示)。证明相应的简单代数同构于 \({\mathbf{M}}_{2}\left( \mathrm{\;K}\right)\) ，其中 \(\mathrm{K}\) 为立方根单位的域。验证练习 13.3 的条件，并推得 \(\rho\) 的特征不是字符 \({\mathrm{l}}_{\mathbf{H}}^{\mathbf{G}},\mathrm{H} \subset  \mathbf{G}\) 的线性组合。

13.5. 设 \(\mathrm{X}\) 与 \(\mathrm{Y}\) 为两个有限集合，群 \(\Gamma\) 在其上作用。若 \(\mathrm{H}\) 是 \(\Gamma\) 的子群，记由 \({X}^{H}\) (分别 \({Y}^{H}\) ) 固定的 \(X\) (分别 \(Y\) ) 元素的集合。证明当且仅当对 \(\Gamma\) 的每一子群 \(H\) 都有 \(\operatorname{Card}\left( {X}^{H}\right) \; = \operatorname{Card}\left( {Y}^{H}\right)\) 时， \(\Gamma\) -集 \(X\) 与 \(Y\) 同构。接着证明下列性质彼此等价:

(i) 与 \(X\) 与 \(Y\) 关联的(线性)置换表示 \({\rho }_{X}\) 与 \({\rho }_{Y}\) 同构。

(ii) 对于 \(\Gamma\) 的每一个循环子群 \(H\) ，都有 \(\operatorname{Card}\left( {X}^{H}\right)  = \operatorname{Card}\left( {Y}^{H}\right)\) 。

(iii) 对于 \(\Gamma\) 的每一个子群 \(H\) ，都有 \(\operatorname{Card}\left( {X/H}\right)  = \operatorname{Card}\left( {Y/H}\right)\) 。

(iv) 对于 \(\Gamma\) 的每一个循环子群 \(H\) ，都有 \(\operatorname{Card}\left( {X/H}\right)  = \operatorname{Card}\left( {Y/H}\right)\) 。

当这些性质成立时，我们称 \(X\) 与 \(Y\) 是弱同构的。

[(i)与(ii)等价可通过计算 \({\rho }_{X}\) 与 \({\rho }_{Y}\) 的特征得到。(i)与(iii)、(iv) 的等价来自于 \(\operatorname{Card}\left( {X/H}\right)\) 为 \({\rho }_{X}\) 的特征与 \({1}_{\mathrm{H}}^{1}\) 的特征的内积这一事实。]

证明若 \(\Gamma\) 为循环群，则 \(\Gamma\) -集 \(\mathbf{X}\) 与 \(\mathbf{Y}\) 同构当且仅当它们弱同构。给出一般情形下弱同构但不同构的集合的例子(取 \(\Gamma\) 为两个 2 阶群的直积)。

13.6. 设 \(\mathbf{X}\) 是 \(\mathbf{G}\) 在 \(\mathbf{Q}\left( m\right)\) 上的不可约特征集， \(\mathbf{Y}\) 是 \(G\) 的共轭类集。设群 \({\Gamma }_{0} = {\left( \mathbf{Z}/m\mathbf{Z}\right) }^{ * }\) 通过 \(\chi  \mapsto  {\sigma }_{t}\left( \chi \right)\) 在 \(X\) 上作用，并通过 \(x \mapsto  {x}^{t}\) 在 \(\mathbf{Y}\) 上作用。

(a) 证明 \({\Gamma }_{\mathrm{K}}\) -集 \(\mathbf{X}\) 与 \(\mathbf{Y}\) 弱同构(参见习题 13.5)。

(b) 证明可将 \(X\) (分别 \(Y\) )识别为从 Q-代数 Cent. Q[G](分别 Q ⊗ R(G))到 Q(m) 的同态集。由此推得:当且仅当 \(\mathrm{Q}\left\lbrack  \mathrm{G}\right\rbrack\) 的中心同构于 \(\mathbf{Q} \otimes  \mathbf{R}\left( \mathbf{G}\right)\) 时， \({\Gamma }_{\mathrm{K}}\) -集 \(\mathrm{X}\) 与 \(\mathrm{Y}\) 同构。

(c) 在下列各种情形中证明 \(\mathbf{Q}\left\lbrack  \mathbf{G}\right\rbrack\) 的中心同构于 \(\mathbf{Q} \otimes  \mathbf{R}\left( \mathbf{G}\right)\) :

\(\left( {c}_{1}\right)\) G 是阿贝尔群(利用从 G 到其对偶群 G 的同构，并注意 \(\mathbf{Q}\left\lbrack  G\right\rbrack   = \mathbf{R}\left( \widehat{G}\right)\) )。

\(\left( {c}_{2}\right) \mathrm{G}\) 是一个 \(p\) -群且 \(p \neq  2\) (利用 \({\Gamma }_{\mathbf{O}}\) 为循环群这一事实)。

(关于存在一个群 \(G\) 使得 \(X\) 与 \(Y\) 非 \({\Gamma }_{{Q}^{ - }}\) 同构的例子，见 J. Thompson, J. of Algebra, 14, 1970, pp. 1-4。)

13.7. 设 \(p\) 为素数 \(\neq  2\) 。令 \(\mathrm{G}\) 为 \({\mathrm{{GL}}}_{3}\left( {\mathbf{F}}_{p}\right)\) 的一个 Sylow \(p\) -子群，令 \({\mathrm{G}}^{\prime }\) 为 \(\mathbf{Z}/p\mathbf{Z}\) 与 \(\mathbf{Z}/{p}^{2}\mathbf{Z}\) 的一个非阿贝尔半直积。因此 \(\operatorname{Card}\left( \mathrm{G}\right)  = \operatorname{Card}\left( {\mathrm{G}}^{\prime }\right)  = {p}^{3}.\)

(a) 证明 \(\mathbf{G}\) 与 \({\mathbf{G}}^{\prime }\) 不同构。

(b) 构造 \(\mathbf{G}\) 与 \({\mathbf{G}}^{\prime }\) 的不可约表示。证明 \(\mathbf{Q}\left\lbrack  \mathrm{G}\right\rbrack\) 与 \(\mathbf{Q}\left\lbrack  {\mathrm{G}}^{\prime }\right\rbrack\) 是由 \(\mathbf{Q},p + 1\) 份域 \(\mathbf{Q}\left( p\right)\) 与矩阵代数 \({\mathbf{M}}_{p}\left( {\mathbf{Q}\left( p\right) }\right)\) 的乘积。特别地， \(\mathbf{Q}\left\lbrack  \mathrm{G}\right\rbrack\) 与 \(\mathbf{Q}\left\lbrack  {\mathrm{G}}^{\prime }\right\rbrack\) 同构。

(c) 证明 \({\mathbf{F}}_{p}\left\lbrack  \mathrm{G}\right\rbrack\) 与 \({\mathbf{F}}_{p}\left\lbrack  {\mathrm{G}}^{\prime }\right\rbrack\) 不同构。

13.8. 设 \(\left\{  {{C}_{1},\ldots ,{C}_{d}}\right\}\) 为 G 的循环子群共轭类的一组代表。证明特征 \({1}_{{C}_{1}}^{G},\ldots ,{1}_{{C}_{d}}^{G}\) 构成 \(\mathbf{Q} \otimes  {\mathbf{R}}_{\mathbf{Q}}\left( \mathrm{G}\right)\) 的一组基。

\section*{13.2 域 \(\mathbf{R}\)}

我们保持前述记号，并取实数域 \(\mathbf{R}\) 为底域 \(\mathrm{K}\) 。对应的群 \({\Gamma }_{\mathbf{R}}\) 是 \({\left( \mathbf{Z}/m\mathbf{Z}\right) }^{ * }\) 的子群 \(\{  \pm  1\}\) ；两个元素 \(x,y\) 与 \(\mathrm{G}\) 在 \({\Gamma }_{\mathbf{R}}\) 下共轭当且仅当 \(y\) 与 \(x\) 或與 \({x}^{-1}\) 共轭。与 \({\Gamma }_{\mathbf{R}}\) 中元素 -1 对应的自同构 \({\sigma }_{-1}\) 就是复共轭 \(z \mapsto  {z}^{ * }\) 。若 \(\chi\) 是在 \(\mathrm{C}\) 上的 G 的一个特征，这里一般公式 \({\sigma }_{t}\left( \chi \right)  = {\Psi }^{t}\left( \chi \right)\) 化简为标准公式

\[
\chi {\left( s\right) }^{ * } = \chi \left( {s}^{-1}\right) ,\;\text{ cf. }{2.1},\text{ prop. }1.
\]

定理 31(Frobenius–Schur)。设 \(\rho  : \mathbf{G} \rightarrow  \mathbf{{GL}}\left( \mathbf{V}\right)\) 是在域 \(\mathbf{C}\) 上、具有特征 \(\chi\) 的线性表示 \(\mathbf{G}\) 。要使得 \(\chi\) 取值于 \(\mathbf{R}\) (或使 \(\rho\) 可在 \(\mathbf{R}\) 上实现)，当且仅当 \(\mathbf{V}\) 存在一个在 G 下不变的非退化双线性型(或对称双线性型)。

群 \(\mathbf{G}\) 在 \(\mathbf{V}\) 的对偶空间 \({\mathbf{V}}^{\prime }\) 上具有自然作用，且容易看出相应的特征 \({\chi }^{\prime }\) 由下式给出

\[
{\chi }^{\prime }\left( s\right)  = \chi {\left( s\right) }^{ * } = \chi \left( {s}^{-1}\right) .
\]

要使 \(\chi\) 取实值，必要且充分的条件是 \(\chi  = {\chi }^{\prime }\) ，即 \(G\) 在 \(V\) 与 \({V}^{\prime }\) 中的表示同构。由 \(V\) 到 \({V}^{\prime }\) 的同构对应于一个在 G 下不变的非退化双线性型。因此，存在这样的型对 \(\chi\) 取实值是必要且充分的。

现假设 \(\rho\) 可在 \(\mathbf{R}\) 上实现。这等价于可以将 \(\mathrm{V}\) 写成如下形式

\[
\mathbf{V} = {\mathbf{V}}_{0} \oplus  i{\mathbf{V}}_{0} = \mathbf{C}{ \otimes  }_{\mathbf{R}}{\mathbf{V}}_{0},
\]

其中 \({\mathrm{V}}_{0}\) 是 \(\mathrm{V}\) 的一个 \(\mathbf{R}\) 子空间，且在所有 \({\rho }_{s}\) 作用下保持稳定。我们知道，在 \({V}_{0}\) 上存在一个正定二次型 \({Q}_{0}\) ，它在 \(G\) 作用下保持不变(取任意正定形式的变换之和)。通过标量扩张， \({Q}_{0}\) 在 \(V\) 上定义了一个二次型，并且相关的双线性型是非退化的、对称的，且在 G 作用下保持不变。

反过来，设 V 带有这样的型 \(\mathrm{B}\left( {x,y}\right)\) 。在 V 上选取一个在 G 下不变的正定埃尔米特内积 \(\left( {x \mid  y}\right)\) ；上述论证表明这样的内积存在(参见 1.3)。对于每个 \(x \in  \mathbf{V}\) ，存在唯一的元素 \(\varphi \left( x\right)\) 于 \(\mathbf{V}\) 使得

\[
\mathrm{B}\left( {x,y}\right)  = {\left( \varphi \left( x\right)  \mid  y\right) }^{ * }\;\text{ for all }y \in  \mathrm{V}.
\]

由此定义的映射 \(\varphi  : \mathrm{V} \rightarrow  \mathrm{V}\) 是反线性的且双射。其平方 \({\varphi }^{2}\) 是 \(\mathrm{V}\) 的自同构。对于 \(x,y \in  \mathrm{V}\) ，我们有

\[
\left( {{\varphi }^{2}\left( x\right)  \mid  y}\right)  = \mathrm{B}{\left( \varphi \left( x\right)  \mid  y\right) }^{ * } = \mathrm{B}{\left( y,\varphi \left( x\right) \right) }^{ * } = \left( {\varphi \left( y\right)  \mid  \varphi \left( x\right) }\right) .
\]

由于 \(\left( {\varphi \left( y\right)  \mid  \varphi \left( x\right) }\right)  = {\left( \varphi \left( x\right)  \mid  \varphi \left( y\right) \right) }^{ * }\) ，我们得到

\[
\left( {{\varphi }^{2}\left( x\right)  \mid  y}\right)  = {\left( {\varphi }^{2}\left( y\right)  \mid  x\right) }^{ * } = \left( {x \mid  {\varphi }^{2}\left( y\right) }\right) ,
\]

这表示 \({\varphi }^{2}\) 为埃尔米特型。此外，公式

\[
\left( {{\varphi }^{2}\left( x\right)  \mid  x}\right)  = \left( {\varphi \left( x\right)  \mid  \varphi \left( x\right) }\right)
\]

证明了 \({\varphi }^{2}\) 为正定。然而我们知道，每当 \(u\) 为厄米正定时，存在唯一的厄米正定 \(v\) 使得 \(u = {v}^{2}\) ，且 \(v\) 可写成 \(\mathrm{P}\left( u\right)\) 的形式，其中 \(\mathrm{P}\) 是实系数多项式(若 \(u\) 的特征值为 \({\lambda }_{1},\ldots ,{\lambda }_{n}\) ，则选择 \(\mathrm{P}\) 使得对所有 i 有 \(\mathrm{P}\left( {\lambda }_{i}\right)  = \sqrt{{\lambda }_{i}}\) )。将此应用于 \(u = {\varphi }^{2}\) ，并设 \(\sigma  = \varphi {v}^{-1}\) 。由于 \(v = \mathrm{P}\left( {\varphi }^{2}\right) ,\varphi\) 与 \(v\) 交换，并且我们有 \({\sigma }^{2} = {\varphi }^{2}{v}^{-2} = 1\) 。令 \(\mathrm{V} \; = {\mathbf{V}}_{0} \oplus  {\mathbf{V}}_{1}\) 为相对于 \(\sigma\) 的特征值 +1 与 -1 的 \(\mathbf{V}\) 的分解。由于 \(\sigma\) 是反线性的，乘以 \(i\) 将 \({V}_{0}\) 映为 \({V}_{1}\) 。因此 \(V = {V}_{0} \oplus  i{V}_{0}\) 。另一方面， \(B\left( {x,y}\right)\) 与 \(\left( {x \mid  y}\right)\) 在 \(G\) 下不变这一事实意味着 \(\varphi ,v\) ，且 \(\sigma\) 与所有 \({\rho }_{s}\) 互相交换。由此 \({V}_{0}\) 与 \({V}_{1}\) 在 \({\rho }_{s}\) 下不变，并且我们得到一个在 R 上的 \(V\) 的实现，从而证明了定理 31。

\section*{备注}

(1) 定理 31 同样适用于紧群的表示，参见第 4 章。本节的其它结果亦同。

(2) 记 \({\mathbf{O}}_{n}\left( \mathbf{C}\right)\) (分别记为 \({\mathbf{O}}_{n}\left( \mathbf{R}\right)\) )为在 \(n\) 个变量上的复(分别实)正交群。上面证明的最后部分实际上表明，每个有限(甚至紧)子群都共轭于包含在 \({\mathbf{O}}_{n}\left( \mathbf{R}\right)\) 中的一个子群；这是李群最大紧子群一般定理的一个特例。

\section*{\(\mathrm{G}\) 的三类不可约表示}

设 \(\rho  : G \rightarrow\) GL(V) 是在 \(C\) 上、次数为 \(n\) 的 \(G\) 的一个不可约表示，且 \(\chi\) 为其特征。可能出现三种情形(互不相容):

(1) \(\chi\) 的某个取值非实。按标量限制， \(\rho\) 定义了一个在 \(\mathbf{R}\) 上、次数为 \({2n}\) 的不可约表示，其特征为 \(\chi  + \overline{\chi }\) 。该表示的对易代数为 \(\mathbf{C}\) 。对应的 \(\mathbf{R}\left\lbrack  \mathbf{G}\right\rbrack\) 的单简单分量同构于 \({\mathbf{M}}_{n}\left( \mathbf{C}\right)\) ；其 Schur 指数为 1。

(2) 所有 \(\chi\) 的取值都是实数，且 \(\rho\) 能由一个在 R 上的表示 \({\rho }_{0}\) 实现。该表示 \({\rho }_{0}\) 是不可约的(甚至是绝对不可约的)，其特征为 \(\chi\) 。它的对易代数是 \(\mathbf{R}\) 。对应的 \(\mathbf{R}\left\lbrack  \mathrm{G}\right\rbrack\) 的简单分量同构于 \({\mathbf{M}}_{n}\left( \mathbf{R}\right)\) ；其 Schur 指数为 1。

(3) 所有 \(\chi\) 的取值都是实数，但 \(\rho\) 无法在 \(\mathbf{R}\) 上实现。通过限制标量， \(\rho\) 定义了一个在 \(\mathbf{R}\) 上的不可约表示，次数为 \({2n}\) ，特征为 \({2\chi }\) 。其对易代数在 \(\mathbf{R}\) 上的次数为 4；它同构于四元数域 \(\mathbf{H}\) 。对应的 \(\mathbf{R}\left\lbrack  \mathbf{G}\right\rbrack\) 的简单分量同构于 \({\mathbf{M}}_{n}\left( \mathbf{H}\right)\) ；其 Schur 指数为 2。

此外， \(\mathbf{G}\) 在 \(\mathbf{R}\) 上的每个不可约表示都可由以上三种方法之一得到:这可以通过将 R[G] 分解为若干简单分量的乘积来证明，并注意到这样的分量形如 \({\mathbf{M}}_{n}\left( \mathbf{R}\right) ,{\mathbf{M}}_{n}\left( \mathbf{C}\right)\) 或 \({\mathbf{M}}_{n}\left( \mathbf{H}\right)\) 。 ( \(\mathbf{R}\left\lbrack  \mathrm{G}\right\rbrack\) 是群代数这一事实在此并不重要:对任何实半单代数同样成立。)

类型 1、2 和 3 可由多种方式刻画:

命题 38。

(a) 若在 \(\mathbf{V}\) 上 \(\mathbf{G}\) 不存在非零不变双线性型，则 \(\rho\) 属于类型 1。

(b) 若存在这样的型，则它在相似变换下唯一，非退化，并且要么为对称型要么为交替型。若为对称型，则 \(\rho\) 属于类型 2；若为交替型，则属于类型 3。

在 \(V\) 上的不变双线性型 \(B \neq  0\) 对应一个将 \(V\) 映入其对偶 \({V}^{\prime }\) 的 G-同态 \(b \neq  0\) 。由于 \(V\) 与 \({V}^{\prime }\) 都是不可约的， \(b\) 是同构的，这表明 B 是非退化的。由定理 31，B 的存在意味着 \(\rho\) 属于类型 2 或 3。此外，Schur 引理表明 \(B\) 在相似变换下是唯一的。若我们将 \(B\) 写成 \(B = {B}_{ + } + {B}_{ - }\) 的形式，且 \({B}_{ + }\) 为对称、 \({B}_{ - }\) 为交替，则 \({B}_{ + }\) 和 \({B}_{ - }\) 亦受 G 不变。由于 \(B\) 是唯一的，我们要么有 \({B}_{ - } = 0\) (且 \(B\) 为对称)，要么有 \({B}_{ + } = 0\) (且 \(B\) 为交替)。由定理 31，第一种情况对应类型 2；因此第二种对应类型 3。

命题 39。要使 \(\rho\) 属于类型 1、2 或 3 的充要条件是数目

\[
\left\langle  {1,{\Psi }^{2}\left( \chi \right) }\right\rangle   = \frac{1}{g}\mathop{\sum }\limits_{{s \in  \mathrm{G}}}\chi \left( {s}^{2}\right) ,\;\text{ where }g = \operatorname{Card}\left( \mathrm{G}\right) ,
\]

分别等于 \(0, + 1\) 或 -1。

令 \({\chi }_{\sigma }^{2}\) (分别为 \({\chi }_{\lambda }^{2}\) )为 V 的对称平方(分别交错平方)的特征。则

\[
{\chi }_{\sigma }^{2} = \frac{1}{2}\left( {{\chi }^{2} + {\Psi }^{2}\chi }\right) ,\;{\chi }_{\lambda }^{2} = \frac{1}{2}\left( {{\chi }^{2} - {\Psi }^{2}\chi }\right) ,
\]

参见 2.1，命题 3。令 \({a}_{ + }\) 和 \({a}_{ - }\) 表示 \(\rho\) 的对称平方与交错平方包含单位表示的次数。则

\[
{a}_{ + } = \left\langle  {1,{\chi }_{\sigma }^{2}}\right\rangle  \;\text{ and }\;{a}_{ - } = \left\langle  {1,{\chi }_{\lambda }^{2}}\right\rangle  .
\]

另一方面，V 的对称(分别交错)平方的对偶可被识别为 V 上的对称(分别交错)双线性形式空间。由于对偶表示包含单位表示的次数相同，由命题 38 可得:

\[
{a}_{ + } = {a}_{ - } = 0\;\text{ in case }1,
\]

\[
{a}_{ + } = 1,{a}_{ - } = 0\;\text{ in case }2,
\]

\[
{a}_{ + } = 0,{a}_{ - } = 1\;\text{ in case }3.
\]

由于 \(\left\langle  {1,{\Psi }^{2}\left( \chi \right) }\right\rangle   = {a}_{ + } - {a}_{ - }\) ，我们确实得到 \(0, + 1\) ，以及在三种相应情形下为 -1。命题得证。

\section*{练习}

13.9. 若 \(c\) 为 \(G\) 的一个共轭类，令 \({c}^{-1}\) 表示由所有 \({x}^{-1}\) 为 \(x \in  c\) 的类组成的类。若 \(c = {c}^{-1}\) ，则称 \(c\) 为偶类。

(a) 证明在域 \(\mathbf{C}\) 上取值为实数的不可约特征的数量等于群 \(\mathbf{G}\) 的偶类数量。

(b) 证明若 G 的阶为奇数，则唯一的偶类是单位元的共轭类。由此推得 \(G\) 唯一的实值不可约特征是单位特征(Burnside)。

13.10. 证明 R-代数 (Cent. R[G]) 与 \(\mathbf{R} \otimes  \mathbf{R}\left( G\right)\) 同构。

13.11. 令 \({X}_{2}\) 与 \({X}_{3}\) 分别表示类型 2 与类型 3 的不可约特征集合。证明整数

\[
\mathop{\sum }\limits_{{\chi  \in  {\mathrm{X}}_{2}}}\chi \left( 1\right)  - \mathop{\sum }\limits_{{\chi  \in  {\mathrm{X}}_{3}}}\chi \left( 1\right)
\]

等于满足 \({s}^{2} = 1\) 的 \(s \in  \mathbf{G}\) 元素的数目。(注意该整数等于 \(\sum \chi \left( 1\right) \left\langle  {1,{\Psi }^{2}\left( \chi \right) }\right\rangle   = \left\langle  {1,{\Psi }^{2}\left( {r}_{\mathrm{G}}\right) }\right\rangle\) ，其中 \({r}_{\mathrm{G}}\) 为 G 的正规表示的特征。)

由此得出，若 G 的阶为偶数，则至少有两个不可约特征为类型 2。

13.12.(Burnside)设 G 的阶为奇数。令 \(h\) 为 G 的共轭类数。证明 \(g \equiv  h\) (模 16)。

[使用公式 \(g = \mathop{\sum }\limits_{{i = 1}}^{h}\chi {\left( 1\right) }^{2}\) ，并注意 \({\chi }_{i} \neq  1\) 成对共轭(参见习题 12.9)，且 \({\chi }_{i}\left( 1\right)\) 为奇数。]

若 \(g\) 的每个素因子都同余于 1 (mod.4)，用同样的方法证明 \(g \equiv  h\) (mod. 32)。

参考书目:第二部分

关于半单代数的一般理论，见:

[8] N. Bourbaki. Algèbre，第 VIII 章，Hermann，巴黎，1958 年。

[9] C. Curtis 与 I. Reiner. Representation Theory of Finite Groups and Associative Algebras. Interscience Publishers，纽约，1962 年。

[10] S. Lang. Algebra. Addison-Wesley，纽约，1965 年。

关于诱导表示、Brauer 定理与有理性问题，见 [9]、[10] 及:

[11] W. Feit. Characters of Finite Groups. W. A. Benjamin Publishers，纽约，1967 年。

[12] R. Brauer 与 J. Tate. On the characters of finite groups. Ann. of Math., 62 (1955), 第 1-7 页。

[13] B. Huppert. Endlichen Gruppen I，第 V 章。Springer-Verlag，1967 年。

关于有限群结构的应用，见 [11]、[13]，以及经典著作:

[14] W. Burnside. Theory of Groups of Finite Order，第二版。剑桥，1911；由 Dover Publishers 于 1955 年再版。

有限“代数”群的特征特别有趣；见:

[15] A. Borel 等. Seminar on Algebraic Groups and Related Finite Groups. Lecture Notes in Mathematics，第 131 卷，Springer-Verlag，1970 年。

[16] P. Deligne 与 G. Lusztig. Representations of reductive groups over finite fields, Ann. of Math., 103 (1976), 第 103-161 页。

可以在下列文献中找到表示论中一些主要问题的清单:

[17] R. Brauer. 有限群的表示。现代数学讲义，第 I 卷，T. Saaty 编。John Wiley and Sons，纽约，1963 年。

\section*{布劳尔理论导论}

这里我们关心的是将有限群在特征 \(p\) 下的表示与特征零下的表示相比较。结果本质上归功于 Brauer，可用“格罗滕迪克群”来最方便地描述；这种方法由 Swan 引入(参见 [21]，[22])，他还得到了一些此处未讨论的结果。

第 14 和 15 章为预备。第 16 章给出主要定理的表述；它们在第 17 章中得证。第 18 章我们用“模特征”来表达这些结果。第 19 章包含对 Artin 表示的应用。附录中收集了一些标准定义:格罗滕迪克群、投射模等。

下述论述仅作入门；特别是，不涉及块理论。有兴趣的读者可参考 Curtis-Reiner [9] 和 Feit 的讲义 [20]，以及 Brauer、Feit、Green、Osima、Suzuki 和 Thompson 的原始论文。

第 14 章 群 \({\mathrm{R}}_{\mathrm{K}}\left( \mathrm{G}\right) ,{\mathrm{R}}_{k}\left( \mathrm{G}\right)\) 和 \({\mathrm{P}}_{k}\left( \mathrm{G}\right)\)

\section*{记号}

在第三部分中，G 表示一有限群，且 \(m\) 是 G 中元素阶的最小公倍数。若一个域包含 \(m\) 次单位根(参见 12.3 定理 24)，则称该域对 G 足够大。

所有讨论的模均假定为有限生成。

我们记 \(\mathrm{K}\) 为相对于离散赋值 \(v\) 完备的域(参见附录)，其赋值环为 A，极大理想为 \(\mathfrak{m}\) ，剩余域为 \(k = \mathbf{A}/\mathfrak{m}\) 。我们假定 \(\mathbf{K}\) 的特征为零且 \(k\) 的特征为 \(p > 0\) (因此“模 m 约降”是从特征零到特征 \(p\) )。

\section*{14.1 环 \({R}_{K}\left( G\right)\) 与 \({R}_{k}\left( G\right)\)}

若 \(L\) 为一域，我们记 \({R}_{L}\left( G\right)\) 为有限生成 L[G]-模范畴的格罗滕迪克群(参见附录)。相对于外张量积(以 L 为基)它是带单位的交换环。若 \(E\) 为一个 \(L\left\lbrack  G\right\rbrack\) -模，我们记 \(\left\lbrack  E\right\rbrack\) 为其在 \({R}_{L}\left( G\right)\) 中的像；所有 \(\left\lbrack  \mathrm{E}\right\rbrack\) 的集合记为 \({\mathrm{R}}_{\mathrm{L}}^{ + }\left( \mathrm{G}\right)\) 。

记 \({\mathrm{S}}_{\mathrm{L}}\) 为简单 L[G]-模的同构类集合(即 \(\mathrm{G}\) 在 \(\mathrm{L}\) 上的不可约表示)。

命题 40. 以所有元素 [E](其中 \(\mathrm{E} \in  {\mathrm{S}}_{\mathrm{L}}\) )构成的族为群 \({\mathrm{R}}_{\mathrm{L}}\left( \mathrm{G}\right)\) 的一组基。

设 \(\mathrm{R}\) 为基为 \({\mathrm{S}}_{\mathrm{L}}\) 的自由 \(\mathrm{Z}\) -模。各个 \(\left\lbrack  \mathrm{E}\right\rbrack\) 、 \(\mathrm{E} \in  {\mathrm{S}}_{\mathrm{L}}\) 的族定义了一个同态 \(\alpha  : \mathrm{R} \rightarrow  {\mathrm{R}}_{\mathrm{L}}\left( \mathrm{G}\right)\) 。另一方面，若 \(\mathrm{F}\) 为一个 \(\mathrm{L}\left\lbrack  \mathrm{G}\right\rbrack\) -模，且若 \(\mathrm{E} \in  {\mathrm{S}}_{\mathrm{L}}\) ，令 \({l}_{\mathrm{E}}\left( \mathrm{F}\right)\) 表示 \(\mathrm{E}\) 在 \(\mathrm{F}\) 的组成列中出现的次数；显然 \({l}_{\mathrm{E}}\) 是 \(F\) 的加法函数。因此存在一个同态 \({\beta }_{E} : {R}_{L}\left( G\right)  \rightarrow  Z\) 使得 对所有 \(\mathrm{F}\) 有 \({\beta }_{\mathrm{E}}\left( \left\lbrack  \mathrm{F}\right\rbrack  \right)  = {l}_{\mathrm{E}}\left( \mathrm{F}\right)\) 。 \({\beta }_{\mathrm{E}}\) 定义了一个同态

\[
\beta  : {\mathrm{R}}_{\mathrm{L}}\left( \mathrm{G}\right)  \rightarrow  \mathrm{R}\text{ , }
\]

并且显然 \(\alpha\) 与 \(\beta\) 互为逆。该命题得证。

更一般地，相同的论证适用于任意环上有限长度模的范畴。

还注意 \({\mathbf{R}}_{\mathbf{L}}^{ + }\left( \mathbf{G}\right)\) 的元素正是基 \({\left( \left\lbrack  E\right\rbrack  \right) }_{E \in  {S}_{I}}\) 的元素以非负整数系数的线性组合。前述讨论尤其适用于域 \(\mathrm{K}\) 与 \(k\) 。由于 \(\mathrm{K}\) 特征为零， \(\mathrm{K}\left\lbrack  \mathrm{G}\right\rbrack\) -模 \(\mathrm{E}\) 的特征 \({\chi }_{\mathrm{E}}\) 已可定义；它是 E 的加法函数。由线性性，我们得到从 \({\mathrm{R}}_{\mathrm{K}}\left( \mathrm{G}\right)\) 到 \(\mathrm{G}\) 的类函数环的线性映射 \(x \mapsto  {\chi }_{x}\) 。该映射事实上将 \({\mathrm{R}}_{\mathrm{K}}\left( \mathrm{G}\right)\) 同构到 \(\mathrm{G}\) 在 K 上的虚特征群，我们常将二者认同(这解释了 12.1 中的记法)。我们也说 \({\chi }_{x}\) 是元素 \(x \in  {\mathrm{R}}_{\mathrm{K}}\left( \mathrm{G}\right)\) 的特征(或虚特征)。

我们将在第 18 章看到，对 \(k\) 有一个类似的关于布劳尔模特征的结果。

备注。如果 \(\mathrm{E}\) 和 \({\mathrm{E}}^{\prime }\) 是两个 \(\mathrm{K}\left\lbrack  \mathrm{G}\right\rbrack\) -模 且 \(\left\lbrack  \mathrm{E}\right\rbrack   = \left\lbrack  {\mathrm{E}}^{\prime }\right\rbrack\) 在 \({\mathrm{R}}_{\mathrm{K}}\left( \mathrm{G}\right)\) 中，则 \(\mathrm{E}\) 与 \({\mathrm{E}}^{\prime }\) 同构:此由 \(\mathrm{E}\) 与 \({\mathrm{E}}^{\prime }\) 为半单模推出。当 \(p\) 整除 \(G\) 的阶时，类似结论对 \(k\left\lbrack  \mathrm{G}\right\rbrack\) -模 不成立，因为存在非半单的模。

\section*{14.2 群 \({\mathrm{P}}_{k}\left( \mathrm{G}\right)\) 与 \({\mathrm{P}}_{\mathrm{A}}\left( \mathrm{G}\right)\)}

这些被定义为 \(k\left\lbrack  \mathrm{G}\right\rbrack\) -模(分别为A[G]-模)范畴的格罗滕迪克群，这些模是射影的(参见附录)。对 \({P}_{k}^{ + }\left( G\right)\) 和 \({P}_{A}^{ + }\left( G\right)\) 也作类似定义。

若 \(E\) (分别 \(F\) )是一个 \(k\left\lbrack  G\right\rbrack\) -模(分别一个投射的 \(k\left\lbrack  G\right\rbrack\) -模)，则 \(\mathrm{E}{ \otimes  }_{k}\mathrm{\;F}\) 是一个投射的 \(k\left\lbrack  \mathrm{\;G}\right\rbrack\) -模(只需在 \(\mathrm{F}\) 为自由时检验，在该情形显而易见)。由此我们得到 \({\mathrm{P}}_{k}\left( \mathrm{G}\right)\) 上的 \({\mathrm{R}}_{k}\left( \mathrm{G}\right)\) -模 结构。

\section*{14.3 \({\mathrm{P}}_{k}\left( \mathrm{G}\right)\) 的结构}

由于 \(k\left\lbrack  \mathrm{G}\right\rbrack\) 是阿廷环，我们可以谈论一个 \(k\left\lbrack  \mathrm{G}\right\rbrack\) -模 M 的投射包(参见 Gabriel [23] 或 Giorgiutti [24])。现简要回顾其含义:

若满足 \(f\left( {\mathrm{M}}^{\prime }\right)  = \mathrm{M}\) 则称模同态 \(f : {\mathrm{M}}^{\prime } \rightarrow  \mathrm{M}\) 为极小本质(essential)；并且若对 \({\mathrm{M}}^{\prime }\) 的所有真子模 \({\mathrm{M}}^{\prime \prime }\) 都有 \(f\left( {\mathrm{M}}^{\prime \prime }\right)  \neq  \mathrm{M}\) 。 \(\mathrm{M}\) 的一个投射包(projective envelope)是一个投射模 \(\mathrm{P}\) ，配备有一个本质同态 \(f : \mathrm{P} \rightarrow  \mathrm{M}\) 。

\section*{命题 41。}

(a) 每个模 \(\mathrm{M}\) 都有一个投射包，且在同构下唯一。

(b) 若 \({\mathrm{P}}_{i}\) 是 \({\mathrm{M}}_{i}\left( {i = 1,\ldots ,n}\right)\) 的投射包，则这些 \({\mathrm{P}}_{i}\) 的直和是对应 \({\mathrm{M}}_{i}\) 的直和的投射包。

(c) 如果 \(\mathrm{P}\) 是一个投射模，且 \(\mathrm{E}\) 是它最大的半单商模，则 \(\mathrm{P}\) 是 \(\mathrm{E}\) 的投射包络。

我们证明 (a)。将 \(M\) 写成 \(L/R\) 的形式，其中 \(L\) 是射影的且 \(R\) 是 \(L\) 的子模(例如可取 \(L\) 为自由模)。对于 \(N \subset  R\) ，令 \({f}_{N}\) 为将 \(L/N\) 映到 \(M = L/R\) 的自然同态。现在在 \(R\) 中取使得 \({f}_{N}\) 为本质(essential)的最小子模 \(N\) ；这样的子模存在，因为 \({f}_{R}\) 是本质的，且 \(k\left\lbrack  \mathrm{G}\right\rbrack\) 是 Artinian 的。置 \(\mathrm{P} = \mathrm{L}/\mathrm{N}\) ，并在那些投影 \(Q \rightarrow  P\) 为满的 L 的子模中取最小的一个 \(\mathrm{Q}\) 。由于 \(L\) 是射影的，投影 \(p : L \rightarrow  P = L/N\) 可以上升为 \(q : L \rightarrow  Q\) ，且 \(Q\) 的最小性表明 \(q\left( L\right)  = Q\) 。令 \({N}^{\prime }\) 为 \(q\) 的核。投影 \({f}_{{\mathrm{N}}^{\prime }} : \mathrm{L}/{\mathrm{N}}^{\prime } \rightarrow  \mathrm{L}/\mathrm{R}\) 分解为 \(\mathrm{L}/{\mathrm{N}}^{\prime } = \mathrm{Q} \rightarrow  \mathrm{L}/\mathrm{N} \rightarrow  \mathrm{L}/\mathrm{R}\) ，且这两个因子都是本质的。既然 \({\mathrm{N}}^{\prime }\) 包含于 \(\mathrm{N}\) ， \(\mathrm{N}\) 的最小性意味着 \({\mathrm{N}}^{\prime } = \mathrm{N}\) ，即 \(\mathrm{Q} \rightarrow  \mathrm{P}\) 是同构。于是模 \(\mathrm{L}\) 为直和 \(\mathrm{L} = \mathrm{N} \oplus  \mathrm{Q}\) ，这表明 \(\mathrm{P} = \mathrm{L}/\mathrm{N}\) 是射影的。然后显然 \(\mathrm{P} \rightarrow  \mathrm{M}\) 是 \(\mathrm{M}\) 的射影包。

设 \({\mathrm{P}}^{\prime } \rightarrow  \mathrm{M}\) 为另一个 \(\mathrm{M}\) 的投射包络。利用 \(\mathrm{P}\) 为投射模，存在 \(g : \mathrm{P} \rightarrow  {\mathrm{P}}^{\prime }\) 使得该三角形

\begin{center}
\includegraphics[max width=0.2\textwidth]{images/bo_d6kn08k601uc73e1e9r0_122_652_1380_219_162_0.jpg}
\end{center}
\hspace*{3em} 

交换。 \(g\left( \mathrm{P}\right)\) 在 \(\mathrm{M}\) 中的像是 \(\mathrm{M}\) ；因 \({\mathrm{P}}^{\prime } \rightarrow  \mathrm{M}\) 为本质，这蕴含 \(g\left( \mathrm{P}\right)  = {\mathrm{P}}^{\prime }\) ，故 \(g\) 为满射。由于 \({\mathrm{P}}^{\prime }\) 为投射模， \(\mathrm{P} \rightarrow  {\mathrm{P}}^{\prime }\) 的核 \(\mathrm{S}\) 在 \(\mathrm{P}\) 中是直和因子，这表明 \(\mathrm{P}\) 分解为 \(S \oplus  {P}^{\prime }\) 。利用 \(P \rightarrow  M\) 为本质，我们得出 \(S = 0\) ，即 \(P \rightarrow  {P}^{\prime }\) 为同构。由此完成了 (a) 的证明。(b) 和 (c) 显然，留给读者(详见 [23],[24])。

注意，在情况 (c) 中， \(E\) 是 \(P\) 对 \({rP}\) 的商，其中 \(r\) 是 \(k\left\lbrack  \mathrm{G}\right\rbrack\) 的激根(最大幂零理想)；这是因为半单 \(k\left\lbrack  \mathrm{G}\right\rbrack\) -模恰好是被 \(\mathrm{r}\) 扰灭的那些模。此外，由 (b) 可知，E 作为简单模直和的每一种分解都对应于 P 的相应分解。因此我们得到:

推论 1. 每个投射 \(k\left\lbrack  \mathrm{G}\right\rbrack\) -模都是投射不可约 \(k\left\lbrack  \mathrm{G}\right\rbrack\) -模的直和；此分解在同构意义下唯一。投射不可约 \(k\left\lbrack  \mathrm{G}\right\rbrack\) -模是简单 \(k\left\lbrack  \mathrm{\;G}\right\rbrack\) -模的投射包。

推论 2. 对每个 \(\mathrm{E} \in  {\mathrm{S}}_{k}\) ，令 \({\mathrm{P}}_{\mathrm{E}}\) 是 \(\mathrm{E}\) 的一个投射包。则 \(\left\lbrack  {\mathrm{P}}_{\mathrm{E}}\right\rbrack  ,\mathrm{E} \in  {\mathrm{S}}_{k}\) 构成 \({\mathrm{P}}_{k}\left( \mathrm{G}\right)\) 的一组基。

推论 3. 两个投射 \(k\left\lbrack  \mathrm{G}\right\rbrack\) -模 \(\mathrm{P}\) 与 \({\mathrm{P}}^{\prime }\) 同构当且仅当它们在 \({\mathrm{P}}_{k}\left( \mathrm{G}\right)\) 中的类 \(\left\lbrack  \mathrm{P}\right\rbrack\) 与 \(\left\lbrack  {\mathrm{P}}^{\prime }\right\rbrack\) 相等。

更确切地，若 \(\left\lbrack  \mathrm{P}\right\rbrack   = \mathop{\sum }\limits_{{\mathrm{E} \in  {\mathrm{S}}_{L}}}{n}_{\mathrm{E}}\left\lbrack  {\mathrm{P}}_{\mathrm{E}}\right\rbrack\) ，则模 \(\mathrm{P}\) 同构于 \(\Pi {\left( {\mathrm{P}}_{\mathrm{E}}\right) }^{{n}_{\mathrm{E}}}\) 。

\section*{练习}

14.1. 证明 \(k\left\lbrack  \mathrm{\;G}\right\rbrack\) 是一个注入的 \(k\left\lbrack  \mathrm{G}\right\rbrack\) -模。由此得出:一个 \(k\left\lbrack  \mathrm{G}\right\rbrack\) -模当且仅当它是注入的才是投射的，且投射不可约的 \(k\left\lbrack  \mathrm{G}\right\rbrack\) -模是简单 \(k\left\lbrack  \mathrm{G}\right\rbrack\) -模的注入包(参见练习 14.6)。

\section*{14.4 \({P}_{A}\left( G\right)\) 的结构}

以下结论是众所周知的:

引理 20. 设 \(\Lambda\) 为交换环， \(\mathrm{P}\) 为 \(\Lambda \left\lbrack  \mathrm{G}\right\rbrack\) -模。为了使 P 在 \(\Lambda \left\lbrack  \mathrm{G}\right\rbrack\) 上为投射模，当且仅当它在 \(\Lambda\) 上为投射模，并且存在一个 \(\Lambda\) -端同态 \(u\) 于 \(\mathrm{P}\) 使得

\[
\mathop{\sum }\limits_{{s \in  \mathrm{G}}}s \cdot  u\left( {{s}^{-1}x}\right)  = x\text{ for all }x \in  \mathrm{P}.
\]

如果 \(\mathrm{P}\) 在 \(\Lambda \left\lbrack  \mathrm{G}\right\rbrack\) 上是射影的，则在 \(\Lambda\) 上也是射影的:这是因为 \(\Lambda \left\lbrack  \mathbf{G}\right\rbrack\) 是 \(\Lambda\) -自由的。反过来，假设作为 \(\Lambda\) -模的底层 \({\mathrm{P}}_{0}\) 是射影的，并设 \(\mathrm{Q} = \Lambda \left\lbrack  \mathrm{G}\right\rbrack  { \otimes  }_{\Lambda }{\mathrm{P}}_{0}\) 。 \(\Lambda \left\lbrack  \mathrm{G}\right\rbrack\) -模 \(\mathrm{Q}\) 是射影的。此外，恒等映射 \({\mathrm{P}}_{0} \rightarrow  \mathrm{P}\) 可扩张为一个满的 \(\Lambda \left\lbrack  \mathrm{G}\right\rbrack\) -同态 \(q : \mathbf{Q} \rightarrow  \mathbf{P}\) 。因此 \(\mathbf{P}\) 为射影当且仅当存在一个 \(\Lambda \left\lbrack  \mathrm{G}\right\rbrack\) -同态 \(v : \mathrm{P} \rightarrow  \mathrm{Q}\) 使得 \(q \circ  v = 1\) 。容易看出每个 \(\Lambda \left\lbrack  \mathrm{G}\right\rbrack\) -同态 \(v : \mathrm{P} \rightarrow  \mathrm{Q}\) 具有形如

\[
x \mapsto  \mathop{\sum }\limits_{{s \in  \mathrm{G}}}s \otimes  u\left( {{s}^{-1}x}\right)
\]

带有 \(u \in  {\operatorname{End}}_{\Lambda }\left( {\mathrm{P}}_{0}\right)\) 。要得到 \(q \circ  v = 1\) 的充要条件是对所有 \(x \in  \mathrm{P}\) 都有 \(\mathop{\sum }\limits_{{s \in  \mathrm{G}}}s \cdot  u\left( {{s}^{-1}x}\right)  = x\) 。这证明了引理。

引理 21. 设 \(\Lambda\) 为局部环，剩余域为 \({k}_{\Lambda } = \Lambda /{\mathfrak{m}}_{\Lambda }\) 。

(a) 设 \(\mathrm{P}\) 为 \(\Lambda\) -自由的 \(\Lambda \left\lbrack  \mathrm{G}\right\rbrack\) -模。为了使 \(\mathrm{P}\) 在 \(\Lambda \left\lbrack  \mathrm{G}\right\rbrack\) 上为射影的，必要且充分的条件是 \({k}_{\Lambda }\left\lbrack  \mathrm{G}\right\rbrack\) -模 \(\overline{\mathrm{P}} = \mathrm{P} \otimes  {k}_{\Lambda }\) 为射影的。

(b) 两个射影的 \(\Lambda \left\lbrack  \mathrm{G}\right\rbrack\) -模 \(\mathrm{P}\) 和 \({\mathrm{P}}^{\prime }\) 同构当且仅当相应的 \({k}_{\Lambda }\left\lbrack  \mathrm{G}\right\rbrack\) -模 \(\overline{\mathrm{P}}\) 和 \(\overline{{\mathrm{P}}^{\prime }}\) 同构。

如果 \(\mathrm{P}\) 是 \(\Lambda \left\lbrack  \mathrm{G}\right\rbrack\) -投射的，则 \(\overline{\mathrm{P}}\) 是 \({k}_{\Lambda }\left\lbrack  \mathrm{G}\right\rbrack\) -投射的。反之，若满足此条件，则前述引理表明存在一个 \({k}_{{\Lambda }^{ - }}\) 端自同态 \(\bar{u}\) 于 \(\bar{P}\) 使得 \(\mathop{\sum }\limits_{{s \in  \mathrm{G}}}s \cdot  \bar{u} \cdot  {s}^{-1} = 1\) 。通过提升 \(\bar{u}\) ，我们得到一个作用在 \(P\) 上的 \(\Lambda\) -端自同态 \(u\) ，使得 \({u}^{\prime } \equiv  1\left( {{\;\operatorname{mod}\;.}{\mathfrak{m}}_{\Lambda }P}\right)\) ，其中 \({u}^{\prime } = \mathop{\sum }\limits_{{s \in  \mathrm{G}}}s \cdot  u \cdot  {s}^{-1}\) 。因此 \({u}^{\prime }\) 是 \(\mathrm{P}\) 的自同构，且与 G 交换。于是 \(\mathop{\sum }\limits_{{s \in  \mathrm{G}}}s \cdot  \left( {u{u}^{\prime  - 1}}\right)  \cdot  {s}^{-1} = 1\) ，这表明 \(\mathrm{P}\) 在 \(\Lambda \left\lbrack  \mathrm{G}\right\rbrack\) 上是投射的，并证明了 (a)。

若 \(\mathrm{P}\) 与 \({\mathrm{P}}^{\prime }\) 都是投射的，且 \(\bar{w} : \overline{\mathrm{P}} \rightarrow  {\overline{\mathrm{P}}}^{\prime }\) 是一个 \({k}_{\Lambda }\left\lbrack  \mathrm{G}\right\rbrack\) -同态， \(\mathrm{P}\) 的投射性表明存在一个 \(\Lambda \left\lbrack  \mathrm{G}\right\rbrack\) -同态 \(w : \mathrm{P} \rightarrow  {\mathrm{P}}^{\prime }\) 提升了 \(\bar{w}\) 。若且仅若 \(\bar{w}\) 还是一个同构，则 Nakayama 引理(或一个基本的行列式论证)表明 \(w\) 为同构。这证得了 (b)。

我们现在回到环 A:

命题 42。

(a) 设 E 为一个 A[G]-模。E 为一个投射的 A[G]-模的充要条件是 \(\mathrm{E}\) 在 \(\mathrm{A}\) 上为自由，并且 \(\mathrm{E}\) 的约化 \(\overline{\mathrm{E}} = \mathrm{E}/\mathrm{{mE}}\) 为一个投射的 \(k\left\lbrack  \mathrm{G}\right\rbrack\) -模。

(b) 若 \(\mathrm{F}\) 为一个投射的 \(k\left\lbrack  \mathrm{G}\right\rbrack\) -模，则存在唯一(至同构)一个投射的 \(\mathrm{A}\left\lbrack  \mathrm{G}\right\rbrack\) -模，其约化 \({\;\operatorname{mod}\;.}\mathfrak{m}\) 与 F 同构。

(a) 部分及 (b) 中的唯一性来自引理 20 和 21。剩下的就是证明 (b) 中的存在性:

设 \(\mathrm{F}\) 为投射的 \(k\left\lbrack  \mathrm{G}\right\rbrack\) -模。若 \(n \geq  1\) 为整数，则令 \({\mathrm{A}}_{n}\) 表示环 \(A/{\mathfrak{m}}^{n}\) ；因此 \({A}_{1} = k\) 和 \(A\) 为 \({A}_{n}\) 的投射极限。环 \({A}_{n}\) 与 \({A}_{n}\left\lbrack  G\right\rbrack\) 是 Artinian 的。前节的论证表明 \({A}_{n}\left\lbrack  G\right\rbrack\) -模 \(F\) 有一个投射包 \({P}_{n}\) ，且 \({P}_{n}\) 在 \({A}_{n}\) 上是自由的。投影 \({P}_{n} \rightarrow  F\) 经由 \({P}_{n}/m{P}_{n} \rightarrow  F\) 分解，且 \({P}_{n}/m{P}_{n} \rightarrow  F\) 是满的。由于 \(F\) 为 \(k\left\lbrack  G\right\rbrack\) -投射，存在一个 \(k\left\lbrack  G\right\rbrack\) -子模 \({F}^{\prime }\) ⊂ \({\mathrm{P}}_{n}/\mathrm{m}{\mathrm{P}}_{n}\) ，其映为与 \(\mathrm{F}\) 同构。 \({\mathrm{F}}^{\prime }\) 在 \({\mathrm{P}}_{n}\) 中的逆像 \({\mathrm{P}}^{\prime }\) 的像为 \(F\) 。因 \({P}_{n} \rightarrow  F\) 为本质的，遂得 \({P}^{\prime } = {P}_{n}\) ，即 \({\mathrm{P}}_{n}/\mathfrak{m}{\mathrm{P}}_{n} \rightarrow  \mathrm{F}\) 是同构。并且这些 \({\mathrm{P}}_{n}\) 构成一个投射系统。它们的投射极限 \(\mathrm{P}\) 是一个 A-自由的 A[G]-模，使得 \(\overline{\mathbf{P}} = \mathbf{P}/\mathfrak{m}\mathbf{P}\) 与 \(\mathbf{F}\) 同构。结合 (a)，此即完成 (b) 的证明。

推论 1. 每个投射的 A[G]-模都可分解为投射不可约 A[G]-模的直和；该分解在同构意义下唯一。投射不可约的 A[G]-模由其约化 \({\;\operatorname{mod.m}}\) 唯一确定至同构，且该约化为投射不可约的 \(k\left\lbrack  \mathrm{G}\right\rbrack\) -模(即某单 \(k\left\lbrack  \mathrm{G}\right\rbrack\) -模的投射包)。

此由前述命题及关于投射 \(k\left\lbrack  \mathrm{\;G}\right\rbrack\) -模的已知结果得到。由此我们得到:

推论 2. 两个投射的 A[G]-模当且仅当其在 \({\mathrm{P}}_{\mathrm{A}}\left( \mathrm{G}\right)\) 中的 \(\left\lbrack  \mathrm{P}\right\rbrack   = \left\lbrack  \mathrm{Q}\right\rbrack\) 相同时同构。

推论 3. 模 \(\mathfrak{m}\) 的约化给出从 \({\mathrm{P}}_{\mathrm{A}}\left( \mathrm{G}\right)\) 到 \({\mathrm{P}}_{k}\left( \mathrm{G}\right)\) 的同构；该同构将 \({\mathrm{P}}_{\mathrm{A}}^{ + }\left( \mathrm{G}\right)\) 映为 \({\mathrm{P}}_{k}^{ + }\left( \mathrm{G}\right)\) 。

将 \({\mathrm{P}}_{\mathrm{A}}\left( \mathrm{G}\right)\) 与 \({\mathrm{P}}_{k}\left( \mathrm{G}\right)\) 认同。

因此我们可以识别出 \({\mathrm{P}}_{\mathrm{A}}\left( \mathrm{G}\right)\) 和 \({\mathrm{P}}_{k}\left( \mathrm{G}\right)\) 。

关于“准局小”范畴中射影包的通论，请参见 Demazure-Gabriel [23]。

\section*{练习}

14.2. 设 \(\Lambda\) 为交换环， \(\mathrm{P}\) 为在 \(\Lambda\) 上射影的 \(\Lambda \left\lbrack  \mathrm{G}\right\rbrack\) -模。证明下列性质等价:

(i) \(\mathrm{P}\) 是一个射影的 \(\Lambda \left\lbrack  \mathrm{G}\right\rbrack\) -模。

(ii) 对于 \(\Lambda\) 的每个极大理想 \(\mathfrak{p}\) ， \(\left( {\Lambda /\mathfrak{p}}\right) \left\lbrack  \mathrm{G}\right\rbrack\) -模 \(\mathrm{P}/\mathfrak{p}\mathrm{P}\) 是射影的。

14.3. (a) 设 \(B\) 为一 A-代数，作为 \(A\) -模自由且有限秩，且令 \(\bar{u}\) 为 \(\overline{\mathbf{B}} = \mathbf{B}/m\mathbf{B}\) 的一幂等元。证明存在 \(\mathbf{B}\) 的一幂等元，其模 \(\mathfrak{m}B\) 约化等于 \(\bar{u}\) 。

(b) 设 \(\mathrm{P}\) 为射影的 \(\mathrm{A}\left\lbrack  \mathrm{G}\right\rbrack\) -模，且设 \(\mathrm{B} = {\operatorname{End}}^{\mathrm{G}}\left( \mathrm{P}\right)\) 。证明 \(\mathrm{B}\) 是 A-自由的，且 \(\overline{\mathbf{B}}\) 可被识别为 \(\overline{\mathbf{P}} = \mathbf{P}/\mathfrak{m}\mathbf{P}\) 的 G-自同构代数。由此与 (a) 推得，每个将 \(\overline{\mathbf{P}}\) 分解为若干 \(k\left\lbrack  \mathrm{\;G}\right\rbrack\) -模直和的分解都能提升为 P 的相应分解。

(c) 用 (b) 给出命题 42(b) 中存在性的另一证明。[把 F 写为自由模的直因子 \(\overline{\mathbf{P}}\) ，把 \(\overline{\mathbf{P}}\) 提升到自由模，并应用 (b)。]

\section*{14.5 对偶}

\({\mathrm{R}}_{\mathrm{K}}\left( \mathrm{G}\right)\) 与 \({\mathrm{R}}_{\mathrm{K}}\left( \mathrm{G}\right)\) 之间的对偶

令 \(\mathrm{E}\) 与 \(\mathrm{F}\) 为 \(\mathrm{K}\left\lbrack  \mathrm{G}\right\rbrack\) -模，并设

\[
\langle \mathrm{E},\mathrm{F}\rangle  = \dim {\operatorname{Hom}}^{\mathrm{G}}\left( {\mathrm{E},\mathrm{F}}\right) ,\;\text{ cf. 7.1. }
\]

映射 \(\left( {E,F}\right)  \mapsto  \langle E,F\rangle\) 在(关于正合列)意义下是“二线性的”，因此定义了一个双线性形式

\[
{\mathrm{R}}_{\mathrm{K}}\left( \mathrm{G}\right)  \times  {\mathrm{R}}_{\mathrm{K}}\left( \mathrm{G}\right)  \rightarrow  \mathbf{Z}
\]

我们记为 \(\langle e,f\rangle\) 或 \(\langle e,f{\rangle }_{\mathrm{K}}\) 。单模的同类 [E] \(E \in  {S}_{K}\) 彼此正交，且 \(\langle E,E\rangle\) 等于自同构域 \({\operatorname{End}}^{G}\left( E\right)\) 的维数 \({d}_{E}\) ，因此 \({d}_{E} \geq  1\) ，当且仅当 \(E\) 绝对简单(即相应的表示是绝对不可约)时取等号，参见 12.1。

当 K 足够大时，由定理 24 可知每个简单的 K[G]-模在代数闭包上仍为简单。因此上述双线性型在 \(\mathbf{Z}\) 上是非退化的，意即它给出了 \({\mathbf{R}}_{\mathbf{K}}\left( \mathbf{G}\right)\) 到其对偶的同构。

\({\mathrm{R}}_{k}\left( \mathrm{G}\right)\) 与 \({\mathrm{P}}_{k}\left( \mathrm{G}\right)\) 之间的对偶性

若 \(\mathrm{E}\) 为投射的 \(k\left\lbrack  \mathrm{G}\right\rbrack\) -模且 \(\mathrm{F}\) 为任意 \(k\left\lbrack  \mathrm{G}\right\rbrack\) -模，则置

\[
\langle \mathrm{E},\mathrm{F}\rangle  = \dim {\mathrm{{Hom}}}^{\mathrm{G}}\left( {\mathrm{E},\mathrm{F}}\right) .
\]

因此我们得到 E 与 F 的一个双线性映射(得益于假设 E 为投射模)，从而得到一个双线性型

\[
{\mathrm{P}}_{k}\left( \mathrm{G}\right)  \times  {\mathrm{R}}_{k}\left( \mathrm{G}\right)  \rightarrow  \mathbf{Z}
\]

表示为 \(\langle e,f\rangle\) 或 \(\langle e,f{\rangle }_{k}\) 。如果 \(\mathrm{E},{\mathrm{E}}^{\prime } \in  {\mathrm{S}}_{k}\) ，我们有

\[
{\mathrm{{Hom}}}^{\mathrm{G}}\left( {{\mathrm{P}}_{\mathrm{E}},{\mathrm{E}}^{\prime }}\right)  = {\mathrm{{Hom}}}^{\mathrm{G}}\left( {\mathrm{E},{\mathrm{E}}^{\prime }}\right) ,
\]

其中 \({\mathrm{P}}_{\mathrm{E}}\) 表示 \(\mathrm{E}\) 的投射包。若 \(\mathrm{E} \neq  {\mathrm{E}}^{\prime }\) ，可见 \(\left\lbrack  {\mathrm{P}}_{\mathrm{E}}\right\rbrack\) 与 \(\left\lbrack  {\mathrm{E}}^{\prime }\right\rbrack\) 正交；对于 \(\mathrm{E} = {\mathrm{E}}^{\prime }\) 我们有

\[
\left\langle  {{\mathrm{P}}_{\mathrm{E}},\mathrm{E}}\right\rangle   = \dim {\operatorname{End}}^{\mathrm{G}}\left( \mathrm{E}\right) .
\]

如前所述， \({d}_{\mathrm{E}} = 1\) 当且仅当 \(\mathrm{E}\) 绝对单。

假设 \(\mathrm{K}\) 足够大，使得 \(k\) 包含第 \(m\) 个单位根。则对于每个 \(\mathrm{E} \in  {\mathrm{S}}_{k}\) 我们都有 \({d}_{\mathrm{E}} = 1\) (见下文)。因此双线性型 \(\langle\) ， \({\rangle }_{k}\) 在 \(\mathbf{Z}\) 上非退化，且基 \(\left\lbrack  \mathrm{E}\right\rbrack\) 与 \(\left\lbrack  {\mathrm{P}}_{\mathrm{E}}\right\rbrack \; \left( {\mathrm{E} \in  {\mathrm{S}}_{k}}\right)\) 关于此型互为对偶。

\section*{备注}

若 \(\mathrm{K}\) 足够大则 \({d}_{\mathrm{E}} = 1\) 的这一事实可用多种方法证明:

(1) 一旦我们知道同态 \(d : {\mathrm{R}}_{\mathrm{K}}\left( \mathrm{G}\right)  \rightarrow  {\mathrm{R}}_{k}\left( \mathrm{G}\right)\) 是满的(参见第 16 章，第 33 定理)，便可通过“模 m 约减”从定理 24 得到此结论。

(2) 也可利用在 \(k\) 上的舒尔指标等于 1(参见 14.6)。这将证明归约到表示在 \(\mathbf{G}\) 的扩张上其特征值总在 \(k\) 中，而这又来自它们是第 \(m\) 个单位根之和这一事实。

\section*{练习}

14.4. 若 \(E\) 是一个 \(k\left\lbrack  G\right\rbrack\) -模，我们以 \({E}^{\prime }\) 表示其对偶。令 \({H}^{0}\left( {G,E}\right)\) 为 \(E\) 中被 \(G\) 固定的元素所成的子空间， \({H}_{0}\left( {G,E}\right)\) 为 \(E\) 除以由 \({sx} - x\) 生成的子空间所得的商，且有 \(x \in  E\) 和 \(s \in  \mathbf{G}\) 。

(a) 证明:若 \(\mathbf{E}\) 为投射的，则映射 \(x \mapsto  \mathop{\sum }\limits_{{s \in  \mathbf{G}}}{sx}\) 通过取商诱导出一个将 \({\mathrm{H}}_{0}\left( {\mathrm{G},\mathrm{E}}\right)\) 同构到 \({\mathrm{H}}^{0}\left( {\mathrm{G},\mathrm{E}}\right)\) 的同构。

(b) 证明 \({\mathrm{H}}^{0}\left( {\mathrm{G},\mathrm{E}}\right)\) 是 \({\mathrm{H}}_{0}\left( {\mathrm{G},{\mathrm{E}}^{\prime }}\right)\) 的对偶。推断如果 \(\mathrm{E}\) 是投射的，则 \({\mathrm{H}}^{0}\left( {\mathrm{G},\mathrm{E}}\right)\) 与 \({\mathrm{H}}^{0}\left( {\mathrm{G},{\mathrm{E}}^{\prime }}\right)\) 维数相同。

14.5. 设 \(\mathrm{E}\) 与 \(\mathrm{F}\) 为两个 \(k\left\lbrack  \mathrm{G}\right\rbrack\) -模，且 \(\mathrm{E}\) 为投射模。证明

\[
\dim {\operatorname{Hom}}^{\mathbf{G}}\left( {\mathbf{E},\mathbf{F}}\right)  = \dim {\operatorname{Hom}}^{\mathbf{G}}\left( {\mathbf{F},\mathbf{E}}\right) .
\]

[将习题 14.4(b) 应用于投射的 \(k\left\lbrack  \mathrm{\;G}\right\rbrack\) -模 \(\operatorname{Hom}\left( {\mathrm{E},\mathrm{F}}\right)\) ，并注意其对偶同构于 \(\operatorname{Hom}\left( {\mathrm{F},\mathrm{E}}\right)\) 。]

14.6. 设 \(\mathrm{S}\) 为一个单 \(k\left\lbrack  \mathrm{G}\right\rbrack\) -模， \({\mathrm{P}}_{\mathrm{S}}\) 为其投射包。证明 \({\mathrm{P}}_{\mathrm{S}}\) 包含一个与 S 同构的子模。[将习题 14.5 与 \(E = {P}_{S},F = S\) 联合使用。] 推断 \({P}_{S}\) 同构于 S 的注入包，参见习题 14.1。特别地，若 S 非投射，则 S 至少在 \({\mathbf{P}}_{\mathbf{S}}\) 的一个组成列中出现两次。

14.7. 设 \(E\) 为半单的 \(k\left\lbrack  G\right\rbrack\) -模， \({P}_{E}\) 为其投射包。证明 \(E\) 的对偶的投射包同构于 \({P}_{E}\) 的对偶 [化简到单模情形并使用习题 14.6]。

\section*{14.6 标量扩张}

若 \({\mathrm{K}}^{\prime }\) 是 \(\mathrm{K}\) 的扩张，每个 \(\mathrm{K}\left\lbrack  \mathrm{G}\right\rbrack\) -模 \(\mathrm{E}\) 通过标量扩张定义出一个 \({\mathrm{K}}^{\prime }\left\lbrack  \mathrm{G}\right\rbrack\) -模 \({\mathrm{K}}^{\prime }{ \otimes  }_{\mathrm{K}}\mathrm{E}\) 。由此我们得到一个同态

\[
{\mathrm{R}}_{\mathrm{K}}\left( \mathrm{G}\right)  \rightarrow  {\mathrm{R}}_{{\mathrm{K}}^{\prime }}\left( \mathrm{G}\right)
\]

这个同态是单射。这一点可以通过确定 \({R}_{K}\left( G\right)\) 的标准基 \(\{ \left\lbrack  E\right\rbrack  \} \left( {E \in  {S}_{K}}\right)\) 的像来看出:若 \({D}_{E}\) 是 \(E\) 的自同态(斜)域，则张量积 \({K}^{\prime } \otimes  {D}_{E}\) 分解为矩阵代数 \({\mathrm{M}}_{{s}_{i}}\left( {\mathrm{D}}_{i}\right)\) 的乘积，其中 \({\mathrm{D}}_{i}\) 是域。每个 \({\mathrm{D}}_{i}\) 对应一个单 \({\mathrm{K}}^{\prime }\left\lbrack  \mathrm{G}\right\rbrack\) -模 \({\mathrm{E}}_{i}^{\prime }\) ，并且 \(\left\lbrack  \mathrm{E}\right\rbrack\) 在 \({\mathrm{R}}_{{\mathrm{K}}^{\prime }}\left( \mathrm{G}\right)\) 中的像等于 \(\sum {s}_{i}\left\lbrack  {\mathrm{E}}_{i}^{\prime }\right\rbrack\) 。此外，每个单 \({\mathrm{K}}^{\prime }\left\lbrack  \mathrm{G}\right\rbrack\) -模同构于唯一的 \({\mathrm{E}}_{i}^{\prime }\) 。这种对 \({\mathrm{R}}_{\mathrm{K}}\left( \mathrm{G}\right)  \rightarrow  {\mathrm{R}}_{{\mathrm{K}}^{\prime }}\left( \mathrm{G}\right)\) 的描述推广了 12.2 中的描述，特别地表明:

若所有 \({\mathbf{D}}_{\mathbf{E}}\) 为交换的，则 \({s}_{i}\) 等于 1，且同态 \({R}_{K}\left( G\right)  \rightarrow  {R}_{{K}^{\prime }}\left( G\right)\) 将第一群识别为第二群的直因子，即为分裂单射。若所有 \(E \in  {S}_{K}\) 绝对单，则 \({R}_{K}\left( G\right)  \rightarrow  {R}_{{K}^{\prime }}\left( G\right)\) 是同构。

类似的结论也适用于下列同态

\[
{\mathrm{R}}_{k}\left( \mathrm{G}\right)  \rightarrow  {\mathrm{R}}_{{k}^{\prime }}\left( \mathrm{G}\right) ,\;{\mathrm{P}}_{k}\left( \mathrm{G}\right)  \rightarrow  {\mathrm{P}}_{{k}^{\prime }}\left( \mathrm{G}\right)
\]

由标量扩张从 \(k\) 到 \({k}^{\prime }\) 定义。情况更为简单:单一 \(k\left\lbrack  \mathrm{G}\right\rbrack\) -模的端同态域总是对 \(k\) 交换且可分的。(当 \(k\) 有限时这是显然的，般情形由标量扩张推出。)因此 \({\mathrm{R}}_{k}\left( \mathrm{G}\right)  \rightarrow  {\mathrm{R}}_{{k}^{\prime }}\left( \mathrm{G}\right)\) 为分裂单射。对 \({\mathrm{P}}_{k}\left( \mathrm{G}\right)  \rightarrow  {\mathrm{P}}_{{k}^{\prime }}\left( \mathrm{G}\right)\) 亦然:因为“标量扩张”函子将投射包保持为投射包。

现在假设 \({\mathrm{K}}^{\prime }\) 是 \(\mathrm{K}\) 的有限扩张。令 \({\mathrm{A}}^{\prime }\) 为 \({\mathrm{K}}^{\prime }\) 的整数环(即 \(\mathrm{A}\) 在 \({\mathrm{K}}^{\prime }\) 中的整闭包)， \({k}^{\prime }\) 为其剩余域。如果 \(E\) 是一个投射的 \(A\left\lbrack  G\right\rbrack\) -模，则 \({E}^{\prime } = {A}^{\prime }{ \otimes  }_{A}E\) 是一个投射的 \({A}^{\prime }\left\lbrack  G\right\rbrack\) -模；此外， \({\mathrm{E}}^{\prime }\) 的约化 \({k}^{\prime }{ \otimes  }_{{\mathrm{A}}^{\prime }}{\mathrm{E}}^{\prime }\) 同构于

\[
{k}^{\prime }{ \otimes  }_{\mathrm{A}}\mathrm{E} = {k}^{\prime }{ \otimes  }_{k}\left( {k{ \otimes  }_{\mathrm{A}}\mathrm{E}}\right) .
\]

该图

\[
{\mathrm{P}}_{\mathrm{A}}\left( \mathrm{G}\right)  \rightarrow  {\mathrm{P}}_{{\mathrm{A}}^{\prime }}\left( \mathrm{G}\right)
\]

\[
\downarrow
\]

\[
{\mathrm{P}}_{k}\left( \mathrm{G}\right)  \rightarrow  {\mathrm{P}}_{{k}^{\prime }}\left( \mathrm{G}\right)
\]

因此是可交换的。由于两条纵箭头都是同构，由上述可得同态 \({P}_{A}\left( G\right)  \rightarrow  {P}_{{A}^{\prime }}\left( G\right)\) 是一个分裂单射。

注。注入映射 \({\mathrm{R}}_{\mathrm{K}}\left( \mathrm{G}\right)  \rightarrow  {\mathrm{R}}_{{\mathrm{K}}^{\prime }}\left( \mathrm{G}\right) ,{\mathrm{R}}_{k}\left( \mathrm{G}\right)  \rightarrow  {\mathrm{R}}_{{k}^{\prime }}\left( \mathrm{G}\right)\) 等与前一节的双线性形式相容。此外，它们与下一章中定义的同态 \(c,d,e\) 交换。

第15章 cde 三角形

我们将定义同态 \(c,d\) 和 \(e\) ，它们构成一个可交换三角形:

\begin{center}
\includegraphics[max width=0.3\textwidth]{images/bo_d6kn08k601uc73e1e9r0_129_591_1158_392_210_0.jpg}
\end{center}
\hspace*{3em} 

15.1 \(c : {\mathrm{P}}_{k}\left( \mathrm{G}\right)  \rightarrow  {\mathrm{R}}_{k}\left( \mathrm{G}\right)\) 的定义

将每个投射的 \(k\left\lbrack  \mathrm{G}\right\rbrack\) -模 \(\mathrm{P}\) 关联到其在群 \({\mathrm{R}}_{k}\left( \mathrm{G}\right)\) 中的类。该类是关于 \(\mathrm{P}\) 的加法函数，因此我们得到一个同态

\[
c : {\mathrm{P}}_{k}\left( \mathrm{G}\right)  \rightarrow  {\mathrm{R}}_{k}\left( \mathrm{G}\right)
\]

称为 Cartan 同态。如果我们用 \({\mathrm{P}}_{k}\left( \mathrm{G}\right)\) 和 \({\mathrm{R}}_{k}\left( \mathrm{G}\right)\) 的标准基 \(\left\lbrack  {\mathrm{P}}_{\mathrm{S}}\right\rbrack\) 和 \(\left\lbrack  \mathrm{S}\right\rbrack  \left( {\mathrm{S} \in  {\mathrm{S}}_{k}}\right)\) 来表示 \(c\) ，就得到一个方阵 C，类型为 \({\mathrm{S}}_{k} \times  {\mathrm{S}}_{k}\) ，称为相对于 \(k\) 的 Cartan 矩阵 \(\mathrm{G}\) 。 \(\mathrm{C}\) 的 (S, T) 系数 \({\mathrm{C}}_{\mathrm{{ST}}}\) 是简单模 S 在投射包覆 \({\mathrm{P}}_{\mathrm{T}}\) 的组成序列中出现的次数，其中 \(\mathrm{T}\) :我们有

\[
\left\lbrack  {\mathrm{P}}_{\mathrm{T}}\right\rbrack   = \mathop{\sum }\limits_{{\mathrm{S} \in  {\mathrm{S}}_{k}}}{\mathrm{C}}_{\mathrm{{ST}}}\left\lbrack  \mathrm{S}\right\rbrack  \text{ in }{\mathrm{R}}_{k}\left( \mathrm{G}\right) .
\]

练习

15.1. 证明当 \(x \in  {\mathrm{R}}_{k}\left( \mathrm{G}\right) ,y \in  {\mathrm{P}}_{k}\left( \mathrm{G}\right)\) 时 \(c\left( {x \cdot  y}\right)  = x \cdot  c\left( y\right)\) 。

\section*{15.2 \(d : {\mathrm{R}}_{\mathrm{K}}\left( \mathrm{G}\right)  \rightarrow  {\mathrm{R}}_{k}\left( \mathrm{G}\right)\) 的定义}

设 \(\mathrm{E}\) 为一个 \(\mathrm{K}\left\lbrack  \mathrm{G}\right\rbrack\) -模。选取 \(\mathrm{E}\) 中的一个格子 \({\mathrm{E}}_{1}\) (即生成 \(\mathrm{E}\) 作为 K-模的一个有限生成的 A-子模)；将 \({\mathrm{E}}_{1}\) 替换为在 \(G\) 各元作用下像的和后，我们可以假定 \({E}_{1}\) 在 \(G\) 作用下稳定。 \({E}_{1}\) 的约化 \({\bar{E}}_{1} = {E}_{1}/m{E}_{1}\) 因此是一个 \(k\left\lbrack  G\right\rbrack\) -模。

定理 32. \({\bar{E}}_{1}\) 在 \({R}_{k}\left( G\right)\) 中的像不依赖于所选的稳定格子 \({E}_{1}\) 。

(由稳定格子 \({\mathrm{E}}_{1}\) 和 \({\mathrm{E}}_{2}\) 的约化得到的两个 \(k\left\lbrack  \mathrm{G}\right\rbrack\) -模 \({\overline{\mathrm{E}}}_{1}\) 与 \({\overline{\mathrm{E}}}_{2}\) 未必同构，参见例 15.1。上述定理的意思是它们具有相同的组成因子。)

设 \({\mathrm{E}}_{2}\) 是在 \(\mathrm{G}\) 作用下稳定的 \(\mathrm{E}\) 的一个格子。我们必须证明 \(\left\lbrack  {\overline{\mathrm{E}}}_{1}\right\rbrack   = \left\lbrack  {\overline{\mathrm{E}}}_{2}\right\rbrack\) 在 \({\mathrm{R}}_{k}\left( \mathrm{G}\right)\) 中。先从一个特例开始:

我们有 \({\mathrm{{mE}}}_{1} \subset  {\mathrm{E}}_{2} \subset  {\mathrm{E}}_{1}\) 。设 \(\mathrm{T}\) 为 \(k\left\lbrack  \mathrm{G}\right\rbrack\) -模 \({\mathrm{E}}_{1}/{\mathrm{E}}_{2}\) 。则存在一条正合列

\[
0 \rightarrow  \mathrm{T} \rightarrow  {\overline{\mathrm{E}}}_{2} \rightarrow  {\overline{\mathrm{E}}}_{1} \rightarrow  \mathrm{T} \rightarrow  0,
\]

其中同态 \(\mathrm{T} \rightarrow  {\overline{\mathrm{E}}}_{2}\) 来自于用理想 \(\mathfrak{m}\) 的一个生成元 \(\pi\) 乘法。通过取 \({\mathrm{R}}_{k}\left( \mathrm{G}\right)\) ，我们得到

\[
\left\lbrack  \mathrm{T}\right\rbrack   - \left\lbrack  {\overline{\mathrm{E}}}_{2}\right\rbrack   + \left\lbrack  {\overline{\mathrm{E}}}_{1}\right\rbrack   - \left\lbrack  \mathrm{T}\right\rbrack   = 0.
\]

因此 \(\left\lbrack  {\overline{\mathrm{E}}}_{1}\right\rbrack   = \left\lbrack  {\overline{\mathrm{E}}}_{2}\right\rbrack\) ，这在此情形下证明了定理。

一般情况。将 \({\mathrm{E}}_{2}\) 乘以一个标量(这不影响 \({\overline{\mathrm{E}}}_{2}\) )后，我们可以假定 \({\mathrm{E}}_{2}\) 包含于 \({\mathrm{E}}_{1}\) 。因此存在一个整数 \(n \geq  0\) 使得

\[
{m}^{n}{E}_{1} \subset  {E}_{2} \subset  {E}_{1}
\]

并且我们对 \(n\) 归纳。令 \({\mathrm{E}}_{3} = {\mathrm{m}}^{n - 1}{\mathrm{E}}_{1} + {\mathrm{E}}_{2}\) 。则

\[
{\mathfrak{m}}^{n - 1}{E}_{1} \subset  {E}_{3} \subset  {E}_{1}\text{ and }\mathfrak{m}{E}_{3} \subset  {E}_{2} \subset  {E}_{3}.
\]

由 (a) 和归纳我们得到

\[
\left\lbrack  {\overline{\mathrm{E}}}_{1}\right\rbrack   = \left\lbrack  {\overline{\mathrm{E}}}_{3}\right\rbrack   = \left\lbrack  {\overline{\mathrm{E}}}_{2}\right\rbrack
\]

这证明了该定理。

现在显然映射 \(E \mapsto  \left\lbrack  {\bar{E}}_{1}\right\rbrack\) 可以扩展为一个环同态

\[
d : {\mathrm{R}}_{\mathrm{K}}\left( \mathrm{G}\right)  \rightarrow  {\mathrm{R}}_{k}\left( \mathrm{G}\right)
\]

称为分解同态。它把 \({R}_{K}^{ + }\left( G\right)\) 送入 \({R}_{k}^{ + }\left( G\right)\) 。相应的矩阵 \(D\) (相对于 \({R}_{K}\left( G\right)\) 和 \({\mathrm{R}}_{k}\left( \mathrm{G}\right)\) 的典范基)称为分解矩阵。它是一个类型为 \({\mathrm{S}}_{k} \times  {\mathrm{S}}_{\mathrm{K}}\) 的矩阵，其系数为非负整数。对于 \(F \in  {S}_{k}\) 和 \(E \in  {S}_{K}\) ， \(D\) 的相应系数 \({D}_{FE}\) 是在模 \(\mathfrak{m}\) 意义下，某个 \(E\) 的稳定格 \({E}_{1}\) 的约化中 \(F\) 出现的次数:因此

\[
\left\lbrack  {\overline{\mathrm{E}}}_{1}\right\rbrack   = \mathop{\sum }\limits_{\mathrm{F}}\mathrm{D}\mathrm{{FE}}\left\lbrack  \mathrm{F}\right\rbrack  \;\text{ in }{\mathrm{R}}_{k}\left( \mathrm{G}\right) .
\]

\section*{Remarks}

(1)假定 \(\mathrm{K}\) 完备在定理 32 的证明及同态 \(d\) 的定义中并不发挥作用。

(2)关于代数群有类似结果，参见 I.H.E.S. Publ. Sci. 第 34 期，1968 年，第 37–52 页。

\section*{Exercises}

15.2. 取 \(p = 2\) 与 \(G\) 为 2 阶。设 \(E = K\left\lbrack  G\right\rbrack\) 。证明 \(E\) 存在其约化为半单的(同构于 \(k \oplus  k\) )稳定格，也存在其约化非半单的(同构于 \(k\left\lbrack  \mathrm{G}\right\rbrack\) )稳定格。

15.3. 设 \(\mathrm{E}\) 为非零的 \(\mathrm{K}\left\lbrack  \mathrm{G}\right\rbrack\) -模，且 \({\mathrm{E}}_{1}\) 为在 \(\mathrm{G}\) 下不变的格。证明下列各项等价:

(i) \({\mathrm{E}}_{1}\) 的约化 \({\overline{\mathrm{E}}}_{1}\) 是一个简单的 \(k\left\lbrack  \mathrm{G}\right\rbrack\) -模。

(ii) 在 \(\mathrm{G}\) 下不变的每个格均为形如 \(a\mathrm{E}\) 且 \(a \in  {\mathrm{K}}^{ * }\) 的格。证明这些条件蕴含 \(\mathrm{E}\) 为一个简单的 \(\mathrm{K}\left\lbrack  \mathrm{G}\right\rbrack\) -模。

15.4.(J. Thompson 之后。)设 \(E\) 为一个 \(Z\) -自由的 \(Z\left\lbrack  G\right\rbrack\) -模，秩为 \(n \geq  2\) 。假设对每一素数 \(p\) ， \(\mathrm{E}\) 的约化 \(\mathrm{E}/p\mathrm{E}\) 是一个简单的 \(\left( {\mathbf{Z}/p\mathbf{Z}}\right)\) [G]-模。

(a)证明存在一个在 \(E\) 上取值于 \(\mathbf{Z}\) 的双线性型 \(B\left( {x,y}\right)\) ，使得对所有 \(x \neq  0\) 都有 \(\mathrm{B}\left( {x,x}\right)  > 0\) 。

(b) 令 B 如 (a) 所述并把它按线性扩张到 Q 向量空间 \(\mathbf{Q} \otimes  \mathbf{E}\) 。证明集合 \({\mathbf{E}}^{\prime }\) ——由满足对所有 \(y \in  \mathbf{E}\) 都有 \(\mathbf{B}\left( {x,y}\right)  \in  \mathbf{Z}\) 的 \(x \in  \mathbf{Q} \otimes  \mathbf{E}\) 组成——具有形式 \({\mathbf{E}}^{\prime } = a\mathbf{E}\) ，其中 \(a \in  {\mathbf{Q}}^{ * }\) (与例 15.3 的论证相同)。由此推出 B 可被选为对 \(\mathbf{Z}\) 非退化，即 使得 \({\mathrm{E}}^{\prime } = \mathrm{E}\) 。若 \(\left( {{e}_{i},\ldots ,{e}_{n}}\right)\) 是 \(\mathrm{E}\) 的一组基，则 \(\mathrm{B}\left( {{e}_{i},{e}_{j}}\right)\) 的矩阵的行列式等于 1 。

(c) 假定已按 (b) 选择了 \(\mathrm{B}\) 。证明存在 \(x \in  \mathrm{E}\) 使得对所有 \(y \in  \mathrm{E}\) 都有 \(\mathrm{B}\left( {y,y}\right)  \equiv  \mathrm{B}\left( {x,y}\right) \left( {\;\operatorname{mod}\;{.2}}\right)\) ，并且这样的 \(x\) 在模 2E 意义下关于 G 不变。由此得出 \(x \equiv  0\) (模 2E)，即二次型 \(\mathrm{B}\left( {x,x}\right)\) 仅取偶值。

(d) 从 (c) 得出同余 \(n \equiv  0\) (模 8)。[使用事实*:任一正定整值二次型，若为偶型且判别式为 1，则其秩为 8 的倍数。]

(e) 证明柯西特群(Coxeter group)型 \({\mathrm{E}}_{8}\) 的反射表示具有上述性质(参见 Bourbaki, Gr. et Alg. de Lie, Ch. VI, §4, no. 10)。

\section*{15.3 \(e : {\mathrm{P}}_{k}\left( \mathrm{G}\right)  \rightarrow  {\mathrm{R}}_{\mathrm{K}}\left( \mathrm{G}\right)\) 的定义}

函子“与 \(\mathrm{K}\) 张量”给出从 \({P}_{A}\left( G\right)\) 到 \({R}_{K}\left( G\right)\) 的同态。将其与规范同构 \({P}_{A}\left( G\right)  \rightarrow  {P}_{k}\left( G\right)\) 的逆结合，参见 14.4，我们得到一个同态

\[
e : {\mathrm{P}}_{k}\left( \mathrm{G}\right)  \rightarrow  {\mathrm{R}}_{\mathrm{K}}\left( \mathrm{G}\right) .
\]

其矩阵记为 \(\mathrm{E}\) ；它是类型 \({\mathrm{S}}_{\mathrm{K}} \times  {\mathrm{S}}_{k}\) 的。

EXERCISE

15.5. 若 \(x \in  {\mathrm{R}}_{\mathrm{K}}\left( \mathrm{G}\right) ,y \in  {\mathrm{P}}_{k}\left( \mathrm{G}\right)\) 则有 \(e\left( {d\left( x\right)  \cdot  y}\right)  = x \cdot  e\left( y\right)\) 。

\section*{15.4 cde 三角形的基本性质}

(a) 它是交换的，即 \(c = d \circ  e\) ，或等价地 \(\mathbf{C} = \mathbf{D} \cdot  \mathbf{E}\) 。此点显然。

(b) 同态 \(d\) 与 \(e\) 关于 14.5 中的双线性型互为伴随:

\[
\langle x,d\left( y\right) {\rangle }_{k} = \langle e\left( x\right) ,y{\rangle }_{\mathrm{K}}\;\text{ if }x \in  {\mathrm{P}}_{k}\left( \mathrm{G}\right) ,y \in  {\mathrm{R}}_{\mathrm{K}}\left( \mathrm{G}\right) .
\]

确实，我们可以假定 \(x = \left\lbrack  \bar{X}\right\rbrack\) ，其中 \(X\) 是一个投射的 \(A\left\lbrack  G\right\rbrack\) -模，并且 \(y = \left\lbrack  {\mathrm{K}{ \otimes  }_{\mathrm{A}}\mathrm{Y}}\right\rbrack\) ，其中 \(\mathrm{Y}\) 是一个 \(\mathrm{A}\left\lbrack  \mathrm{G}\right\rbrack\) -模且为 \(\mathrm{A}\) -自由。于是 A-模 \({\operatorname{Hom}}^{\mathrm{G}}\left( {\mathrm{X},\mathrm{Y}}\right)\) 是自由的；设 \(r\) 为其秩。于是我们有自然同构:

\[
\mathrm{K} \otimes  {\mathrm{{Hom}}}^{\mathrm{G}}\left( {\mathrm{X},\mathrm{Y}}\right)  = {\mathrm{{Hom}}}^{\mathrm{G}}\left( {\mathrm{K} \otimes  \mathrm{X},\mathrm{K} \otimes  \mathrm{Y}}\right)
\]

以及

\[
k \otimes  {\operatorname{Hom}}^{\mathrm{G}}\left( {\mathrm{X},\mathrm{Y}}\right)  = {\operatorname{Hom}}^{\mathrm{G}}\left( {k \otimes  \mathrm{X},k \otimes  \mathrm{Y}}\right) .
\]

这表明 \(\langle e\left( x\right) ,y\rangle  = r = \langle x,d\left( y\right) \rangle\) 。

(c) 假定 \(\mathrm{K}\) 足够大。根据 14.5， \({P}_{k}\left( G\right)\) (或 \({R}_{K}\left( G\right)\) )以及 \({R}_{k}\left( G\right)\) (或 \({R}_{K}\left( G\right)\) )的规范基关于双线性形式 \(\langle a,b{\rangle }_{k}\) (或形式 \(\langle a,b{\rangle }_{\mathrm{K}}\) )互为对偶。这意味着 \(e\) 可被识别为 \(d\) 的转置；特别地我们有 \(\mathrm{E} = {}^{t}\mathrm{D}\) 。由于 \(\mathrm{C} = \mathrm{D} \cdot  \mathrm{E} = \mathrm{D} \cdot  {}^{t}\mathrm{D}\) ，可见 \(\mathrm{C}\) 是对称矩阵。

习题

15.6. 设有 \(S,T \in  {S}_{k}\) 并设 \({P}_{S},{P}_{T}\) 为它们的投射包。我们定义

\[
{d}_{\mathrm{S}} = \dim {\operatorname{End}}^{\mathrm{G}}\left( \mathrm{S}\right) ,\;{d}_{\mathrm{T}} = \dim {\operatorname{End}}^{\mathrm{G}}\left( \mathrm{T}\right) ,
\]

并设 \({C}_{ST}\) (或 \({C}_{TS}\) )为在 \({P}_{T}\) (或 \({P}_{S}\) )的组成列中 \(S\) (或 \(T\) )的重数，参见 15.1。

(a) 证明 \({\mathrm{C}}_{\mathrm{{ST}}}{d}_{\mathrm{S}} = \dim {\operatorname{Hom}}^{\mathrm{G}}\left( {{\mathrm{P}}_{\mathrm{S}},{\mathrm{P}}_{\mathrm{T}}}\right)\) 。

(b) 证明 \({\mathrm{C}}_{\mathrm{{ST}}}{d}_{\mathrm{S}} = {\mathrm{C}}_{\mathrm{{TS}}}{d}_{\mathrm{T}}\) [应用习题 14.5]。当 \(\mathrm{K}\) 足够大时， \({d}_{\mathrm{S}}\) 等于 1，并且我们再次得到矩阵 \(\mathbf{C} = \left( {\mathbf{C}}_{\mathbf{{ST}}}\right)\) 为对称的结论。

15.7. 保持习题 15.6 的记号。证明要么 \(S\) 是投射的，要么 \({P}_{S} \cong  S\) 与 \({\mathrm{C}}_{\mathrm{{SS}}} = 1\) ，或 \({\mathrm{C}}_{\mathrm{{SS}}} \geq  2\) [参用习题 14.6]。

15.8. 若 \(x \in  {\mathrm{P}}_{k}\left( \mathrm{G}\right)\) ，则有 \(\langle x,c\left( x\right) {\rangle }_{k} = \langle e\left( x\right) ,e\left( x\right) {\rangle }_{\mathrm{K}}\) 。由此得出:若 \(\mathrm{K}\) 足够大，由 Cartan 矩阵 \(C\) 定义的二次型是正定的。

\section*{15.5 例: \({p}^{\prime }\) -群}

命题 43。假设 \(\mathrm{G}\) 的阶与 p 互素。则:

(i)每个 \(k\left\lbrack  \mathrm{G}\right\rbrack\) -模(或每个 A-自由的 \(\mathrm{A}\left\lbrack  \mathrm{G}\right\rbrack\) -模)均为投射模。

(ii)对 \(\mathfrak{m}\) 取模的约化运算在 \({\mathrm{S}}_{\mathrm{K}}\) 与 \({\mathrm{S}}_{k}\) 之间给出双射。

(iii)如果按(ii)将 \({\mathrm{S}}_{\mathrm{K}}\) 与 \({\mathrm{S}}_{k}\) 识别，则 \(\mathrm{C},\mathrm{D},\mathrm{E}\) 的矩阵全是单位矩阵。

(简言之:群 G 的表示论在 \(k\) 上与在 K 上“相同”。)

设 E 为一 A[G]-模且在 A 上自由。可将 E 写成自由 A[G]-模 L 的商 L/R。由于 E 在 A 上自由，存在 A 线性投影 \(\pi\) 将 \(\mathrm{L}\) 映到 \(\mathrm{R}\) ；又因 \(\mathrm{G}\) 的阶 \(g\) 在 \(\mathrm{A}\) 中可逆，可将 \(\pi\) 替换为其共轭的平均值 \(\left( {1/g}\right) \mathop{\sum }\limits_{{s \in  \mathrm{G}}}{s\pi }{s}^{-1}\) ，由此得到的投影为 A[G]-线性。这说明 E 为 A[G]-投射模。同样的论证适用于 \(k\left\lbrack  \mathrm{G}\right\rbrack\) -模。这证明了(i)，并且卡坦矩阵为单位矩阵。

如果 \(E \in  {S}_{k}\) ，相对于 \(A\left\lbrack  G\right\rbrack\) 的 \(E\) 的投射包络 \({E}_{1}\) 是一个投射的 \(\mathrm{A}\left\lbrack  \mathrm{G}\right\rbrack\) -模，其约化 \({\overline{\mathrm{E}}}_{1} = {\mathrm{E}}_{1}/m{\mathrm{E}}_{1}\) 为 \(\mathrm{E}\) 。若置 \(\mathrm{F} = \mathrm{K} \otimes  {\mathrm{E}}_{1}\) ，则 \(d\left( \left\lbrack  \mathrm{\;F}\right\rbrack  \right)  = \left\lbrack  \mathrm{E}\right\rbrack\) 。因 \(\mathrm{E}\) 为单模，故 \(\mathrm{F}\) 亦为单模，因而与 \({S}_{K}\) 中某一元素同构。由此得到从 \({S}_{k}\) 到 \({S}_{K}\) 的映射 \(E \mapsto  F\) ，且显然该映射是 \(d\) 的逆映射。这证明了(ii)和(iii)。

备注。 \(\mathrm{D}\) 为单位矩阵表明 \(d\) 将 \({\mathrm{R}}_{\mathrm{K}}^{ + }\left( \mathrm{G}\right)\) 映到 \({\mathrm{R}}_{k}^{ + }\left( \mathrm{G}\right)\) ；换言之， \(\mathrm{G}\) 在 \(\mathrm{K}\) 上的任一线性表示都可提升为在 A 上的表示，该结果亦可直接验证(参见下文练习 15.9)。

练习

15.9. 设 \(g\) 与 \(p\) 互素。令 \(\mathrm{E}\) 为一 A-自由模。

(a)令 \(n \geq  1\) 为整数，且设

\[
{\rho }_{n} : \mathrm{G} \rightarrow  \mathrm{{GL}}\left( {\mathrm{E}/{\mathrm{m}}^{n}\mathrm{E}}\right)
\]

为把 \(G\) 同构到 \(E/{\mathfrak{m}}^{n}E\) 的自同构群的同态。证明 \({\rho }_{n}\) 可以被提升为

\[
{\rho }_{n + 1} : \mathrm{G} \rightarrow  \mathrm{{GL}}\left( {\mathrm{E}/{\mathfrak{m}}^{n + 1}\mathrm{E}}\right)
\]

并且这种提升是唯一的，至多相差一个与 1 式模 \(.{\mathfrak{m}}^{n}\) 同余的 \(\mathrm{E}/{\mathfrak{m}}^{n + 1}\mathrm{E}\) 自同构的共轭。〔利用 \(\mathrm{G}\) 的 1、2 维上同调群以 \(\operatorname{End}\left( {\mathrm{E}/\mathfrak{m}\mathrm{E}}\right)\) 为系数为零的事实。〕

(b) 从 (a) 推出每一个线性表示

\[
{\rho }_{1} : \mathrm{G} \rightarrow  \mathrm{{GL}}\left( {\mathrm{E}/\mathrm{{mE}}}\right)
\]

G 在 \(k\) 上的表示可以以本质唯一的方式提升为 G 在 A 上的表示。

\section*{15.6 例: \(p\) -群}

假设 \(\mathrm{G}\) 是一个 \(p\) -群，阶为 \({p}^{n}\) 。我们已见(8.3，命题26的推论)在特征为 \(p\) 时 \(\mathrm{G}\) 唯一的不可约表示是平凡表示。由此可知艾蒂安环 \(k\left\lbrack  \mathrm{G}\right\rbrack\) 是一个剩余域为 \(k\) 的局部环。简单 \(k\left\lbrack  \mathrm{G}\right\rbrack\) -模 \(k\) 的投射包是 \(k\left\lbrack  \mathrm{G}\right\rbrack\) ，即 \(\mathrm{G}\) 的正则表示。群 \({\mathrm{R}}_{k}\left( \mathrm{G}\right)\) 与 \({\mathrm{P}}_{k}\left( \mathrm{G}\right)\) 均可鉴别为 \(\mathbf{Z}\) ，Cartan 同态 \(c : \mathbf{Z} \rightarrow  \mathbf{Z}\) 为乘以 \({p}^{n}\) 。同态 \(d : {\mathrm{R}}_{\mathrm{K}}\left( \mathrm{G}\right)  \rightarrow  \mathbf{Z}\) 对应于 K-秩；同态 \(e : \mathbf{Z} \rightarrow  {\mathrm{R}}_{\mathrm{K}}\left( \mathrm{G}\right)\) 将整数 \(n\) 映为 \(n\) 倍的 \(\mathrm{G}\) 的正则表示类。

\section*{15.7 例: \({p}^{\prime }\) -群与 \(p\) -群的直积}

假设 \(\mathrm{G} = \mathrm{S} \times  \mathrm{P}\) ，其间 \(\mathrm{S}\) 的阶与 \(p\) 互素，且 \(\mathrm{P}\) 是一个 \(p\) -群。我们有 \(k\left\lbrack  \mathrm{G}\right\rbrack   = k\left\lbrack  \mathrm{\;S}\right\rbrack   \otimes  k\left\lbrack  \mathrm{P}\right\rbrack\) 。此外:

(a) \({Ak}\left\lbrack  \mathrm{G}\right\rbrack\) -模 \(\mathrm{E}\) 当且仅当 \(\mathrm{P}\) 在 \(\mathrm{E}\) 上平凡作用时为半单。

充分性由每个 \(k\left\lbrack  \mathrm{\;S}\right\rbrack\) -模皆为半单(参见15.5)可得。为证必要性，可假设 E 为单模。由15.6，子空间 \({E}^{\prime }\) (由被 \(P\) 固定的元素构成)非零。因 \(P\) 在 \(G\) 中为正规子群，该子空间在 \(G\) 下稳定，遂等于 \(E\) ，即 \(\mathrm{P}\) 作平凡作用。

(b) \({Ak}\left\lbrack  \mathrm{\;F}\right\rbrack\) -模 \(\mathrm{E}\) 当且仅当同构于 \(\mathrm{F} \otimes  k\left\lbrack  \mathrm{P}\right\rbrack\) (其中 \(\mathrm{F}\) 为一个 \(k\left\lbrack  \mathrm{\;S}\right\rbrack\) -模)时为投射模。

由于 \(\mathrm{F}\) 是一个投射的 \(k\left\lbrack  \mathrm{\;S}\right\rbrack\) -模(参见 15.5)， \(\mathrm{F} \otimes  k\left\lbrack  \mathrm{P}\right\rbrack\) 是一个投射的 \(k\left\lbrack  \mathrm{G}\right\rbrack\) -模。此外，很明显 \(\mathrm{F}\) 是 \(\mathrm{F} \otimes  k\left\lbrack  \mathrm{P}\right\rbrack\) 在 \(P\) 上平凡作用时的最大商。由 (a) 这意味着 \(F \otimes  k\left\lbrack  P\right\rbrack\) 是 F 的投射包。然而，每个投射模都是其最大半单商的投射包。因此我们看到每个投射模都具有形如 \(\mathrm{F} \otimes  k\left\lbrack  \mathrm{P}\right\rbrack\) 的形式。

(c) \({An}\mathrm{\;A}\left\lbrack  \mathrm{G}\right\rbrack\) -模 \(\widetilde{\mathrm{E}}\) 当且仅当它同构于 \(\widetilde{\mathrm{F}} \otimes  \mathrm{A}\left\lbrack  \mathrm{P}\right\rbrack\) 时是投射的，其中 \(\widetilde{\mathrm{F}}\) 是一个关于 A 自由的 \(\mathrm{A}\left\lbrack  \mathrm{S}\right\rbrack\) -模。

显然形如 \(\widetilde{F} \otimes  A\left\lbrack  P\right\rbrack\) 的模是投射的。逆命题通过将 (b) 应用于 \(E = \widetilde{E}/\mathfrak{m}\widetilde{E}\) 来证明:若 \(\widetilde{E}\) 是投射的，则有 \(E \; \simeq  \mathrm{F} \otimes  k\left\lbrack  \mathrm{P}\right\rbrack\) ，且我们知道 \(\mathrm{F}\) 可以提升为一个对 \(\mathrm{A}\left\lbrack  \mathrm{S}\right\rbrack\) 的模 \(\widetilde{\mathrm{F}}\) ，该模在 \(A\) 上是自由的(甚至是 \(A\left\lbrack  S\right\rbrack\) -投射的，参见上文)。模 \(\widetilde{F} \otimes  A\left\lbrack  P\right\rbrack\) 是 \(\mathrm{F} \otimes  k\left\lbrack  \mathrm{P}\right\rbrack\) 的投射包，因此同构于 \(\widetilde{\mathrm{E}}\) 。

性质 (a) 与 (b) 特别表明 \(\mathbf{G}\) 的卡坦矩阵是标量矩阵 \({p}^{n}\) ，其中 \({p}^{n} = \operatorname{Card}\left( \mathrm{P}\right)\) 。

\section*{第16章 定理}

16.1 cde 三角的性质

主要结果如下*:

定理 33. 同态 \(d : {\mathrm{R}}_{\mathrm{K}}\left( \mathrm{G}\right)  \rightarrow  {\mathrm{R}}_{k}\left( \mathrm{G}\right)\) 是满射。

证明将在 17.3 给出。

\section*{备注}

(1) 这特别适用于 \(k = \mathbf{Z}/p\mathbf{Z}\) ，取 \(\mathrm{K}\) 为 \(p\) -进域 \({\mathbf{Q}}_{p}\) ；此时环 \(\mathbf{A}\) 就是 \(p\) -进整数的环 \({\mathbf{Z}}_{p}\) 。

(2) 粗略地说，该定理断言:若接受“虚表示”，即格罗滕迪克群中的元素 \({R}_{K}\left( G\right)\) ，则群 G 在 \(k\) 上的任一线性表示都可升至特征 0。这对许多应用极其有用。

定理 34. 同态 \(e : {\mathrm{P}}_{k}\left( \mathrm{G}\right)  \rightarrow  {\mathrm{R}}_{\mathrm{K}}\left( \mathrm{G}\right)\) 是分裂单射。

当 \(\mathrm{K}\) 足够大时， \(e\) 是 \(d\) 的转置(参见 15.4)，而 \(d\) 的满射性意味着 \(e\) 为分裂单射。一般情况下，令 \({\mathrm{K}}^{\prime }\) 为 \(\mathrm{K}\) 的一个有限且足够大的扩张， \({k}^{\prime }\) 为其剩余域。考虑下列图:

\begin{center}
\includegraphics[max width=0.4\textwidth]{images/bo_d6kn08k601uc73e1e9r0_136_536_1841_439_168_0.jpg}
\end{center}
\hspace*{3em} 

\HRule

— 在本书第一版法文版中，定理 33 仅在域 K 足够大时给出。Claude Chevalley 与 Andreas Dress 独立地指出该命题对一般情形也成立。

\HRule

正如我们刚才所见， \({e}^{\prime }\) 是分裂单射。根据 14.6，对 \({\mathrm{P}}_{k}\left( \mathrm{G}\right)  \rightarrow  {\mathrm{P}}_{{k}^{\prime }}\left( \mathrm{G}\right)\) 也成立。它们的复合同样是分裂单射，因此 \(e\) 也具有同样性质。

同时我们已证明:

推论 1. 对于每个有限扩张 \({\mathrm{K}}^{\prime }\) 的 \(\mathrm{K}\) ，同态

\[
{\mathrm{P}}_{k}\left( \mathrm{G}\right) \overset{e}{ \rightarrow  }{\mathrm{R}}_{\mathrm{K}}\left( \mathrm{G}\right)  \rightarrow  {\mathrm{R}}_{{\mathrm{K}}^{\prime }}\left( \mathrm{G}\right)
\]

是分裂单射。

\(e\) 的单射性等价于:

推论 2. 设 \(\mathrm{P}\) 与 \({\mathrm{P}}^{\prime }\) 为投射的 \(\mathrm{A}\left\lbrack  \mathrm{G}\right\rbrack\) -模。若 \(\mathrm{K}\left\lbrack  \mathrm{G}\right\rbrack\) -模 \(\mathrm{K} \otimes  \mathrm{P}\) 与 \(\mathrm{K} \otimes  {\mathrm{P}}^{\prime }\) 同构，则 \(\mathrm{P}\) 与 \({\mathrm{P}}^{\prime }\) 为 \(\mathrm{A}\left\lbrack  \mathrm{G}\right\rbrack\) -同构。

(事实上我们知道，在 \({\mathrm{R}}_{\mathrm{A}}\left( \mathrm{G}\right)  \simeq  {\mathrm{R}}_{k}\left( \mathrm{G}\right)\) 中等式 \(\left\lbrack  \mathrm{P}\right\rbrack   = \left\lbrack  {\mathrm{P}}^{\prime }\right\rbrack\) 等价于 \(\mathrm{P} \simeq  {\mathrm{P}}^{\prime }\) 。)

定理 35. 令 \({p}^{n}\) 为整除 |G| 的最大 \(p\) 幂。则凡被 \({p}^{n}\) 整除的 \({\mathrm{R}}_{k}\left( \mathrm{G}\right)\) 中元素，均属于 Cartan 映射 \(c : {\mathrm{P}}_{k}\left( \mathrm{G}\right)  \rightarrow  {\mathrm{R}}_{k}\left( \mathrm{G}\right)\) 的像。

证明将在 17.4 给出。

推论 1. 映射 \(c : {\mathrm{P}}_{k}\left( \mathrm{G}\right)  \rightarrow  {\mathrm{R}}_{k}\left( \mathrm{G}\right)\) 是单射，且其余商是有限的 p 群。

第二项断言由定理 35 立即得出；第一项随后成立，因为 \({P}_{k}\left( G\right)\) 与 \({R}_{k}\left( G\right)\) 是秩相同的自由 \(Z\) -模，即秩为 \(\operatorname{Card}\left( {\mathrm{S}}_{k}\right)\) 。

推论 2. 若两个投射的 \(k\left( \mathrm{G}\right)\) -模具有相同的组成因子(计重)，则它们同构。

这就是对 \(c\) 单射性的重述。

推论 3. 设 \(\mathrm{K}\) 足够大。则卡坦矩阵 \(\mathrm{C}\) 为对称，且相应的二次型正定。 \(\mathbf{C}\) 的行列式是 \(p\) 的幂。

所讨论的二次型为

\[
x \mapsto  \langle x,c\left( x\right) {\rangle }_{k} = \langle x,d\left( {e\left( x\right) }\right) {\rangle }_{k} = \langle e\left( x\right) ,e\left( x\right) {\rangle }_{\mathrm{K}},\;x \in  {\mathrm{P}}_{k}\left( \mathrm{G}\right) .
\]

由于形式 \(\langle a,b{\rangle }_{\mathrm{K}}\) 显然正定，且 \(e\) 是单射(定理 34)，可见上述形式亦正定。于是 \(\mathbf{C}\) 的行列式为 \(> 0\) 。这意味着 \(\det \left( \mathbf{C}\right)\) 是 \(p\) 的幂，因为 \(c\) 的余核是一个 \(p\) -群。

备注

(1) 上述论证表明 \(c\) 的单射性可由 \(e\) 的单射性推出。

(2) 定理 35 等价于存在一个同态 \({c}^{\prime } : {\mathrm{R}}_{k}\left( \mathrm{G}\right)  \rightarrow  {\mathrm{P}}_{k}\left( \mathrm{G}\right)\) 使得 \(c \circ  {c}^{\prime } = {p}^{n}\) (这暗含 \(\left. {{c}^{\prime } \circ  c = {p}^{n}}\right)\) )。

(3) 定理 35 中的指数 \(n\) 是最优的，参见例 16.3。

练习

16.1. 证明当 \(\mathrm{K}\) 不完备时，若 \(\mathrm{K}\) 足够大，定理 33 仍然成立。(设 \(\widehat{\mathrm{K}}\) 为 \(\mathrm{K}\) 的完备化，注意同态 \({R}_{K}\left( G\right)  \rightarrow  {R}_{\widehat{K}}\left( G\right)\) 是同构，并将定理 33 应用于 K。)

16.2. 证明若 \(\mathrm{G}\) 为 4 阶循环群，则 \(d : {\mathrm{R}}_{\mathbf{Q}}\left( \mathrm{G}\right)  \rightarrow  {\mathrm{R}}_{\mathbf{Z}/5\mathbf{Z}}\left( \mathrm{G}\right)\) 非满射。

16.3. 设 \(\mathrm{H}\) 为 \(p\) 在 \(\mathrm{G}\) 中的 Sylow-子群。证明若 \(\mathrm{E}\) 是一个投射的 \(k\left\lbrack  \mathrm{G}\right\rbrack\) -模，则 \(\mathrm{E}\) 为自由的 \(k\left\lbrack  \mathrm{H}\right\rbrack\) -模(参见 15.6)，因此 dim \(\mathrm{E}\) 可被 \({p}^{n}\) 整除。由此得出映射 \(\left\lbrack  \mathrm{E}\right\rbrack   \mapsto  \dim \mathrm{E}\) 通过取商诱导出一个满的同态 Coker \(\left( c\right)  \rightarrow  \mathbf{Z}/{p}^{n}\mathbf{Z}\) 。特别地， \({\mathrm{R}}_{k}\left( \mathrm{G}\right)\) 中的元素 \({p}^{n - 1}\) 不在 \(c\) 的像中。

16.2 描述 \(e\) 像的特征

一个 \(\mathrm{G}\) 的元素称为 p-奇异的，当且仅当它不是 \(p\) -正则的(参见 10.1)，即当其阶可被 \(p\) 整除。还记得每个 \({\mathbf{R}}_{\mathbf{K}}\left( \mathbf{G}\right)\) 的元素都可被识别为定义在 \(\mathbf{G}\) 上的类函数，即其特征(参见 12.1 和 14.1)。

定理 36. \(e : {\mathrm{P}}_{k}\left( \mathrm{G}\right)  \rightarrow  {\mathrm{R}}_{\mathrm{K}}\left( \mathrm{G}\right)\) 的像由那些在 \(\mathrm{G}\) 的 p-奇异元素上其特征为零的 \({\mathrm{R}}_{\mathrm{K}}\left( \mathrm{G}\right)\) 的元素组成。

我们甚至有更精确的结果:

定理 37. 设 \({\mathrm{K}}^{\prime }\) 是 \(\mathrm{K}\) 的有限扩张。要使一个 \({\mathrm{R}}_{{\mathrm{K}}^{\prime }}\left( \mathrm{G}\right)\) 的元素属于在 \(e\) 下 \({\mathrm{P}}_{\mathrm{A}}\left( \mathrm{G}\right)  = {\mathrm{P}}_{k}\left( \mathrm{G}\right)\) 的像，当且仅当其特征取值于 \(\mathrm{K}\) 并且在 \(\mathbf{G}\) 的 p-奇异元素上为零。

证明见 17.5。

\section*{练习}

16.4.(Swan)设 \(\Lambda\) 为商域为 F 的 Dedekind 范畴。假定对每个除 \(\mathbf{G}\) 的阶的素数 \(p\) ，存在 \(\Lambda\) 的素理想 \(\mathfrak{p}\) 使得 \(\Lambda /\mathfrak{p}\) 的特征为 \(p\) 。设 \(\mathrm{P}\) 为一个投射的 \(\Lambda \left\lbrack  \mathrm{G}\right\rbrack\) -模。证明 \(\mathrm{F} \otimes  \mathrm{P}\) 为自由的 \(\mathrm{F}\left\lbrack  \mathrm{G}\right\rbrack\) -模。[对在此类素理想 \(\mathfrak{p}\) 处完成后从 \(\mathrm{P}\) 得到的模套用定理 36。由此推得 \(\mathrm{F} \otimes  \mathrm{P}\) 的特征在 \(\mathrm{G}\) 的非单位元上为零。]

此练习特别适用于 \(\Lambda\) 为代数数域的整数环的情形。

\section*{16.3 通过特征对投射的 A[G]-模的刻画}

这样的刻画等同于确定那些在 K 上的 G 表示，其含有一个在 G 下不变且作为 A[G]-模为投射的格。换言之，即刻画 \({\mathbf{P}}_{\mathbf{K}}^{ + }\left( \mathbf{G}\right)  = {\mathbf{P}}_{\mathbf{A}}^{ + }\left( \mathbf{G}\right)\) 在 \(e\) 下的像。仅知部分结果。首先，

引理 22. 设 \(x \in  {\mathrm{P}}_{\mathrm{A}}\left( \mathrm{G}\right)\) ，并设 \(n \geq  1\) 为一整数。如果 \({nx} \in  {\mathrm{P}}_{\mathrm{A}}^{ + }\left( \mathrm{G}\right)\) ，我们有 \(x \in  {\mathrm{P}}_{\mathrm{A}}^{ + }\left( \mathrm{G}\right)\) 。

这很清楚:如果 \(r = \operatorname{Card}\left( {\mathrm{S}}_{k}\right)\) ，那么可将 \({\mathrm{P}}_{\mathrm{A}}\left( \mathrm{G}\right)\) 识别为 \({\mathbf{Z}}^{r}\) ，将 \({\mathrm{P}}_{\mathrm{A}}^{ + }\left( \mathrm{G}\right)\) 识别为 \({\mathrm{N}}^{r}\) ，参见 14.3 和 14.4。

命题 44。设 \({\mathrm{K}}^{\prime }\) 为 \(\mathrm{K}\) 的有限扩张，设 \({\mathrm{A}}^{\prime }\) 为 \({\mathrm{K}}^{\prime }\) 的整环。对 \({\mathrm{R}}_{{\mathrm{K}}^{\prime }}\left( \mathrm{G}\right)\) 中的元素 \(x\) 假定下列两条条件:

(a) \(x\) 的特征取值于 \(\mathrm{K}\) 。

(b) 存在整数 \(n \geq  1\) ，使得 nx 通过标量扩张来源于一个投射的 \({\mathrm{A}}^{\prime }\left\lbrack  \mathrm{G}\right\rbrack\) -模。

则 \(x\) 来源于一个投射的 \(\mathrm{A}\left\lbrack  \mathrm{G}\right\rbrack\) -模，且此模在同构下唯一确定。

设 \(\mathrm{N} = \left\lbrack  {{\mathrm{K}}^{\prime } : \mathrm{K}}\right\rbrack   = \left\lbrack  {{\mathrm{A}}^{\prime } : \mathrm{A}}\right\rbrack\) 。令 \({\mathrm{E}}^{\prime }\) 为一个投射的 \({\mathrm{A}}^{\prime }\left\lbrack  \mathrm{G}\right\rbrack\) -模，其在 \({\mathrm{R}}_{{\mathrm{K}}^{\prime }}\left( \mathrm{G}\right)\) 中的像为 \({nx}\) ，并令 \(\mathrm{E}\) 为由将 \({\mathrm{E}}^{\prime }\) 限制到 \(\mathrm{A}\left\lbrack  \mathrm{G}\right\rbrack\) 得到的 \(\mathrm{A}\left\lbrack  \mathrm{G}\right\rbrack\) -模。容易验证 \(\mathrm{K} \otimes  \mathrm{E}\) 的特征等于 \(n\mathrm{\;N}\) 倍的 \(x\) 的特征。

因此

\[
e\left( \left\lbrack  \mathrm{E}\right\rbrack  \right)  = n\mathrm{\;N} \cdot  x\;\text{ in }{\mathrm{R}}_{{\mathrm{K}}^{\prime }}\left( \mathrm{G}\right) .
\]

由定理 36， \(e\left( \left\lbrack  \mathrm{\;E}\right\rbrack  \right)\) 的特征在 \(\mathrm{G}\) 的 \(p\) -奇异元上为零；因此 \(x\) 亦然。于是由定理 37，我们有 \(x = e\left( y\right)\) ，并且 \(y \in  {\mathrm{P}}_{\mathrm{A}}\left( \mathrm{G}\right)\) 。由于 \(e\) 是单射(定理 34)，这蕴含 \(\left\lbrack  \mathrm{E}\right\rbrack   = n\mathrm{\;N} \cdot  y\) ，引理 22 表明 \(y\) 属于 \({\mathrm{P}}_{\mathrm{A}}^{ + }\left( \mathrm{G}\right)\) 。因此存在一个投射的 \(\mathrm{A}\left\lbrack  \mathrm{G}\right\rbrack\) -模 \(\mathrm{Y}\) 使得 \(\left\lbrack  {\mathrm{K} \otimes  \mathrm{Y}}\right\rbrack   = x\) 在 \({\mathrm{R}}_{\mathrm{K}}\left( \mathrm{G}\right)\) 中； \(\mathrm{Y}\) 的唯一性(同构下)由定理 34 的推论 2 得出。

有人可以问是否 \(e\left( {{\mathrm{P}}_{\mathrm{A}}^{ + }\left( \mathrm{G}\right) }\right)  = e\left( {{\mathrm{P}}_{\mathrm{A}}\left( \mathrm{G}\right) }\right)  \cap  {\mathrm{R}}_{\mathrm{K}}^{ + }\left( \mathrm{G}\right)\) 。这在一般情况下不成立(参见例 16.5 和 16.7)。然而，我们有如下判准:

命题 45。假设满足下列条件:

(R) 存在 \(\mathrm{K}\) 的有限扩张 \({\mathrm{K}}^{\prime }\) ，其剩余域为 \({k}^{\prime }\) ，使得 \(d\left( {{\mathrm{\;R}}_{{\mathrm{K}}^{\prime }}^{ + }\left( \mathrm{G}\right) }\right)  = {\mathrm{R}}_{{k}^{\prime }}^{ + }\left( \mathrm{G}\right) .\)

则我们有 \(e\left( {{\mathrm{P}}_{\mathrm{A}}^{ + }\left( \mathrm{G}\right) }\right)  = e\left( {{\mathrm{P}}_{\mathrm{A}}\left( \mathrm{G}\right) }\right)  \cap  {\mathrm{R}}_{\mathrm{K}}^{ + }\left( \mathrm{G}\right)\) 。

由命题 44 足以证明的是

\[
e\left( {{\mathrm{P}}_{\mathrm{A}}^{ + }\left( \mathrm{G}\right) }\right)  = e\left( {{\mathrm{P}}_{\mathrm{A}}\left( \mathrm{G}\right) }\right)  \cap  {\mathrm{R}}_{\mathrm{K}}^{ + }\left( \mathrm{G}\right)
\]

当 \(\mathrm{K}\) 足够大时，此时条件 \(\left( \mathrm{R}\right)\) 仅意味着 \(d\) 将 \({\mathrm{R}}_{\mathrm{K}}^{ + }\left( \mathrm{G}\right)\) 映到 \({\mathrm{R}}_{k}^{ + }\left( \mathrm{G}\right)\) 。现在令

\[
x \in  e\left( {{\mathrm{P}}_{\mathrm{A}}\left( \mathrm{G}\right) }\right)  \cap  {\mathrm{R}}_{\mathrm{K}}^{ + }\left( \mathrm{G}\right) .
\]

由于 \(x \in  e\left( {{\mathrm{P}}_{\mathrm{A}}\left( \mathrm{G}\right) }\right)\) ，我们可以将 \(x\) 写成

\[
x = \mathop{\sum }\limits_{{\mathrm{E} \in  {\mathrm{S}}_{k}}}{n}_{\mathrm{E}}e\left( \left\lbrack  {\widetilde{\mathrm{P}}}_{\mathrm{E}}\right\rbrack  \right)
\]

其中 \({\widetilde{P}}_{E}\) 表示一个 A[G]-投射模，其模 \(\mathfrak{m}\) 意义下的化简是 \(E\) 的投射包 \({P}_{E}\) (参见 14.4)。我们必须证明整数 \({n}_{E}\) 非负。由 \(\left( \mathrm{R}\right)\) ，对每个 \(\mathrm{E} \in  {\mathrm{S}}_{k}\) 存在 \({z}_{\mathrm{E}} \in  {\mathrm{R}}_{\mathrm{K}}^{ + }\left( \mathrm{G}\right)\) 使得 \(d\left( {z}_{\mathrm{E}}\right)  = \left\lbrack  \mathrm{E}\right\rbrack\) 。由于 \(x \in  {\mathrm{R}}_{k}^{ + }\left( \mathrm{G}\right)\) ，我们有 \({\left\langle  x,{z}_{\mathrm{E}}\right\rangle  }_{\mathrm{K}} \geq  0\) 。另一方面， \(d\) 与 \(e\) 为伴随函子这一事实表明 \({\left\langle  x,{z}_{\mathrm{E}}\right\rangle  }_{\mathrm{K}} = {n}_{\mathrm{E}}\) 。尤其 \({n}_{\mathrm{E}}\) 非负，证毕。

结合命题 45 与定理 36，得到:

推论。假设 G 满足命题 45 的条件 (R)。一个定义在 \(\mathrm{K}\) 上的线性表示来自一个投射的 \(\mathrm{A}\left\lbrack  \mathrm{G}\right\rbrack\) -模，当且仅当其特征在 \(\mathrm{G}\) 的 p-奇异元素上消失。

注。条件 (R) 等价于下列条件:

(R’) 若 \(\mathrm{K}\) 足够大，则每个简单的 \(k\left\lbrack  \mathrm{G}\right\rbrack\) -模都是某个 K[G]-模(必然为简单模)对模 \(\mathfrak{m}\) 的化简。

(换言之， \(\mathbf{G}\) 在 \(k\) 上的每个不可约线性表示都可提升到 K。)

定理 38. (Fong–Swan)。假设 \(\mathbf{G}\) 是 p-可解的，即存在一个组成列，其商要么是 p-群，要么是与 p 互素阶的群。则 \(\mathrm{G}\) 满足上述条件 \(\left( \mathrm{R}\right)\) 与 \(\left( {\mathrm{R}}^{\prime }\right)\) 。

证明见 17.6。

练习

16.5. 采用命题 44 中的记法，证明

\[
{\mathrm{P}}_{\mathrm{A}}^{ + }\left( \mathrm{G}\right)  = {\mathrm{P}}_{{\mathrm{A}}^{\prime }}^{ + }\left( \mathrm{G}\right)  \cap  {\mathrm{P}}_{\mathrm{A}}\left( \mathrm{G}\right)  = {\mathrm{P}}_{{\mathrm{A}}^{\prime }}^{ + }\left( \mathrm{G}\right)  \cap  {\mathrm{R}}_{\mathrm{K}}\left( \mathrm{G}\right) .
\]

16.6. 证明:当 K 足够大时，条件 (R) 等价于条件 \(e\left( {{\mathrm{P}}_{\mathrm{A}}^{ + }\left( \mathrm{G}\right) }\right)  = e\left( {{\mathrm{P}}_{\mathrm{A}}\left( \mathrm{G}\right) }\right)  \cap  {\mathrm{R}}_{\mathrm{K}}^{ + }\left( \mathrm{G}\right)\) 。(注意， \({P}_{k}\left( G\right)\) 中的元素 \(x\) 属于 \({P}_{k}^{ + }\left( G\right)\) 当且仅当对于所有 \(y \in  {R}_{k}^{ + }\left( G\right)\) 都有 \(\langle x,y{\rangle }_{k} \geq  0\) 。)

16.7. 取 \(G\) 为群 \({SL}\left( V\right)\) ，其中 \(V\) 是维数为 2 的 \({\mathbf{F}}_{p} = \mathbf{Z}/p\mathbf{Z}\) 上的向量空间。证明对 \(i < p\) ， \(\mathbf{G}\) 在 \(\mathbf{V}\) 的第 \(i\) 次对称幂 \({\mathbf{V}}_{i}\) 中的自然表示是绝对不可约的。(由于 \(G\) 的 \(p\) -正则共轭类数为 \(p\) ，故这些表示(至同构)构成 G 的全部不可约表示，参见 18.2，定理 42 的推论 2。)给出这些表示即便在充分大的特征零域 K 上也无法提升的例子。(对于 \(p = 7,i = 4\) ，我们有 \(\dim {V}_{i} = 5\) ，且 5 不整除 \(G\) 的阶；因此 \({V}_{i}\) 无法提升。)

\section*{16.4 投影 A[G]-模的例子:缺陷为零的不可约表示}

本节中我们假设 \(\mathbf{K}\) 足够大。

命题 46. 设 \(\mathrm{E}\) 为一个简单的 \(\mathrm{K}\left\lbrack  \mathrm{G}\right\rbrack\) -模，且 \(\mathrm{P}\) 是 \(\mathrm{E}\) 中一个在 G 下不变的格。假设 \(\mathrm{E}\) 的维数 \(\mathrm{N}\) 可被最大的幂 \({p}^{n}\) (该幂为 \(p\) 的某次幂)整除，而这一幂整除群 \(\mathrm{G}\) 的阶 \(g\) 。则:

(a) \(\mathrm{P}\) 是一个投影的 \(\mathrm{A}\left\lbrack  \mathrm{G}\right\rbrack\) -模。

(b) 自然映射 \(\mathrm{A}\left\lbrack  \mathrm{G}\right\rbrack   \rightarrow  {\operatorname{End}}_{\mathrm{A}}\left( \mathrm{P}\right)\) 是满射，其核作为双边理想在 \(\mathrm{A}\left\lbrack  \mathrm{G}\right\rbrack\) 中是一个直和因子。

(c) \(\mathrm{P}\) 的约化 \(\overline{\mathrm{P}} = \mathrm{P}/\mathrm{{mP}}\) 是一个简单且投影的 \(\mathrm{A}\left\lbrack  \mathrm{G}\right\rbrack\) -模。

注意 (a) 蕴含(参见定理 37):

推论。群 \(\mathrm{E}\) 的特征 \({\chi }_{\mathrm{E}}\) 在 \(\mathrm{G}\) 的 p-奇元素上为零。

首先，既然 \(\mathrm{N}\) 可被 \({p}^{n}\) 整除，商 \(\mathrm{N}/g\) 属于环 A。这样我们可以在不引入任何“分母”的情况下(即在环 A 内)应用傅里叶反演 (6.2., 命题 11)。更精确地，设 \({s}_{\mathrm{P}}\) 为由 \(s \in  \mathrm{G}\) 定义的 \(\mathrm{P}\) 的自同态；如果 \(\phi  \in  {\operatorname{End}}_{\mathrm{A}}\left( \mathrm{P}\right)\) ，则 \({s}_{\mathrm{P}}^{-1}\phi\) 的迹 \(\operatorname{Tr}\left( {{s}_{\mathrm{P}}^{-1}\phi }\right)\) 属于 \(\mathrm{A}\) ，因此我们可以定义该元素

\[
{u}_{\phi } = \frac{\mathrm{N}}{g}\mathop{\sum }\limits_{{s \in  \mathrm{G}}}\operatorname{Tr}\left( {{s}_{\mathrm{P}}^{-1}\phi }\right) s\;\text{ of the ring }\mathrm{A}\left\lbrack  \mathrm{G}\right\rbrack  .
\]

由命题 11 得到 \({u}_{\phi }\) 在 \({\operatorname{End}}_{\mathrm{K}}\left( \mathrm{E}\right)\) 中的像为 \(1 \otimes  \phi\) ，并且对每个不与 \(E\) 同构的单 \(K\left\lbrack  G\right\rbrack\) -模 \({E}^{\prime }\) 在 \({\operatorname{End}}_{K}\left( {E}^{\prime }\right)\) 中的像为 0。特别地， \({u}_{\phi }\) 在 \({\operatorname{End}}_{A}\left( P\right)\) 中的像为 \(\phi\) ，这证明了 (b)。断言 (a) 随即由下面这一基本事实得出: \(P\) 在环 \({\operatorname{End}}_{A}\left( P\right)\) 上是投射的；同样的论证也适用于 (c)。

注。在模块块理论的表述中(参见 [9], [20])，命题 46 是具有唯一不可约特征的块(或缺陷零)的情形。

例。若 \(\mathbf{G}\) 是特征为 \(p\) 的有限域上的一个半单线性群，则存在一个 \(\mathbf{G}\) 的线性不可约表示(在 \(\mathbf{Q}\) 上)，其次数等于 \({p}^{n}\) ；这是由 \(\mathrm{R}\) 发现的 \(\mathrm{G}\) 的特殊表示。Steinberg(参见 Canad. J. of Math., 8, 1956, p. 580-591 和 9, 1957, p. 347-351)。根据 Solomon–Tits 的结果，它可实现为与 \({\mathrm{G}}^{ * }\) 关联的 Tits 建筑的最高维同调表示。

\section*{练习}

16.8. 取 \(G = {\mathfrak{A}}_{4}\) ,参见 5.7。证明:对于 \(p = 2\) ，群 \(G\) 不存在命题 46 所述类型的不可约表示，但对于 \(p = 3\) 存在这样的表示。对于 \({\mathfrak{S}}_{4}\) 也问同样的问题。

16.9. 设 \(\mathrm{S} \in  {\mathrm{S}}_{k}\) 。证明下列性质等价:

(i) \(\mathrm{S}\) 是一个投射的 \(k\left\lbrack  \mathrm{G}\right\rbrack\) -模。

(ii) \(\mathrm{S}\) 同构于满足命题 46 条件的模 \(\mathrm{P}\) 的约化 \({\;\operatorname{mod.}}\mathfrak{m}\) 。

(iii) \({\mathrm{C}}_{\mathrm{{SS}}}\) 为 \(\mathrm{G}\) 的卡坦矩阵的对角元等于 1。(关于 (i) 与 (iii) 的等价性，参见习题 15.7。)

\HRule

* 参见 L. Solomon, The Steinberg character of a finite group with a BN-pair. Theory of Finite Groups, edited by R. Brauer and C.-H Sah. W. A. Benjamin, New York, 1969, p. 213-221.

\HRule

第 17 章

证明

\section*{17.1 群的变换}

设 \(\mathrm{H}\) 为 \(\mathrm{G}\) 的一个子群。我们已相对于 \({\mathbf{R}}_{\mathbf{K}}\) 定义了限制与诱导同态:

\[
{\operatorname{Res}}_{\mathrm{H}}^{\mathrm{G}} : {\mathrm{R}}_{\mathrm{K}}\left( \mathrm{G}\right)  \rightarrow  {\mathrm{R}}_{\mathrm{K}}\left( \mathrm{H}\right) \;\text{ and }\;{\operatorname{Ind}}_{\mathrm{H}}^{\mathrm{G}} : {\mathrm{R}}_{\mathrm{K}}\left( \mathrm{H}\right)  \rightarrow  {\mathrm{R}}_{\mathrm{K}}\left( \mathrm{G}\right) .
\]

同样的定义也适用于 \({\mathrm{R}}_{k}\) 与 \({\mathrm{P}}_{k}\) :通过限制，每个 \(k\) [G]-模给出一个 \(k\left\lbrack  \mathrm{H}\right\rbrack\) -模，若原模是投射的则该模亦为投射。通过传到格罗滕迪克群，我们得到同态

\[
{\operatorname{Res}}_{\mathrm{H}}^{\mathrm{G}} : {\mathrm{R}}_{k}\left( \mathrm{G}\right)  \rightarrow  {\mathrm{R}}_{k}\left( \mathrm{H}\right) \;\text{ and }\;{\operatorname{Res}}_{\mathrm{H}}^{\mathrm{G}} : {\mathrm{P}}_{k}\left( \mathrm{G}\right)  \rightarrow  {\mathrm{P}}_{k}\left( \mathrm{H}\right) .
\]

另一方面，若 \(\mathrm{E}\) 是一个 \(k\left\lbrack  \mathrm{H}\right\rbrack\) -模，则 Ind \(\mathrm{E} = k\left\lbrack  \mathrm{G}\right\rbrack  { \otimes  }_{k\left\lbrack  \mathrm{H}\right\rbrack  }\mathrm{E}\) 是一个 \(k\left\lbrack  \mathrm{G}\right\rbrack\) -模(称由 \(\mathrm{E}\) 诱导)，若 \(\mathrm{E}\) 是投射的则诱导后亦为投射。因此我们有同态

\[
{\operatorname{Ind}}_{\mathrm{H}}^{\mathrm{G}} : {\mathrm{R}}_{k}\left( \mathrm{H}\right)  \rightarrow  {\mathrm{R}}_{k}\left( \mathrm{G}\right) \;\text{ and }\;{\operatorname{Ind}}_{\mathrm{H}}^{\mathrm{G}} : {\mathrm{P}}_{k}\left( \mathrm{H}\right)  \rightarrow  {\mathrm{P}}_{k}\left( \mathrm{G}\right) .
\]

利用张量积的结合律，我们容易得到公式

(*)
\[
{\operatorname{Ind}}_{\mathrm{H}}^{\mathrm{G}}\left( {x \cdot  {\operatorname{Res}}_{\mathrm{H}}^{\mathrm{G}}y}\right)  = {\operatorname{Ind}}_{\mathrm{H}}^{\mathrm{G}}\left( x\right)  \cdot  y
\]

在下列每种情形中:

(a) \(x \in  {\mathrm{R}}_{\mathrm{K}}\left( \mathrm{H}\right) ,y \in  {\mathrm{R}}_{\mathrm{K}}\left( \mathrm{G}\right) \;\) 与 \(\;{\operatorname{Ind}}_{\mathrm{H}}^{\mathrm{G}}\left( x\right)  \cdot  y \in  {\mathrm{R}}_{\mathrm{K}}\left( \mathrm{G}\right)\) ，

(b) \(x \in  {\mathrm{R}}_{k}\left( \mathrm{H}\right) ,y \in  {\mathrm{R}}_{k}\left( \mathrm{G}\right) \;\) 与 \(\;{\operatorname{Ind}}_{\mathrm{H}}^{\mathrm{G}}\left( x\right)  \cdot  y \in  {\mathrm{R}}_{k}\left( \mathrm{G}\right)\) ,

(c) \(x \in  {\mathrm{R}}_{k}\left( \mathrm{H}\right) ,y \in  {\mathrm{P}}_{k}\left( \mathrm{G}\right) \;\) 与 \(\;{\operatorname{Ind}}_{\mathrm{H}}^{\mathrm{G}}\left( x\right)  \cdot  y \in  {\mathrm{P}}_{k}\left( \mathrm{G}\right)\) 。

[情况 (c) 有道理，因为 \({\mathrm{P}}_{k}\left( \mathrm{G}\right)\) 是 \({\mathrm{R}}_{k}\left( \mathrm{G}\right)\) 的一个模。]

此外，第15章中的同态 \(c,d,e\) 与同态 \({\operatorname{Res}}_{\mathrm{H}}^{\mathrm{G}}\) 和 \({\operatorname{Ind}}_{\mathrm{H}}^{\mathrm{G}}\) 互相交换。

习题

17.1. 将 \({\operatorname{Res}}_{\mathrm{H}}^{\mathrm{G}}\) 与 \({\operatorname{Ind}}_{\mathrm{H}}^{\mathrm{G}}\) 的定义推广到同态 \(\mathrm{H} \rightarrow  \mathrm{G}\) 的情形，其核的阶与 \(p\) 互素(参见习题 7.1)。

\section*{17.2 模情形下的布劳尔定理}

定理 39。设 \(\mathbf{X}\) 为 \(\mathrm{G}\) 的所有 \({\Gamma }_{\mathbf{K}}\) -初等子群的集合(参见 12.6)。同态

\[
\text{ Ind: }\mathop{\oplus }\limits_{{\mathrm{H} \in  \mathrm{X}}}{\mathrm{R}}_{k}\left( \mathrm{H}\right)  \rightarrow  {\mathrm{R}}_{k}\left( \mathrm{G}\right)
\]

和

\[
\text{ Ind: }\mathop{\oplus }\limits_{{\mathrm{H} \in  \mathrm{X}}}{\mathrm{P}}_{k}\left( \mathrm{H}\right)  \rightarrow  {\mathrm{P}}_{k}\left( \mathrm{G}\right)
\]

由 \({\operatorname{Ind}}_{\mathrm{H}}^{\mathrm{G}}\) 为 \(\mathrm{H} \in  \mathrm{X}\) 定义，是满射。

(换言之，定理27 对 \({\mathrm{R}}_{k}\) 与 \({\mathrm{P}}_{k}\) 成立。)

令 \({l}_{K}\) (分别 \({l}_{k}\) )为环 \({R}_{K}\left( G\right)\) (分别 \({R}_{k}\left( G\right)\) )的单位元。我们有 \(d\left( {l}_{K}\right)  = {l}_{k}\) 。由定理27 我们可以将 \({l}_{K}\) 写成

\[
{\mathrm{l}}_{\mathrm{K}} = \mathop{\sum }\limits_{{\mathrm{H} \in  \mathrm{X}}}{\operatorname{Ind}}_{\mathrm{H}}\left( {x}_{\mathrm{H}}\right) \;\text{ with }{x}_{\mathrm{H}} \in  {\mathrm{R}}_{\mathrm{K}}\left( \mathrm{H}\right) .
\]

作用 \(d\) ，并利用 \(d\) 与 \({\operatorname{Ind}}_{\mathrm{H}}^{\mathrm{G}}\) 交换的事实，我们得到对 \({1}_{k}\) 的类似公式:

\[
{\mathrm{l}}_{k} = \mathop{\sum }\limits_{{\mathrm{H} \in  \mathrm{X}}}{\operatorname{Ind}}_{\mathrm{H}}\left( {x}_{\mathrm{H}}^{\prime }\right) ,\;\text{ with }{x}_{\mathrm{H}}^{\prime } = d\left( {x}_{\mathrm{H}}\right)  \in  {\mathrm{R}}_{k}\left( \mathrm{H}\right) .
\]

对 \(y \in  {\mathrm{R}}_{k}\left( \mathrm{G}\right)\) (分别 \({\mathrm{P}}_{k}\left( \mathrm{G}\right)\) )，通过乘法我们得到:

\[
y = {\mathrm{l}}_{k} \cdot  y = \mathop{\sum }\limits_{{\mathrm{H} \in  \mathrm{X}}}{\operatorname{Ind}}_{\mathrm{H}}\left( {x}_{\mathrm{H}}^{\prime }\right)  \cdot  y = \mathop{\sum }\limits_{{\mathrm{H} \in  \mathrm{X}}}{\operatorname{Ind}}_{\mathrm{H}}^{\mathrm{G}}\left( {{x}_{\mathrm{H}}^{\prime } \cdot  {\operatorname{Res}}_{\mathrm{H}}^{\mathrm{G}}y}\right) ,
\]

从而证明了该定理。

推论. 若 \(\mathrm{K}\) 足够大，则 \({\mathrm{R}}_{k}\left( \mathrm{G}\right)\) 的每个元素(分别 \({\mathrm{P}}_{k}\left( \mathrm{G}\right)\) 的元素)均可表示为形如 \({\operatorname{Ind}}_{\mathrm{H}}\left( {y}_{\mathrm{H}}\right)\) 的元素之和，其中 \(\mathrm{H}\) 为 \(\mathrm{G}\) 的一个初等子群，且 \({y}_{\mathrm{H}}\) 属于 \({\mathrm{R}}_{k}\left( \mathrm{H}\right)\) (分别 属于 \({\mathrm{P}}_{k}\left( \mathrm{H}\right)\) )。

确实，当 \(\mathrm{K}\) 足够大时， \(\mathrm{X}\) 仅为 \(\mathbf{G}\) 的所有基元子群的集合。

注:证明定理39所用的论证可适用于许多其他情形(参见 Swan [21]，§§ 3,4)。例如，它给出了 Artin 定理的以下类比(参见定理26):

定理40。设 \(\mathrm{T}\) 为 \(\mathrm{G}\) 的所有循环子群的集合。群同态

\[
\mathrm{Q} \otimes  \text{ Ind: }\underset{\mathrm{H} \in  \mathrm{T}}{ \oplus  }\mathrm{Q} \otimes  {\mathrm{R}}_{k}\left( \mathrm{H}\right)  \rightarrow  \mathrm{Q} \otimes  {\mathrm{R}}_{k}\left( \mathrm{G}\right)
\]

和

\[
\mathbf{Q} \otimes  \text{ Ind: }\underset{\mathrm{H} \in  \mathrm{T}}{ \oplus  }\mathbf{Q} \otimes  {\mathrm{P}}_{k}\left( \mathrm{H}\right)  \rightarrow  \mathbf{Q} \otimes  {\mathrm{P}}_{k}\left( \mathrm{G}\right)
\]

是满射。

\section*{17.3 定理33的证明}

我们须证明 \(d : {\mathrm{R}}_{\mathrm{K}}\left( \mathrm{G}\right)  \rightarrow  {\mathrm{R}}_{k}\left( \mathrm{G}\right)\) 是满的。由定理 \({39},{\mathrm{R}}_{k}\left( \mathrm{G}\right)\) 可知它由各个 \({\operatorname{Ind}}_{\mathrm{H}}^{\mathrm{G}}\left( {{\mathrm{R}}_{k}\left( \mathrm{H}\right) }\right)\) 生成，其中 \(\mathrm{H}\) 为 \({\Gamma }_{\mathrm{K}}\) -元的。既然 \(d\) 与 \({\operatorname{Ind}}_{\mathrm{H}}^{\mathrm{G}}\) 交换，便足以证明 \({\mathrm{R}}_{k}\left( \mathrm{H}\right)  = d\left( {{\mathrm{R}}_{\mathrm{K}}\left( \mathrm{H}\right) }\right)\) 。因此我们归约到 \(\mathrm{G}\) 为 \({\Gamma }_{\mathrm{K}}\) -元的情形。在此情形下我们有更精确的结论:

定理41。令 \(l\) 为素数。设 \(\mathrm{G}\) 为一个 l-群 \(\mathrm{P}\) 与一循环正规子群 \(\mathrm{C}\) (阶与 \(l\) 互素)的半直积。则每个单 \(k\left\lbrack  \mathrm{G}\right\rbrack\) -模 \(\mathrm{E}\) 都可提升(即是某个 \(\mathrm{A}\) -自由 A[G]-模的约化 \({\;\operatorname{mod}\;\mathfrak{m}}\) )。

(换言之， \(d\) 将 \({\mathrm{R}}_{\mathrm{K}}^{ + }\left( \mathrm{G}\right)\) 映到 \({\mathrm{R}}_{k}^{ + }\left( \mathrm{G}\right)\) 。)

假设 \(l \neq  p\) 。令 \({\mathrm{C}}_{p}\) 为 \(\mathrm{C}\) 的 \(p\) -Sylow 子群，令 \({\mathrm{E}}^{\prime }\) 为 \(\mathrm{E}\) 中被 \({\mathrm{C}}_{p}\) 固定的向量子空间。由于 \({\mathrm{C}}_{p}\) 是 \(p\) -群，我们有 \({\mathrm{E}}^{\prime } \neq  0\) ，参见 8.3，命题26。因 \({\mathrm{C}}_{p}\) 在 \(\mathrm{G}\) 中为正规子群，故 \({\mathrm{E}}^{\prime }\) 在 \(\mathrm{G}\) 下稳定。于是 \({\mathrm{E}}^{\prime } = \mathrm{E}\) ，这意味着 \({C}_{p}\) 在 \(\mathrm{E}\) 上作平凡作用，且 \(\mathrm{G}\) 在 \(\mathrm{E}\) 上的表示来源于 \(\mathrm{G}/{\mathrm{C}}_{p}\) 的表示。由于 \(\mathrm{G}/{\mathrm{C}}_{p}\) 的阶与 \(p\) 互素，此类表示可立即被提升(参见15.5)。

假设现在 \(l = p\) 。我们对 \(\mathrm{G}\) 的阶按归纳法证明。由于 \(\mathrm{C}\) 的阶与 \(p\) 互素， \(\mathrm{C}\) 在 \(\mathrm{E}\) 中的表示是半单的。将其分解为同型型 \(k\left\lbrack  \mathrm{C}\right\rbrack\) -模的直和(参见 8.1 命题 24):

\[
\mathrm{E} = \underset{\alpha }{ \oplus  }{\mathrm{E}}_{\alpha }
\]

群 \(\mathrm{G}\) 对 \({\mathrm{E}}_{\alpha }\) 们作置换；因 \(\mathrm{E}\) 为单群， \(\mathrm{G}\) 对非零的 \({\mathrm{E}}_{\alpha }\) 作传递置换。令 \({\mathrm{E}}_{\beta }\) 为其中之一，令 \({\mathrm{G}}_{\beta }\) 为 \(\mathrm{G}\) 的子群，包含那些使得 \(s{\mathrm{E}}_{\beta } = {\mathrm{E}}_{\beta }\) 的元素 \(s\) 。显然 \({\mathrm{E}}_{\beta }\) 为一个 \(k\left\lbrack  {\mathrm{G}}_{\beta }\right\rbrack\) -模，且 \(\mathrm{E}\) 与相应的诱导模 \({\operatorname{Ind}}_{{G}_{\beta }}^{G}\left( {E}_{\beta }\right)\) 同构。此外， \({G}_{\beta }\) 是 \(P\) 的一个子群与群 \(C\) 的半直积。如果 \({E}_{\beta } \neq  E\) ，则有 \({G}_{\beta } \neq  G\) ，并且对 \({G}_{\beta }\) 应用归纳假设表明 \({E}_{\beta }\) 可以被提升；同样的结论也适用于 E。

因此我们可假定 \(\mathrm{E}\) 是一 个同型型的 \(k\left\lbrack  \mathrm{C}\right\rbrack\) -模。令 \(\rho\) 表示从 \(k\left\lbrack  \mathrm{G}\right\rbrack\) 到 \({\operatorname{End}}_{k}\left( \mathrm{E}\right)\) 的同态，该同态决定了 E 的 \(k\left\lbrack  \mathrm{G}\right\rbrack\) -模结构。 \(\mathrm{E}\) 为 \(k\left\lbrack  \mathrm{C}\right\rbrack\) -同型型等价于 \(k\left\lbrack  \mathrm{C}\right\rbrack\) 在 \(\rho\) 下的像是一个域 \({k}^{\prime }\) ，它是 \(k\) 的有限扩张。将 \(\rho\) 限制到 \(\mathbf{C}\) 得到同态 \(\phi  : \mathbf{C} \rightarrow  {k}^{\prime  * }\) ，且 \({k}^{\prime }\) 由 \(\phi \left( \mathrm{C}\right)\) 在 \(k\) 上生成。因此 \(\mathrm{E}\) 具有 \({k}^{\prime }\) -向量空间的结构。现在选取一个在 \(\mathrm{P}\) 下不变的 \(v \neq  0\) 属于 \(\mathrm{E}\) ；这又是可能的，因为 \(\mathrm{P}\) 是一个 \(p\) -群，参见 8.3 命题 26。对 \(x \in  \mathrm{C},s \in  \mathrm{P}\) ，置 \({}^{s}x = {sx}{s}^{-1}\) 。我们有

\[
\rho \left( s\right) \left( {\phi \left( x\right)  \cdot  v}\right)  = \rho \left( {{sx}{s}^{-1}}\right) \rho \left( s\right)  \cdot  v = \phi \left( {{}^{s}x}\right)  \cdot  v.
\]

因此由 \(\phi \left( x\right)  \cdot  v,x \in  \mathbf{C}\) 生成的子空间 \({k}^{\prime }v\) 在 \(\mathrm{C}\) 和 \(\mathrm{P}\) 下不变，因此等于 \(\mathrm{E}\) 。因此 \({\dim }_{{k}^{\prime }}\mathrm{E} = 1\) 。这使我们能以使 \(v\) 成为 \({k}^{\prime }\) 的单位元的方式将E与 \({k}^{\prime }\) 识别。对所有 \(t \in  \mathrm{G} \; \rho \left( t\right)\) 都是 \(k\) -向量空间 \({k}^{\prime }\) 的端om \({\sigma }_{t}\) 。对于 \(s \in  \mathbf{P}\) ，我们按构造有 \({\sigma }_{s}\left( 1\right)  = 1\) 。此外，上述公式表明

\[
{\sigma }_{s}\left( {\phi \left( x\right) }\right)  = \phi \left( {{}^{s}x}\right) \;\text{ for all }x \in  \mathrm{C},
\]

因此

\[
{\sigma }_{s}\left( {\phi \left( x\right) \phi \left( {x}^{\prime }\right) }\right)  = {\sigma }_{s}\left( {\phi \left( x\right) }\right) {\sigma }_{s}\left( {\phi \left( {x}^{\prime }\right) }\right) \;\text{ for all }x,{x}^{\prime } \in  \mathrm{C}.
\]

由于 \({k}^{\prime }\) 由 \(\phi \left( x\right)\) 生成，我们得到

\[
{\sigma }_{s}\left( {a{a}^{\prime }}\right)  = {\sigma }_{s}\left( a\right) {\sigma }_{s}\left( {a}^{\prime }\right) \;\text{ if }a,{a}^{\prime } \in  {k}^{\prime };
\]

换言之， \({\sigma }_{s}\) 是域 \({k}^{\prime }\) 的自同构，映射 \(s \mapsto  {\sigma }_{s}\) 是一个同态 \(\sigma  : \mathrm{P} \rightarrow  \operatorname{Gal}\left( {{k}^{\prime }/k}\right)\) ，后者表示 \({k}^{\prime }/k\) 的伽罗瓦群。现在易于定义 \(\mathrm{E}\) 的提升:令 \({\mathrm{K}}^{\prime }\) 为与残量扩张 \({k}^{\prime }/k\) 对应的不分歧扩张 \(\mathrm{K}\) ，并令 \({\mathrm{A}}^{\prime }\) 为 \({\mathrm{K}}^{\prime }\) 的整数环。典范同构

\[
\mathrm{{Gal}}\left( {{\mathrm{K}}^{\prime }/\mathrm{K}}\right)  \rightarrow  \mathrm{{Gal}}\left( {{k}^{\prime }/k}\right)
\]

给出 \(P\) 在 \({K}^{\prime }\) 和 \({A}^{\prime }\) 上的作用(使用 \(\sigma\) )。另一方面，同态 \(\phi  : \mathbf{C} \rightarrow  {k}^{\prime  * }\) 可唯一提升(例如使用乘法表示)为同态 \(\widetilde{\phi } : C \rightarrow  {A}^{\prime  * }\) ，它通过乘法在 \({A}^{\prime }\) 上给出 \(C\) 的作用。由唯一性立刻(可见)我们仍然有

\[
{\sigma }_{s}\left( {\widetilde{\phi }\left( x\right) }\right)  = \widetilde{\phi }\left( {{}^{s}x}\right) \;\text{ for }x \in  \mathrm{C},s \in  \mathrm{P}.
\]

这意味着 \(\mathrm{C}\) 和 \(\mathrm{P}\) 在 \({\mathrm{A}}^{\prime }\) 上的作用结合成了G的作用。赋予这样的A[G]-模结构后， \({\mathrm{A}}^{\prime }\) 就是所需的提升。

注:当 \(\mathrm{K}\) 足够大时，我们仅在 \(\mathrm{G}\) 为初等群即与 \(\mathrm{C}\) 直积为 \(\mathrm{P}\) 的情形需要定理41。上面的证明变得更简单:群 \(\mathrm{P}\) 在简单模 \(\mathrm{E}\) 上作平凡作用，因而可将其视为简单的 \(k\left\lbrack  \mathrm{C}\right\rbrack\) -模并无难度地提升。

\section*{17.4 定理35的证明}

令 \({p}^{n}\) 为除尽G的阶的最大 \(p\) 冪。我们要证明 \(c : {\mathrm{P}}_{k}\left( \mathrm{G}\right)  \rightarrow  {\mathrm{R}}_{k}\left( \mathrm{G}\right)\) 的余核被 \({p}^{n}\) 消去。我们区分两种情形:

\section*{(a) \(\mathrm{K}\) 足够大}

根据推论， \({39},{\mathrm{R}}_{k}\left( \mathrm{G}\right)\) 由具有 \(\mathrm{H}\) 元素的 \({\operatorname{Ind}}_{\mathrm{H}}^{\mathrm{G}}\left( {{\mathrm{R}}_{k}\left( \mathrm{H}\right) }\right)\) 所生成。因此我们归结到 \(\mathrm{G}\) 为元的情况，故分解为乘积 \(\mathrm{S} \times  \mathrm{P}\) ，其中 \(\mathrm{S}\) 的阶与 \(p\) 互素，且 \(\mathrm{P}\) 是一个 \(p\) -群。我们在 15.7 中已见此类群的卡坦矩阵为数乘矩阵 \({p}^{n}\) 。在此情形下定理成立。

\section*{(b) 一般情形}

设 \({\mathrm{K}}^{\prime }\) 为 \(\mathrm{K}\) 的一有限足够大的扩张，其剩余域为 \({k}^{\prime }\) 。从 \(k\) 到 \({k}^{\prime }\) 的标量扩张给出交换图:

\(\downarrow  c\; \downarrow  {c}^{\prime }\; \downarrow  \gamma\)

\[
0 \rightarrow  \;{\mathrm{P}}_{k}\left( \mathrm{G}\right) \; \rightarrow  \;{\mathrm{P}}_{{k}^{\prime }}\left( \mathrm{G}\right) \; \rightarrow  \;\mathrm{P}\; \rightarrow  0
\]

\[
0 \rightarrow  {\mathrm{P}}_{k}\left( \mathrm{G}\right) \; \rightarrow  \;{\mathrm{R}}_{{k}^{\prime }}\left( \mathrm{G}\right) \; \rightarrow  \;\mathrm{R}\; \rightarrow  0,
\]

其中 \(P = {P}_{{k}^{\prime }}\left( G\right) /{P}_{k}\left( G\right)\) 与 \(R = {R}_{{k}^{\prime }}\left( G\right) /{R}_{k}\left( G\right)\) 。由此得正合列:

\[
0 \rightarrow  \operatorname{Ker}\left( c\right)  \rightarrow  \operatorname{Ker}\left( {c}^{\prime }\right)  \rightarrow  \operatorname{Ker}\left( \gamma \right)  \rightarrow  \operatorname{Coker}\left( c\right)  \rightarrow  \operatorname{Coker}\left( {c}^{\prime }\right) .
\]

由 (a) 知，Coker \(\left( {c}^{\prime }\right)\) 被 \({p}^{n}\) 杀死。因 \({\mathrm{P}}_{{k}^{\prime }}\left( \mathrm{G}\right)\) 与 \({\mathrm{R}}_{{k}^{\prime }}\left( \mathrm{G}\right)\) 秩相同，故 \({c}^{\prime }\) 为单射，因而 \(c\) 亦然，故 Coker \(\left( c\right)\) 有限。但我们已知(参见 14.6) \({\mathrm{P}}_{k}\left( \mathrm{G}\right)  \rightarrow  {\mathrm{P}}_{{k}^{\prime }}\left( \mathrm{G}\right)\) 是分裂单射。故群 \(\mathrm{P}\) 为 \(\mathbf{Z}\) -自由， \(\operatorname{Ker}\left( \gamma \right)\) 亦是。既然 \(\operatorname{Ker}\left( {c}^{\prime }\right)  = 0\) ，且 \(\operatorname{Coker}\left( c\right)\) 有限，上述正合列显示 \(\operatorname{Ker}\left( \gamma \right)  = 0\) ；因此 \(\operatorname{Coker}\left( c\right)\) 嵌入于 \(\operatorname{Coker}\left( {c}^{\prime }\right)\) 。既然后者被 \({p}^{n}\) 杀死，Coker \(\left( c\right)\) 亦然，从而证得本定理。

\section*{17.5 定理 37 的证明}

如有必要，通过扩张 \({\mathrm{K}}^{\prime }\) ，我们可假设 \({\mathrm{K}}^{\prime }\) 足够大。

\section*{(i) 必要性}

令 \(\mathrm{E}\) 为投射的 \(\mathrm{A}\left\lbrack  \mathrm{G}\right\rbrack\) -模，令 \(\chi\) 为 \({\mathrm{K}}^{\prime }\left\lbrack  \mathrm{G}\right\rbrack\) -模 \({\mathrm{K}}^{\prime }{ \otimes  }_{{\mathrm{A}}^{\prime }}\mathrm{E}\) 的特征。若 \(s \in  \mathrm{G}\) 为 \(p\) -奇异，需证明 \(\chi \left( s\right)  = 0\) 。将 \(G\) 替换为由 \(s\) 生成的循环子群，可假定 \(G\) 为循环群，因此形如 \(S \times  P\) ，其中 \(S\) 的阶与 \(p\) 互素，且 \(P\) 为一 \(p\) -群。由 15.7， \(E\) 同构于 \(F \otimes  {A}^{\prime }\left\lbrack  P\right\rbrack\) ，其中 \(F\) 为一个无 \({A}^{\prime }\) 的 \({A}^{\prime }\left\lbrack  S\right\rbrack\) -模。于是 \({\mathrm{K}}^{\prime } \otimes  \mathrm{E}\) 的特征 \(\chi\) 等于 \(\psi  \otimes  {r}_{\mathrm{p}}\) ，其中 \(\psi\) 为 \(\mathrm{S}\) 的一个特征， \({r}_{\mathrm{p}}\) 为 \(\mathrm{P}\) 的正则表示的特征。这样的特征在 \(S\) 外显然为零，因而特别有 \(\chi \left( s\right)  = 0\) 。

\section*{(ii) 充要性(第一部分)}

设 \(y \in  {\mathrm{R}}_{{k}^{\prime }}\left( \mathrm{G}\right)\) ，令 \(\chi\) 为相应的虚特征，并假定对 \(G\) 的每个 \(p\) -奇异元 \(s\) 都有 \(\chi \left( s\right)  = 0\) 。

我们将证明 \(y\) 属于 \({\mathrm{P}}_{{\mathrm{A}}^{\prime }}\left( \mathrm{G}\right)\) (此群通过 \(e\) 与 \({\mathrm{R}}_{{k}^{\prime }}\left( \mathrm{G}\right)\) 的一个子群相认同)。

由定理 39 的推论，我们有

\[
1 = \sum \operatorname{Ind}\left( {x}_{\mathrm{H}}\right) ,\;\text{ with }{x}_{\mathrm{H}} \in  {\mathrm{R}}_{{\mathrm{K}}^{\prime }}\left( \mathrm{G}\right) ,
\]

其中 \(\mathrm{H}\) 取遍 G 的所有元胞子群。两边乘以 \(y\) ，得:

\[
y = \sum \operatorname{Ind}\left( {y}_{\mathrm{H}}\right) ,\;\text{ with }{y}_{\mathrm{H}} = {x}_{\mathrm{H}} \cdot  {\operatorname{Res}}_{\mathrm{H}}\left( y\right)  \in  {\mathrm{R}}_{{\mathrm{K}}^{\prime }}\left( \mathrm{H}\right) .
\]

\({y}_{\mathrm{H}}\) 的特征在 \(\mathrm{H}\) 的 \(p\) -奇异元上为零。若我们知道 \({y}_{\mathrm{H}}\) 属于 \({\mathrm{P}}_{{\mathrm{A}}^{\prime }}\left( \mathrm{H}\right)\) ，则可推出 \(y\) 属于 \({\mathrm{P}}_{{\mathrm{A}}^{\prime }}\left( \mathrm{G}\right)\) 。因此，我们归约到 \(\mathbf{G}\) 为元胞群的情形。

现在按上述分解 \(G = S \times  P\) 。我们有

\[
{\mathrm{R}}_{{\mathrm{K}}^{\prime }}\left( \mathrm{G}\right)  = {\mathrm{R}}_{{\mathrm{K}}^{\prime }}\left( \mathrm{S}\right)  \otimes  {\mathrm{R}}_{{\mathrm{K}}^{\prime }}\left( \mathrm{P}\right) .
\]

由于 \(\chi\) 在 \(S\) 之外为零，我们可将 \(\chi\) 写成 \(f \otimes  {r}_{\mathrm{P}}\) 的形式，其中 \(f\) 是定义在 \(S\) 上的类函数， \({r}_{p}\) 是 \(P\) 的正则表示的特征。若 \(\rho\) 是 \(S\) 的一特征，则

\[
\left\langle  {f \otimes  {r}_{\mathbf{P}},\rho  \otimes  1}\right\rangle   = \left\langle  {f,\rho }\right\rangle   \cdot  \left\langle  {{r}_{\mathbf{P}},1}\right\rangle   = \langle f,\rho \rangle .
\]

因左边等于 \(\langle \chi ,\rho  \otimes  1\rangle\) ，故为整数；因此对所有 \(\rho\) 有 \(\langle f,\rho \rangle  \in  \mathbf{Z}\) ，这证明了 \(f\) 是 \(\mathbf{S}\) 的虚特征。于是我们可以把 \(y\) 写成

\[
y = {y}_{\mathrm{S}} \otimes  {y}_{\mathrm{P}}
\]

有 \({y}_{\mathrm{S}} \in  {\mathrm{R}}_{{\mathrm{K}}^{\prime }}\left( \mathrm{S}\right)\) ，且 \({y}_{\mathrm{P}}\) 为正则表示的类 \(\mathrm{P}\) 。由于 \({y}_{\mathrm{S}} \in  {\mathrm{P}}_{{\mathrm{A}}^{\prime }}\left( \mathrm{S}\right)\) 与 \({y}_{\mathrm{P}} \in  {\mathrm{P}}_{{\mathrm{A}}^{\prime }}\left( \mathrm{P}\right)\) ，我们确实有 \(y \in  {\mathrm{P}}_{{\mathrm{A}}^{\prime }}\left( \mathrm{G}\right)\) 。

\section*{(iii) 充分性(第二部分)}

保留记号 (ii)，并且额外假设 \(y\) 的特征 \(\chi\) 的取值在 \(\mathbf{K}\) 中。我们必须证明 \(y\) 属于 \({\mathrm{P}}_{\mathrm{A}}\left( \mathrm{G}\right)\) 。由 (ii)，我们至少知道 \(y \in  {\mathbf{P}}_{{\mathbf{A}}^{\prime }}\left( \mathrm{G}\right)\) 。

设 \(r\) 为扩张 \({\mathrm{K}}^{\prime }/\mathrm{K}\) 的次数。每个 \({\mathrm{A}}^{\prime }\left\lbrack  \mathrm{G}\right\rbrack\) -模通过限制定义为一个 A[G]-模，且当原模为投射模时此模也是投射的。因此我们有一个同态

\[
\pi  : {\mathrm{P}}_{{\mathrm{A}}^{\prime }}\left( \mathrm{G}\right)  \rightarrow  {\mathrm{P}}_{\mathrm{A}}\left( \mathrm{G}\right) .
\]

令 \(z = \pi \left( y\right)\) 。则 \(z = r \cdot  y\) 。事实上，只需在 \({\mathbf{R}}_{{\mathbf{K}}^{\prime }}\left( \mathbf{G}\right)\) 中验证此等式，而为此足以证明与 \(z\) 关联的特征 \({\chi }_{z}\) 等于 \(r \cdot  \chi\) 。但我们有

\[
{\chi }_{z} = {\operatorname{Tr}}_{{\mathbf{K}}^{\prime }/\mathbf{K}}\left( \chi \right)
\]

并且因为 \(\chi\) 的取值在 \(\mathbf{K}\) 中，我们得到 \({\chi }_{z} = r \cdot  \chi\) 。

因此 \(y \in  {\mathbf{P}}_{{\mathbf{A}}^{\prime }}\left( \mathrm{G}\right)\) 与 \(r \cdot  y \in  {\mathbf{P}}_{\mathbf{A}}\left( \mathrm{G}\right)\) 。但包含映射 \({\mathbf{P}}_{\mathbf{A}}\left( \mathrm{G}\right)  \rightarrow  {\mathbf{P}}_{{\mathbf{A}}^{\prime }}\left( \mathrm{G}\right)\) 是一个分裂单射，参见 14.6。由于 \(r \cdot  y\) 在 \({\mathbf{P}}_{{\mathbf{A}}^{\prime }}\left( \mathrm{G}\right)\) 中可被 \(r\) 整除，在 \({P}_{A}\left( G\right)\) 中亦然，这意味着 \(y \in  {P}_{A}\left( G\right)\) ，并完成证明。

\section*{17.6 定理 38 的证明}

我们称群 \(G\) 为高度为 \(h\) 的 \(p\) -可解，若其是由 \(h\) 个群逐次扩张所得，这些群要么阶与 \(p\) 互素，要么阶为 \(p\) 的幂。我们要证明，如果 \(\mathrm{K}\) 足够大，则每个单 \(k\left\lbrack  \mathrm{\;G}\right\rbrack\) -模都可以提升为一个 A 自由的 A[G]-模。

我们对 \(h\) 归纳(当 \(h = 0\) 的情形平凡)，并且对高度为 \(h\) 的群按群阶归纳。

设 I 为 \(\mathbf{G}\) 的正规子群，其阶要么与 \(p\) 互素，要么为 \(p\) 的幂，且 \(\mathbf{G}/\mathbf{I}\) 的高度为 \(h - 1\) 。设 \(\mathbf{E}\) 为一个单(因而绝对单) \(k\left\lbrack  \mathrm{\;G}\right\rbrack\) -模。若 I 为 \(p\) -群，则由 \(I\) 不动的 \(\mathrm{E}\) 的所有元素所成子空间 \({\mathrm{E}}^{\mathrm{I}}\) 为 \(\neq  0\) ，因此等于 \(E\) ；于是 \(E\) 为简单的 \(k\left\lbrack  {\mathrm{G}/\mathrm{I}}\right\rbrack\) -模。由归纳可把它提升为一个 A-自由的 \(\mathrm{A}\left\lbrack  {\mathrm{G}/\mathrm{I}}\right\rbrack\) 模，此情形下结论成立。

现在假设 I 的阶与 \(p\) 互素。将 \(\mathrm{E}\) 分解为同型模的直和，即同构单模的直和，作为 \(k\) [I]-模:

\[
\mathrm{E} =  \oplus  {\mathrm{E}}_{\alpha },
\]

其中 \({\mathrm{E}}_{\alpha }\) 为类型为 \({\overline{\mathrm{S}}}_{\alpha }\) 的同型 \(k\left\lbrack  \mathrm{I}\right\rbrack\) -模。

群 \(\mathrm{G}\) 对 \({\mathrm{E}}_{\alpha }\) 作置换；因 \(\mathrm{E}\) 为简单模，它在非零的那些分块上作传递置换。取其中之一记为 \({\mathrm{E}}_{\alpha }\) ，令 \({\mathrm{G}}_{\alpha }\) 为 \(G\) 的子群，由所有满足 \(s\left( {E}_{\alpha }\right)  = {E}_{\alpha }\) 的 \(s \in  G\) 构成。则 \({E}_{\alpha }\) 为一个 \(k\left\lbrack  {G}_{\alpha }\right\rbrack\) -模，且显然 E 为相应的诱导模。若 \({\mathrm{E}}_{\alpha } \neq  \mathrm{E}\) ，则有 \({\mathrm{G}}_{\alpha } \neq  \mathrm{G}\) ，对 \({\mathrm{G}}_{\alpha }\) 应用归纳假设表明 \({\mathrm{E}}_{\alpha }\) 可提升；因此 \(\mathrm{E}\) 也可提升。

现在我们可以假设 \(\mathbf{E}\) 为类型为 \(\overline{\mathbf{S}}\) 的同型 \(k\left\lbrack  \mathrm{I}\right\rbrack\) -模，其中 \(\overline{\mathbf{S}}\) 是一个简单的 \(k\left\lbrack  \mathrm{I}\right\rbrack\) -模。由于 \(\mathrm{I}\) 的阶与 \(p\) 互素，我们可以以本质唯一的方式将 \(\overline{\mathbf{S}}\) 提升为一个 A-自由的 A[I]-模，记为 S，且显然 \(\mathrm{K} \otimes  \mathrm{S}\) 是绝对简单的。由 6.5 中命题 16 的推论 2 得出 dim S 除以 I 的阶；特别地，dim S 与 p 互素。

现在设 \(s \in  \mathrm{G}\) ，并记由 \(x \mapsto  {sx}{s}^{-1}\) 诱发的 I 的自同构为 \({i}_{s}\) 。由于 E 为类型为 \(\bar{S}\) 的同型模，故 \(\bar{S}\) (因而 \(S\) )与其被 \({i}_{s}\) 变换后的模同构。可表述如下:

设 \(\rho  : I \rightarrow  \operatorname{Aut}\left( S\right)\) 为定义 \(S\) 的 I-模结构的同态，令 \({U}_{s}\) 为满足 … 的 \(t \in  \operatorname{Aut}\left( S\right)\) 的集合

\[
{t\rho }\left( x\right) {t}^{-1} = \rho \left( {{sx}{s}^{-1}}\right) \;\text{ for all }x \in  \mathrm{I}.
\]

则 \({\mathrm{U}}_{s}\) 非空。

设 \({\mathrm{G}}_{1}\) 为所有满足 \(s \in  \mathrm{G},t \in  {\mathrm{U}}_{s}\) 的对 \(\left( {s,t}\right)\) 组成的群。映射 \(\left( {s,t}\right)  \mapsto  s\) 是一个满同态 \({\mathrm{G}}_{1} \rightarrow  \mathrm{G}\) ；其核等于 \({\mathrm{U}}_{1}\) ，即 \(A\) 的乘法群 \({A}^{ * }\) 。因此群 \({G}_{1}\) 是 \(\mathbf{G}\) 被 \(\mathbf{A}\) 所中心扩张的群*；它通过同态 \(\left( {s,t}\right)  \mapsto  t\) 在 \(\mathbf{S}\) 上作用。

现在我们把 \({G}_{1}\) 换成一个有限群。设 \(d = \dim S\) 。若 \(s \in  G\) ，则元素 \(\det \left( t\right) ,t \in  {\mathrm{U}}_{s}\) 构成模 \({\mathrm{A}}^{d}\) 的 \({\mathrm{A}}^{ * }\) 的一个陪集。通过扩大 \(\mathrm{K}\) (这是可以的，因为这不改变 \({R}_{K}\left( G\right)\) )，我们可以假定这些陪集都平凡，换言之每个 \({U}_{s}\) 包含行列式为 1 的元素。做到这点后，令 \(C\) 为由所有 \(\det \left( {\rho \left( x\right) }\right) ,x \in  \mathrm{I}\) 组成的 \({A}^{ * }\) 的子群，令 \({\mathrm{G}}_{2}\) 为由所有满足 \(t \in  {\mathrm{U}}_{s}\) 与 \(\det \left( t\right)  \in  \mathbf{C}\) 的 \(\left( {s,t}\right)\) 组成的 \({\mathrm{G}}_{1}\) 的子群。群 \({\mathrm{G}}_{2}\) 映到 \(\mathbf{G}\) ；映射 \({\mathrm{G}}_{2} \rightarrow  \mathrm{G}\) 的核 \(\mathrm{N}\) 与由所有满足 \({a}^{d} \in  \mathrm{C}\) 的 \(a\) 组成的 \({\mathrm{A}}^{ * }\) 的子群同构。由于 \(d\) 与 Card (C) 与 \(p\) 互素，我们得出 \(\mathrm{N}\) 是一个与 \(p\) 互素阶的循环群。

记 \({\rho }_{2} : {\mathrm{G}}_{2} \rightarrow\) Aut(S) 为 \({\mathrm{G}}_{2}\) 的表示 \(\left( {s,t}\right)  \mapsto  t\) 。若借助 \(x \mapsto  \left( {x,\rho \left( x\right) }\right)\) 将 I 识别为 \({\mathrm{G}}_{2}\) 的一个子群，则可见 \({\rho }_{2}\) 限制到 I 等于 \(\rho\) 。于是我们已把 \(\rho\) 扩张，不是扩到 \(\mathrm{G}\) 本身，但至少扩到 \(G\) 的一个中心扩张(按 Schur 的意义我们得到了 \(G\) 的“投影”表示)。注意 \(I\) 在 \({G}_{2}\) 中是正规子群，且 \(I \cap  N = \{ 1\}\) 。

现在回到原来的 \(k\left\lbrack  \mathrm{G}\right\rbrack\) -模 E。设 \(\mathrm{F} = {\operatorname{Hom}}^{1}\left( {\overline{\mathrm{S}},\mathrm{E}}\right)\) 并设 \(u : \overline{\mathrm{S}} \otimes  \mathrm{F} \rightarrow  \mathrm{E}\) 是把 \(a \otimes  t(a \in  \overline{\mathrm{S}}\) 、 \(b \in  \mathrm{F}\) 关联到 E 中元素 \(b\left( a\right)\) 的同态。由 \(\mathrm{E}\) 为类型 \(\overline{\mathrm{S}}\) 的同型体可轻易推得 \(u\) 将 \(\overline{\mathrm{S}} \otimes  \mathrm{F}\) 同构到 \(\mathrm{E}\) 上。

群 \({\mathrm{G}}_{2}\) 通过约化的 \({\rho }_{2}\) 在 \(\overline{\mathrm{S}}\) 上作用；它也通过 \({\mathrm{G}}_{2} \rightarrow  \mathrm{G}\) 在 \(\mathrm{E}\) 上作用；因此它在 \(\mathrm{F}\) 上作用。该同构

\[
u : \overline{\mathrm{S}} \otimes  \mathrm{F} \rightarrow  \mathrm{E}
\]

与 \({\mathrm{G}}_{2}\) 的此作用相容。因此将 \(\mathrm{E}\) 视为一个 \(k\left\lbrack  {\mathrm{G}}_{2}\right\rbrack\) -模时，可被识别为 \(k\left\lbrack  {\mathrm{G}}_{2}\right\rbrack\) -模 \(\overline{\mathbf{S}}\) 与 \(\mathbf{F}\) 的张量积。为了把 E 提升，只需分别提升 \(\overline{\mathbf{S}}\) 与 \(\mathbf{F}\) 并取这些提升的张量积。这样我们就得到一个 A-自由的 A[G]-模 E。因 N 的阶与 \(p\) 互素且在 \(\widetilde{\mathrm{E}}\) 的降阶 \(\mathrm{E}\) 上平凡作用，故 \(\mathrm{N}\) 在 E 上也平凡作用(参见 15.5)，且 E 可被视为一个 A[G]-模——确实，我们已将 E 提升。

因此剩下要证明 F 可被提升( \(\bar{S}\) 的情形已解决)。但 \(\mathrm{F}\) 是一个简单的 \(k\left\lbrack  {\mathrm{G}}_{2}\right\rbrack\) -模(因 \(\mathrm{E}\) 是)，按构造 I 在其上平凡作用。故我们可以将其视为一个简单的 \(k\left\lbrack  \mathrm{H}\right\rbrack\) -模，而 \(\mathrm{H} = {\mathrm{G}}_{2}/\mathrm{I}\) 。

群 \(\mathrm{H}\) 是 \(\mathrm{G}/\mathrm{I}\) 的一个中心扩张(该 \(\mathrm{G}/\mathrm{I}\) 是高度为 \(\leq  h - 1\) 的 \(p\) -可解群)，扩张核为循环群 \(\mathrm{N}\) ，其阶与 \(p\) 互素。若 \(h = 1\) ，则有 \(H = N\) ，且 \(F\) 的提升是显然的(15.5)。若 \(h \geq  2\) ，则 H/N 包含一个满足下列两条的正规子群 M/N:

(a) \(\mathrm{H}/\mathrm{M} = \left( {\mathrm{H}/\mathrm{N}}\right) /\left( {\mathrm{M}/\mathrm{N}}\right)\) 的高度为 \(\leq  h - 2\) 。

(b) \(\mathrm{M}/\mathrm{N}\) 要么是一个 \(p\) -群，要么是阶与 \(p\) 互素的群。

若 \(\mathrm{M}/\mathrm{N}\) 是一个 \(p\) -群，则由于 \(\mathrm{N}\) 的阶与 \(p,\mathrm{M}\) 互素，可写成乘积 \(\mathrm{N} \times  \mathrm{P}\) ，其中 \(\mathrm{P}\) 是一个 \(p\) -群。证明开头的论证说明 P 在 F 上平凡作用，故 F 可视为一个 \(k\left\lbrack  {\mathrm{H}/\mathrm{P}}\right\rbrack\) -模。但显然 \(\mathrm{H}/\mathrm{P}\) 的高度为 \(\leq  h - 1\) ，所以可由归纳法提升 \(\mathrm{F}\) 。剩下的情形是 \(M/N\) 的阶与 \(p\) 互素。此时 \(M\) 的阶亦与 \(p\) 互素，且由于 \(H/M\) 的高度为 \(\leq  h - 2\) ， \(H\) 的高度为 \(\leq  h - 1\) ，再次可用归纳法。证毕。

\section*{第18章 模拟特征}

我们所讨论的结果大多归功于 R. Brauer。他用稍有不同的语言来表述——模拟特征的语言，下面我们将加以描述。

为简便起见，我们假定 \(\mathrm{K}\) 足够大。

\section*{18.1 表示的模拟特征}

设 \({\mathrm{G}}_{\text{ reg }}\) 为 \(\mathrm{G}\) 中 \(p\) -正则元素构成的集合，设 \({m}^{\prime }\) 为 \({G}_{\text{ reg }}\) 中元素阶的最小公倍数。根据假设， \(K\) 包含 \({m}^{\prime }\) 次单位根组成的群 \({\mu }_{K}\) ；此外，由于 \({m}^{\prime }\) 与 \(p\) 互质，模 \(\mathfrak{m}\) 约化是从 \({\mu }_{\mathrm{K}}\) 到剩余域 \(k\) 的 \({m}^{\prime }\) 次单位根组成的群 \({\mu }_{k}\) 的同构。对于 \(\lambda  \in  {\mu }_{k}\) ，我们用 \(\widetilde{\lambda }\) 表示 \({\mu }_{\mathrm{K}}\) 中模 \(\mathfrak{m}\) 约化后为 \(\lambda\) 的元素。

令 \(\mathrm{E}\) 为维数为 \(n\) 的 \(k\left\lbrack  \mathrm{G}\right\rbrack\) -模，令 \(s \in  {\mathrm{G}}_{\text{ reg }}\) ，并令 \({s}_{\mathrm{E}}\) 为由 \(s\) 定义的 \(\mathrm{E}\) 上的自同态。因 \(s\) 的阶与 \(p,{s}_{\mathrm{E}}\) 互素，故其可对角化，且其特征值 \(\left( {{\lambda }_{1},\ldots ,{\lambda }_{n}}\right)\) 属于 \({\mu }_{k}\) 。设

\[
{\phi }_{\mathrm{E}}\left( s\right)  = \mathop{\sum }\limits_{{i = 1}}^{{i = n}}{\widetilde{\lambda }}_{i}
\]

如此定义的函数 \({\phi }_{\mathrm{E}} : {\mathrm{G}}_{\text{ reg }} \rightarrow  \mathrm{A}\) 称为 E 的模特征(或 Brauer 特征)。下列性质显然:

(i) 我们有 \({\phi }_{\mathrm{E}}\left( 1\right)  = n = \dim \mathrm{E}\) 。

(ii) \({\phi }_{\mathrm{E}}\) 是作用在 \({\mathrm{G}}_{\text{ reg }}\) 上的类函数，即，

\[
{\phi }_{\mathrm{E}}\left( {{ts}{t}^{-1}}\right)  = {\phi }_{\mathrm{E}}\left( s\right) \;\text{ if }s \in  {\mathrm{G}}_{\text{ reg }}\text{ and }t \in  \mathrm{G}.
\]

(iii) 若 \(0 \rightarrow  \mathrm{E} \rightarrow  {\mathrm{E}}^{\prime } \rightarrow  {\mathrm{E}}^{\prime \prime } \rightarrow  0\) 为一列确切的 \(k\left\lbrack  \mathrm{G}\right\rbrack\) -模序列，则有

\[
{\phi }_{{\mathrm{E}}^{\prime }} = {\phi }_{\mathrm{E}} + {\phi }_{{\mathrm{E}}^{\prime \prime }}.
\]

(iv) 我们有

\[
{\phi }_{{\mathrm{E}}_{1} \otimes  {\mathrm{E}}_{2}} = {\phi }_{{\mathrm{E}}_{1}} \cdot  {\phi }_{{\mathrm{E}}_{2}}.
\]

(v) 若 \(t \in  \mathrm{G}\) 具有 \({p}^{\prime }\) -分量 \(s \in  {\mathrm{G}}_{\text{ reg }}\) ，则端同态 \({t}_{\mathrm{E}}\) 在 \(\mathrm{E}\) 上的迹是 \({\phi }_{\mathrm{E}}\left( s\right)\) 关于模 \(\mathfrak{m}\) 的约化:我们有

\[
\operatorname{Tr}\left( {t}_{\mathrm{E}}\right)  = \overline{{\phi }_{\mathrm{E}}\left( s\right) },
\]

其中上横线表示模 m 约化。(可由观察 \({\left( {t}^{-1}s\right) }_{\mathrm{E}}\) 的特征值为 \({p}^{n}\) 次单位根，因 \(k\) 的特征为 \(p\) ，故这些根等于 1。由此可知 \({t}_{\mathrm{E}}\) 的特征值与 \({s}_{\mathrm{E}}\) 的相同，从而得出所需公式。)

(vi) 设 \(F\) 为一 \(K\left\lbrack  G\right\rbrack\) -模，具有特征 \(\chi\) ，令 \({E}_{1}\) 为在 \(G\) 下不变的 \(F\) 的格子，且令 \(E = {E}_{1}/m{E}_{1}\) 为其对 \(\mathfrak{m}\) 的化简。则 \({\phi }_{E}\) 为 \(\chi\) 限制到 \({\mathrm{G}}_{\text{ reg }}\) 的结果。(只需在 \(\mathrm{G}\) 为与 \(p\) 互素阶的循环群时验证。此外，定理 32 显示 \({\phi }_{\mathrm{E}}\) 不依赖于所选的稳定格子 \({E}_{1}\) 。因此可以约化到 \({E}_{1}\) 由 \(G\) 的特征向量生成的情形，此时结论显然。)

(vii) 若 \(F\) 为射影的 \(k\left\lbrack  G\right\rbrack\) -模，且 \(\widetilde{F}\) 为其化简为 \(F\) 的射影 \(A\left\lbrack  G\right\rbrack\) -模，我们将把 \(\widetilde{F}\) 的特征(即 \(K\left\lbrack  G\right\rbrack\) -模 \(\mathrm{K} \otimes  \widetilde{\mathrm{F}}\) 的特征)记作 \({\Phi }_{\mathrm{F}}\) 。若 \(\mathrm{E}\) 为任一 \(k\left\lbrack  \mathrm{G}\right\rbrack\) -模，我们知道 \(\mathrm{E} \otimes  \mathrm{F}\) 为射影的 \(k\left\lbrack  \mathrm{\;G}\right\rbrack\) -模，因此 \({\Phi }_{\mathrm{E} \otimes  \mathrm{F}}\) 有意义。我们有

\[
{\Phi }_{\mathrm{E} \otimes  \mathrm{F}}\left( s\right)  = \left\{  \begin{array}{ll} {\phi }_{\mathrm{E}}\left( s\right) {\Phi }_{\mathrm{F}}\left( s\right) & \text{ if }s \in  {\mathrm{G}}_{\text{ reg }} \\  0 & \text{ otherwise, } \end{array}\right.
\]

一个更简洁地写作 \({\Phi }_{\mathrm{E} \otimes  \mathrm{F}} = {\phi }_{\mathrm{E}} \cdot  {\Phi }_{\mathrm{F}}\) 的公式，尽管 \({\phi }_{\mathrm{E}}\) 在 \({\mathrm{G}}_{\text{ reg }}\) 之外未定义。(我们知道若 \({\Phi }_{\mathrm{E} \otimes  \mathrm{F}}\left( s\right)  = 0\) 则 \(s \notin  {\mathrm{G}}_{\text{ reg }}\) ，参见定理 36。并且由 (vi) \({\Phi }_{\mathrm{E} \otimes  \mathrm{F}}\) 限制到 \({\mathrm{G}}_{\text{ reg }}\) 等于 \(\mathrm{E} \otimes  \mathrm{F}\) 的模特征，而按 (iv) 该模特征即 \({\phi }_{\mathrm{E}} \cdot  {\Phi }_{\mathrm{F}}\) 。)

(viii) 在与 (vii) 相同的假设下，我们有

\[
\langle \mathrm{E},\mathrm{F}{\rangle }_{k} = \frac{1}{g}\mathop{\sum }\limits_{{s \in  {\mathbf{G}}_{\mathrm{{reg}}}}}{\Phi }_{\mathrm{F}}\left( {s}^{-1}\right) {\phi }_{\mathrm{E}}\left( s\right)  = \left\langle  {{\phi }_{\mathrm{E}},{\Phi }_{\mathrm{F}}}\right\rangle  ,
\]

其中 \(g = \operatorname{Card}\left( \mathrm{G}\right)\) 。(根据定义， \(\langle \mathrm{E},\mathrm{F}{\rangle }_{k}\) 是 \(\mathrm{H} = \operatorname{Hom}\left( {\mathrm{F},\mathrm{E}}\right)\) 中被 \(\mathrm{G}\) 固定的最大子空间 \({\mathrm{H}}^{\mathrm{G}}\) 的维数。然而， \(\mathrm{H}\) 是一个射影 \(k\left\lbrack  \mathrm{\;G}\right\rbrack\) -模，所以如果 \(\widetilde{\mathrm{H}}\) 是对应的射影 \(\mathrm{A}\left\lbrack  \mathrm{G}\right\rbrack\) -模，我们很容易看出 \({\dim }_{k}{\mathrm{H}}^{\mathrm{G}} = {\operatorname{rank}}_{\mathrm{A}}{\widetilde{\mathrm{H}}}^{\mathrm{G}}\) 。如果 \({\Phi }_{\mathrm{H}}\) 是 \(\mathrm{K} \otimes  \widetilde{\mathrm{H}}\) 的特征标，我们有

\[
\langle \mathrm{E},\mathrm{F}{\rangle }_{k} = \left\langle  {1,{\Phi }_{\mathrm{H}}}\right\rangle   = \frac{1}{g}\mathop{\sum }\limits_{{s \in  \mathrm{G}}}{\Phi }_{\mathrm{H}}\left( s\right) .
\]

但 \(\mathrm{H}\) 与 \(\mathrm{E}\) 与 \(\mathrm{F}\) 的对偶的张量积同构。由 (vii) 对于 \(s \in  {\mathrm{G}}_{\text{ reg }}\) 我们有 \({\Phi }_{\mathrm{H}}\left( s\right)  = {\Phi }_{\mathrm{F}}\left( {s}^{-1}\right) {\phi }_{\mathrm{E}}\left( s\right)\) ，否则有 \({\Phi }_{\mathrm{H}}\left( s\right)  = 0\) 。由此得证。)

我们注意到特殊情况: \(\mathrm{E}\) 为平凡表示时:

(ix) 由 \(G\) 不变的元素构成的子空间 \({F}^{G}\) 的维数为

\[
\left\langle  {1,{\Phi }_{\mathrm{F}}}\right\rangle   = \frac{1}{g}\mathop{\sum }\limits_{{s \in  {\mathrm{G}}_{\text{ reg }}}}{\Phi }_{\mathrm{F}}\left( s\right) .
\]

注:性质 (iii) 使我们能够为任一 \({\mathrm{R}}_{k}\left( \mathrm{G}\right)\) 的元素 \(x\) 定义虚拟模特征 \({\phi }_{x}\) 。由 (vi)，若 \(x = d\left( y\right)\) 伴随 \(y \in  {\mathrm{R}}_{\mathrm{G}}\left( \mathrm{G}\right)\) ，则 \({\phi }_{x}\) 只是 \(y\) 的虚拟特征 \({\chi }_{y}\) 限制到 \({\mathrm{G}}_{\text{ reg }}\) 。

可以对定义在域为 \(k\) 的任一线性代数群 \(\mathrm{G}\) (此处为简便假设 \(k\) 代数闭合)给出类似定义。集合 \({\mathrm{G}}_{\text{ reg }}\) 定义为 \(\mathrm{G}\) 的半单元素的集合。若 \(\mathrm{E}\) 是 \(G\) 的线性表示，且若 \(s \in  {G}_{\text{ reg }}\) ，则将 \({\phi }_{E}\left( s\right)\) 定义为 \({s}_{E}\) 的特征值的乘法代表元之和；由此定义的模特征 \({\phi }_{\mathrm{E}}\) 是在 \({\mathrm{G}}_{\text{ reg }}\) 上取值于 \(\mathrm{A}\) 的类函数。

\section*{18.2 模特征的独立性}

回想 \({S}_{k}\) 表示简单 \(k\left\lbrack  \mathrm{G}\right\rbrack\) -模同构类的集合。对应于 \({\mathrm{S}}_{k}\) 的元素 \(\mathrm{E}\) 的各个 \({\phi }_{\mathrm{E}}\) 被称为群 \(\mathrm{G}\) 的不可约模特征。

定理 42。(R. Brauer)。不可约模特征 \({\phi }_{\mathrm{E}}\left( {\mathrm{E} \in  {\mathrm{S}}_{k}}\right)\) 构成以 \(\mathrm{K}\) 为系数、取值于 K 的在 \({\mathrm{G}}_{\text{ reg }}\) 上的类函数向量空间的一组基。

这可以用以下等价的形式表述:

定理 42’。映射 \(x \mapsto  {\phi }_{x}\) 可扩展为从 \(\mathrm{K} \otimes  {\mathrm{R}}_{k}\left( \mathrm{G}\right)\) 到取值于 \(\mathrm{K}\) 的在 \({\mathrm{G}}_{\text{ reg }}\) 上的类函数代数的同构。

这些定理立即给出:

推论 1。设 \(\mathrm{F}\) 与 \({\mathrm{F}}^{\prime }\) 为具有相同模特征的两个 \(k\left\lbrack  \mathrm{\;G}\right\rbrack\) -模。则 \(\left\lbrack  \mathrm{F}\right\rbrack   = \left\lbrack  {\mathrm{F}}^{\prime }\right\rbrack\) 属于 \({\mathrm{R}}_{k}\left( \mathrm{G}\right)\) ；若 \(\mathrm{F}\) 与 \({\mathrm{F}}^{\prime }\) 半单，则它们同构。

推论 2。同态 \(d : {\mathrm{R}}_{\mathrm{K}}\left( \mathrm{G}\right)  \rightarrow  {\mathrm{R}}_{k}\left( \mathrm{G}\right)\) 的核由那些其虚特征 \({\chi }_{x}\) 在 \({\mathrm{G}}_{\text{ reg }}\) 上为零的元素 \(x\) 组成。

(由于 \(d\) 是满射，这给出了将 \({\mathbf{R}}_{k}\left( \mathbf{G}\right)\) 明确描述为 \({R}_{K}\left( G\right)\) 的一个商的方式。)

推论 3。简单 \(k\left\lbrack  \mathrm{G}\right\rbrack\) -模的同类数等于 \(\mathbf{G}\) 的 p-正规共轭类的数目。

定理 42 的证明。

(a)我们首先证明 \({\phi }_{\mathbf{E}}\left( {\mathbf{E} \in  {\mathbf{S}}_{k}}\right)\) 在 \(\mathbf{K}\) 上线性无关。确实，假设存在关系 \(\sum {a}_{\mathrm{E}}{\phi }_{\mathrm{E}} = 0\) ，其中 \({a}_{\mathrm{E}} \in  \mathrm{K}\) ，不全为零。将这些 \({a}_{\mathrm{E}}\) 乘以 \(\mathrm{K}\) 的某个元素，可假定它们都属于环 A，且至少有一个不属于 \(\mathfrak{m}\) 。对 \(\mathfrak{m}\) 取模约化后，我们得到

\[
\mathop{\sum }\limits_{{\mathrm{E} \in  {\mathrm{S}}_{k}}}{\bar{a}}_{\mathrm{E}}\overline{{\phi }_{\mathrm{E}}\left( s\right) } = 0\;\text{ for all }s \in  {\mathrm{G}}_{\mathrm{{reg}}},
\]

其中上横线表示对 \(\mathfrak{m}\) 的模约化，且其中一个 \({\bar{a}}_{\mathrm{E}}\) 非零。由前一节的公式 (v)，我们得到

\[
\sum {\bar{a}}_{\mathrm{E}}\operatorname{Tr}\left( {t}_{\mathrm{E}}\right)  = 0\;\text{ for all }t \in  \mathrm{G},
\]

因此对所有 \(t \in  k\left\lbrack  \mathrm{G}\right\rbrack\) 亦然。然而，由于 \(\mathrm{K}\) 足够大，模 \(E\) 是绝对单的，于是由密度定理([8],§4, no. 2)，同态 \(k\left\lbrack  \mathrm{G}\right\rbrack   \rightarrow  { \oplus  }_{\mathrm{E} \in  {\mathrm{S}}_{k}}{\operatorname{End}}_{k}\left( \mathrm{E}\right)\) 是满的。现在取 \(\mathrm{E} \in  {\mathrm{S}}_{k}\) 使得 \({\bar{a}}_{\mathrm{E}} \neq  0\) ，取 \(u \in  {\operatorname{End}}_{k}\left( \mathrm{E}\right)\) 具有迹 1(例如投影到一条直线)，并取 \(t\) 为 \(k\left\lbrack  \mathrm{G}\right\rbrack\) 中的一个元素，其在 \({\operatorname{End}}_{k}\left( \mathrm{E}\right)\) 的像为 \(u\) ，在为 \({\mathrm{E}}^{\prime } \neq  \mathrm{E}\) 的 \({\operatorname{End}}_{k}\left( {\mathrm{E}}^{\prime }\right)\) 中为 0。于是我们发现 \({\bar{a}}_{\mathrm{E}} \cdot  1 = 0\) ，矛盾。

这部分证明同样适用于线性代数群。

(b) 我们需证明 \({\phi }_{\mathrm{E}}\) 生成了 \({G}_{\text{ reg }}\) 上类函数的向量空间。设 \(f\) 为此类函数，并将其扩展为群 G 上的类函数 \({f}^{\prime }\) 。我们知道 \({f}^{\prime }\) 可以写成 \(\sum {\lambda }_{i}{\chi }_{i}\) 的形式，且 \({\lambda }_{i} \in  \mathrm{K}\) 与 \({\chi }_{i} \in  {\mathrm{R}}_{\mathrm{K}}\left( \mathrm{G}\right)\) 。因此 \(f = \sum {\lambda }_{i}d\left( {\chi }_{i}\right)\) ，其中 \(d\left( {\chi }_{i}\right)\) 是 \({\chi }_{i}\) 限制到 \({\mathrm{G}}_{\text{ reg }}\) 的结果。由于每个 \(d\left( {\chi }_{i}\right)\) 都是 \({\phi }_{\mathrm{E}}\) 的线性组合，因此得出所需结论。

问题(Brauer)。设 \(\mathrm{E} \in  {\mathrm{S}}_{k}\) ，令 \({p}^{e}\) 为能整除 dim E 的 \(p\) 的最大幂。是否有 \({p}^{e}\) 整除 G 的阶？

\section*{练习}

18.1.(在本练习中我们不假设 \(G\) 有限或 \(k\) 特征为 \(\neq  0\) 。)设 \(E\) 与 \({E}^{\prime }\) 为半单 \(k\left\lbrack  G\right\rbrack\) -模。假定对每个 \(s \in  \mathrm{G}\) ，多项式 \(\det \left( {1 + {s}_{\mathrm{E}}\mathrm{T}}\right)\) 与 \(\det \left( {1 + {s}_{{\mathrm{E}}^{\prime }}\mathrm{T}}\right)\) 相等。证明 \(\mathrm{E}\) 与 \({\mathrm{E}}^{\prime }\) 同构。[约化到 \(k\) 代数闭的情形，并按定理 42 证明的 (a) 部分的论法进行。] 作为推论，证明若 \(\mathrm{E}\) 为半单且所有 \({s}_{\mathrm{E}}\) 都是单一的，则 \(\mathrm{G}\) 在 E 上做平凡作用(Kolchin 定理)。

18.2. 设 \(\mathrm{H}\) 为 \(\mathrm{G}\) 的一个子群，设 \(\mathrm{F}\) 为一个 \(k\left\lbrack  \mathrm{H}\right\rbrack\) -模，且设 \(\mathrm{E} = {\operatorname{Ind}}_{\mathrm{H}}^{\mathrm{G}}\mathrm{F}\) 。证明 \(\mathrm{E}\) 的模特征 \({\phi }_{\mathrm{E}}\) 由与特征零情形相同的公式由 \({\phi }_{\mathrm{F}}\) 得到。

18.3. 环 \({R}_{k}\left( G\right)\) 的谱是什么？

18.4. 证明不可约模特征构成以 A 为系数的类函数模在 \({G}_{\text{ reg }}\) 上的基。［使用 10.3 的引理 8 证明每个取值于 A 的 \({G}_{\text{ reg }}\) 上的类函数可扩张为取值于 A 且属于 \(A \otimes  {R}_{K}\left( G\right)\) 的 \(G\) 上的类函数。］

\section*{18.3 等价表述}

我们刚看到 \(x \mapsto  {\phi }_{x}\) 将 \(K \otimes  {R}_{K}\left( G\right)\) 同构到 \({G}_{\text{ reg }}\) 上的类函数空间。另一方面，映射 \(\mathrm{K} \otimes  e : \mathrm{K} \otimes  {\mathrm{P}}_{k}\left( \mathrm{G}\right)  \rightarrow  \mathrm{K} \otimes  {\mathrm{R}}_{\mathrm{K}}\left( \mathrm{G}\right)\) 将 \(\mathrm{K} \otimes  {\mathrm{P}}_{k}\left( \mathrm{G}\right)\) 同一化为在 \(G\) 上且在 \({G}_{\text{ reg }}\) 外为零的类函数的向量空间(例如通过比较两空间的维数可检验)。与 K 张量后，cde 三角形变为:

\begin{center}
\includegraphics[max width=0.7\textwidth]{images/bo_d6kn08k601uc73e1e9r0_156_371_896_790_253_0.jpg}
\end{center}
\hspace*{3em} 

在 \(\mathbf{G}\) 上的类函数，

映射 K \(\otimes  c,\) K \(\otimes  d,\) K \(\otimes  e\) 为显然的映射:限制、限制、包含。注意 \(\mathrm{K} \otimes  c\) 是一个同构，这与定理 35 的推论 1 相符。

矩阵 \(C\) 与 \(D\) 可按如下方式解释:若 \(F \in  {S}_{K}\) ，记 \({\chi }_{F}\) 为 \(F\) 的特征；若 \(E \in  {S}_{k}\) ，记 \({\phi }_{E}\) 为 \(\mathrm{E}\) 的模特征，且记 \({\Phi }_{\mathrm{E}}\) 为 E 的投影包络的特征。则

\[
{\chi }_{\mathrm{F}} = \mathop{\sum }\limits_{{\mathrm{E} \in  {\mathrm{S}}_{k}}}{\mathrm{D}}_{\mathrm{{EF}}}{\phi }_{\mathrm{E}}\;\text{ on }{\mathrm{G}}_{\text{ reg }}
\]

\[
{\Phi }_{\mathrm{E}} = \mathop{\sum }\limits_{{\mathrm{F} \in  {\mathrm{S}}_{\mathrm{K}}}}{\mathrm{D}}_{\mathrm{{EF}}}{\chi }_{\mathrm{F}}\;\text{ on }\mathrm{G}
\]

\[
{\Phi }_{\mathrm{E}} = \mathop{\sum }\limits_{{{\mathrm{E}}^{\prime } \in  {\mathrm{S}}_{k}}}{\mathrm{C}}_{{\mathrm{E}}^{\prime }\mathrm{E}}{\phi }_{{\mathrm{E}}^{\prime }}\;\text{ on }{\mathrm{G}}_{\text{ reg }}
\]

且我们有正交关系

\[
\left\langle  {{\Phi }_{\mathrm{E}},{\phi }_{{\mathrm{E}}^{\prime }}}\right\rangle   = {\delta }_{{\mathrm{{EE}}}^{\prime }},\;\text{ where }\left\langle  {{\Phi }_{\mathrm{E}},{\phi }_{{\mathrm{E}}^{\prime }}}\right\rangle   = \frac{1}{g}\mathop{\sum }\limits_{{s \in  {\mathrm{G}}_{\mathrm{{reg}}}}}{\Phi }_{\mathrm{E}}\left( {s}^{-1}\right) {\phi }_{{\mathrm{E}}^{\prime }}\left( s\right) .
\]

我们还提到定理 35 的以下版本:

定理 35’. 设 \({p}^{n}\) 为除尽 G 阶的最大的 \(p\) 的幂。若 \(\phi\) 是 \(\mathbf{G}\) 的一个模特征，且若通过公式定义 \(\Phi\) ，

\[
\Phi \left( s\right)  = {p}^{n}\phi \left( s\right) \;\text{ if }s \in  {\mathbf{G}}_{\text{ reg }}
\]

\[
\Phi \left( s\right)  = 0\;\text{ if }s \notin  {\mathrm{G}}_{\text{ reg }}
\]

则 \(\Phi\) 是 \(\mathbf{G}\) 的一个虚特征。

留给读者将此类命题作进一步等价表述的任务。

练习

18.5. 若 \(s \in  {\mathrm{G}}_{\text{ reg }}\) ，记 \({p}^{z\left( s\right) }\) 为在 \(\mathbf{G}\) 中对 \(s\) 的中心化子的一个 \(p\) -Sylow 子群的阶。

(a) 设 \(\Phi\) 是在 \(G\) 上取值于 \(K\) 的类函数。证明 \(\Phi  \in  A \otimes  {P}_{k}\left( G\right)\) 当且仅当 \(\Phi\) 在 \({G}_{\text{ reg }}\) 之外为 0 且 对每个 \(s \in  {G}_{\text{ reg }}\) 有 \(\Phi \left( s\right)  \in  {p}^{z\left( s\right) }A\) (使用习题 18.4 以及正交关系 \(\left\langle  {{\Phi }_{\mathbf{E}},{\phi }_{{\mathbf{E}}^{\prime }}}\right\rangle   = {\delta }_{\left. \mathbf{E}{\mathbf{E}}^{\prime }\right) }\) )。

(b) 用 (a) 证明

\[
\operatorname{Coker}\left( c\right)  \simeq  \prod \mathbf{Z}/{p}^{z\left( s\right) }\mathbf{Z}\;\text{ and }\;\det \left( \mathrm{C}\right)  = {p}^{\sum z\left( s\right) },
\]

其中 \(s\) 运行于 G 的一套 \(p\) -规则类的代表上。

18.6. 设 \(G\) 为 \(p\) -可解(参见 16.3)。若 \(F \in  {S}_{K}\) ，令 \({\phi }_{F}\) 表示 \({\chi }_{F}\) 限制到 \({G}_{\text{ reg }}\) 。证明函数 \(\phi\) 在 \({G}_{\text{ reg }}\) 上当且仅当满足下列两条时，才是一个简单 \(k\left\lbrack  \mathrm{\;G}\right\rbrack\) -模的模特征:

(a) 存在 \(F \in  {S}_{K}\) 使得 \(\phi  = {\phi }_{F}\) 。

(b) 若 \({\left( {n}_{\mathrm{F}}\right) }_{\mathrm{F} \in  {\mathrm{S}}_{\mathrm{K}}}\) 是一族整数 \(\geq  0\) 且 \(\phi  = \sum {n}_{\mathrm{F}}{\phi }_{\mathrm{F}}\) ，则其中一个 \({n}_{\mathrm{F}}\) 等于 1，其余均为 0。[使用 Fong–Swan 定理。]

\section*{18.4 对 \(d\) 的截面}

同态 \(d : {\mathrm{R}}_{\mathrm{K}}\left( \mathrm{G}\right)  \rightarrow  {\mathrm{R}}_{k}\left( \mathrm{G}\right)\) 是满射(定理 33)。我们现在将描述 \(d\) 的一个截面，即一个同态

\[
\sigma  : {\mathrm{R}}_{k}\left( \mathrm{G}\right)  \rightarrow  {\mathrm{R}}_{\mathrm{K}}\left( \mathrm{G}\right)
\]

使得 \(d \circ  \sigma  = 1\) 。

对于 \(s \in  \mathrm{G}\) ，令 \({s}^{\prime }\) 表示 \(s\) 的 \({p}^{\prime }\) 分量。若 \(f\) 是在 \({\mathrm{G}}_{\text{ reg }}\) 上的类函数，按下式在 \(\mathrm{G}\) 上定义类函数 \({f}^{\prime }\) :

\[
{f}^{\prime }\left( s\right)  = f\left( {s}^{\prime }\right) .
\]

\section*{定理 43。}

(i) 若 \(f\) 是 \(\mathbf{G},{f}^{\prime }\) 的模特征，则它是 \(\mathbf{G}\) 的虚特征。

(ii) 映射 \(f \mapsto  {f}^{\prime }\) 将定义一个从 \({\mathrm{R}}_{k}\left( \mathrm{G}\right)\) 到 \({\mathrm{R}}_{\mathrm{K}}\left( \mathrm{G}\right)\) 的同态，该同态是 \(d\) 的一个截面。

为了证明 \({f}^{\prime }\) 是 \(G\) 的虚特征(即属于 \({R}_{K}\left( G\right)\) )，只需证明对 \(G\) 的每个初等子群 \(H\) ， \({f}^{\prime }\) 限制到 \(H\) 属于 \({R}_{K}\left( H\right)\) (参见 11.1，第 21 题)。因此我们化简为 \(G\) 为初等群的情形，于是可分解为 \(G = S \times  P\) ，其中 \(S\) 的阶与 \(p\) 互素，且 \(P\) 是一个 \(p\) -群。此外，我们可以假定 \(f\) 是简单 \(k\left\lbrack  \mathrm{G}\right\rbrack\) -模 E 的模特征。按 15.7 的讨论，E 甚至是一个简单的 \(k\left\lbrack  \mathrm{\;S}\right\rbrack\) -模，并且我们可以将其提升(lift)为一个 \(\mathrm{K}\left\lbrack  \mathrm{S}\right\rbrack\) -简单模，在其上 \(\mathrm{P}\) 起平凡作用。该模的特征显然是 \({f}^{\prime }\) ，这证明了 (i)。

断言 (ii) 可由 (i) 推出，因为将 \({f}^{\prime }\) 限制到 \({\mathrm{G}}_{\text{ reg }}\) 等于 \(f\) 。

练习

18.7. 设 \(m\) 为 G 中元素阶的最小公倍数。把 \(m\) 写成 \({p}^{n}{m}^{\prime }\) 的形式，且 \(\left( {p,{m}^{\prime }}\right)  = 1\) (参见 18.1)，并选择一个整数 \(q\) 使得 \(q \equiv  0 \; \left( {{\;\operatorname{mod}\;.}{p}^{n}}\right)\) 与 \(q \equiv  1\left( {{\;\operatorname{mod}\;.}{m}^{\prime }}\right)\) 。

(a) 证明，如果 \(s \in  \mathrm{G}\) ，则 \(s\) 的 \({p}^{\prime }\) 分量 \({s}^{\prime }\) 等于 \({s}^{q}\) 。

(b) 设 \(f\) 是 \(G\) 的一个模特征，且令 \(\phi\) 为 \({R}_{K}\left( G\right)\) 中的一个元素，其限制到 \({G}_{\text{ reg }}\) 为 \(f\) (此类元素由第 33 题存在)。按第 43 题的记法，证明 \({f}^{\prime } = {\Psi }^{q}\phi\) ，其中 \({\Psi }^{q}\) 是在练习 9.3 中定义的算子。由此再推导出 \({f}^{\prime }\) 属于 \({\mathrm{R}}_{\mathrm{K}}\left( \mathrm{G}\right)\) 的另一证明［注意 \({\mathrm{R}}_{\mathrm{K}}\left( \mathrm{G}\right)\) 在 \({\Psi }^{q}\) 下不变］。

18.8. 在不假定 \(K\) 足够大的情况下证明第 43 题［使用前一道练习的方法］。

18.5 例:对称群 \({\mathfrak{S}}_{4}\) 的模特征

群 \({\mathfrak{S}}_{4}\) 是对 \(\{ a,b,c,d\}\) 的置换群。回顾其特征表(参见 5.8):

\begin{center}
\adjustbox{max width=\textwidth}{
\begin{tabular}{|c|c|c|c|c|c|}
\hline
\phantom{X} & 1 & (ab) & (ab)(cd) & (abc) & (abcd) \\
\cline{1-6}
\({\chi }_{1}\) & 1 & 1 & 1 & 1 & 1 \\
\cline{1-6}
\({\chi }_{2}\) & 1 & -1 & 1 & 1 & -1 \\
\cline{1-6}
\({\chi }_{3}\) & 2 & 0 & 2 & -1 & 0 \\
\cline{1-6}
\({\chi }_{4}\) & 3 & 1 & -1 & 0 & -1 \\
\cline{1-6}
\({\chi }_{5}\) & 3 & -1 & -1 & 0 & 1 \\
\cline{1-6}
\hline
\end{tabular}
}
\end{center}

我们将确定其在特征 \(p\) 下的不可约模表示的特征。我们可以假定 \(p\) 整除 \(G\) 的阶，即 \(p = 2\) 或 \(p = 3\) 。

(a) 情况 \(p = 2\)

有两个 \(p\) -正则类:单位元所在的类和 (abc) 所在的类。由定理 42 的推论 3，在特征 2 下(同构意义下)有两个不可约表示。唯一的 1 维表示是恒等表示，其模特征为 \({\phi }_{1} = 1\) 。另一方面， \({\mathfrak{S}}_{4}\) 的 2 维不可约表示模 2 约化后给出表示 \({\rho }_{2}\) ，其模特征 \({\phi }_{2}\) 在元素 (abc) 处取值 -1。因此， \({\rho }_{2}\) 不能是两个 1 维表示的扩张(否则我们会有 \({\phi }_{2} = 2{\phi }_{1} = 2\) )，因而它是不可约的。 \({\mathfrak{S}}_{4}\) 的不可约模特征因此为 \({\phi }_{1}\) 和 \({\phi }_{2}\) :

\begin{center}
\adjustbox{max width=\textwidth}{
\begin{tabular}{|c|c|c|}
\hline
\phantom{X} & 1 & (abc) \\
\cline{1-3}
\({\phi }_{1}\) & 1 & 1 \\
\cline{1-3}
\({\phi }_{2}\) & 2 & -1 \\
\cline{1-3}
\hline
\end{tabular}
}
\end{center}

分解矩阵 \(\mathbf{D}\) 是通过将特征 \({\chi }_{1},\ldots ,{\chi }_{5}\) 限制到 \({\mathrm{G}}_{\text{ reg }}\) 后，以 \({\phi }_{1}\) 和 \({\phi }_{2}\) 为基表示得到的。我们得到

\[
{\chi }_{1} = {\phi }_{1}
\]

\[
{\chi }_{2} = {\phi }_{1}
\]

\[
\text{ on }{\mathrm{G}}_{\text{ reg }}
\]

\[
{\chi }_{3} = {\phi }_{2}
\]

\[
\text{ on }{\mathrm{G}}_{\mathrm{{reg}}}
\]

\[
{\chi }_{4} = {\phi }_{1} + {\phi }_{2}
\]

\[
\text{ on }{\mathrm{G}}_{\text{ reg }}
\]

\[
{\chi }_{5} = {\phi }_{1} + {\phi }_{2}
\]

因此

\[
\mathrm{D} = \left( \begin{array}{lllll} 1 & 1 & 0 & 1 & 1 \\  0 & 0 & 1 & 1 & 1 \end{array}\right)
\]

对应于 \({\phi }_{1}\) 和 \({\phi }_{2}\) 的投射不可约模的特征 \({\Phi }_{1}\) 和 \({\Phi }_{2}\) 是由 D 的转置矩阵得到的:

\[
{\Phi }_{1} = {\chi }_{1} + {\chi }_{2} + {\chi }_{4} + {\chi }_{5}
\]

\[
{\Phi }_{2} = {\chi }_{3} + {\chi }_{4} + {\chi }_{5}
\]

对应的表示度为 8。Cartan 矩阵 \(C = D \cdot  {}^{t}D\) 是行列式为 8 的矩阵 \(\left( \begin{array}{ll} 4 & 2 \\  2 & 3 \end{array}\right)\) 。它表示了 \({\Phi }_{1}\) 和 \({\Phi }_{2}\) 在 \({G}_{\text{ reg }}\) 上的以下分解:

\[
{\Phi }_{1} = 4{\phi }_{1} + 2{\phi }_{2}\;\text{ on }{\mathrm{G}}_{\text{ reg }}
\]

\[
{\Phi }_{2} = 2{\phi }_{1} + 3{\phi }_{2}\;\text{ on }{\mathrm{G}}_{\text{ reg }}.
\]

(b) 情况 \(p = 3\)

有四个 \(p\) -正则类: \(1,\left( {ab}\right) ,\left( {ab}\right) \left( {cd}\right) ,\left( {abcd}\right)\) ，因此在特征 \(p = 3\) 下有四个不可约表示。另一方面，特征 \({\chi }_{1},{\chi }_{2},{\chi }_{4}\) 和 \({\chi }_{5}\) 的约化是不可约的:对于前两个显然如此，它们的度为 1；对于后两个，则由其度是整除群阶的最大 \(p\) 幂次这一事实得出(参见 16.4，命题 46)。由于它们的模特征彼此不同，它们构成了该 \({\mathfrak{S}}_{4}\) 的全部不可约模特征。若记它们为 \({\phi }_{1},{\phi }_{2}\) 、 \({\phi }_{3},{\phi }_{4}\) ，则有下表:

\begin{center}
\adjustbox{max width=\textwidth}{
\begin{tabular}{|c|c|c|c|c|}
\hline
\phantom{X} & 1 & (肛门) & (肛门)(阴道) & (肛门阴道) \\
\cline{1-5}
\({\phi }_{1}\) & 1 & 1 & 1 & 1 \\
\cline{1-5}
\({\phi }_{2}\) & 1 & -1 & 1 & -1 \\
\cline{1-5}
\({\phi }_{3}\) & 3 & 1 & -1 & -1 \\
\cline{1-5}
\({\phi }_{4}\) & 3 & -1 & -1 & 1 \\
\cline{1-5}
\hline
\end{tabular}
}
\end{center}

自从 \({\chi }_{3} = {\phi }_{1} + {\phi }_{2}\) 于 \({\mathrm{G}}_{\text{ reg }}\) 起我们得到如下分解矩阵 D 与卡坦矩阵 C:

\[
\mathrm{D} = \left( \begin{array}{lllll} 1 & 0 & 1 & 0 & 0 \\  0 & 1 & 1 & 0 & 0 \\  0 & 0 & 0 & 1 & 0 \\  0 & 0 & 0 & 0 & 1 \end{array}\right) ,\;\mathrm{C} = \mathrm{D} \cdot  {}^{\prime }\mathrm{D} = \left( \begin{array}{llll} 2 & 1 & 0 & 0 \\  1 & 2 & 0 & 0 \\  0 & 0 & 1 & 0 \\  0 & 0 & 0 & 1 \end{array}\right) ,\det \left( \mathrm{C}\right)  = 3.
\]

投射不可约模的特征 \({\Phi }_{1},\ldots ,{\Phi }_{4}\) 为:

\[
{\Phi }_{1} = {\chi }_{1} + {\chi }_{3}
\]

\[
{\Phi }_{2} = {\chi }_{2} + {\chi }_{3}
\]

\[
{\Phi }_{3} = {\chi }_{4}
\]

\[
{\Phi }_{4} = {\chi }_{5}
\]

(注意 \({\Phi }_{3}\) 与 \({\Phi }_{4}\) 的简单表示，参见命题 46。)

\section*{练习}

18.9. 验证 Fong–Swan 定理对 \({\mathfrak{S}}_{4}\) 的适用性【检验每个 \({\phi }_{i}\) 是否为某个 \({\chi }_{j}\) 限制到 \({\mathrm{G}}_{\text{ reg }}\) 的结果】。

18.10. 证明 \({\mathfrak{S}}_{4}\) 的不可约表示可在素域上实现(任意特征)。

18.11. 群 \({\mathfrak{S}}_{4}\) 有一正规子群 \(\mathrm{N}\) ，阶为 4，使得 \({\mathfrak{S}}_{4}/\mathrm{N}\) 同构于 \({\mathfrak{S}}_{3}\) 。证明在特征 2 下 \(\mathrm{N}\) 在每个 \({\mathfrak{S}}_{4}\) 的不可约表示中作平凡作用。利用此分类此类表示。

\section*{18.6 示例:交替群 \({\mathfrak{A}}_{5}\) 的模特征}

群 \({\mathfrak{A}}_{5}\) 是 \(\{ a,b,c,d,e\}\) 的偶置换群。它有 60 个元素，分为 5 个共轭类:

恒等元 1，

15 个与 \(\left( {ab}\right) \left( {cd}\right)\) 共轭的元素，阶为 2，

20 个与 \(\left( {abc}\right)\) 共轭的元素，阶为 3，

12 个与 \(s = \left( {abcde}\right)\) 共轭的元素，阶为 5，

另有 12 个与 \({s}^{2}\) 共轭的元素，阶为 5。

共有 5 个不可约特征，列于下表:

\begin{center}
\adjustbox{max width=\textwidth}{
\begin{tabular}{|c|c|c|c|c|c|}
\hline
\phantom{X} & 1 & (ab)(cd) & (abc) & \(S\) & \({s}^{2}\) \\
\cline{1-6}
\({\chi }_{1}\) & 1 & 1 & 1 & 1 & 1 \\
\cline{1-6}
\({\chi }_{2}\) & 3 & -1 & 0 & \(z = \frac{1 + \sqrt{5}}{2}\) & \({z}^{\prime }\) \\
\cline{1-6}
\({\chi }_{3}\) & 3 & -1 & 0 & \({z}^{\prime } = \frac{1 - \sqrt{5}}{2}\) & \(z\) \\
\cline{1-6}
\({\chi }_{4}\) & 4 & 0 & 1 & -1 & \phantom{X} \\
\cline{1-6}
\({\chi }_{5}\) & 5 & 1 & -1 & 0 & 0 \\
\cline{1-6}
\hline
\end{tabular}
}
\end{center}

相应的表示为:

\({\chi }_{1}\) :单位表示

\({\chi }_{2}\) 和 \({\chi }_{3}\) :两个 3 次表示，可在域 \(\mathbf{Q}\left( \sqrt{5}\right)\) 上实现，并在 \(\mathbf{Q}\) 上共轭。它们可由下述观察得到: \(\{  \pm  1\}  \times  {\mathfrak{A}}_{5}\) 是带图 o\_3\_o\_5 o 的“考克斯特群”，然后考虑该群的反射表示(参见 Bourbaki, Gr. et Alg. de Lie, Ch. VI, p. 231, ex. 11)。

\({\chi }_{4}\) :一个 4 次表示，可在 \(\mathbf{Q}\) 上实现，由 \({\mathfrak{A}}_{5}\) 在 \(\{ a,b,c,d,e\}\) 上的置换表示中去掉单位表示得到，参见 ex. 2.6。

\({\chi }_{5}\) :一个 5 次表示，可在 \(\mathbf{Q}\) 上实现，由 \({\mathfrak{A}}_{5}\) 在其 6 个 5 阶子群集合上的置换表示中去掉单位表示得到。

我们确定 \(p = 2,3,5\) 的 \({\mathfrak{A}}_{5}\) 的模不可约特征:

(a) 情形 \(p = 2\)

有四个 \(p\) -正规类，因此有 4 个模不可约特征。其中两个显然:单位特征和 \({\chi }_{4}\) 的限制(参见命题 46)。另一方面，我们有

\[
{\chi }_{2} + {\chi }_{3} = 1 + {\chi }_{5}\text{ on }{\mathrm{G}}_{\text{ reg }},
\]

这表明两个 3 次不可约表示的约化都不是不可约的(由于 \(\sqrt{5} \notin  {\mathbf{Q}}_{2}\) ，它们的特征在 2-进数域 \({\mathbf{Q}}_{2}\) 上是共轭的)。每个在 \({\mathbf{R}}_{k}\left( \mathbf{G}\right)\) 中必须分解为单位表示与一個 2 次表示之和，该 2 次表示必然不可约。因此，不可约模特征 \({\phi }_{1},{\phi }_{2},{\phi }_{3},{\phi }_{4}\) 由下表给出:

\begin{center}
\adjustbox{max width=\textwidth}{
\begin{tabular}{|c|c|c|c|c|}
\hline
\phantom{X} & 1 & (abc) & \(S\) & \({s}^{2}\) \\
\cline{1-5}
\({\phi }_{1}\) & 1 & 1 & 1 & 1 \\
\cline{1-5}
\({\phi }_{2}\) & 2 & -1 & \(z - 1\) & \({z}^{\prime } - 1\) \\
\cline{1-5}
\({\phi }_{3}\) & 2 & -1 & \({z}^{\prime } - 1\) & \(z - 1\) \\
\cline{1-5}
\({\phi }_{4}\) & 4 & 1 & -1 & -1 \\
\cline{1-5}
\hline
\end{tabular}
}
\end{center}

我们有

\[
{\chi }_{1} = {\phi }_{1}
\]

\[
\text{ on }{\mathrm{G}}_{\mathrm{{reg}}}
\]

\[
{\chi }_{2} = {\phi }_{1} + {\phi }_{2}
\]

\[
\text{ on }{\mathrm{G}}_{\mathrm{{reg}}}
\]

\[
{\chi }_{3} = {\phi }_{1} + {\phi }_{3}
\]

\[
\text{ on }{\mathrm{G}}_{\mathrm{{reg}}}
\]

\[
{\chi }_{4} = {\phi }_{4}
\]

\[
\text{ on }{\mathrm{G}}_{\text{ reg }}
\]

\[
{\chi }_{5} = {\phi }_{1} + {\phi }_{2} + {\phi }_{3}
\]

\[
\text{ on }{\mathrm{G}}_{\text{ reg }}\text{ . }
\]

矩阵 \(\mathrm{D}\) 与 \(\mathrm{C}\) 的由来:

\[
\mathrm{D} = \left( \begin{array}{lllll} 1 & 1 & 1 & 0 & 1 \\  0 & 1 & 0 & 0 & 1 \\  0 & 0 & 1 & 0 & 1 \\  0 & 0 & 0 & 1 & 0 \end{array}\right) ,\;\mathrm{C} = \left( \begin{array}{llll} 4 & 2 & 2 & 0 \\  2 & 2 & 1 & 0 \\  2 & 1 & 2 & 0 \\  0 & 0 & 0 & 1 \end{array}\right) ,\;\det \left( \mathrm{C}\right)  = 4.
\]

(b) 情形 \(p = 3\)

在特征 3 中可找到 4 个不可约表示，分别是度数为 1、3 和 4 的不可约表示的约化(度数为 3 的有两个)。此外，我们在 \({\mathrm{G}}_{\text{ reg }}\) 上有 \({\chi }_{5} = 1 + {\chi }_{4}\) 。因此:

\[
\mathrm{D} = \left( \begin{array}{lllll} 1 & 0 & 0 & 0 & 1 \\  0 & 1 & 0 & 0 & 0 \\  0 & 0 & 1 & 0 & 0 \\  0 & 0 & 0 & 1 & 1 \end{array}\right) ,\;\mathrm{C} = \left( \begin{array}{llll} 2 & 0 & 0 & 1 \\  0 & 1 & 0 & 0 \\  0 & 0 & 1 & 0 \\  1 & 0 & 0 & 2 \end{array}\right) ,\;\det \left( \mathrm{C}\right)  = 3.
\]

(c) 情形 \(p = 5\)

在特征 5 中有 3 个不可约表示，即度数为 1、3 和 5 的不可约表示的约化(注意两个度数为 3 的表示的约化同构)。此外，我们在 \({\mathrm{G}}_{\text{ reg }}\) 上有 \({\chi }_{4} = {\chi }_{1} + {\chi }_{3}\) 。因此

\[
\mathrm{D} = \left( \begin{array}{lllll} 1 & 0 & 0 & 1 & 0 \\  0 & 1 & 1 & 1 & 0 \\  0 & 0 & 0 & 0 & 1 \end{array}\right) ,\;\mathrm{C} = \left( \begin{array}{lll} 2 & 1 & 0 \\  1 & 3 & 0 \\  0 & 0 & 1 \end{array}\right) ,\;\det \left( \mathrm{C}\right)  = 5.
\]

练习

18.12. 核实断言 (b) 与 (c)。

18.13. 证明 \({\mathfrak{A}}_{5}\) 在特征 2 下的度数 2 的不可约表示可在有 4 个元素的域 \({\mathbf{F}}_{4}\) 上实现；由此得到 \({\mathfrak{A}}_{5}\) 与群 \({\mathbf{{SL}}}_{2}\left( {\mathbf{F}}_{4}\right)\) 的一个同构。

18.14. 证明 \({\mathfrak{A}}_{5}\) 同构于 \({\mathbf{{SL}}}_{2}\left( {\mathbf{F}}_{5}\right) /\{  \pm  1\}\) ，并利用该同构得到 \({\mathfrak{A}}_{5}\) 在特征 5 下的不可约表示列表。

18.15. 证明 \({\chi }_{5}\) 是单项表示(monomial)，而 \({\chi }_{2},{\chi }_{3},{\chi }_{4}\) 不是。

\section*{第 19 章 Artin 表示的应用}

\section*{19.1 Artin 与 Swan 表示}

设 E 是关于离散赋值完备的域，令 F/E 是 E 的有限 Galois 扩张，其 Galois 群为 G，并为简化起见假设 E 与 F 有相同的剩余域。若 \(s \neq  1\) 是 \(G\) 的一个元素且 \(\pi\) 是 \(F\) 的一个素元素，则置

\[
{i}_{\mathrm{G}}\left( s\right)  = {v}_{\mathrm{F}}\left( {s\left( \pi \right)  - \pi }\right) ,
\]

其中 \({v}_{\mathrm{F}}\) 表示对 \(\mathrm{F}\) 的赋值，并归一化使得 \({v}_{\mathrm{F}}\left( \pi \right)  = 1\) 。

置

\[
{a}_{\mathrm{G}}\left( s\right)  =  - {i}_{\mathrm{G}}\left( s\right) \text{ if }s \neq  1
\]

\[
{a}_{\mathrm{G}}\left( 1\right)  = \mathop{\sum }\limits_{{s \neq  1}}{i}_{\mathrm{G}}\left( s\right)
\]

显然 \({v}_{\mathrm{F}}\) 是 \(\mathrm{G}\) 上取整数值的类函数。此外:

定理。函数 \({a}_{\mathrm{G}}\) 是 \(\mathrm{G}\) 的一个表示的特征标(在足够大的域上)。

换言之，若 \(\chi\) 是 \(G\) 的任意特征标，则该数

\[
f\left( \chi \right)  = \left\langle  {{a}_{\mathrm{G}},\chi }\right\rangle
\]

是一个非负整数。

利用 \({a}_{\mathrm{G}}\) 的形式性质(参见 [25], 第六章)，我们看到 \(f\left( \chi \right)  \geq  0\) ，并可容易地将整性问题化归到 \(G\) 为循环群的情形(若愿意，甚至化归到 G 为余域特征的幂阶循环群的情形)。然后我们可以用几种方法进行证明:

(i) 若 \(\chi\) 是 \(G\) 的一维特征标，可证明 \(f\left( \chi \right)\) 与在局部类域论意义下 \(\chi\) 的导出子(导数)模的值相吻合，而该值显然为整数。此方法在有限余域(最初由 Artin 处理)或代数闭余域(利用局部类域论的“几何”类似物)情形均可适用；此外，一般情形亦可由代数闭余域情形容易推得。

(ii) 断言 \(f\left( \chi \right)\) 为整数等价于扩张 F/E 的“分歧数”的某些同余性质。这些性质可直接证明，参见 [25], 第五章 §7，以及 S. Sen, Ann. of Math., 90, 1969, 第33–46页。

现在令 \({r}_{\mathrm{G}}\) 为 \(\mathrm{G}\) 的正则表示的特征标，并取 \({u}_{\mathrm{G}} = {r}_{\mathrm{G}} - 1\) 。令 \(s{w}_{\mathrm{G}} = {a}_{\mathrm{G}} - {u}_{\mathrm{G}}\) 。则

\[
s{w}_{\mathrm{G}}\left( s\right)  = 1 - {i}_{\mathrm{G}}\left( s\right) \text{ if }s \neq  1
\]

\[
s{w}_{\mathrm{G}}\left( 1\right)  = \mathop{\sum }\limits_{{s \neq  1}}\left( {{i}_{\mathrm{G}}\left( s\right)  - 1}\right) .
\]

易检验若 \(\chi\) 为 \(G\) 的特征标，则数乘积 \(\left\langle  {s{w}_{\mathrm{G}},\chi }\right\rangle\) 等于 \(\geq  0\) 。由上述定理，得 \(\left\langle  {s{w}_{\mathrm{G}},\chi }\right\rangle\) 对所有 \(\chi\) 均为非负整数，即 \(s{w}_{\mathrm{G}}\) 为 \(\mathrm{G}\) 的特征标。

特征标 \({a}_{\mathrm{G}}\) (分别为 \(s{w}_{\mathrm{G}}\) )称为 Galois 群 G 的 Artin(分别 Swan)特征标；对应的表示称为 G 的 Artin(分别 Swan)表示。这些表示的显式构造尚未知。尽管如此，我们可以给出 \(g \cdot  {a}_{\mathrm{G}}\) 与 \(g \cdot  s{w}_{\mathrm{G}}\) 的一个简单描述，其中 \(g = \operatorname{Card}\left( \mathrm{G}\right)\) :

令 \({\mathrm{G}}_{i}\left( {i = 0,1,\ldots }\right)\) 表示 \(\mathrm{G}\) 的分歧子群；因此当且仅当 \(s \in  {\mathrm{G}}_{i}\) 或 \({i}_{\mathrm{G}}\left( s\right)  \geq  i + 1\) 或 \(s = 1\) 时成立。置 \(\operatorname{Card}\left( {\mathrm{G}}_{i}\right)  = {g}_{i}\) 。然后可检验

\[
g \cdot  {a}_{\mathrm{G}} = \mathop{\sum }\limits_{{i = 0}}^{\infty }{g}_{i} \cdot  {\operatorname{Ind}}_{{\mathrm{G}}_{i}}^{\mathrm{G}}\left( {u}_{{\mathrm{G}}_{i}}\right)
\]

以及

\[
g \cdot  s{w}_{\mathrm{G}} = \mathop{\sum }\limits_{{i = 1}}^{\infty }{g}_{i} \cdot  {\operatorname{Ind}}_{{\mathrm{G}}_{i}}^{\mathrm{G}}\left( {u}_{{\mathrm{G}}_{i}}\right)
\]

并且 \({u}_{{\mathrm{G}}_{i}} = {r}_{{\mathrm{G}}_{i}} - 1\) 。

特别地，我们有 \(s{w}_{\mathrm{G}} = 0\) 当且仅当 \({\mathrm{G}}_{1} = \{ 1\}\) ，即 \(\mathrm{G}\) 的阶与 \(\mathrm{E}\) 的残特征互素。(换言之， \(s{w}_{\mathrm{G}} = 0\) 当且仅当 \(\mathrm{F}/\mathrm{E}\) 是温和分岔的。)

\section*{19.2 阿廷表示与斯旺表示的有理性}

尽管 \({a}_{\mathrm{G}}\) 与 \(s{w}_{\mathrm{G}}\) 在 \(\mathbf{Z}\) 取值，但可以给出相应表示无法在 \(\mathbf{Q}\) 上实现，甚至无法在 R 上实现的例子(参见 [26]，§4 与 §5)。不过:

定理 44. 设 \(l\) 为与 E 的残特征不同的素数。

(i) 阿廷表示与斯旺表示可在 l-进数域 \({\mathbf{Q}}_{l}\) 上实现。

(ii) 存在一个射影 \({\mathbf{Z}}_{l}\left\lbrack  \mathrm{G}\right\rbrack\) -模 \({\mathrm{{Sw}}}_{\mathrm{G}}\) ，其同构类唯一，使得 \({\mathbf{Q}}_{l} \otimes  {\mathrm{{Sw}}}_{\mathrm{G}}\) 的表征为 \({\mathrm{{sw}}}_{\mathrm{G}}\) 。

只需证明 (ii)；由此可得 (i)，因为 \({a}_{\mathrm{G}}\) 由向 \(s{w}_{\mathrm{G}}\) 添加入可在任意域上实现的 \({u}_{\mathrm{G}}\) 得到。

为此，我们运用命题 44，取 \(p = l,\mathrm{\;K} = {\mathbf{Q}}_{l},n = g = \operatorname{Card}\left( \mathrm{G}\right)\) ，并为 \({\mathrm{K}}^{\prime }\) 选择一个足够大的 \({\mathbf{Q}}_{l}\) 的有限扩张。该命题的条件 (a) 得到满足，参见 19.1。

为验证 (b)，我们使用上述公式

\[
g \cdot  s{w}_{\mathrm{G}} = \mathop{\sum }\limits_{{i \geq  1}}{g}_{i} \cdot  {\operatorname{Ind}}_{{\mathrm{G}}_{i}}^{\mathrm{G}}\left( {u}_{{\mathrm{G}}_{i}}\right)
\]

。按分岔理论，这些 \({\mathrm{G}}_{i}\left( {i \geq  1}\right)\) 的阶与 \(l\) 互素；因此每个 \({\mathrm{A}}^{\prime }\left\lbrack  {\mathrm{G}}_{i}\right\rbrack\) -模都是射影的(参见 15.5)，其中 \({\mathrm{A}}^{\prime }\) 表示 \({\mathrm{K}}^{\prime }\) 的整环。因此 \({u}_{{\mathrm{G}}_{i}}\) 由一个射影 \({A}^{\prime }\left\lbrack  {\mathrm{G}}_{i}\right\rbrack\) -模给予(若愿意可取为射影 \({\mathbf{Z}}_{l}\left\lbrack  {\mathrm{G}}_{i}\right\rbrack\) -模)，相应的诱导 \({\mathrm{A}}^{\prime }\left\lbrack  \mathrm{G}\right\rbrack\) -模亦为射影。取这些模的直和(每个重复 \({g}_{i}\) 次)，我们得到一个表征为 \(g \cdot  s{w}_{\mathrm{G}}\) 的射影 \({\mathrm{A}}^{\prime }\left\lbrack  \mathrm{G}\right\rbrack\) -模。故命题 44 的所有条件均满足，定理得证。

\section*{注记}

(1) 定理 44 的 (i) 在 [26] 中以稍复杂的方法证明，且该方法给出更强的结论:代数 \({\mathbf{Q}}_{l}\left\lbrack  \mathrm{G}\right\rbrack\) 是准分裂的(参见 12.2)。

(2) 也可由 (i) 与 Fong–Swan 定理(定理 38)结合得到 (ii)。

(3) 有些例子中，Artin 与 Swan 表示不能在 \({\mathbf{Q}}_{p}\) 上实现，其中 \(p\) 为 E 的残值特征。然而，J.-M. Fontaine 已指出(参见 [27])，这些表示可在 \({e}_{0}\) 的 Witt 矢量域上实现， \({e}_{0}\) 表示 E 的残值域中最大属于素域的代数子域。

\section*{19.3 一个不变量}

设 \(l\) 为不等于 E 残值特征的素数。记作 \(k = \mathbf{Z}/l\mathbf{Z}\) ，并令 \(\mathbf{M}\) 为一個 \(k\left\lbrack  \mathrm{G}\right\rbrack\) 模。我们通过下式定义 \(b\left( \mathbf{M}\right)\) 为 \(\mathbf{M}\) 的一个不变量

\[
b\left( \mathrm{M}\right)  = \langle \overline{{\mathrm{{Sw}}}_{\mathrm{G}}},\mathrm{M}{\rangle }_{k} = \dim {\operatorname{Hom}}^{\mathrm{G}}\overline{{\mathrm{{Sw}}}_{\mathrm{G}}},\mathrm{M}) = \dim {\operatorname{Hom}}_{{\mathbf{Z}}_{l}\left\lbrack  \mathrm{G}\right\rbrack  }\left( {{\mathrm{{Sw}}}_{\mathrm{G}},\mathrm{M}}\right) ,
\]

其中 \(\overline{{\mathrm{{Sw}}}_{\mathrm{G}}} = {\mathrm{{Sw}}}_{\mathrm{G}}/l \cdot  {\mathrm{{Sw}}}_{\mathrm{G}}\) 表示由定理 44 定义的 \({\mathrm{Z}}_{l}\left\lbrack  \mathrm{G}\right\rbrack\) -模 \({\mathrm{{Sw}}}_{\mathrm{G}}\) 对 \(l\) 的模化约。标量积 \({\left\langle  \overline{{\mathrm{{Sw}}}_{\mathrm{G}}},\mathrm{M}\right\rangle  }_{k}\) 有意义，因为 \(\overline{{\mathrm{{Sw}}}_{\mathrm{G}}}\) 是投射的，参见 14.5。

不变量 \(b\left( \mathrm{M}\right)\) 有下列性质:

(i) 若 \(0 \rightarrow  {\mathrm{M}}^{\prime } \rightarrow  \mathrm{M} \rightarrow  {\mathrm{M}}^{\prime \prime } \rightarrow  0\) 为一段 \(k\left\lbrack  \mathrm{G}\right\rbrack\) -模的正合列，则 \(b\left( \mathrm{M}\right)  = b\left( {\mathrm{M}}^{\prime }\right)  + b\left( {\mathrm{M}}^{\prime \prime }\right)\) 。

(ii) 若 \({\phi }_{M}\) 表示 \(M\) 的模特征，则

\[
b\left( \mathbf{M}\right)  = \left\langle  {s{w}_{\mathrm{G}},{\phi }_{\mathbf{M}}}\right\rangle   = \frac{1}{g}\mathop{\sum }\limits_{{s \in  {\mathrm{G}}_{\text{ reg }}}}s{w}_{\mathrm{G}}\left( {s}^{-1}\right) {\phi }_{\mathbf{M}}\left( s\right) ,
\]

参见 18.1，公式 (viii)。

(iii)

\[
b\left( \mathrm{\;M}\right)  = \mathop{\sum }\limits_{{i = 1}}^{\infty }\frac{{g}_{i}}{g}{\dim }_{k}\left( {\mathrm{\;M}/{\mathrm{M}}^{{\mathrm{G}}_{i}}}\right)
\]

其中 \({\mathrm{M}}^{{\mathrm{G}}_{i}}\) 表示被第 \(i\) 层分馏群 \({\mathrm{G}}_{i}\) 所固定的 \(\mathrm{M}\) 的最大子空间。

(这可由公式 \(g \cdot  s{w}_{\mathrm{G}} = \mathop{\sum }\limits_{{i \geq  l}}{g}_{i}{\operatorname{Ind}}_{{\mathrm{G}}_{i}}^{\mathrm{G}}\left( {u}_{{\mathrm{G}}_{i}}\right)\) 推出，注意到 \(\left\langle  {{\operatorname{Ind}}_{{\mathrm{G}}_{i}}^{\mathrm{G}}\left( {u}_{{\mathrm{G}}_{i}}\right) ,{\phi }_{\mathrm{M}}}\right\rangle\) 等于 \(\left. {{\dim }_{k}\left( {\mathrm{M}/{\mathrm{M}}^{{\mathrm{G}}_{i}}}\right) \text{ if }i \geq  1}\right)\) 。)

(iv) 当且仅当 \({\mathrm{G}}_{1}\) 在 \(\mathrm{M}\) 上平凡作用时有 \(b\left( \mathrm{M}\right)  = 0\) ，即 \(\mathrm{G}\) 在 \(\mathrm{M}\) 上的作用是“温和”的。[此句从 (iii) 得出。]

因此 \(b\left( \mathrm{M}\right)\) 度量了模 M 的“野分馏”。该不变量出现在许多问题中:代数曲线的上同调、ζ 函数的局部因子、椭圆曲线的导子(参见 [28], [29], [30])。

\section*{附录}

\section*{Artinian 环}

若环 A 满足下列等价条件，则称其为阿廷环(参见 Bourbaki, Alg. Ch. VIII, §2):

(a) A 的每个左理想递减序列最终稳定。

(b) 作为左 A-模的 A 具有有限长度。

(c) 每个有限生成的左 A-模具有有限长度。

若 \(A\) 为阿廷环，则其根 \(r\) 为幂零，商环 \(S = A/r\) 是半单的。环 \(S\) 可分解为简单环的直积 \(\prod {S}_{i}\) ；每个 \({S}_{i}\) 同构于 (矩阵) 代数 \({\mathbf{M}}_{{n}_{i}}\left( {D}_{i}\right)\) (上)于(非)交换体 \({\mathrm{D}}_{i}\) ，并拥有唯一的简单模 \({\mathrm{E}}_{i}\) ，该模为维数为 \({n}_{i}\) 的 \({\mathrm{D}}_{i}\) -向量空间。每个半单 A-模都被 \(\mathfrak{r}\) 消去，因此可被视为 S-模；若该模为单模，则同构于上述某个 \({\mathrm{E}}_{i}\) 。

例:在域 \(k\) 上具有有限维的代数是阿廷环；这尤其适用于有限群 \(\mathrm{G}\) 的代数 \(k\left\lbrack  \mathrm{G}\right\rbrack\) 。

\section*{Grothendieck 群}

设 A 为一环， \(\mathcal{F}\) 为左 A-模的一个范畴。 \(\mathfrak{F}\) 的 Grothendieck 群，记为 \(\mathrm{K}\left( \mathfrak{F}\right)\) ，是由下列生成元与关系定义的阿贝尔群:

生成元。对每个 E 属于 \(\in  \mathcal{F}\) ，对应一个生成元 [E]。

关系。对每个短正合列，关联一条关系 \(\left\lbrack  E\right\rbrack   = \left\lbrack  {E}^{\prime }\right\rbrack   + \left\lbrack  {E}^{\prime \prime }\right\rbrack\) 。

\[
0 \rightarrow  \mathrm{E} \rightarrow  {\mathrm{E}}^{\prime } \rightarrow  {\mathrm{E}}^{\prime \prime } \rightarrow  0\text{ where }\mathrm{E},{\mathrm{E}}^{\prime },{\mathrm{E}}^{\prime \prime } \in  \mathcal{F}.
\]

若 \(H\) 是阿贝尔群，则同态 \(f : K\left( \mathcal{F}\right)  \rightarrow  H\) 与满足“加法性”的映射 \(\phi  : \mathcal{F} \rightarrow  H\) 一一对应，即对上述每个正合列有 \(\phi \left( \mathrm{E}\right)  = \phi \left( {\mathrm{E}}^{\prime }\right)  + \phi \left( {\mathrm{E}}^{\prime \prime }\right)\) 。

最常见的两个例子是: \(\mathcal{F}\) 为所有有限生成 A-模的范畴，或为所有有限生成投射 A-模的范畴。

\section*{投射模}

设 A 为一环，P 为左 A-模。若 P 满足下列等价条件，则称其为投射模(参见 Bourbaki, Alg., Ch. II, §2):

(a) 存在一个自由 A-模使得 \(\mathrm{P}\) 为其直和因子。

(b) 对任意满同态 \(f : \mathbf{E} \rightarrow  {\mathbf{E}}^{\prime }\) (左 A-模之间)及任意同态 \({g}^{\prime } : \mathrm{P} \rightarrow  {\mathrm{E}}^{\prime }\) ，存在同态 \(g : \mathrm{P} \rightarrow  \mathrm{E}\) 使得 \({g}^{\prime } = f \circ  g\) 。

(c) 函子 \(E \mapsto  {\operatorname{Hom}}_{A}\left( {P,E}\right)\) 是正合的。

为了使 \(A\) 的左理想 \(a\) 作为模成为 \(A\) 的直和因子，当且仅当存在 \(e \in  A\) ，使得 \({e}^{2} = e\) 和 \(a = {Ae}\) 成立；这样的理想是 A-投射模。

\section*{离散赋值}

设 \(\mathrm{K}\) 为域，令 \({\mathrm{K}}^{ * }\) 为 K 的非零元素乘法群。K 的一个离散赋值(参见 [24])是一个满同态 \(v : {\mathrm{K}}^{ * } \rightarrow  \mathbf{Z}\) ，满足

\[
v\left( {x + y}\right)  \geq  \operatorname{Inf}\left( {v\left( x\right) ,v\left( y\right) }\right) \;\text{ for }x,y \in  {\mathrm{K}}^{ * }.
\]

这里通过规定 \(v\left( 0\right)  =  + \infty\) 将 \(v\) 扩张到 \(\mathrm{K}\) 。

集合 \(\mathrm{A}\) := \{ \(x \in  \mathrm{K}\) | \(v\left( x\right)  \geq  0\) \} 是 \(\mathrm{K}\) 的子环，称为 \(v\) 的赋值环(或 \(\mathrm{K}\) 的整数环)。它有唯一的极大理想，即所有满足 \(v\left( x\right)  \geq  1\) 的元素构成的集合 \(\mathfrak{m}\) 。域 \(k = \mathrm{A}/\mathfrak{m}\) 被称为 \(\mathrm{A}\) (或 \(v\) )的剩余域。

为了使 \(\mathrm{K}\) 在由 \(\mathfrak{m}\) 的幂定义的拓扑下完备，当且仅当 A 到各 \(\mathrm{A}/{\mathfrak{m}}^{n}\) 的乘极限的自然映射为同构。

参考书目:第三部分

关于模表示，见 Curtis 与 Reiner [9] 以及:

[18] R. Brauer. Über die Darstellung von Gruppen in Galoisschen Feldern. Act. Sci. Ind., 195 (1935).

[19] R. Brauer. Zur Darstellungstheorie der Gruppen endlicher Ordnung. Math. Zeit., 63 (1956), p. 406-444.

[20] W. Feit. Representations of Finite Groups I. 耶鲁大学，影印讲义，1969。

关于格罗滕迪克群及其在有限群表示中的应用，见:

[21] R. Swan. Induced representations and projective modules. Ann. of Math., 71 (1960), p. 552-578.

[22] R. Swan. The Grothendieck group of a finite group. Topology, (1963), p. 85-110.

关于投射包(projective envelopes)，见:

[23] M. Demazure 和 P. Gabriel。Groupes algébriques，Tome I，Chapter V，§2，no. 4。Masson and North-Holland，1970。

[24] I. Giorgiutti。Groupes de Grothendieck。Ann. Fac. Sci. Univ. Toulouse，26 (1962)，p. 151-207。

关于局部域以及 Artin 和 Swan 表示，见:

[25] J.-P. Serre。Corps Locaux。Act. Sci. Ind.，1296 (1962)。

[26] J.-P. Serre。Sur la rationalité des représentations d'Artin。Ann. of Math.，72 (1960)，p. 406-420。

[27] J.-M. Fontaine。Groupes de ramification et représentations d'Artin。Ann. Sci. E.N.S.，4 (1971)，p. 337-391。

从 Swan 表示得到的不变量被用于:

[28] M. Raynaud. 关于一个层的欧拉–彭加莱特征与阿贝尔簇的上同调。Bourbaki 讨论会，报告 286，1964/65，W. A. Benjamin 出版社，纽约，1966。

[29] A. P. Ogg. 椭圆曲线与狂野分歧。Amer. J. of Math., 89 (1967), 页 1-21。

[30] J.-P. Serre. 代数簇的 ζ 函数的局部因子(定义与猜想)。Delange-Pisot-Poitou 讨论会，巴黎，1969/70，报告 19。

符号索引

V, GL(V): 1.1

\(\rho ,{\rho }_{s} = \rho \left( s\right)  : {1.1}\)

\({\mathbf{C}}^{ * } = \mathbf{C} - \{ 0\}  : {1.2}\)

\(V = W \oplus  {W}^{\prime } : {1.3}\)

\(g =\) 群的阶:1.3,2.2

\({\mathrm{V}}_{1} \otimes  {\mathrm{V}}_{2},{\rho }_{1} \otimes  {\rho }_{2},{\operatorname{Sym}}^{2}\left( \mathrm{\;V}\right) , \; {\text{ Alt }}^{2}\left( \mathrm{\;V}\right)  : {1.5}\)

\(\operatorname{Tr}\left( a\right)  = \sum {a}_{ii},{\chi }_{\rho }\left( s\right)  = \operatorname{Tr}\left( {\rho }_{s}\right)  : {2.1}\)

\({z}^{ * } = \bar{z} = x - {iy} : {2.1}\)

\({\chi }_{\sigma }^{2},{\chi }_{\alpha }^{2} : {2.1}\)

\({\delta }_{i,j}\left(  = \right. 1\) 若 \(i = j, = 0\) 否则 \() :\) 2.2

\(\langle \varphi ,\psi \rangle  = \left( {1/g}\right) \mathop{\sum }\limits_{{t \in  \mathrm{G}}}\varphi \left( {t}^{-1}\right) \psi \left( t\right)  :\) 2.2

\(\check{\varphi }\left( t\right)  = \varphi {\left( {t}^{-1}\right) }^{ * } : {2.3}\)

\(\left( {\varphi  \mid  \psi }\right)  = \langle \varphi ,\check{\psi }\rangle\)

\(= \left( {1/g}\right) \mathop{\sum }\limits_{{t \in  \mathrm{G}}}\varphi \left( t\right) \psi {\left( t\right) }^{ * } : {2.3}\)

\({\chi }_{1},\ldots ,{\chi }_{h};{n}_{1},\ldots ,{n}_{h};{\mathrm{W}}_{1},\ldots , \; {\mathrm{W}}_{h} : {2.4}\)

\({\mathrm{C}}_{1},\ldots ,{\mathrm{C}}_{k};{c}_{s} : {2.5}\)

\(\mathrm{V} = {\mathrm{V}}_{1} \oplus  \cdots  \oplus  {\mathrm{V}}_{h}\) (正则分解):2.6

\({p}_{i}\) (到 \({\mathrm{V}}_{i}\) 的典范投影):2.6

\({p}_{\alpha \beta } : {2.7}\)

\(\mathrm{G} = {\mathrm{G}}_{1} \times  {\mathrm{G}}_{2} : {3.2}\)

\(\rho ,\theta ,{\chi }_{\rho },{\chi }_{\theta } : {3.3}\)

\(\mathrm{G}/\mathrm{H},s\mathrm{H},\mathrm{R} : {3.3}\)

\({\int }_{\mathrm{G}}f\left( t\right) {dt} : {4.2}\)

\(\left( {\varphi  \mid  \psi }\right)  = {\int }_{\mathrm{G}}\varphi \left( t\right) \psi {\left( t\right) }^{ * }{dt} : {4.2}\)

\({\mathrm{C}}_{n} : {5.1}\)

\({\mathrm{C}}_{\infty } : {5.2}\)

\({\mathrm{D}}_{n},{\mathrm{C}}_{nv} : {5.3}\)

\(\mathrm{I} = \{ 1,\iota \} ;{\mathrm{D}}_{nh} = {\mathrm{D}}_{n} \times  \mathrm{I} : {5.4}\)

\({\chi }_{g},{\chi }_{u} : {5.4}\)

\({\mathrm{D}}_{\infty } : {5.5}\)

\({\mathrm{D}}_{\infty h} = {\mathrm{D}}_{\infty } \times  \mathrm{I} : {5.6}\)

\({\mathfrak{A}}_{4} = \mathrm{H} \cdot  \mathrm{K} : {5.7}\)

\({\mathfrak{S}}_{4} = \mathrm{H} \cdot  \mathrm{L} : {5.8}\)

\(\mathrm{G} = {\mathfrak{S}}_{3} \cdot  \mathrm{M} = {\mathfrak{S}}_{4} \times  \mathrm{I} : {5.9}\)

符号索引

K[G]: 6.1

中心 C[G], \({\omega }_{i}\) : 6.3

\({\operatorname{Ind}}_{\mathbf{H}}^{\mathbf{G}}\left( \mathbf{W}\right)\) , Ind W: 7.1

\({f}^{\prime } = \operatorname{Ind}f = {\operatorname{Ind}}_{\mathbf{H}}^{\mathbf{G}}f : {7.2}\)

Res φ, Res V: 7.2

\(\mathrm{K} \smallsetminus  \mathrm{G}/\mathrm{H},{\mathrm{W}}_{s},{\rho }^{s} : {7.3}\)

\({\theta }_{i,\rho } : {8.2}\)

\({\mathbf{R}}^{ + }\left( \mathbf{G}\right) ,\mathbf{R}\left( \mathbf{G}\right) ,{\mathbf{F}}_{\mathbf{C}\left( \mathbf{G}\right)  : {9.1}}\)

\({\operatorname{Res}}_{\mathrm{H}}^{\mathrm{G}},\operatorname{Res},{\operatorname{Ind}}_{\mathrm{H}}^{\mathrm{G}},\operatorname{Ind} : {9.1}\)

\({\Psi }^{k}\left( f\right) ,{\chi }_{\sigma }^{k},{\chi }_{\lambda }^{k},{\sigma }_{\mathrm{T}}\left( \chi \right) ,{\lambda }_{\mathrm{T}}\left( \chi \right)  : {9.1},\) 例 3

\({\theta }_{\mathrm{A}} : {9.4}\)

\(x = {x}_{r} \cdot  {x}_{u},\mathrm{H} = \mathrm{C} \cdot  \mathrm{P} : {10.1}\)

\(g = {p}^{n}l : {10.2}\)

\({\mathrm{V}}_{p}\) , Ind, A : 10.2

A, \(g,{\Psi }^{n} : {11.2}\)

Spec, \(\mathrm{{Cl}}\left( \mathrm{G}\right) ,{\mathrm{M}}_{c},{\mathrm{P}}_{\mathrm{M},c} : {11.4}\)

\(\mathrm{K},\mathrm{C},{\mathrm{R}}_{\mathrm{K}}\left( \mathrm{G}\right) ,{\overline{\mathrm{R}}}_{\mathrm{K}}\left( \mathrm{G}\right)  : {12.1}\)

\({\mathrm{A}}_{i},{\mathrm{V}}_{i},{\rho }_{i},{\chi }_{i},{\varphi }_{i},{\psi }_{i},{m}_{i} : {12.2}\)

\({\Gamma }_{\mathrm{C}},{\sigma }_{t},{\Psi }^{t} : {12.4}\)

\({\mathrm{X}}_{\mathbf{K}},{\mathrm{X}}_{\mathbf{K}}\left( p\right) ,g = {p}^{n}l,{\mathrm{\;V}}_{\mathbf{K},p} : {12.6},\) 12.7

A, \({\mathfrak{p}}_{i},\mathrm{\;N}\left( x\right)  : {12.7}\)

\(\mathbf{Q}\left( m\right) ,{\mathbf{l}}_{\mathrm{H}}^{\mathrm{G}} : {13.1}\)

\(\mathbf{K},\mathbf{A},\mathfrak{m},p,\mathbf{G},m : {14},\) 符号约定

\({\mathbf{S}}_{\mathbf{K}},{\mathbf{S}}_{k},{\mathbf{R}}_{\mathbf{K}}\left( \mathbf{G}\right) ,\;{\mathbf{R}}_{\mathbf{K}}^{ + }\left( \mathbf{G}\right) ,{\mathbf{R}}_{k}\left( \mathbf{G}\right) , \; {\mathbf{R}}_{k}^{ + }\left( \mathrm{G}\right)  : {14.1}\)

\({\mathbf{P}}_{k}\left( \mathrm{G}\right) ,{\mathbf{P}}_{k}^{ + }\left( \mathrm{G}\right) ,{\mathbf{P}}_{\mathrm{A}}\left( \mathrm{G}\right) ,{\mathbf{P}}_{\mathrm{A}}^{ + }\left( \mathrm{G}\right)  : {14.2}\)

\({\mathrm{P}}_{\mathrm{E}} : {14.3}\)

\(\overline{\mathrm{P}} = \mathrm{P}/\mathrm{{mP}} : {14.4}\)

\(\langle e,f{\rangle }_{\mathbf{K}},\langle e,f{\rangle }_{k} : {14.5}\)

\(c,\mathbf{C},{\mathbf{C}}_{\mathbf{{ST}}} : {15.1}\)

\(d,D,{D}_{FE} : {15.2}\)

\(e,E : {15.3}\)

\({\operatorname{Res}}_{\mathrm{H}}^{\mathrm{G}},{\operatorname{Ind}}_{\mathrm{H}}^{\mathrm{G}} : {17.1}\)

\({\mathrm{G}}_{\text{ reg }},{\mu }_{\mathrm{K}},{\mu }_{k},\widetilde{\lambda },{\varphi }_{\mathrm{E}},{\varphi }_{x},{s}_{\mathrm{E}},{\Phi }_{\mathrm{F}} : {18.1}\)

\({\chi }_{\mathrm{F}}\left( {\mathrm{F} \in  {s}_{\mathrm{K}}}\right) ,{\varphi }_{\mathrm{E}},{\Phi }_{\mathrm{E}}\left( {\mathrm{E} \in  {\mathrm{S}}_{k}}\right)  : {18.3}\)

\({a}_{\mathrm{G}},{i}_{\mathrm{G}},s{w}_{\mathrm{G}},{r}_{\mathrm{G}},{u}_{\mathrm{G}} : {19.1}\)

\({\mathrm{{Sw}}}_{\mathrm{G}} : {19.2}\)

\(b\left( \mathrm{M}\right)  : {19.3}\)

术语索引

绝对不可约(表示):12.1

代数(有限群的):6.1

阿尔廷(表示):19.1

阿尔廷定理:9.2, 12.5, 17.2

阿蒂安(环):附录

关联(将 \(p\) -元素子群 … 与 \({p}^{\prime }\) -元关联):10.1

布劳尔定理(关于提供表示的域):12.3

布劳尔定理(关于诱导特征):10.1, 12.6, 17.2

布劳尔定理(关于模特征):18.2

中心(群代数的):6.3

特征(表示的):2.1

特征(模的):18.1

类函数:2.1, 2.5

\({\Gamma }_{\mathrm{K}}\) -类:12.6

紧致(群):4.1

补空间(向量空间的):1.3

共轭类:2.5

共轭(元素):2.5

\({\Gamma }_{\mathrm{K}}\) -共轭(元素):12.4

分解(表示的规范…):2.6

分解(同态、…矩阵):15.3

次数(表示的):1.1

二面体(群):5.3

直和(两个表示的):1.3

双陪集:7.3

初等(子群):10.5

\({\Gamma }_{\mathrm{K}}\) -初等(子群):12.6

包络(模的射影…):14.3

Fong–Swan(定理):16.3, 17.6

傅里叶(反演公式):6.2

Frobenius(互反公式):7.2

Frobenius(子群):练习 7.3

Frobenius(定理):11.2

Grothendieck(群):附录

Haar(测度):4.2

Higman(定理):练习 6.3

子群的指数:3.1, 3.3

诱导函数:7.2

诱导表示:3.3, 7.1, 17.1

Z-上的整元:6.4

不可约特征:2.3

不可约(模)特征:18.2

不可约表示:1.4

同型分型模/表示:8.1

\({\Gamma }_{\mathrm{K}}\) 类:12.6

\({\Gamma }_{\mathrm{K}}\) 共轭(元素):12.4

\({\Gamma }_{\mathrm{K}}\) -初等， \({\Gamma }_{\mathrm{K}}\) - \(p\) -初等(子群):12.6

Kronecker(直积):1.5

格(K-向量空间的):15.2

左陪集(子群的):3.3

Mackey(不可约性判据):7.4

表示的矩阵形式:2.1

单项式表示:7.1

幂零群:8.3

在 \(\mathbf{Z}\) 上非退化(二次或双线性式):14.5

正交关系(关于特征):2.3

正交关系(关于系数):2.2

\(p\) -分量与 \({p}^{\prime }\) -分量(元素的):10.1

\(p\) -元素， \({p}^{\prime }\) -元素:10.1

\(p\) -初等(子群):10.1

\({\Gamma }_{\mathrm{K}}\) - \(p\) -初等(子群):12.6

p-群:8.3

Plancherel(公式):例 6.2

p-regular(元):10.1

p-regular(共轭类):11.4

p-solvable(群):16.3

直积(两个群的...):3.2

标量乘积:1.3

标量乘积(两个函数的...):2.3

半直积(两个群的...):8.2

张量积(两个表示的...):1.5, 3.2

投影:1.3

射影(模):附录

p-singular(元):16.2

p-unipotent(元):10.1

拟分裂代数:12.2

四元数(群):8.5, 例 8.11

有理(K 上的表示):12.1

约化(模 m):14.4，15.2

表示:1.1，6.1

表示(置换):1.2

表示(正则):1.2

表示(空间):11.1

表示(幺元):1.2

限制(表示的):7.2，9.1，17.1

Schur(指标):12.2

Schur 引理:2.2

单(表示):1.3

可解(群):8.3

谱(交换环的):11.4

分裂(单射):11.1

子表示:1.3

充分大的(域):14，记号

超可解(群):8.3

天鹅(表示论):19.1

西洛(定理):8.4

西洛子群:8.4

对称平方与交替平方(表示的):1.5

迹(endomorphism 的):2.1

赋值(域的离散……):附录

虚(特征):9.1

\section*{Graduate Texts in Mathematics}

每卷至第14卷有精软/精装两种版本，第15卷起仅有精装

1 TAKEUTI/ZARING. 公理集合论导论。vii, 250页。1971。

2 OxtOBY. 测度与范畴。viii, 95页。1971。

3 Schaeffer. 拓扑向量空间。xi, 294页。1971。

4 HILTON/STAMMBACH. 同调代数课程。ix, 338页。1971。(仅精装)

5 MACLANE. 实用数学家的范畴论。ix, 262页。1972。

6 Hughes/PIPER. 投影平面。xii, 291页。1973。

7 SERRE. 算术课程。x, 115页。1973。

8 TAKEUTI/ZARING. 公理化集合论。viii, 238 页。1973。

9 HUMPHREYS. 李代数与表示论导论。xiv, 169 页。1972。

10 COHEN. 单纯同伦理论课程。xii, 114 页。1973。

11 Conway. 单复变函数。第二次更正再版。xiii, 313 页。1975。(仅精装版。)

12 BEALS. 高等数学分析。xi, 230 页。1973。

13 Anderson/Fuller. 环与模的范畴。ix, 339 页。1974。

14 GOLUBITSKY/GUILLEMIN. 稳定映射及其奇点。x, 211 页。1974。

15 BERBERIAN. 泛函分析与算子理论讲义。x, 356 页。1974。

16 WINTER. 域的结构。xiii, 205 页。1974。

17 ROSENBLATT. 随机过程。第二版。x, 228 页。1974。

18 HALMOS. 测度论。xi, 304 页。1974。

19 HALMOS. 希尔伯特空间问题集。xvii, 365 页。1974。

20 Husemoller. 纤维丛。第二版。xvi, 344 页。1975。

21 Humphreys. 线性代数群。xiv, 272 页。1975。

22 BARNES/MACK. 数理逻辑的代数引论。x, 137 页。1975。

23 GREUB. 线性代数。第四版。xvii, 451 页。1975。

24 Holmes。《几何泛函分析及其应用》。x, 246 页。1975。

25 HEWITT/STROMBERG。《实变与抽象分析》。第3次印刷。viii, 476 页。1975。

26 Manes。《代数理论》。x, 356 页。1976。

27 KELLEY。《一般拓扑》。xiv, 298 页。1975。

28 Zariski/Samuel。《交换代数 I》。xi, 329 页。1975。

29 ZARISKI/SAMUEL。《交换代数 II》。x, 414 页。1976。

30 Jacosson。《抽象代数讲义 I:基本概念》。xii, 205 页。1976。

31 JACOBSON。《抽象代数讲义 II:线性代数》。xii, 280 页。1975。

32 JACOBSON。《抽象代数讲义 III:域论与伽罗瓦理论》。ix, 324 页。1976。

33 Hirsch。《微分拓扑》。x, 222 页。1976。

34 SPITZER。《随机游动原理》。第2版。xiii, 408 页。1976。

35 Wermer。《巴拿赫代数与若干复变量》。第2版。xiv, 162 页。1976。

36 KELLEY/NAMIOKA。《线性拓扑空间》。xv, 256 页。1976。

37 Monκ。《数理逻辑》。x, 531 页。1976。

8 GRAUERT/FRITZSCHE。《若干复变量》。viii, 207 页。1976。

39 ARVESON。《C*-代数导引》。x, 106 页。1976。

40 KEMENY/SNELL/KNAPP. 可数马尔可夫链。第2版。xii, 484页。1976。

41 APOSTOL. 模形式与数论中的狄利克雷级数。x, 198页。1976。

42 SERRE. 有限群的线性表示。约176页。1977。

GILLMAN/JERISON. 连续函数环。xiii, 300页。1976。

KENDIG. 初等代数几何。约304页。1977。

45 Loève. 概率论。第4版。第1卷。xvii, 425页。1977。

46 Loève. 概率论。第4版。第2卷。准备中。1977。
\end{document}